\documentclass[12pt]{book}
\usepackage{fontspec}
\usepackage{microtype}
\usepackage{geometry}
\usepackage{hyperref}
\usepackage{graphicx}
\geometry{paperwidth=6.25in,paperheight=9.25in,inner=1.0in,outer=0.6in,top=0.8in,bottom=0.8in}
\setmainfont{Baskerville}
\setlength{\parindent}{0pt}
\setlength{\parskip}{0.6em}
\linespread{1.1}
\raggedbottom
\sloppy
\hfuzz=2pt
\setcounter{secnumdepth}{1}
\setcounter{tocdepth}{1}
\usepackage{titlesec}
\titleformat{\chapter}[display]
  {\normalfont\bfseries}
  {\small\chaptertitlename\ \thechapter}
  {0.5em}
  {\huge}
\usepackage{tocloft}
\renewcommand{\cftchapfont}{\normalfont\small}
\renewcommand{\cftchappagefont}{\normalfont\small}
\renewcommand{\cftchapleader}{\cftdotfill{\cftdotsep}}
\setlength{\cftbeforechapskip}{0.5em}
\usepackage{fancyhdr}
\setlength{\headheight}{20pt}
\pagestyle{fancy}
\fancyhf{}
\fancyhead[LE,RO]{\small\textit{\thepage}}
\fancyhead[RE,LO]{}
\renewcommand{\headrulewidth}{0pt}
\renewcommand{\chaptermark}[1]{}
\makeatletter
\@addtoreset{footnote}{chapter}
\makeatother
\hypersetup{colorlinks=true,linkcolor=black,urlcolor=black}
\providecommand{\pandocbounded}[1]{#1}
\let\oldincludegraphics\includegraphics
\renewcommand{\includegraphics}[2][]{\oldincludegraphics[width=\textwidth,height=4cm,keepaspectratio,#1]{#2}}
\title{Essays, IV}
\author{Paul Graham}
\date{}
\begin{document}
\frontmatter
\maketitle
\tableofcontents
\clearpage
\mainmatter

\chapter{How You Know}

December 2014

I've read Villehardouin's chronicle of the Fourth Crusade at least two
times, maybe three. And yet if I had to write down everything I remember
from it, I doubt it would amount to much more than a page. Multiply this
times several hundred, and I get an uneasy feeling when I look at my
bookshelves. What use is it to read all these books if I remember so
little from them?

A few months ago, as I was reading Constance Reid's excellent biography
of Hilbert, I figured out if not the answer to this question, at least
something that made me feel better about it. She writes:

\begin{quote}
Hilbert had no patience with mathematical lectures which filled the
students with facts but did not teach them how to frame a problem and
solve it. He often used to tell them that ``a perfect formulation of a
problem is already half its solution.''
\end{quote}

That has always seemed to me an important point, and I was even more
convinced of it after hearing it confirmed by Hilbert.

But how had I come to believe in this idea in the first place? A
combination of my own experience and other things I'd read. None of
which I could at that moment remember! And eventually I'd forget that
Hilbert had confirmed it too. But my increased belief in the importance
of this idea would remain something I'd learned from this book, even
after I'd forgotten I'd learned it.

Reading and experience train your model of the world. And even if you
forget the experience or what you read, its effect on your model of the
world persists. Your mind is like a compiled program you've lost the
source of. It works, but you don't know why.

The place to look for what I learned from Villehardouin's chronicle is
not what I remember from it, but my mental models of the crusades,
Venice, medieval culture, siege warfare, and so on. Which doesn't mean I
couldn't have read more attentively, but at least the harvest of reading
is not so miserably small as it might seem.

This is one of those things that seem obvious in retrospect. But it was
a surprise to me and presumably would be to anyone else who felt uneasy
about (apparently) forgetting so much they'd read.

Realizing it does more than make you feel a little better about
forgetting, though. There are specific implications.

For example, reading and experience are usually ``compiled'' at the time
they happen, using the state of your brain at that time. The same book
would get compiled differently at different points in your life. Which
means it is very much worth reading important books multiple times. I
always used to feel some misgivings about rereading books. I
unconsciously lumped reading together with work like carpentry, where
having to do something again is a sign you did it wrong the first time.
Whereas now the phrase ``already read'' seems almost ill-formed.

Intriguingly, this implication isn't limited to books. Technology will
increasingly make it possible to relive our experiences. When people do
that today it's usually to enjoy them again (e.g.~when looking at
pictures of a trip) or to find the origin of some bug in their compiled
code (e.g.~when Stephen Fry succeeded in remembering the childhood
trauma that prevented him from singing). But as technologies for
recording and playing back your life improve, it may become common for
people to relive experiences without any goal in mind, simply to learn
from them again as one might when rereading a book.

Eventually we may be able not just to play back experiences but also to
index and even edit them. So although not knowing how you know things
may seem part of being human, it may not be.

\textbf{Thanks} to Sam Altman, Jessica Livingston, and Robert Morris for
reading drafts of this.
\chapter{How to Be an Expert in a Changing World}

December 2014

If the world were static, we could have monotonically increasing
confidence in our beliefs. The more (and more varied) experience a
belief survived, the less likely it would be false. Most people
implicitly believe something like this about their opinions. And they're
justified in doing so with opinions about things that don't change much,
like human nature. But you can't trust your opinions in the same way
about things that change, which could include practically everything
else.

When experts are wrong, it's often because they're experts on an earlier
version of the world.

Is it possible to avoid that? Can you protect yourself against obsolete
beliefs? To some extent, yes. I spent almost a decade investing in early
stage startups, and curiously enough protecting yourself against
obsolete beliefs is exactly what you have to do to succeed as a startup
investor. Most really good startup ideas look like bad ideas at first,
and many of those look bad specifically because some change in the world
just switched them from bad to good. I spent a lot of time learning to
recognize such ideas, and the techniques I used may be applicable to
ideas in general.

The first step is to have an explicit belief in change. People who fall
victim to a monotonically increasing confidence in their opinions are
implicitly concluding the world is static. If you consciously remind
yourself it isn't, you start to look for change.

Where should one look for it? Beyond the moderately useful
generalization that human nature doesn't change much, the unfortunate
fact is that change is hard to predict. This is largely a tautology but
worth remembering all the same: change that matters usually comes from
an unforeseen quarter.

So I don't even try to predict it. When I get asked in interviews to
predict the future, I always have to struggle to come up with something
plausible-sounding on the fly, like a student who hasn't prepared for an
exam. \footnote{My usual trick is to talk about aspects of the present that most
people haven't noticed yet.} But it's not out of laziness that I
haven't prepared. It seems to me that beliefs about the future are so
rarely correct that they usually aren't worth the extra rigidity they
impose, and that the best strategy is simply to be aggressively
open-minded. Instead of trying to point yourself in the right direction,
admit you have no idea what the right direction is, and try instead to
be super sensitive to the winds of change.

It's ok to have working hypotheses, even though they may constrain you a
bit, because they also motivate you. It's exciting to chase things and
exciting to try to guess answers. But you have to be disciplined about
not letting your hypotheses harden into anything more.
\footnote{Especially if they become well enough known that people start to
identify them with you. You have to be extra skeptical about things you
want to believe, and once a hypothesis starts to be identified with you,
it will almost certainly start to be in that category.}

I believe this passive m.o. works not just for evaluating new ideas but
also for having them. The way to come up with new ideas is not to try
explicitly to, but to try to solve problems and simply not discount
weird hunches you have in the process.

The winds of change originate in the unconscious minds of domain
experts. If you're sufficiently expert in a field, any weird idea or
apparently irrelevant question that occurs to you is ipso facto worth
exploring. \footnote{In practice ``sufficiently expert'' doesn't require one to be
recognized as an expert---which is a trailing indicator in any case. In
many fields a year of focused work plus caring a lot would be enough.} Within Y Combinator, when an idea is
described as crazy, it's a compliment---in fact, on average probably a
higher compliment than when an idea is described as good.

Startup investors have extraordinary incentives for correcting obsolete
beliefs. If they can realize before other investors that some apparently
unpromising startup isn't, they can make a huge amount of money. But the
incentives are more than just financial. Investors' opinions are
explicitly tested: startups come to them and they have to say yes or no,
and then, fairly quickly, they learn whether they guessed right. The
investors who say no to a Google (and there were several) will remember
it for the rest of their lives.

Anyone who must in some sense bet on ideas rather than merely commenting
on them has similar incentives. Which means anyone who wants such
incentives can have them, by turning their comments into bets: if you
write about a topic in some fairly durable and public form, you'll find
you worry much more about getting things right than most people would in
a casual conversation. \footnote{Though they are public and persist indefinitely, comments on
e.g.~forums and places like Twitter seem empirically to work like casual
conversation. The threshold may be whether what you write has a title.}

Another trick I've found to protect myself against obsolete beliefs is
to focus initially on people rather than ideas. Though the nature of
future discoveries is hard to predict, I've found I can predict quite
well what sort of people will make them. Good new ideas come from
earnest, energetic, independent-minded people.

Betting on people over ideas saved me countless times as an investor. We
thought Airbnb was a bad idea, for example. But we could tell the
founders were earnest, energetic, and independent-minded. (Indeed,
almost pathologically so.) So we suspended disbelief and funded them.

This too seems a technique that should be generally applicable. Surround
yourself with the sort of people new ideas come from. If you want to
notice quickly when your beliefs become obsolete, you can't do better
than to be friends with the people whose discoveries will make them so.

It's hard enough already not to become the prisoner of your own
expertise, but it will only get harder, because change is accelerating.
That's not a recent trend; change has been accelerating since the
paleolithic era. Ideas beget ideas. I don't expect that to change. But I
could be wrong.


\textbf{Thanks} to Sam Altman, Patrick Collison, and Robert Morris for
reading drafts of this.
\chapter{Let the Other 95\% of Great Programmers In}

December 2014

American technology companies want the government to make immigration
easier because they say they can't find enough programmers in the US.
Anti-immigration people say that instead of letting foreigners take
these jobs, we should train more Americans to be programmers. Who's
right?

The technology companies are right. What the anti-immigration people
don't understand is that there is a huge variation in ability between
competent programmers and exceptional ones, and while you can train
people to be competent, you can't train them to be exceptional.
Exceptional programmers have an aptitude for and
\href{genius.html}{interest in} programming that is not merely the
product of training. \footnote{How much better is a great programmer than an ordinary one? So
much better that you can't even measure the difference directly. A great
programmer doesn't merely do the same work faster. A great programmer
will invent things an ordinary programmer would never even think of.
This doesn't mean a great programmer is infinitely more valuable,
because any invention has a finite market value. But it's easy to
imagine cases where a great programmer might invent things worth 100x or
even 1000x an average programmer's salary.}

The US has less than 5\% of the world's population. Which means if the
qualities that make someone a great programmer are evenly distributed,
95\% of great programmers are born outside the US.

The anti-immigration people have to invent some explanation to account
for all the effort technology companies have expended trying to make
immigration easier. So they claim it's because they want to drive down
salaries. But if you talk to startups, you find practically every one
over a certain size has gone through legal contortions to get
programmers into the US, where they then paid them the same as they'd
have paid an American. Why would they go to extra trouble to get
programmers for the same price? The only explanation is that they're
telling the truth: there are just not enough great programmers to go
around. \footnote{There are a handful of consulting firms that rent out big pools
of foreign programmers they bring in on H1-B visas. By all means crack
down on these. It should be easy to write legislation that distinguishes
them, because they are so different from technology companies. But it is
dishonest of the anti-immigration people to claim that companies like
Google and Facebook are driven by the same motives. An influx of
inexpensive but mediocre programmers is the last thing they'd want; it
would destroy them.}

I asked the CEO of a startup with about 70 programmers how many more
he'd hire if he could get all the great programmers he wanted. He said
``We'd hire 30 tomorrow morning.'' And this is one of the hot startups
that always win recruiting battles. It's the same all over Silicon
Valley. Startups are that constrained for talent.

It would be great if more Americans were trained as programmers, but no
amount of training can flip a ratio as overwhelming as 95 to 5.
Especially since programmers are being trained in other countries too.
Barring some cataclysm, it will always be true that most great
programmers are born outside the US. It will always be true that most
people who are great at anything are born outside the US.
\footnote{Though this essay talks about programmers, the group of people
we need to import is broader, ranging from designers to programmers to
electrical engineers. The best one could do as a general term might be
``digital talent.'' It seemed better to make the argument a little too
narrow than to confuse everyone with a neologism.}

Exceptional performance implies immigration. A country with only a few
percent of the world's population will be exceptional in some field only
if there are a lot of immigrants working in it.

But this whole discussion has taken something for granted: that if we
let more great programmers into the US, they'll want to come. That's
true now, and we don't realize how lucky we are that it is. If we want
to keep this option open, the best way to do it is to take advantage of
it: the more of the world's great programmers are here, the more the
rest will want to come here.

And if we don't, the US could be seriously fucked. I realize that's
strong language, but the people dithering about this don't seem to
realize the power of the forces at work here. Technology gives the best
programmers huge leverage. The world market in programmers seems to be
becoming dramatically more liquid. And since good people like good
colleagues, that means the best programmers could collect in just a few
hubs. Maybe mostly in one hub.

What if most of the great programmers collected in one hub, and it
wasn't here? That scenario may seem unlikely now, but it won't be if
things change as much in the next 50 years as they did in the last 50.

We have the potential to ensure that the US remains a technology
superpower just by letting in a few thousand great programmers a year.
What a colossal mistake it would be to let that opportunity slip. It
could easily be the defining mistake this generation of American
politicians later become famous for. And unlike other potential mistakes
on that scale, it costs nothing to fix.

So please, get on with it.


\textbf{Thanks} to Sam Altman, John Collison, Patrick Collison, Jessica
Livingston, Geoff Ralston, Fred Wilson, and Qasar Younis for reading
drafts of this.
\chapter{Don't Talk to Corp Dev}

January 2015

Corporate Development, aka corp dev, is the group within companies that
buys other companies. If you're talking to someone from corp dev, that's
why, whether you realize it yet or not.

It's usually a mistake to talk to corp dev unless (a) you want to sell
your company right now and (b) you're sufficiently likely to get an
offer at an acceptable price. In practice that means startups should
only talk to corp dev when they're either doing really well or really
badly. If you're doing really badly, meaning the company is about to
die, you may as well talk to them, because you have nothing to lose. And
if you're doing really well, you can safely talk to them, because you
both know the price will have to be high, and if they show the slightest
sign of wasting your time, you'll be confident enough to tell them to
get lost.

The danger is to companies in the middle. Particularly to young
companies that are growing fast, but haven't been doing it for long
enough to have grown big yet. It's usually a mistake for a promising
company less than a year old even to talk to corp dev.

But it's a mistake founders constantly make. When someone from corp dev
wants to meet, the founders tell themselves they should at least find
out what they want. Besides, they don't want to offend Big Company by
refusing to meet.

Well, I'll tell you what they want. They want to talk about buying you.
That's what the title ``corp dev'' means. So before agreeing to meet
with someone from corp dev, ask yourselves, ``Do we want to sell the
company right now?'' And if the answer is no, tell them ``Sorry, but
we're focusing on growing the company.'' They won't be offended. And
certainly the founders of Big Company won't be offended. If anything
they'll think more highly of you. You'll remind them of themselves. They
didn't sell either; that's why they're in a position now to buy other
companies. \footnote{I'm not saying you should never sell. I'm saying you should be
clear in your own mind about whether you want to sell or not, and not be
led by manipulation or wishful thinking into trying to sell earlier than
you otherwise would have.}

Most founders who get contacted by corp dev already know what it means.
And yet even when they know what corp dev does and know they don't want
to sell, they take the meeting. Why do they do it? The same mix of
denial and wishful thinking that underlies most mistakes founders make.
It's flattering to talk to someone who wants to buy you. And who knows,
maybe their offer will be surprisingly high. You should at least see
what it is, right?

No.~If they were going to send you an offer immediately by email, sure,
you might as well open it. But that is not how conversations with corp
dev work. If you get an offer at all, it will be at the end of a long
and unbelievably distracting process. And if the offer is surprising, it
will be surprisingly low.

Distractions are the thing you can least afford in a startup. And
conversations with corp dev are the worst sort of distraction, because
as well as consuming your \href{top.html}{attention} they undermine your
morale. One of the tricks to surviving a grueling process is not to stop
and think how tired you are. Instead you get into a sort of flow.
\footnote{In a startup, as in most competitive sports, the task at hand
almost does this for you; you're too busy to feel tired. But when you
lose that protection, e.g.~at the final whistle, the fatigue hits you
like a wave. To talk to corp dev is to let yourself feel it mid-game.} Imagine what it would do to you if at mile 20 of
a marathon, someone ran up beside you and said ``You must feel really
tired. Would you like to stop and take a rest?'' Conversations with corp
dev are like that but worse, because the suggestion of stopping gets
combined in your mind with the imaginary high price you think they'll
offer.

And then you're really in trouble. If they can, corp dev people like to
turn the tables on you. They like to get you to the point where you're
trying to convince them to buy instead of them trying to convince you to
sell. And surprisingly often they succeed.

This is a very slippery slope, greased with some of the most powerful
forces that can work on founders' minds, and attended by an experienced
professional whose full time job is to push you down it.

Their tactics in pushing you down that slope are usually fairly brutal.
Corp dev people's whole job is to buy companies, and they don't even get
to choose which. The only way their performance is measured is by how
cheaply they can buy you, and the more ambitious ones will stop at
nothing to achieve that. For example, they'll almost always start with a
lowball offer, just to see if you'll take it. Even if you don't, a low
initial offer will demoralize you and make you easier to manipulate.

And that is the most innocent of their tactics. Just wait till you've
agreed on a price and think you have a done deal, and then they come
back and say their boss has vetoed the deal and won't do it for more
than half the agreed upon price. Happens all the time. If you think
investors can behave badly, it's nothing compared to what corp dev
people can do. Even corp dev people at companies that are otherwise
benevolent.

I remember once complaining to a friend at Google about some nasty trick
their corp dev people had pulled on a YC startup.

``What happened to Don't be Evil?'' I asked.

``I don't think corp dev got the memo,'' he replied.

The tactics you encounter in M\&A conversations can be like nothing
you've experienced in the otherwise comparatively
\href{mean.html}{upstanding} world of Silicon Valley. It's as if a chunk
of genetic material from the old-fashioned robber baron business world
got incorporated into the startup world. \footnote{To be fair, the apparent misdeeds of corp dev people are
magnified by the fact that they function as the face of a large
organization that often doesn't know its own mind. Acquirers can be
surprisingly indecisive about acquisitions, and their flakiness is
indistinguishable from dishonesty by the time it filters down to you.}

The simplest way to protect yourself is to use the trick that John D.
Rockefeller, whose grandfather was an alcoholic, used to protect himself
from becoming one. He once told a Sunday school class

\begin{quote}
Boys, do you know why I never became a drunkard? Because I never took
the first drink.
\end{quote}

Do you want to sell your company right now? Not eventually, right now.
If not, just don't take the first meeting. They won't be offended. And
you in turn will be guaranteed to be spared one of the worst experiences
that can happen to a startup.

If you do want to sell, there's another set of
\href{https://justinkan.com/the-founders-guide-to-selling-your-company-a1b2025c9481}{techniques}
for doing that. But the biggest mistake founders make in dealing with
corp dev is not doing a bad job of talking to them when they're ready
to, but talking to them before they are. So if you remember only the
title of this essay, you already know most of what you need to know
about M\&A in the first year.


\textbf{Thanks} to Marc Andreessen, Jessica Livingston, Geoff Ralston,
and Qasar Younis for reading drafts of this.
\chapter{What Doesn't Seem Like Work?}

January 2015

My father is a mathematician. For most of my childhood he worked for
Westinghouse, modelling nuclear reactors.

He was one of those lucky people who know early on what they want to do.
When you talk to him about his childhood, there's a clear watershed at
about age 12, when he ``got interested in maths.''

He grew up in the small Welsh seacoast town of
\href{https://goo.gl/maps/rkzUm}{Pwllheli}. As we retraced his walk to
school on Google Street View, he said that it had been nice growing up
in the country.

``Didn't it get boring when you got to be about 15?'' I asked.

``No,'' he said, ``by then I was interested in maths.''

In another conversation he told me that what he really liked was solving
problems. To me the exercises at the end of each chapter in a math
textbook represent work, or at best a way to reinforce what you learned
in that chapter. To him the problems were the reward. The text of each
chapter was just some advice about solving them. He said that as soon as
he got a new textbook he'd immediately work out all the problems --- to
the slight annoyance of his teacher, since the class was supposed to
work through the book gradually.

Few people know so early or so certainly what they want to work on. But
talking to my father reminded me of a heuristic the rest of us can use.
If something that seems like work to other people doesn't seem like work
to you, that's something you're well suited for. For example, a lot of
programmers I know, including me, actually like debugging. It's not
something people tend to volunteer; one likes it the way one likes
popping zits. But you may have to like debugging to like programming,
considering the degree to which programming consists of it.

The stranger your tastes seem to other people, the stronger evidence
they probably are of what you should do. When I was in college I used to
write papers for my friends. It was quite interesting to write a paper
for a class I wasn't taking. Plus they were always so relieved.

It seemed curious that the same task could be painful to one person and
pleasant to another, but I didn't realize at the time what this
imbalance implied, because I wasn't looking for it. I didn't realize how
hard it can be to decide what you should work on, and that you sometimes
have to \href{love.html}{figure it out} from subtle clues, like a
detective solving a case in a mystery novel. So I bet it would help a
lot of people to ask themselves about this explicitly. What seems like
work to other people that doesn't seem like work to you?

\textbf{Thanks} to Sam Altman, Trevor Blackwell, Jessica Livingston,
Robert Morris, and my father for reading drafts of this.
\chapter{The Ronco Principle}

January 2015

No one, VC or angel, has invested in more of the top startups than Ron
Conway. He knows what happened in every deal in the Valley, half the
time because he arranged it.

And yet he's a super nice guy. In fact, nice is not the word. Ronco is
good. I know of zero instances in which he has behaved badly. It's hard
even to imagine.

When I first came to Silicon Valley I thought ``How lucky that someone
so powerful is so benevolent.'' But gradually I realized it wasn't luck.
It was by being benevolent that Ronco became so powerful. All the deals
he gets to invest in come to him through referrals. Google did. Facebook
did. Twitter was a referral from Evan Williams himself. And the reason
so many people refer deals to him is that he's proven himself to be a
good guy.

Good does not mean being a pushover. I would not want to face an angry
Ronco. But if Ron's angry at you, it's because you did something wrong.
Ron is so old school he's Old Testament. He will smite you in his just
wrath, but there's no malice in it.

In almost every domain there are advantages to seeming good. It makes
people trust you. But actually being good is an expensive way to seem
good. To an amoral person it might seem to be overkill.

In some fields it might be, but apparently not in the startup world.
Though plenty of investors are jerks, there is a clear trend among them:
the most successful investors are also the most upstanding.
\footnote{I'm not saying that if you sort investors by benevolence you've
also sorted them by returns, but rather that if you do a scatterplot
with benevolence on the x axis and returns on the y, you'd see a clear
upward trend.}

It was not always this way. I would not feel confident saying that about
investors twenty years ago.

What changed? The startup world became more transparent and more
unpredictable. Both make it harder to seem good without actually being
good.

It's obvious why transparency has that effect. When an investor
maltreats a founder now, it gets out. Maybe not all the way to the
press, but other founders hear about it, and that investor starts to
lose deals. \footnote{Y Combinator in particular, because it aggregates data from so
many startups, has a pretty comprehensive view of investor behavior.}

The effect of unpredictability is more subtle. It increases the work of
being inconsistent. If you're going to be two-faced, you have to know
who you should be nice to and who you can get away with being nasty to.
In the startup world, things change so rapidly that you can't tell. The
random college kid you talk to today might in a couple years be the CEO
of the hottest startup in the Valley. If you can't tell who to be nice
to, you have to be nice to everyone. And probably the only people who
can manage that are the people who are genuinely good.

In a sufficiently connected and unpredictable world, you can't seem good
without being good.

As often happens, Ron discovered how to be the investor of the future by
accident. He didn't foresee the future of startup investing, realize it
would pay to be upstanding, and force himself to behave that way. It
would feel unnatural to him to behave any other way. He was already
\href{startupideas.html}{living in the future}.

Fortunately that future is not limited to the startup world. The startup
world is more transparent and unpredictable than most, but almost
everywhere the trend is in that direction.


\textbf{Thanks} to Sam Altman and Jessica Livingston for reading drafts
of this.
\chapter{What Microsoft Is this the Altair Basic of?}

February 2015

One of the most valuable exercises you can try if you want to understand
startups is to look at the most successful companies and explain why
they were not as lame as they seemed when they first launched. Because
they practically all seemed lame at first. Not just small, lame. Not
just the first step up a big mountain. More like the first step into a
swamp.

A Basic interpreter for the Altair? How could that ever grow into a
giant company? People sleeping on airbeds in strangers' apartments? A
web site for college students to stalk one another? A wimpy little
single-board computer for hobbyists that used a TV as a monitor? A new
search engine, when there were already about 10, and they were all
trying to de-emphasize search? These ideas didn't just seem small. They
seemed wrong. They were the kind of ideas you could not merely ignore,
but ridicule.

Often the founders themselves didn't know why their ideas were
promising. They were attracted to these ideas by instinct, because they
were \href{startupideas.html}{living in the future} and they sensed that
something was missing. But they could not have put into words exactly
how their ugly ducklings were going to grow into big, beautiful swans.

Most people's first impulse when they hear about a lame-sounding new
startup idea is to make fun of it. Even a lot of people who should know
better.

When I encounter a startup with a lame-sounding idea, I ask ``What
Microsoft is this the Altair Basic of?'' Now it's a puzzle, and the
burden is on me to solve it. Sometimes I can't think of an answer,
especially when the idea is a made-up one. But it's remarkable how often
there does turn out to be an answer. Often it's one the founders
themselves hadn't seen yet.

Intriguingly, there are sometimes multiple answers. I talked to a
startup a few days ago that could grow into 3 distinct Microsofts.
They'd probably vary in size by orders of magnitude. But you can never
predict how big a Microsoft is going to be, so in cases like that I
encourage founders to follow whichever path is most immediately exciting
to them. Their instincts got them this far. Why stop now?
\chapter{Change Your Name}

August 2015

If you have a US startup called X and you don't have x.com, you should
probably change your name.

The reason is not just that people can't find you. For companies with
mobile apps, especially, having the right domain name is not as critical
as it used to be for getting users. The problem with not having the .com
of your name is that it signals weakness. Unless you're so big that your
reputation precedes you, a marginal domain suggests you're a marginal
company. Whereas (as Stripe shows) having x.com signals strength even if
it has no relation to what you do.

Even good founders can be in denial about this. Their denial derives
from two very powerful forces: identity, and lack of imagination.

X is what we \emph{are}, founders think. There's no other name as good.
Both of which are false.

You can fix the first by stepping back from the problem. Imagine you'd
called your company something else. If you had, surely you'd be just as
attached to that name as you are to your current one. The idea of
switching to your current name would seem repellent.
\footnote{Incidentally, this thought experiment works for
\href{identity.html}{nationality and religion} too.}

There's nothing intrinsically great about your current name. Nearly all
your attachment to it comes from it being attached to you.
\footnote{The liking you have for a name that has become part of your
identity manifests itself not directly, which would be easy to discount,
but as a collection of specious beliefs about its intrinsic qualities.
(This too is true of nationality and religion as well.)}

The way to neutralize the second source of denial, your inability to
think of other potential names, is to acknowledge that you're bad at
naming. Naming is a completely separate skill from those you need to be
a good founder. You can be a great startup founder but hopeless at
thinking of names for your company.

Once you acknowledge that, you stop believing there is nothing else you
could be called. There are lots of other potential names that are as
good or better; you just can't think of them.

How do you find them? One answer is the default way to solve problems
you're bad at: find someone else who can think of names. But with
company names there is another possible approach. It turns out almost
any word or word pair that is not an obviously bad name is a
sufficiently good one, and the number of such domains is so large that
you can find plenty that are cheap or even untaken. So make a list and
try to buy some. That's what
\href{http://www.quora.com/How-did-Stripe-come-up-with-its-name?share=1}{Stripe}
did. (Their search also turned up parse.com, which their friends at
Parse took.)

The reason I know that naming companies is a distinct skill orthogonal
to the others you need in a startup is that I happen to have it. Back
when I was running YC and did more office hours with startups, I would
often help them find new names. 80\% of the time we could find at least
one good name in a 20 minute office hour slot.

Now when I do office hours I have to focus on more important questions,
like what the company is doing. I tell them when they need to change
their name. But I know the power of the forces that have them in their
grip, so I know most won't listen. \footnote{Sometimes founders know it's a problem that they don't have the
.com of their name, but delusion strikes a step later in the belief that
they'll be able to buy it despite having no evidence it's for sale.
Don't believe a domain is for sale unless the owner has already told you
an asking price.}

There are of course examples of startups that have succeeded without
having the .com of their name. There are startups that have succeeded
despite any number of different mistakes. But this mistake is less
excusable than most. It's something that can be fixed in a couple days
if you have sufficient discipline to acknowledge the problem.

100\% of the top 20 YC companies by valuation have the .com of their
name. 94\% of the top 50 do. But only 66\% of companies in the current
batch have the .com of their name. Which suggests there are lessons
ahead for most of the rest, one way or another.


\textbf{Thanks} to Sam Altman, Jessica Livingston, and Geoff Ralston for
reading drafts of this.
\chapter{Why It's Safe for Founders to Be Nice}

August 2015

I recently got an email from a founder that helped me understand
something important: why it's safe for startup founders to be nice
people.

I grew up with a cartoon idea of a very successful businessman (in the
cartoon it was always a man): a rapacious, cigar-smoking, table-thumping
guy in his fifties who wins by exercising power, and isn't too fussy
about how. As I've written before, one of the things that has surprised
me most about startups is \href{mean.html}{how few} of the most
successful founders are like that. Maybe successful people in other
industries are; I don't know; but not startup founders.
\footnote{Many think successful startup founders are driven by money. In
fact the secret weapon of the most successful founders is that they
aren't. If they were, they'd have taken one of the acquisition offers
that every fast-growing startup gets on the way up. What drives the most
successful founders is the same thing that drives most people who make
things: the company is their project.}

I knew this empirically, but I never saw the math of why till I got this
founder's email. In it he said he worried that he was fundamentally
soft-hearted and tended to give away too much for free. He thought
perhaps he needed ``a little dose of sociopath-ness.''

I told him not to worry about it, because so long as he built something
good enough to spread by word of mouth, he'd have a superlinear growth
curve. If he was bad at extracting money from people, at worst this
curve would be some constant multiple less than 1 of what it might have
been. But a constant multiple of any curve is exactly the same shape.
The numbers on the Y axis are smaller, but the curve is just as steep,
and when anything grows at the rate of a successful startup, the Y axis
will take care of itself.

Some examples will make this clear. Suppose your company is making
\$1000 a month now, and you've made something so great that it's growing
at 5\% a week. Two years from now, you'll be making about \$160k a
month.

Now suppose you're so un-rapacious that you only extract half as much
from your users as you could. That means two years later you'll be
making \$80k a month instead of \$160k. How far behind are you? How long
will it take to catch up with where you'd have been if you were
extracting every penny? A mere 15 weeks. After two years, the
un-rapacious founder is only 3.5 months behind the rapacious one.
\footnote{In fact since 2 ≈ 1.05 \^{} 15, the un-rapacious founder is
always 15 weeks behind the rapacious one.}

If you're going to optimize a number, the one to choose is your
\href{growth.html}{growth rate}. Suppose as before that you only extract
half as much from users as you could, but that you're able to grow 6\% a
week instead of 5\%. Now how are you doing compared to the rapacious
founder after two years? You're already ahead---\$214k a month versus
\$160k---and pulling away fast. In another year you'll be making \$4.4
million a month to the rapacious founder's \$2 million.

Obviously one case where it would help to be rapacious is when growth
depends on that. What makes startups different is that usually it
doesn't. Startups usually win by making something so great that people
recommend it to their friends. And being rapacious not only doesn't help
you do that, but probably hurts. \footnote{The other reason it might help to be good at squeezing money out
of customers is that startups usually lose money at first, and making
more per customer makes it easier to get to profitability before your
initial funding runs out. But while it is very common for startups to
\href{pinch.html}{die} from running through their initial funding and
then being unable to raise more, the underlying cause is usually slow
growth or excessive spending rather than insufficient effort to extract
money from existing customers.}

The reason startup founders can safely be nice is that making great
things is compounded, and rapacity isn't.

So if you're a founder, here's a deal you can make with yourself that
will both make you happy and make your company successful. Tell yourself
you can be as nice as you want, so long as you work hard on your growth
rate to compensate. Most successful startups make that tradeoff
unconsciously. Maybe if you do it consciously you'll do it even better.


\textbf{Thanks} to Sam Altman, Harj Taggar, Jessica Livingston, and
Geoff Ralston for reading drafts of this, and to Randall Bennett for
being such a nice guy.
\chapter{Default Alive or Default Dead?}

October 2015

When I talk to a startup that's been operating for more than 8 or 9
months, the first thing I want to know is almost always the same.
Assuming their expenses remain constant and their revenue growth is what
it has been over the last several months, do they make it to
profitability on the money they have left? Or to put it more
dramatically, by default do they live or die?

The startling thing is how often the founders themselves don't know.
Half the founders I talk to don't know whether they're default alive or
default dead.

If you're among that number, Trevor Blackwell has made a handy
\href{http://growth.tlb.org/\#}{calculator} you can use to find out.

The reason I want to know first whether a startup is default alive or
default dead is that the rest of the conversation depends on the answer.
If the company is default alive, we can talk about ambitious new things
they could do. If it's default dead, we probably need to talk about how
to save it. We know the current trajectory ends badly. How can they get
off that trajectory?

Why do so few founders know whether they're default alive or default
dead? Mainly, I think, because they're not used to asking that. It's not
a question that makes sense to ask early on, any more than it makes
sense to ask a 3 year old how he plans to support himself. But as the
company grows older, the question switches from meaningless to critical.
That kind of switch often takes people by surprise.

I propose the following solution: instead of starting to ask too late
whether you're default alive or default dead, start asking too early.
It's hard to say precisely when the question switches polarity. But it's
probably not that dangerous to start worrying too early that you're
default dead, whereas it's very dangerous to start worrying too late.

The reason is a phenomenon I wrote about earlier: the
\href{pinch.html}{fatal pinch}. The fatal pinch is default dead + slow
growth + not enough time to fix it. And the way founders end up in it is
by not realizing that's where they're headed.

There is another reason founders don't ask themselves whether they're
default alive or default dead: they assume it will be easy to raise more
money. But that assumption is often false, and worse still, the more you
depend on it, the falser it becomes.

Maybe it will help to separate facts from hopes. Instead of thinking of
the future with vague optimism, explicitly separate the components. Say
``We're default dead, but we're counting on investors to save us.''
Maybe as you say that, it will set off the same alarms in your head that
it does in mine. And if you set off the alarms sufficiently early, you
may be able to avoid the fatal pinch.

It would be safe to be default dead if you could count on investors
saving you. As a rule their interest is a function of growth. If you
have steep revenue growth, say over 5x a year, you can start to count on
investors being interested even if you're not profitable.
\footnote{Steep usage growth will also interest investors. Revenue will
ultimately be a constant multiple of usage, so x\% usage growth predicts
x\% revenue growth. But in practice investors discount merely predicted
revenue, so if you're measuring usage you need a higher growth rate to
impress investors.} But investors are so fickle that you can never
do more than start to count on them. Sometimes something about your
business will spook investors even if your growth is great. So no matter
how good your growth is, you can never safely treat fundraising as more
than a plan A. You should always have a plan B as well: you should know
(as in write down) precisely what you'll need to do to survive if you
can't raise more money, and precisely when you'll have to switch to plan
B if plan A isn't working.

In any case, growing fast versus operating cheaply is far from the sharp
dichotomy many founders assume it to be. In practice there is
surprisingly little connection between how much a startup spends and how
fast it grows. When a startup grows fast, it's usually because the
product hits a nerve, in the sense of hitting some big need straight on.
When a startup spends a lot, it's usually because the product is
expensive to develop or sell, or simply because they're wasteful.

If you're paying attention, you'll be asking at this point not just how
to avoid the fatal pinch, but how to avoid being default dead. That one
is easy: don't hire too fast. Hiring too fast is by far the biggest
killer of startups that raise money. \footnote{Startups that don't raise money are saved from hiring too fast
because they can't afford to. But that doesn't mean you should avoid
raising money in order to avoid this problem, any more than that total
abstinence is the only way to avoid becoming an alcoholic.}

Founders tell themselves they need to hire in order to grow. But most
err on the side of overestimating this need rather than underestimating
it. Why? Partly because there's so much work to do. Naive founders think
that if they can just hire enough people, it will all get done. Partly
because successful startups have lots of employees, so it seems like
that's what one does in order to be successful. In fact the large staffs
of successful startups are probably more the effect of growth than the
cause. And partly because when founders have slow growth they don't want
to face what is usually the real reason: the product is not appealing
enough.

Plus founders who've just raised money are often encouraged to overhire
by the VCs who funded them. Kill-or-cure strategies are optimal for VCs
because they're protected by the portfolio effect. VCs want to blow you
up, in one sense of the phrase or the other. But as a founder your
incentives are different. You want above all to survive.
\footnote{I would not be surprised if VCs' tendency to push founders to
overhire is not even in their own interest. They don't know how many of
the companies that get killed by overspending might have done well if
they'd survived. My guess is a significant number.}

Here's a common way startups die. They make something moderately
appealing and have decent initial growth. They raise their first round
fairly easily, because the founders seem smart and the idea sounds
plausible. But because the product is only moderately appealing, growth
is ok but not great. The founders convince themselves that hiring a
bunch of people is the way to boost growth. Their investors agree. But
(because the product is only moderately appealing) the growth never
comes. Now they're rapidly running out of runway. They hope further
investment will save them. But because they have high expenses and slow
growth, they're now unappealing to investors. They're unable to raise
more, and the company dies.

What the company should have done is address the fundamental problem:
that the product is only moderately appealing. Hiring people is rarely
the way to fix that. More often than not it makes it harder. At this
early stage, the product needs to evolve more than to be ``built out,''
and that's usually easier with fewer people. \footnote{After reading a draft, Sam Altman wrote:

``I think you should make the hiring point more strongly. I think it's
roughly correct to say that YC's most successful companies have never
been the fastest to hire, and one of the marks of a great founder is
being able to resist this urge.''

Paul Buchheit adds:

``A related problem that I see a lot is premature scaling---founders
take a small business that isn't really working (bad unit economics,
typically) and then scale it up because they want impressive growth
numbers. This is similar to over-hiring in that it makes the business
much harder to fix once it's big, plus they are bleeding cash really
fast.''}

Asking whether you're default alive or default dead may save you from
this. Maybe the alarm bells it sets off will counteract the forces that
push you to overhire. Instead you'll be compelled to seek growth in
other ways. For example, by \href{ds.html}{doing things that don't
scale}, or by redesigning the product in the way only founders can. And
for many if not most startups, these paths to growth will be the ones
that actually work.

Airbnb waited 4 months after raising money at the end of Y~Combinator
before they hired their first employee. In the meantime the founders
were terribly overworked. But they were overworked evolving Airbnb into
the astonishingly successful organism it is now.


\textbf{Thanks} to Sam Altman, Paul Buchheit, Joe Gebbia, Jessica
Livingston, and Geoff Ralston for reading drafts of this.
\chapter{Write Like You Talk}

October 2015

Here's a simple trick for getting more people to read what you write:
write in spoken language.

Something comes over most people when they start writing. They write in
a different language than they'd use if they were talking to a friend.
The sentence structure and even the words are different. No one uses
``pen'' as a verb in spoken English. You'd feel like an idiot using
``pen'' instead of ``write'' in a conversation with a friend.

The last straw for me was a sentence I read a couple days ago:

\begin{quote}
The mercurial Spaniard himself declared: ``After Altamira, all is
decadence.''
\end{quote}

It's from Neil Oliver's \emph{A History of Ancient Britain}. I feel bad
making an example of this book, because it's no worse than lots of
others. But just imagine calling Picasso ``the mercurial Spaniard'' when
talking to a friend. Even one sentence of this would raise eyebrows in
conversation. And yet people write whole books of it.

Ok, so written and spoken language are different. Does that make written
language worse?

If you want people to read and understand what you write, yes. Written
language is more complex, which makes it more work to read. It's also
more formal and distant, which gives the reader's attention permission
to drift. But perhaps worst of all, the complex sentences and fancy
words give you, the writer, the false impression that you're saying more
than you actually are.

You don't need complex sentences to express complex ideas. When
specialists in some abstruse topic talk to one another about ideas in
their field, they don't use sentences any more complex than they do when
talking about what to have for lunch. They use different words,
certainly. But even those they use no more than necessary. And in my
experience, the harder the subject, the more informally experts speak.
Partly, I think, because they have less to prove, and partly because the
harder the ideas you're talking about, the less you can afford to let
language get in the way.

Informal language is the athletic clothing of ideas.

I'm not saying spoken language always works best. Poetry is as much
music as text, so you can say things you wouldn't say in conversation.
And there are a handful of writers who can get away with using fancy
language in prose. And then of course there are cases where writers
don't want to make it easy to understand what they're saying---in
corporate announcements of bad news, for example, or at the more
\href{https://scholar.google.com/scholar?hl=en&as_sdt=1,5&q=transgression+narrative+postmodern+gender}{bogus}
end of the humanities. But for nearly everyone else, spoken language is
better.

It seems to be hard for most people to write in spoken language. So
perhaps the best solution is to write your first draft the way you
usually would, then afterward look at each sentence and ask ``Is this
the way I'd say this if I were talking to a friend?'' If it isn't,
imagine what you would say, and use that instead. After a while this
filter will start to operate as you write. When you write something you
wouldn't say, you'll hear the clank as it hits the page.

Before I publish a new essay, I read it out loud and fix everything that
doesn't sound like conversation. I even fix bits that are phonetically
awkward; I don't know if that's necessary, but it doesn't cost much.

This trick may not always be enough. I've seen writing so far removed
from spoken language that it couldn't be fixed sentence by sentence. For
cases like that there's a more drastic solution. After writing the first
draft, try explaining to a friend what you just wrote. Then replace the
draft with what you said to your friend.

People often tell me how much my essays sound like me talking. The fact
that this seems worthy of comment shows how rarely people manage to
write in spoken language. Otherwise everyone's writing would sound like
them talking.

If you simply manage to write in spoken language, you'll be ahead of
95\% of writers. And it's so easy to do: just don't let a sentence
through unless it's the way you'd say it to a friend.

\textbf{Thanks} to Patrick Collison and Jessica Livingston for reading
drafts of this.
\chapter{A Way to Detect Bias}

October 2015

This will come as a surprise to a lot of people, but in some cases it's
possible to detect bias in a selection process without knowing anything
about the applicant pool. Which is exciting because among other things
it means third parties can use this technique to detect bias whether
those doing the selecting want them to or not.

You can use this technique whenever (a) you have at least a random
sample of the applicants that were selected, (b) their subsequent
performance is measured, and (c) the groups of applicants you're
comparing have roughly equal distribution of ability.

How does it work? Think about what it means to be biased. What it means
for a selection process to be biased against applicants of type x is
that it's harder for them to make it through. Which means applicants of
type x have to be better to get selected than applicants not of type x.
\footnote{This technique wouldn't work if the selection process looked for
different things from different types of applicants---for example, if an
employer hired men based on their ability but women based on their
appearance.} Which means applicants of type x who do make it
through the selection process will outperform other successful
applicants. And if the performance of all the successful applicants is
measured, you'll know if they do.

Of course, the test you use to measure performance must be a valid one.
And in particular it must not be invalidated by the bias you're trying
to measure. But there are some domains where performance can be
measured, and in those detecting bias is straightforward. Want to know
if the selection process was biased against some type of applicant?
Check whether they outperform the others. This is not just a heuristic
for detecting bias. It's what bias means.

For example, many suspect that venture capital firms are biased against
female founders. This would be easy to detect: among their portfolio
companies, do startups with female founders outperform those without? A
couple months ago, one VC firm (almost certainly unintentionally)
published a study showing bias of this type. First Round Capital found
that among its portfolio companies, startups with female founders
\href{http://10years.firstround.com/\#one}{outperformed} those without
by 63\%. \footnote{As Paul Buchheit points out, First Round excluded their most
successful investment, Uber, from the study. And while it makes sense to
exclude outliers from some types of studies, studies of returns from
startup investing, which is all about hitting outliers, are not one of
them.}

The reason I began by saying that this technique would come as a
surprise to many people is that we so rarely see analyses of this type.
I'm sure it will come as a surprise to First Round that they performed
one. I doubt anyone there realized that by limiting their sample to
their own portfolio, they were producing a study not of startup trends
but of their own biases when selecting companies.

I predict we'll see this technique used more in the future. The
information needed to conduct such studies is increasingly available.
Data about who applies for things is usually closely guarded by the
organizations selecting them, but nowadays data about who gets selected
is often publicly available to anyone who takes the trouble to aggregate
it.


\textbf{Thanks} to Sam Altman, Jessica Livingston, and Geoff Ralston for
reading drafts of this.
\chapter{Jessica Livingston}

November 2015

A few months ago an article about Y Combinator said that early on it had
been a ``one-man show.'' It's sadly common to read that sort of thing.
But the problem with that description is not just that it's unfair. It's
also misleading. Much of what's most novel about YC is due to Jessica
Livingston. If you don't understand her, you don't understand YC. So let
me tell you a little about Jessica.

YC had 4 founders. Jessica and I decided one night to start it, and the
next day we recruited my friends Robert Morris and Trevor Blackwell.
Jessica and I ran YC day to day, and Robert and Trevor read applications
and did interviews with us.

Jessica and I were already dating when we started YC. At first we tried
to act ``professional'' about this, meaning we tried to conceal it. In
retrospect that seems ridiculous, and we soon dropped the pretense. And
the fact that Jessica and I were a couple is a big part of what made YC
what it was. YC felt like a family. The founders early on were mostly
young. We all had dinner together once a week, cooked for the first
couple years by me. Our first building had been a private home. The
overall atmosphere was shockingly different from a VC's office on Sand
Hill Road, in a way that was entirely for the better. There was an
authenticity that everyone who walked in could sense. And that didn't
just mean that people trusted us. It was the perfect quality to instill
in startups. Authenticity is one of the most important things YC looks
for in founders, not just because fakers and opportunists are annoying,
but because authenticity is one of the main things that separates the
most successful startups from the rest.

Early YC was a family, and Jessica was its mom. And the culture she
defined was one of YC's most important innovations. Culture is important
in any organization, but at YC culture wasn't just how we behaved when
we built the product. At YC, the culture was the product.

Jessica was also the mom in another sense: she had the last word.
Everything we did as an organization went through her first --- who to
fund, what to say to the public, how to deal with other companies, who
to hire, everything.

Before we had kids, YC was more or less our life. There was no real
distinction between working hours and not. We talked about YC all the
time. And while there might be some businesses that it would be tedious
to let infect your private life, we liked it. We'd started YC because it
was something we were interested in. And some of the problems we were
trying to solve were endlessly difficult. How do you recognize good
founders? You could talk about that for years, and we did; we still do.

I'm better at some things than Jessica, and she's better at some things
than me. One of the things she's best at is judging people. She's one of
those rare individuals with x-ray vision for character. She can see
through any kind of faker almost immediately. Her nickname within YC was
the Social Radar, and this special power of hers was critical in making
YC what it is. The earlier you pick startups, the more you're picking
the founders. Later stage investors get to try products and look at
growth numbers. At the stage where YC invests, there is often neither a
product nor any numbers.

Others thought YC had some special insight about the future of
technology. Mostly we had the same sort of insight Socrates claimed: we
at least knew we knew nothing. What made YC successful was being able to
pick good founders. We thought Airbnb was a bad idea. We funded it
because we liked the founders.

During interviews, Robert and Trevor and I would pepper the applicants
with technical questions. Jessica would mostly watch. A lot of the
applicants probably read her as some kind of secretary, especially early
on, because she was the one who'd go out and get each new group and she
didn't ask many questions. She was ok with that. It was easier for her
to watch people if they didn't notice her. But after the interview, the
three of us would turn to Jessica and ask ``What does the Social Radar
say?'' \footnote{Harj Taggar reminded me that while Jessica didn't ask many
questions, they tended to be important ones:

``She was always good at sniffing out any red flags about the team or
their determination and disarmingly asking the right question, which
usually revealed more than the founders realized.''}

Having the Social Radar at interviews wasn't just how we picked founders
who'd be successful. It was also how we picked founders who were good
people. At first we did this because we couldn't help it. Imagine what
it would feel like to have x-ray vision for character. Being around bad
people would be intolerable. So we'd refuse to fund founders whose
characters we had doubts about even if we thought they'd be successful.

Though we initially did this out of self-indulgence, it turned out to be
very valuable to YC. We didn't realize it in the beginning, but the
people we were picking would become the YC alumni network. And once we
picked them, unless they did something really egregious, they were going
to be part of it for life. Some now think YC's alumni network is its
most valuable feature. I personally think YC's advice is pretty good
too, but the alumni network is certainly among the most valuable
features. The level of trust and helpfulness is remarkable for a group
of such size. And Jessica is the main reason why.

(As we later learned, it probably cost us little to reject people whose
characters we had doubts about, because how good founders are and how
well they do are \href{mean.html}{not orthogonal}. If bad founders
succeed at all, they tend to sell early. The most successful founders
are almost all good.)

If Jessica was so important to YC, why don't more people realize it?
Partly because I'm a writer, and writers always get disproportionate
attention. YC's brand was initially my brand, and our applicants were
people who'd read my essays. But there is another reason: Jessica hates
attention. Talking to reporters makes her nervous. The thought of giving
a talk paralyzes her. She was even uncomfortable at our wedding, because
the bride is always the center of attention. \footnote{Or more precisely, while she likes getting attention in the
sense of getting credit for what she has done, she doesn't like getting
attention in the sense of being watched in real time. Unfortunately, not
just for her but for a lot of people, how much you get of the former
depends a lot on how much you get of the latter.

Incidentally, if you saw Jessica at a public event, you would never
guess she hates attention, because (a) she is very polite and (b) when
she's nervous, she expresses it by smiling more.}

It's not just because she's shy that she hates attention, but because it
throws off the Social Radar. She can't be herself. You can't watch
people when everyone is watching you.

Another reason attention worries her is that she hates bragging. In
anything she does that's publicly visible, her biggest fear (after the
obvious fear that it will be bad) is that it will seem ostentatious. She
says being too modest is a common problem for women. But in her case it
goes beyond that. She has a horror of ostentation so visceral it's
almost a phobia.

She also hates fighting. She can't do it; she just shuts down. And
unfortunately there is a good deal of fighting in being the public face
of an organization.

So although Jessica more than anyone made YC unique, the very qualities
that enabled her to do it mean she tends to get written out of YC's
history. Everyone buys this story that PG started YC and his wife just
kind of helped. Even YC's haters buy it. A couple years ago when people
were attacking us for not funding more female founders (than exist),
they all treated YC as identical with PG. It would have spoiled the
narrative to acknowledge Jessica's central role at YC.

Jessica was boiling mad that people were accusing \emph{her} company of
sexism. I've never seen her angrier about anything. But she did not
contradict them. Not publicly. In private there was a great deal of
profanity. And she wrote three separate essays about the question of
female founders. But she could never bring herself to publish any of
them. She'd seen the level of vitriol in this debate, and she shrank
from engaging. \footnote{The existence of people like Jessica is not just something the
mainstream media needs to learn to acknowledge, but something feminists
need to learn to acknowledge as well. There are successful women who
don't like to fight. Which means if the public conversation about women
consists of fighting, their voices will be silenced.

There's a sort of Gresham's Law of conversations. If a conversation
reaches a certain level of incivility, the more thoughtful people start
to leave. No one understands female founders better than Jessica. But
it's unlikely anyone will ever hear her speak candidly about the topic.
She ventured a toe in that water a while ago, and the reaction was so
violent that she decided ``never again.''}

It wasn't just because she disliked fighting. She's so sensitive to
character that it repels her even to fight with dishonest people. The
idea of mixing it up with linkbait journalists or Twitter trolls would
seem to her not merely frightening, but disgusting.

But Jessica knew her example as a successful female founder would
encourage more women to start companies, so last year she did something
YC had never done before and hired a PR firm to get her some interviews.
At one of the first she did, the reporter brushed aside her insights
about startups and turned it into a sensationalistic story about how
some guy had tried to chat her up as she was waiting outside the bar
where they had arranged to meet. Jessica was mortified, partly because
the guy had done nothing wrong, but more because the story treated her
as a victim significant only for being a woman, rather than one of the
most knowledgeable investors in the Valley.

After that she told the PR firm to stop.

You're not going to be hearing in the press about what Jessica has
achieved. So let me tell you what Jessica has achieved. Y Combinator is
fundamentally a nexus of people, like a university. It doesn't make a
product. What defines it is the people. Jessica more than anyone curated
and nurtured that collection of people. In that sense she literally made
YC.

Jessica knows more about the qualities of startup founders than anyone
else ever has. Her immense data set and x-ray vision are the perfect
storm in that respect. The qualities of the founders are the best
predictor of how a startup will do. And startups are in turn the most
important source of growth in mature economies.

The person who knows the most about the most important factor in the
growth of mature economies --- that is who Jessica Livingston is.
Doesn't that sound like someone who should be better known?


\textbf{Thanks} to Sam Altman, Paul Buchheit, Patrick Collison, Daniel
Gackle, Carolynn Levy, Jon Levy, Kirsty Nathoo, Robert Morris, Geoff
Ralston, and Harj Taggar for reading drafts of this. And yes, Jessica
Livingston, who made me cut surprisingly little.
\chapter{The Refragmentation}

January 2016

One advantage of being old is that you can see change happen in your
lifetime. A lot of the change I've seen is fragmentation. US politics is
much more polarized than it used to be. Culturally we have ever less
common ground. The creative class flocks to a handful of happy cities,
abandoning the rest. And increasing economic inequality means the spread
between rich and poor is growing too. I'd like to propose a hypothesis:
that all these trends are instances of the same phenomenon. And
moreover, that the cause is not some force that's pulling us apart, but
rather the erosion of forces that had been pushing us together.

Worse still, for those who worry about these trends, the forces that
were pushing us together were an anomaly, a one-time combination of
circumstances that's unlikely to be repeated --- and indeed, that we
would not want to repeat.

The two forces were war (above all World War II), and the rise of large
corporations.

The effects of World War II were both economic and social. Economically,
it decreased variation in income. Like all modern armed forces,
America's were socialist economically. From each according to his
ability, to each according to his need. More or less. Higher ranking
members of the military got more (as higher ranking members of socialist
societies always do), but what they got was fixed according to their
rank. And the flattening effect wasn't limited to those under arms,
because the US economy was conscripted too. Between 1942 and 1945 all
wages were set by the National War Labor Board. Like the military, they
defaulted to flatness. And this national standardization of wages was so
pervasive that its effects could still be seen years after the war
ended. \footnote{Lester Thurow, writing in 1975, said the wage differentials
prevailing at the end of World War II had become so embedded that they
``were regarded as `just' even after the egalitarian pressures of World
War II had disappeared. Basically, the same differentials exist to this
day, thirty years later.'' But Goldin and Margo think market forces in
the postwar period also helped preserve the wartime compression of wages
--- specifically increased demand for unskilled workers, and oversupply
of educated ones.

(Oddly enough, the American custom of having employers pay for health
insurance derives from efforts by businesses to circumvent NWLB wage
controls in order to attract workers.)}

Business owners weren't supposed to be making money either. FDR said
``not a single war millionaire'' would be permitted. To ensure that, any
increase in a company's profits over prewar levels was taxed at 85\%.
And when what was left after corporate taxes reached individuals, it was
taxed again at a marginal rate of 93\%. \footnote{As always, tax rates don't tell the whole story. There were lots
of exemptions, especially for individuals. And in World War II the tax
codes were so new that the government had little acquired immunity to
tax avoidance. If the rich paid high taxes during the war it was more
because they wanted to than because they had to.

After the war, federal tax receipts as a percentage of GDP were about
the same as they are now. In fact, for the entire period since the war,
tax receipts have stayed close to 18\% of GDP, despite dramatic changes
in tax rates. The lowest point occurred when marginal income tax rates
were highest: 14.1\% in 1950. Looking at the data, it's hard to avoid
the conclusion that tax rates have had little effect on what people
actually paid.}

Socially too the war tended to decrease variation. Over 16 million men
and women from all sorts of different backgrounds were brought together
in a way of life that was literally uniform. Service rates for men born
in the early 1920s approached 80\%. And working toward a common goal,
often under stress, brought them still closer together.

Though strictly speaking World War II lasted less than 4 years for the
US, its effects lasted longer. Wars make central governments more
powerful, and World War II was an extreme case of this. In the US, as in
all the other Allied countries, the federal government was slow to give
up the new powers it had acquired. Indeed, in some respects the war
didn't end in 1945; the enemy just switched to the Soviet Union. In tax
rates, federal power, defense spending, conscription, and nationalism,
the decades after the war looked more like wartime than prewar
peacetime. \footnote{Though in fact the decade preceding the war had been a time of
unprecedented federal power, in response to the Depression. Which is not
entirely a coincidence, because the Depression was one of the causes of
the war. In many ways the New Deal was a sort of dress rehearsal for the
measures the federal government took during wartime. The wartime
versions were much more drastic and more pervasive though. As Anthony
Badger wrote, ``for many Americans the decisive change in their
experiences came not with the New Deal but with World War II.''} And the social effects lasted too.
The kid pulled into the army from behind a mule team in West Virginia
didn't simply go back to the farm afterward. Something else was waiting
for him, something that looked a lot like the army.

If total war was the big political story of the 20th century, the big
economic story was the rise of a new kind of company. And this too
tended to produce both social and economic cohesion.
\footnote{I don't know enough about the origins of the world wars to say,
but it's not inconceivable they were connected to the rise of big
corporations. If that were the case, 20th century cohesion would have a
single cause.}

The 20th century was the century of the big, national corporation.
General Electric, General Foods, General Motors. Developments in
finance, communications, transportation, and manufacturing enabled a new
type of company whose goal was above all scale. Version 1 of this world
was low-res: a Duplo world of a few giant companies dominating each big
market. \footnote{More precisely, there was a bimodal economy consisting, in
Galbraith's words, of ``the world of the technically dynamic, massively
capitalized and highly organized corporations on the one hand and the
hundreds of thousands of small and traditional proprietors on the
other.'' Money, prestige, and power were concentrated in the former, and
there was near zero crossover.}

The late 19th and early 20th centuries had been a time of consolidation,
led especially by J. P. Morgan. Thousands of companies run by their
founders were merged into a couple hundred giant ones run by
professional managers. Economies of scale ruled the day. It seemed to
people at the time that this was the final state of things. John D.
Rockefeller said in 1880

\begin{quote}
The day of combination is here to stay. Individualism has gone, never to
return.
\end{quote}

He turned out to be mistaken, but he seemed right for the next hundred
years.

The consolidation that began in the late 19th century continued for most
of the 20th. By the end of World War II, as Michael Lind writes, ``the
major sectors of the economy were either organized as government-backed
cartels or dominated by a few oligopolistic corporations.''

For consumers this new world meant the same choices everywhere, but only
a few of them. When I grew up there were only 2 or 3 of most things, and
since they were all aiming at the middle of the market there wasn't much
to differentiate them.

One of the most important instances of this phenomenon was in TV. Here
there were 3 choices: NBC, CBS, and ABC. Plus public TV for eggheads and
communists. The programs that the 3 networks offered were
indistinguishable. In fact, here there was a triple pressure toward the
center. If one show did try something daring, local affiliates in
conservative markets would make them stop. Plus since TVs were
expensive, whole families watched the same shows together, so they had
to be suitable for everyone.

And not only did everyone get the same thing, they got it at the same
time. It's difficult to imagine now, but every night tens of millions of
families would sit down together in front of their TV set watching the
same show, at the same time, as their next door neighbors. What happens
now with the Super Bowl used to happen every night. We were literally in
sync. \footnote{I wonder how much of the decline in families eating together was
due to the decline in families watching TV together afterward.}

In a way mid-century TV culture was good. The view it gave of the world
was like you'd find in a children's book, and it probably had something
of the effect that (parents hope) children's books have in making people
behave better. But, like children's books, TV was also misleading.
Dangerously misleading, for adults. In his autobiography, Robert MacNeil
talks of seeing gruesome images that had just come in from Vietnam and
thinking, we can't show these to families while they're having dinner.

I know how pervasive the common culture was, because I tried to opt out
of it, and it was practically impossible to find alternatives. When I
was 13 I realized, more from internal evidence than any outside source,
that the ideas we were being fed on TV were crap, and I stopped watching
it. \footnote{I know when this happened because it was the season Dallas
premiered. Everyone else was talking about what was happening on Dallas,
and I had no idea what they meant.} But it wasn't just TV. It seemed like
everything around me was crap. The politicians all saying the same
things, the consumer brands making almost identical products with
different labels stuck on to indicate how prestigious they were meant to
be, the balloon-frame houses with fake ``colonial'' skins, the cars with
several feet of gratuitous metal on each end that started to fall apart
after a couple years, the ``red delicious'' apples that were red but
only nominally apples. And in retrospect, it \emph{was} crap.
\footnote{I didn't realize it till I started doing research for this
essay, but the meretriciousness of the products I grew up with is a
well-known byproduct of oligopoly. When companies can't compete on
price, they compete on tailfins.}

But when I went looking for alternatives to fill this void, I found
practically nothing. There was no Internet then. The only place to look
was in the chain bookstore in our local shopping mall.
\footnote{Monroeville Mall was at the time of its completion in 1969 the
largest in the country. In the late 1970s the movie \emph{Dawn of the
Dead} was shot there. Apparently the mall was not just the location of
the movie, but its inspiration; the crowds of shoppers drifting through
this huge mall reminded George Romero of zombies. My first job was
scooping ice cream in the Baskin-Robbins.} There I found a copy of \emph{The Atlantic}. I
wish I could say it became a gateway into a wider world, but in fact I
found it boring and incomprehensible. Like a kid tasting whisky for the
first time and pretending to like it, I preserved that magazine as
carefully as if it had been a book. I'm sure I still have it somewhere.
But though it was evidence that there was, somewhere, a world that
wasn't red delicious, I didn't find it till college.

It wasn't just as consumers that the big companies made us similar. They
did as employers too. Within companies there were powerful forces
pushing people toward a single model of how to look and act. IBM was
particularly notorious for this, but they were only a little more
extreme than other big companies. And the models of how to look and act
varied little between companies. Meaning everyone within this world was
expected to seem more or less the same. And not just those in the
corporate world, but also everyone who aspired to it --- which in the
middle of the 20th century meant most people who weren't already in it.
For most of the 20th century, working-class people tried hard to look
middle class. You can see it in old photos. Few adults aspired to look
dangerous in 1950.

But the rise of national corporations didn't just compress us
culturally. It compressed us economically too, and on both ends.

Along with giant national corporations, we got giant national labor
unions. And in the mid 20th century the corporations cut deals with the
unions where they paid over market price for labor. Partly because the
unions were monopolies. \footnote{Labor unions were exempted from antitrust laws by the Clayton
Antitrust Act in 1914 on the grounds that a person's work is not ``a
commodity or article of commerce.'' I wonder if that means service
companies are also exempt.} Partly because, as
components of oligopolies themselves, the corporations knew they could
safely pass the cost on to their customers, because their competitors
would have to as well. And partly because in mid-century most of the
giant companies were still focused on finding new ways to milk economies
of scale. Just as startups rightly pay AWS a premium over the cost of
running their own servers so they can focus on growth, many of the big
national corporations were willing to pay a premium for labor.
\footnote{The relationships between unions and unionized companies can
even be symbiotic, because unions will exert political pressure to
protect their hosts. According to Michael Lind, when politicians tried
to attack the A\&P supermarket chain because it was putting local
grocery stores out of business, ``A\&P successfully defended itself by
allowing the unionization of its workforce in 1938, thereby gaining
organized labor as a constituency.'' I've seen this phenomenon myself:
hotel unions are responsible for more of the political pressure against
Airbnb than hotel companies.}

As well as pushing incomes up from the bottom, by overpaying unions, the
big companies of the 20th century also pushed incomes down at the top,
by underpaying their top management. Economist J. K. Galbraith wrote in
1967 that ``There are few corporations in which it would be suggested
that executive salaries are at a maximum.'' \footnote{Galbraith was clearly puzzled that corporate executives would
work so hard to make money for other people (the shareholders) instead
of themselves. He devoted much of \emph{The New Industrial State} to
trying to figure this out.

His theory was that professionalism had replaced money as a motive, and
that modern corporate executives were, like (good) scientists, motivated
less by financial rewards than by the desire to do good work and thereby
earn the respect of their peers. There is something in this, though I
think lack of movement between companies combined with self-interest
explains much of observed behavior.}

To some extent this was an illusion. Much of the de facto pay of
executives never showed up on their income tax returns, because it took
the form of perks. The higher the rate of income tax, the more pressure
there was to pay employees upstream of it. (In the UK, where taxes were
even higher than in the US, companies would even pay their kids' private
school tuitions.) One of the most valuable things the big companies of
the mid 20th century gave their employees was job security, and this too
didn't show up in tax returns or income statistics. So the nature of
employment in these organizations tended to yield falsely low numbers
about economic inequality. But even accounting for that, the big
companies paid their best people less than market price. There was no
market; the expectation was that you'd work for the same company for
decades if not your whole career. \footnote{Galbraith (p.~94) says a 1952 study of the 800 highest paid
executives at 300 big corporations found that three quarters of them had
been with their company for more than 20 years.}

Your work was so illiquid there was little chance of getting market
price. But that same illiquidity also encouraged you not to seek it. If
the company promised to employ you till you retired and give you a
pension afterward, you didn't want to extract as much from it this year
as you could. You needed to take care of the company so it could take
care of you. Especially when you'd been working with the same group of
people for decades. If you tried to squeeze the company for more money,
you were squeezing the organization that was going to take care of
\emph{them}. Plus if you didn't put the company first you wouldn't be
promoted, and if you couldn't switch ladders, promotion on this one was
the only way up. \footnote{It seems likely that in the first third of the 20th century
executive salaries were low partly because companies then were more
dependent on banks, who would have disapproved if executives got too
much. This was certainly true in the beginning. The first big company
CEOs were J. P. Morgan's hired hands.

Companies didn't start to finance themselves with retained earnings till
the 1920s. Till then they had to pay out their earnings in dividends,
and so depended on banks for capital for expansion. Bankers continued to
sit on corporate boards till the Glass-Steagall act in 1933.

By mid-century big companies funded 3/4 of their growth from earnings.
But the early years of bank dependence, reinforced by the financial
controls of World War II, must have had a big effect on social
conventions about executive salaries. So it may be that the lack of
movement between companies was as much the effect of low salaries as the
cause.

Incidentally, the switch in the 1920s to financing growth with retained
earnings was one cause of the 1929 crash. The banks now had to find
someone else to lend to, so they made more margin loans.}

To someone who'd spent several formative years in the armed forces, this
situation didn't seem as strange as it does to us now. From their point
of view, as big company executives, they were high-ranking officers.
They got paid a lot more than privates. They got to have expense account
lunches at the best restaurants and fly around on the company's
Gulfstreams. It probably didn't occur to most of them to ask if they
were being paid market price.

The ultimate way to get market price is to work for yourself, by
starting your own company. That seems obvious to any ambitious person
now. But in the mid 20th century it was an alien concept. Not because
starting one's own company seemed too ambitious, but because it didn't
seem ambitious enough. Even as late as the 1970s, when I grew up, the
ambitious plan was to get lots of education at prestigious institutions,
and then join some other prestigious institution and work one's way up
the hierarchy. Your prestige was the prestige of the institution you
belonged to. People did start their own businesses of course, but
educated people rarely did, because in those days there was practically
zero concept of starting what we now call a \href{growth.html}{startup}:
a business that starts small and grows big. That was much harder to do
in the mid 20th century. Starting one's own business meant starting a
business that would start small and stay small. Which in those days of
big companies often meant scurrying around trying to avoid being
trampled by elephants. It was more prestigious to be one of the
executive class riding the elephant.

By the 1970s, no one stopped to wonder where the big prestigious
companies had come from in the first place. It seemed like they'd always
been there, like the chemical elements. And indeed, there was a double
wall between ambitious kids in the 20th century and the origins of the
big companies. Many of the big companies were roll-ups that didn't have
clear founders. And when they did, the founders didn't seem like us.
Nearly all of them had been uneducated, in the sense of not having been
to college. They were what Shakespeare called rude mechanicals. College
trained one to be a member of the professional classes. Its graduates
didn't expect to do the sort of grubby menial work that Andrew Carnegie
or Henry Ford started out doing. \footnote{Even now it's hard to get them to. One of the things I find
hardest to get into the heads of would-be startup founders is how
important it is to do certain kinds of menial work early in the life of
a company. Doing \href{ds.html}{things that don't scale} is to how Henry
Ford got started as a high-fiber diet is to the traditional peasant's
diet: they had no choice but to do the right thing, while we have to
make a conscious effort.}

And in the 20th century there were more and more college graduates. They
increased from about 2\% of the population in 1900 to about 25\% in
2000. In the middle of the century our two big forces intersect, in the
form of the GI Bill, which sent 2.2 million World War II veterans to
college. Few thought of it in these terms, but the result of making
college the canonical path for the ambitious was a world in which it was
socially acceptable to work for Henry Ford, but not to be Henry Ford.
\footnote{Founders weren't celebrated in the press when I was a kid.
``Our founder'' meant a photograph of a severe-looking man with a walrus
mustache and a wing collar who had died decades ago. The thing to be
when I was a kid was an \emph{executive}. If you weren't around then
it's hard to grasp the cachet that term had. The fancy version of
everything was called the ``executive'' model.}

I remember this world well. I came of age just as it was starting to
break up. In my childhood it was still dominant. Not quite so dominant
as it had been. We could see from old TV shows and yearbooks and the way
adults acted that people in the 1950s and 60s had been even more
conformist than us. The mid-century model was already starting to get
old. But that was not how we saw it at the time. We would at most have
said that one could be a bit more daring in 1975 than 1965. And indeed,
things hadn't changed much yet.

But change was coming soon. And when the Duplo economy started to
disintegrate, it disintegrated in several different ways at once.
Vertically integrated companies literally dis-integrated because it was
more efficient to. Incumbents faced new competitors as (a) markets went
global and (b) technical innovation started to trump economies of scale,
turning size from an asset into a liability. Smaller companies were
increasingly able to survive as formerly narrow channels to consumers
broadened. Markets themselves started to change faster, as whole new
categories of products appeared. And last but not least, the federal
government, which had previously smiled upon J. P. Morgan's world as the
natural state of things, began to realize it wasn't the last word after
all.

What J. P. Morgan was to the horizontal axis, Henry Ford was to the
vertical. He wanted to do everything himself. The giant plant he built
at River Rouge between 1917 and 1928 literally took in iron ore at one
end and sent cars out the other. 100,000 people worked there. At the
time it seemed the future. But that is not how car companies operate
today. Now much of the design and manufacturing happens in a long supply
chain, whose products the car companies ultimately assemble and sell.
The reason car companies operate this way is that it works better. Each
company in the supply chain focuses on what they know best. And they
each have to do it well or they can be swapped out for another supplier.

Why didn't Henry Ford realize that networks of cooperating companies
work better than a single big company? One reason is that supplier
networks take a while to evolve. In 1917, doing everything himself
seemed to Ford the only way to get the scale he needed. And the second
reason is that if you want to solve a problem using a network of
cooperating companies, you have to be able to coordinate their efforts,
and you can do that much better with computers. Computers reduce the
transaction costs that Coase argued are the raison d'etre of
corporations. That is a fundamental change.

In the early 20th century, big companies were synonymous with
efficiency. In the late 20th century they were synonymous with
inefficiency. To some extent this was because the companies themselves
had become sclerotic. But it was also because our standards were higher.

It wasn't just within existing industries that change occurred. The
industries themselves changed. It became possible to make lots of new
things, and sometimes the existing companies weren't the ones who did it
best.

Microcomputers are a classic example. The market was pioneered by
upstarts like Apple. When it got big enough, IBM decided it was worth
paying attention to. At the time IBM completely dominated the computer
industry. They assumed that all they had to do, now that this market was
ripe, was to reach out and pick it. Most people at the time would have
agreed with them. But what happened next illustrated how much more
complicated the world had become. IBM did launch a microcomputer. Though
quite successful, it did not crush Apple. But even more importantly, IBM
itself ended up being supplanted by a supplier coming in from the side
--- from software, which didn't even seem to be the same business. IBM's
big mistake was to accept a non-exclusive license for DOS. It must have
seemed a safe move at the time. No other computer manufacturer had ever
been able to outsell them. What difference did it make if other
manufacturers could offer DOS too? The result of that miscalculation was
an explosion of inexpensive PC clones. Microsoft now owned the PC
standard, and the customer. And the microcomputer business ended up
being Apple vs Microsoft.

Basically, Apple bumped IBM and then Microsoft stole its wallet. That
sort of thing did not happen to big companies in mid-century. But it was
going to happen increasingly often in the future.

Change happened mostly by itself in the computer business. In other
industries, legal obstacles had to be removed first. Many of the
mid-century oligopolies had been anointed by the federal government with
policies (and in wartime, large orders) that kept out competitors. This
didn't seem as dubious to government officials at the time as it sounds
to us. They felt a two-party system ensured sufficient competition in
politics. It ought to work for business too.

Gradually the government realized that anti-competitive policies were
doing more harm than good, and during the Carter administration it
started to remove them. The word used for this process was misleadingly
narrow: deregulation. What was really happening was de-oligopolization.
It happened to one industry after another. Two of the most visible to
consumers were air travel and long-distance phone service, which both
became dramatically cheaper after deregulation.

Deregulation also contributed to the wave of hostile takeovers in the
1980s. In the old days the only limit on the inefficiency of companies,
short of actual bankruptcy, was the inefficiency of their competitors.
Now companies had to face absolute rather than relative standards. Any
public company that didn't generate sufficient returns on its assets
risked having its management replaced with one that would. Often the new
managers did this by breaking companies up into components that were
more valuable separately. \footnote{The wave of hostile takeovers in the 1980s was enabled by a
combination of circumstances: court decisions striking down state
anti-takeover laws, starting with the Supreme Court's 1982 decision in
Edgar v. MITE Corp.; the Reagan administration's comparatively
sympathetic attitude toward takeovers; the Depository Institutions Act
of 1982, which allowed banks and savings and loans to buy corporate
bonds; a new SEC rule issued in 1982 (rule 415) that made it possible to
bring corporate bonds to market faster; the creation of the junk bond
business by Michael Milken; a vogue for conglomerates in the preceding
period that caused many companies to be combined that never should have
been; a decade of inflation that left many public companies trading
below the value of their assets; and not least, the increasing
complacency of managements.}

Version 1 of the national economy consisted of a few big blocks whose
relationships were negotiated in back rooms by a handful of executives,
politicians, regulators, and labor leaders. Version 2 was higher
resolution: there were more companies, of more different sizes, making
more different things, and their relationships changed faster. In this
world there were still plenty of back room negotiations, but more was
left to market forces. Which further accelerated the fragmentation.

It's a little misleading to talk of versions when describing a gradual
process, but not as misleading as it might seem. There was a lot of
change in a few decades, and what we ended up with was qualitatively
different. The companies in the S\&P 500 in 1958 had been there an
average of 61 years. By 2012 that number was 18 years.
\footnote{Foster, Richard. ``Creative Destruction Whips through Corporate
America.'' Innosight, February 2012.}

The breakup of the Duplo economy happened simultaneously with the spread
of computing power. To what extent were computers a precondition? It
would take a book to answer that. Obviously the spread of computing
power was a precondition for the rise of startups. I suspect it was for
most of what happened in finance too. But was it a precondition for
globalization or the LBO wave? I don't know, but I wouldn't discount the
possibility. It may be that the refragmentation was driven by computers
in the way the industrial revolution was driven by steam engines.
Whether or not computers were a precondition, they have certainly
accelerated it.

The new fluidity of companies changed people's relationships with their
employers. Why climb a corporate ladder that might be yanked out from
under you? Ambitious people started to think of a career less as
climbing a single ladder than as a series of jobs that might be at
different companies. More movement (or even potential movement) between
companies introduced more competition in salaries. Plus as companies
became smaller it became easier to estimate how much an employee
contributed to the company's revenue. Both changes drove salaries toward
market price. And since people vary dramatically in productivity, paying
market price meant salaries started to diverge.

By no coincidence it was in the early 1980s that the term ``yuppie'' was
coined. That word is not much used now, because the phenomenon it
describes is so taken for granted, but at the time it was a label for
something novel. Yuppies were young professionals who made lots of
money. To someone in their twenties today, this wouldn't seem worth
naming. Why wouldn't young professionals make lots of money? But until
the 1980s, being underpaid early in your career was part of what it
meant to be a professional. Young professionals were paying their dues,
working their way up the ladder. The rewards would come later. What was
novel about yuppies was that they wanted market price for the work they
were doing now.

The first yuppies did not work for startups. That was still in the
future. Nor did they work for big companies. They were professionals
working in fields like law, finance, and consulting. But their example
rapidly inspired their peers. Once they saw that new BMW 325i, they
wanted one too.

Underpaying people at the beginning of their career only works if
everyone does it. Once some employer breaks ranks, everyone else has to,
or they can't get good people. And once started this process spreads
through the whole economy, because at the beginnings of people's careers
they can easily switch not merely employers but industries.

But not all young professionals benefitted. You had to produce to get
paid a lot. It was no coincidence that the first yuppies worked in
fields where it was easy to measure that.

More generally, an idea was returning whose name sounds old-fashioned
precisely because it was so rare for so long: that you could make your
fortune. As in the past there were multiple ways to do it. Some made
their fortunes by creating wealth, and others by playing zero-sum games.
But once it became possible to make one's fortune, the ambitious had to
decide whether or not to. A physicist who chose physics over Wall Street
in 1990 was making a sacrifice that a physicist in 1960 didn't have to
think about.

The idea even flowed back into big companies. CEOs of big companies make
more now than they used to, and I think much of the reason is prestige.
In 1960, corporate CEOs had immense prestige. They were the winners of
the only economic game in town. But if they made as little now as they
did then, in real dollar terms, they'd seem like small fry compared to
professional athletes and whiz kids making millions from startups and
hedge funds. They don't like that idea, so now they try to get as much
as they can, which is more than they had been getting.
\footnote{CEOs of big companies may be overpaid. I don't know enough
about big companies to say. But it is certainly not impossible for a CEO
to make 200x as much difference to a company's revenues as the average
employee. Look at what Steve Jobs did for Apple when he came back as
CEO. It would have been a good deal for the board to give him 95\% of
the company. Apple's market cap the day Steve came back in July 1997 was
1.73 billion. 5\% of Apple now (January 2016) would be worth about 30
billion. And it would not be if Steve hadn't come back; Apple probably
wouldn't even exist anymore.

Merely including Steve in the sample might be enough to answer the
question of whether public company CEOs in the aggregate are overpaid.
And that is not as facile a trick as it might seem, because the broader
your holdings, the more the aggregate is what you care about.}

Meanwhile a similar fragmentation was happening at the other end of the
economic scale. As big companies' oligopolies became less secure, they
were less able to pass costs on to customers and thus less willing to
overpay for labor. And as the Duplo world of a few big blocks fragmented
into many companies of different sizes --- some of them overseas --- it
became harder for unions to enforce their monopolies. As a result
workers' wages also tended toward market price. Which (inevitably, if
unions had been doing their job) tended to be lower. Perhaps
dramatically so, if automation had decreased the need for some kind of
work.

And just as the mid-century model induced social as well as economic
cohesion, its breakup brought social as well as economic fragmentation.
People started to dress and act differently. Those who would later be
called the ``creative class'' became more mobile. People who didn't care
much for religion felt less pressure to go to church for appearances'
sake, while those who liked it a lot opted for increasingly colorful
forms. Some switched from meat loaf to tofu, and others to Hot Pockets.
Some switched from driving Ford sedans to driving small imported cars,
and others to driving SUVs. Kids who went to private schools or wished
they did started to dress ``preppy,'' and kids who wanted to seem
rebellious made a conscious effort to look disreputable. In a hundred
ways people spread apart. \footnote{The late 1960s were famous for social upheaval. But that was
more rebellion (which can happen in any era if people are provoked
sufficiently) than fragmentation. You're not seeing fragmentation unless
you see people breaking off to both left and right.}

Almost four decades later, fragmentation is still increasing. Has it
been net good or bad? I don't know; the question may be unanswerable.
Not entirely bad though. We take for granted the forms of fragmentation
we like, and worry only about the ones we don't. But as someone who
caught the tail end of mid-century
\href{https://books.google.com/ngrams/graph?content=well-adjusted&year_start=1800&year_end=2000&corpus=15&smoothing=3}{conformism},
I can tell you it was no utopia. \footnote{Globally the trend has been in the other direction. While the
US is becoming more fragmented, the world as a whole is becoming less
fragmented, and mostly in good ways.}

My goal here is not to say whether fragmentation has been good or bad,
just to explain why it's happening. With the centripetal forces of total
war and 20th century oligopoly mostly gone, what will happen next? And
more specifically, is it possible to reverse some of the fragmentation
we've seen?

If it is, it will have to happen piecemeal. You can't reproduce
mid-century cohesion the way it was originally produced. It would be
insane to go to war just to induce more national unity. And once you
understand the degree to which the economic history of the 20th century
was a low-res version 1, it's clear you can't reproduce that either.

20th century cohesion was something that happened at least in a sense
naturally. The war was due mostly to external forces, and the Duplo
economy was an evolutionary phase. If you want cohesion now, you'd have
to induce it deliberately. And it's not obvious how. I suspect the best
we'll be able to do is address the symptoms of fragmentation. But that
may be enough.

The form of fragmentation people worry most about lately is
\href{ineq.html}{economic inequality}, and if you want to eliminate that
you're up against a truly formidable headwind that has been in operation
since the stone age. Technology.

Technology is a lever. It magnifies work. And the lever not only grows
increasingly long, but the rate at which it grows is itself increasing.

Which in turn means the variation in the amount of wealth people can
create has not only been increasing, but accelerating. The unusual
conditions that prevailed in the mid 20th century masked this underlying
trend. The ambitious had little choice but to join large organizations
that made them march in step with lots of other people --- literally in
the case of the armed forces, figuratively in the case of big
corporations. Even if the big corporations had wanted to pay people
proportionate to their value, they couldn't have figured out how. But
that constraint has gone now. Ever since it started to erode in the
1970s, we've seen the underlying forces at work again.
\footnote{There were a handful of ways to make a fortune in the mid 20th
century. The main one was drilling for oil, which was open to newcomers
because it was not something big companies could dominate through
economies of scale. How did individuals accumulate large fortunes in an
era of such high taxes? Giant tax loopholes defended by two of the most
powerful men in Congress, Sam Rayburn and Lyndon Johnson.

But becoming a Texas oilman was not in 1950 something one could aspire
to the way starting a startup or going to work on Wall Street were in
2000, because (a) there was a strong local component and (b) success
depended so much on luck.}

Not everyone who gets rich now does it by creating wealth, certainly.
But a significant number do, and the Baumol Effect means all their peers
get dragged along too. \footnote{The Baumol Effect induced by startups is very visible in
Silicon Valley. Google will pay people millions of dollars a year to
keep them from leaving to start or join startups.} And as long as it's
possible to get rich by creating wealth, the default tendency will be
for economic inequality to increase. Even if you eliminate all the other
ways to get rich. You can mitigate this with subsidies at the bottom and
taxes at the top, but unless taxes are high enough to discourage people
from creating wealth, you're always going to be fighting a losing battle
against increasing variation in productivity. \footnote{I'm not claiming variation in productivity is the only cause of
economic inequality in the US. But it's a significant cause, and it will
become as big a cause as it needs to, in the sense that if you ban other
ways to get rich, people who want to get rich will use this route
instead.}

That form of fragmentation, like the others, is here to stay. Or rather,
back to stay. Nothing is forever, but the tendency toward fragmentation
should be more forever than most things, precisely because it's not due
to any particular cause. It's simply a reversion to the mean. When
Rockefeller said individualism was gone, he was right for a hundred
years. It's back now, and that's likely to be true for longer.

I worry that if we don't acknowledge this, we're headed for trouble. If
we think 20th century cohesion disappeared because of few policy tweaks,
we'll be deluded into thinking we can get it back (minus the bad parts,
somehow) with a few countertweaks. And then we'll waste our time trying
to eliminate fragmentation, when we'd be better off thinking about how
to mitigate its consequences.


\textbf{Thanks} to Sam Altman, Trevor Blackwell, Paul Buchheit, Patrick
Collison, Ron Conway, Chris Dixon, Benedict Evans, Richard Florida, Ben
Horowitz, Jessica Livingston, Robert Morris, Tim O'Reilly, Geoff
Ralston, Max Roser, Alexia Tsotsis, and Qasar Younis for reading drafts
of this. Max also told me about several valuable sources.

\textbf{Bibliography}

Allen, Frederick Lewis. \emph{The Big Change}. Harper, 1952.

Averitt, Robert. \emph{The Dual Economy}. Norton, 1968.

Badger, Anthony. \emph{The New Deal}. Hill and Wang, 1989.

Bainbridge, John. \emph{The Super-Americans}. Doubleday, 1961.

Beatty, Jack. \emph{Collossus}. Broadway, 2001.

Brinkley, Douglas. \emph{Wheels for the World}. Viking, 2003.

Brownleee, W. Elliot. \emph{Federal Taxation in America}. Cambridge,
1996.

Chandler, Alfred. \emph{The Visible Hand}. Harvard, 1977.

Chernow, Ron. \emph{The House of Morgan}. Simon \& Schuster, 1990.

Chernow, Ron. \emph{Titan: The Life of John D. Rockefeller}. Random
House, 1998.

Galbraith, John. \emph{The New Industrial State}. Houghton Mifflin,
1967.

Goldin, Claudia and Robert A. Margo. ``The Great Compression: The Wage
Structure in the United States at Mid-Century.'' NBER Working Paper
3817, 1991.

Gordon, John. \emph{An Empire of Wealth}. HarperCollins, 2004.

Klein, Maury. \emph{The Genesis of Industrial America, 1870-1920}.
Cambridge, 2007.

Lind, Michael. \emph{Land of Promise}. HarperCollins, 2012.

Mickelthwaite, John, and Adrian Wooldridge. \emph{The Company}. Modern
Library, 2003.

Nasaw, David. \emph{Andrew Carnegie}. Penguin, 2006.

Sobel, Robert. \emph{The Age of Giant Corporations}. Praeger, 1993.

Thurow, Lester. \emph{Generating Inequality: Mechanisms of
Distribution}. Basic Books, 1975.

Witte, John. \emph{The Politics and Development of the Federal Income
Tax}. Wisconsin, 1985.

\textbf{Related:}
\chapter{Economic Inequality}

January 2016

Since the 1970s, economic inequality in the US has increased
dramatically. And in particular, the rich have gotten a lot richer.
Nearly everyone who writes about the topic says that economic inequality
should be decreased.

I'm interested in this question because I was one of the founders of a
company called Y Combinator that helps people start startups. Almost by
definition, if a startup succeeds, its founders become rich. Which means
by helping startup founders I've been helping to increase economic
inequality. If economic inequality should be decreased, I shouldn't be
helping founders. No one should be.

But that doesn't sound right. What's going on here? What's going on is
that while economic inequality is a single measure (or more precisely,
two: variation in income, and variation in wealth), it has multiple
causes. Many of these causes are bad, like tax loopholes and drug
addiction. But some are good, like Larry Page and Sergey Brin starting
the company you use to find things online.

If you want to understand economic inequality --- and more importantly,
if you actually want to fix the bad aspects of it --- you have to tease
apart the components. And yet the trend in nearly everything written
about the subject is to do the opposite: to squash together all the
aspects of economic inequality as if it were a single phenomenon.

Sometimes this is done for ideological reasons. Sometimes it's because
the writer only has very high-level data and so draws conclusions from
that, like the proverbial drunk who looks for his keys under the
lamppost, instead of where he dropped them, because the light is better
there. Sometimes it's because the writer doesn't understand critical
aspects of inequality, like the role of technology in wealth creation.
Much of the time, perhaps most of the time, writing about economic
inequality combines all three.

\_\_\_

The most common mistake people make about economic inequality is to
treat it as a single phenomenon. The most naive version of which is the
one based on the pie fallacy: that the rich get rich by taking money
from the poor.

Usually this is an assumption people start from rather than a conclusion
they arrive at by examining the evidence. Sometimes the pie fallacy is
stated explicitly:

\begin{quote}
\ldots those at the top are grabbing an increasing fraction of the
nation's income --- so much of a larger share that what's left over for
the rest is diminished\ldots. \footnote{Stiglitz, Joseph. \emph{The Price of Inequality}. Norton, 2012.
p.~32.}
\end{quote}

Other times it's more unconscious. But the unconscious form is very
widespread. I think because we grow up in a world where the pie fallacy
is actually true. To kids, wealth \emph{is} a fixed pie that's shared
out, and if one person gets more, it's at the expense of another. It
takes a conscious effort to remind oneself that the real world doesn't
work that way.

In the real world you can create wealth as well as taking it from
others. A woodworker creates wealth. He makes a chair, and you willingly
give him money in return for it. A high-frequency trader does not. He
makes a dollar only when someone on the other end of a trade loses a
dollar.

If the rich people in a society got that way by taking wealth from the
poor, then you have the degenerate case of economic inequality, where
the cause of poverty is the same as the cause of wealth. But instances
of inequality don't have to be instances of the degenerate case. If one
woodworker makes 5 chairs and another makes none, the second woodworker
will have less money, but not because anyone took anything from him.

Even people sophisticated enough to know about the pie fallacy are led
toward it by the custom of describing economic inequality as a ratio of
one quantile's income or wealth to another's. It's so easy to slip from
talking about income shifting from one quantile to another, as a figure
of speech, into believing that is literally what's happening.

Except in the degenerate case, economic inequality can't be described by
a ratio or even a curve. In the general case it consists of multiple
ways people become poor, and multiple ways people become rich. Which
means to understand economic inequality in a country, you have to go
find individual people who are poor or rich and figure out why.
\footnote{Particularly since economic inequality is a matter of outliers,
and outliers are disproportionately likely to have gotten where they are
by ways that have little do with the sort of things economists usually
think about, like wages and productivity, but rather by, say, ending up
on the wrong side of the ``War on Drugs.''}

If you want to understand \emph{change} in economic inequality, you
should ask what those people would have done when it was different. This
is one way I know the rich aren't all getting richer simply from some
new system for transferring wealth to them from everyone else. When you
use the would-have method with startup founders, you find what most
would have done \href{re.html}{back in 1960}, when economic inequality
was lower, was to join big companies or become professors. Before Mark
Zuckerberg started Facebook, his default expectation was that he'd end
up working at Microsoft. The reason he and most other startup founders
are richer than they would have been in the mid 20th century is not
because of some right turn the country took during the Reagan
administration, but because progress in technology has made it much
easier to start a new company that \href{growth.html}{grows fast}.

Traditional economists seem strangely averse to studying individual
humans. It seems to be a rule with them that everything has to start
with statistics. So they give you very precise numbers about variation
in wealth and income, then follow it with the most naive speculation
about the underlying causes.

But while there are a lot of people who get rich through rent-seeking of
various forms, and a lot who get rich by playing zero-sum games, there
are also a significant number who get rich by creating wealth. And
creating wealth, as a source of economic inequality, is different from
taking it --- not just morally, but also practically, in the sense that
it is harder to eradicate. One reason is that variation in productivity
is accelerating. The rate at which individuals can create wealth depends
on the technology available to them, and that grows exponentially. The
other reason creating wealth is such a tenacious source of inequality is
that it can expand to accommodate a lot of people.

\_\_\_

I'm all for shutting down the crooked ways to get rich. But that won't
eliminate great variations in wealth, because as long as you leave open
the option of getting rich by creating wealth, people who want to get
rich will do that instead.

Most people who get rich tend to be fairly driven. Whatever their other
flaws, laziness is usually not one of them. Suppose new policies make it
hard to make a fortune in finance. Does it seem plausible that the
people who currently go into finance to make their fortunes will
continue to do so, but be content to work for ordinary salaries? The
reason they go into finance is not because they love finance but because
they want to get rich. If the only way left to get rich is to start
startups, they'll start startups. They'll do well at it too, because
determination is the main factor in the success of a startup.
\footnote{Determination is the most important factor in deciding between
success and failure, which in startups tend to be sharply
differentiated. But it takes more than determination to create one of
the hugely successful startups. Though most founders start out excited
about the idea of getting rich, purely mercenary founders will usually
take one of the big acquisition offers most successful startups get on
the way up. The founders who go on to the next stage tend to be driven
by a sense of mission. They have the same attachment to their companies
that an artist or writer has to their work. But it is very hard to
predict at the outset which founders will do that. It's not simply a
function of their initial attitude. Starting a company changes people.} And while it would probably be a good thing for
the world if people who wanted to get rich switched from playing
zero-sum games to creating wealth, that would not only not eliminate
great variations in wealth, but might even exacerbate them. In a
zero-sum game there is at least a limit to the upside. Plus a lot of the
new startups would create new technology that further accelerated
variation in productivity.

Variation in productivity is far from the only source of economic
inequality, but it is the irreducible core of it, in the sense that
you'll have that left when you eliminate all other sources. And if you
do, that core will be big, because it will have expanded to include the
efforts of all the refugees. Plus it will have a large Baumol penumbra
around it: anyone who could get rich by creating wealth on their own
account will have to be paid enough to prevent them from doing it.

You can't prevent great variations in wealth without preventing people
from getting rich, and you can't do that without preventing them from
starting startups.

So let's be clear about that. Eliminating great variations in wealth
would mean eliminating startups. And that doesn't seem a wise move.
Especially since it would only mean you eliminated startups in your own
country. Ambitious people already move halfway around the world to
further their careers, and startups can operate from anywhere nowadays.
So if you made it impossible to get rich by creating wealth in your
country, people who wanted to do that would just leave and do it
somewhere else. Which would certainly get you a lower Gini coefficient,
along with a lesson in being careful what you ask for.
\footnote{After reading a draft of this essay, Richard Florida told me how
he had once talked to a group of Europeans ``who said they wanted to
make Europe more entrepreneurial and more like Silicon Valley. I said by
definition this will give you more inequality. They thought I was insane
--- they could not process it.''}

I think rising economic inequality is the inevitable fate of countries
that don't choose something worse. We had a 40 year stretch in the
middle of the 20th century that convinced some people otherwise. But as
I explained in \href{re.html}{The Refragmentation}, that was an anomaly
--- a unique combination of circumstances that compressed American
society not just economically but culturally too.
\footnote{Economic inequality has been decreasing globally. But this is
mainly due to the erosion of the kleptocracies that formerly dominated
all the poorer countries. Once the playing field is leveler politically,
we'll see economic inequality start to rise again. The US is the
bellwether. The situation we face here, the rest of the world will
sooner or later.}

And while some of the growth in economic inequality we've seen since
then has been due to bad behavior of various kinds, there has
simultaneously been a huge increase in individuals' ability to create
wealth. Startups are almost entirely a product of this period. And even
within the startup world, there has been a qualitative change in the
last 10 years. Technology has decreased the cost of starting a startup
so much that founders now have the upper hand over investors. Founders
get less diluted, and it is now common for them to retain
\href{control.html}{board control} as well. Both further increase
economic inequality, the former because founders own more stock, and the
latter because, as investors have learned, founders tend to be better at
running their companies than investors.

While the surface manifestations change, the underlying forces are very,
very old. The acceleration of productivity we see in Silicon Valley has
been happening for thousands of years. If you look at the history of
stone tools, technology was already accelerating in the Mesolithic. The
acceleration would have been too slow to perceive in one lifetime. Such
is the nature of the leftmost part of an exponential curve. But it was
the same curve.

You do not want to design your society in a way that's incompatible with
this curve. The evolution of technology is one of the most powerful
forces in history.

Louis Brandeis said ``We may have democracy, or we may have wealth
concentrated in the hands of a few, but we can't have both.'' That
sounds plausible. But if I have to choose between ignoring him and
ignoring an exponential curve that has been operating for thousands of
years, I'll bet on the curve. Ignoring any trend that has been operating
for thousands of years is dangerous. But exponential growth, especially,
tends to bite you.

\_\_\_

If accelerating variation in productivity is always going to produce
some baseline growth in economic inequality, it would be a good idea to
spend some time thinking about that future. Can you have a healthy
society with great variation in wealth? What would it look like?

Notice how novel it feels to think about that. The public conversation
so far has been exclusively about the need to decrease economic
inequality. We've barely given a thought to how to live with it.

I'm hopeful we'll be able to. Brandeis was a product of the Gilded Age,
and things have changed since then. It's harder to hide wrongdoing now.
And to get rich now you don't have to buy politicians the way railroad
or oil magnates did. \footnote{Some people still get rich by buying politicians. My point is
that it's no longer a precondition.} The great concentrations of
wealth I see around me in Silicon Valley don't seem to be destroying
democracy.

There are lots of things wrong with the US that have economic inequality
as a symptom. We should fix those things. In the process we may decrease
economic inequality. But we can't start from the symptom and hope to fix
the underlying causes. \footnote{As well as problems that have economic inequality as a symptom,
there are those that have it as a cause. But in most if not all,
economic inequality is not the primary cause. There is usually some
injustice that is allowing economic inequality to turn into other forms
of inequality, and that injustice is what we need to fix. For example,
the police in the US treat the poor worse than the rich. But the
solution is not to make people richer. It's to make the police treat
people more equitably. Otherwise they'll continue to maltreat people who
are weak in other ways.}

The most obvious is poverty. I'm sure most of those who want to decrease
economic inequality want to do it mainly to help the poor, not to hurt
the rich. \footnote{Some who read this essay will say that I'm clueless or even
being deliberately misleading by focusing so much on the richer end of
economic inequality --- that economic inequality is really about
poverty. But that is exactly the point I'm making, though sloppier
language than I'd use to make it. The real problem is poverty, not
economic inequality. And if you conflate them you're aiming at the wrong
target.

Others will say I'm clueless or being misleading by focusing on people
who get rich by creating wealth --- that startups aren't the problem,
but corrupt practices in finance, healthcare, and so on. Once again,
that is exactly my point. The problem is not economic inequality, but
those specific abuses.

It's a strange task to write an essay about why something isn't the
problem, but that's the situation you find yourself in when so many
people mistakenly think it is.} Indeed, a good number are merely being
sloppy by speaking of decreasing economic inequality when what they mean
is decreasing poverty. But this is a situation where it would be good to
be precise about what we want. Poverty and economic inequality are not
identical. When the city is turning off your
\href{http://www.theatlantic.com/business/archive/2014/07/what-happens-when-detroit-shuts-off-the-water-of-100000-people/374548/}{water}
because you can't pay the bill, it doesn't make any difference what
Larry Page's net worth is compared to yours. He might only be a few
times richer than you, and it would still be just as much of a problem
that your water was getting turned off.

Closely related to poverty is lack of social mobility. I've seen this
myself: you don't have to grow up rich or even upper middle class to get
rich as a startup founder, but few successful founders grew up
desperately poor. But again, the problem here is not simply economic
inequality. There is an enormous difference in wealth between the
household Larry Page grew up in and that of a successful startup
founder, but that didn't prevent him from joining their ranks. It's not
economic inequality per se that's blocking social mobility, but some
specific combination of things that go wrong when kids grow up
sufficiently poor.

One of the most important principles in Silicon Valley is that ``you
make what you measure.'' It means that if you pick some number to focus
on, it will tend to improve, but that you have to choose the right
number, because only the one you choose will improve; another that seems
conceptually adjacent might not. For example, if you're a university
president and you decide to focus on graduation rates, then you'll
improve graduation rates. But only graduation rates, not how much
students learn. Students could learn less, if to improve graduation
rates you made classes easier.

Economic inequality is sufficiently far from identical with the various
problems that have it as a symptom that we'll probably only hit
whichever of the two we aim at. If we aim at economic inequality, we
won't fix these problems. So I say let's aim at the problems.

For example, let's attack poverty, and if necessary damage wealth in the
process. That's much more likely to work than attacking wealth in the
hope that you will thereby fix poverty. \footnote{Particularly since many causes of poverty are only partially
driven by people trying to make money from them. For example, America's
abnormally high incarceration rate is a major cause of poverty. But
although
\href{https://www.washingtonpost.com/posteverything/wp/2015/04/28/how-for-profit-prisons-have-become-the-biggest-lobby-no-one-is-talking-about/}{for-profit
prison companies} and
\href{http://mic.com/articles/41531/union-of-the-snake-how-california-s-prison-guards-subvert-democracy}{prison
guard unions} both spend a lot lobbying for harsh sentencing laws, they
are not the original source of them.} And if
there are people getting rich by tricking consumers or lobbying the
government for anti-competitive regulations or tax loopholes, then let's
stop them. Not because it's causing economic inequality, but because
it's stealing. \footnote{Incidentally, tax loopholes are definitely not a product of
some power shift due to recent increases in economic inequality. The
golden age of economic equality in the mid 20th century was also the
golden age of tax avoidance. Indeed, it was so widespread and so
effective that I'm skeptical whether economic inequality was really so
low then as we think. In a period when people are trying to hide wealth
from the government, it will tend to be hidden from statistics too. One
sign of the potential magnitude of the problem is the discrepancy
between government receipts as a percentage of GDP, which have remained
more or less constant during the entire period from the end of World War
II to the present, and tax rates, which have varied dramatically.}

If all you have is statistics, it seems like that's what you need to
fix. But behind a broad statistical measure like economic inequality
there are some things that are good and some that are bad, some that are
historical trends with immense momentum and others that are random
accidents. If we want to fix the world behind the statistics, we have to
understand it, and focus our efforts where they'll do the most good.


\textbf{Thanks} to Sam Altman, Tiffani Ashley Bell, Patrick Collison,
Ron Conway, Richard Florida, Ben Horowitz, Jessica Livingston, Robert
Morris, Tim O'Reilly, Max Roser, and Alexia Tsotsis for reading drafts
of this.

\textbf{Note:} This is a new version from which I removed a pair of
metaphors that made a lot of people mad, essentially by macroexpanding
them. If anyone wants to see the old version, I put it
\href{ineqold.html}{here}.

\textbf{Related:}
\chapter{Life is Short}

January 2016

Life is short, as everyone knows. When I was a kid I used to wonder
about this. Is life actually short, or are we really complaining about
its finiteness? Would we be just as likely to feel life was short if we
lived 10 times as long?

Since there didn't seem any way to answer this question, I stopped
wondering about it. Then I had kids. That gave me a way to answer the
question, and the answer is that life actually is short.

Having kids showed me how to convert a continuous quantity, time, into
discrete quantities. You only get 52 weekends with your 2 year old. If
Christmas-as-magic lasts from say ages 3 to 10, you only get to watch
your child experience it 8 times. And while it's impossible to say what
is a lot or a little of a continuous quantity like time, 8 is not a lot
of something. If you had a handful of 8 peanuts, or a shelf of 8 books
to choose from, the quantity would definitely seem limited, no matter
what your lifespan was.

Ok, so life actually is short. Does it make any difference to know that?

It has for me. It means arguments of the form ``Life is too short for
x'' have great force. It's not just a figure of speech to say that life
is too short for something. It's not just a synonym for annoying. If you
find yourself thinking that life is too short for something, you should
try to eliminate it if you can.

When I ask myself what I've found life is too short for, the word that
pops into my head is ``bullshit.'' I realize that answer is somewhat
tautological. It's almost the definition of bullshit that it's the stuff
that life is too short for. And yet bullshit does have a distinctive
character. There's something fake about it. It's the junk food of
experience. \footnote{At first I didn't like it that the word that came to mind was
one that had other meanings. But then I realized the other meanings are
fairly closely related. Bullshit in the sense of things you waste your
time on is a lot like intellectual bullshit.}

If you ask yourself what you spend your time on that's bullshit, you
probably already know the answer. Unnecessary meetings, pointless
disputes, bureaucracy, posturing, dealing with other people's mistakes,
traffic jams, addictive but unrewarding pastimes.

There are two ways this kind of thing gets into your life: it's either
forced on you, or it tricks you. To some extent you have to put up with
the bullshit forced on you by circumstances. You need to make money, and
making money consists mostly of errands. Indeed, the law of supply and
demand ensures that: the more rewarding some kind of work is, the
cheaper people will do it. It may be that less bullshit is forced on you
than you think, though. There has always been a stream of people who opt
out of the default grind and go live somewhere where opportunities are
fewer in the conventional sense, but life feels more authentic. This
could become more common.

You can do it on a smaller scale without moving. The amount of time you
have to spend on bullshit varies between employers. Most large
organizations (and many small ones) are steeped in it. But if you
consciously prioritize bullshit avoidance over other factors like money
and prestige, you can probably find employers that will waste less of
your time.

If you're a freelancer or a small company, you can do this at the level
of individual customers. If you fire or avoid toxic customers, you can
decrease the amount of bullshit in your life by more than you decrease
your income.

But while some amount of bullshit is inevitably forced on you, the
bullshit that sneaks into your life by tricking you is no one's fault
but your own. And yet the bullshit you choose may be harder to eliminate
than the bullshit that's forced on you. Things that lure you into
wasting your time have to be really good at tricking you. An example
that will be familiar to a lot of people is arguing online. When someone
contradicts you, they're in a sense attacking you. Sometimes pretty
overtly. Your instinct when attacked is to defend yourself. But like a
lot of instincts, this one wasn't designed for the world we now live in.
Counterintuitive as it feels, it's better most of the time not to defend
yourself. Otherwise these people are literally taking your life.
\footnote{I chose this example deliberately as a note to self. I get
attacked a lot online. People tell the craziest lies about me. And I
have so far done a pretty mediocre job of suppressing the natural human
inclination to say ``Hey, that's not true!''}

Arguing online is only incidentally addictive. There are more dangerous
things than that. As I've written before, one byproduct of technical
progress is that things we like tend to become
\href{addiction.html}{more addictive}. Which means we will increasingly
have to make a conscious effort to avoid addictions --- to stand outside
ourselves and ask ``is this how I want to be spending my time?''

As well as avoiding bullshit, one should actively seek out things that
matter. But different things matter to different people, and most have
to learn what matters to them. A few are lucky and realize early on that
they love math or taking care of animals or writing, and then figure out
a way to spend a lot of time doing it. But most people start out with a
life that's a mix of things that matter and things that don't, and only
gradually learn to distinguish between them.

For the young especially, much of this confusion is induced by the
artificial situations they find themselves in. In middle school and high
school, what the other kids think of you seems the most important thing
in the world. But when you ask adults what they got wrong at that age,
nearly all say they cared too much what other kids thought of them.

One heuristic for distinguishing stuff that matters is to ask yourself
whether you'll care about it in the future. Fake stuff that matters
usually has a sharp peak of seeming to matter. That's how it tricks you.
The area under the curve is small, but its shape jabs into your
consciousness like a pin.

The things that matter aren't necessarily the ones people would call
``important.'' Having coffee with a friend matters. You won't feel later
like that was a waste of time.

One great thing about having small children is that they make you spend
time on things that matter: them. They grab your sleeve as you're
staring at your phone and say ``will you play with me?'' And odds are
that is in fact the bullshit-minimizing option.

If life is short, we should expect its shortness to take us by surprise.
And that is just what tends to happen. You take things for granted, and
then they're gone. You think you can always write that book, or climb
that mountain, or whatever, and then you realize the window has closed.
The saddest windows close when other people die. Their lives are short
too. After my mother died, I wished I'd spent more time with her. I
lived as if she'd always be there. And in her typical quiet way she
encouraged that illusion. But an illusion it was. I think a lot of
people make the same mistake I did.

The usual way to avoid being taken by surprise by something is to be
consciously aware of it. Back when life was more precarious, people used
to be aware of death to a degree that would now seem a bit morbid. I'm
not sure why, but it doesn't seem the right answer to be constantly
reminding oneself of the grim reaper hovering at everyone's shoulder.
Perhaps a better solution is to look at the problem from the other end.
Cultivate a habit of impatience about the things you most want to do.
Don't wait before climbing that mountain or writing that book or
visiting your mother. You don't need to be constantly reminding yourself
why you shouldn't wait. Just don't wait.

I can think of two more things one does when one doesn't have much of
something: try to get more of it, and savor what one has. Both make
sense here.

How you live affects how long you live. Most people could do better. Me
among them.

But you can probably get even more effect by paying closer attention to
the time you have. It's easy to let the days rush by. The ``flow'' that
imaginative people love so much has a darker cousin that prevents you
from pausing to savor life amid the daily slurry of errands and alarms.
One of the most striking things I've read was not in a book, but the
title of one: James Salter's \emph{Burning the Days}.

It is possible to slow time somewhat. I've gotten better at it. Kids
help. When you have small children, there are a lot of moments so
perfect that you can't help noticing.

It does help too to feel that you've squeezed everything out of some
experience. The reason I'm sad about my mother is not just that I miss
her but that I think of all the things we could have done that we
didn't. My oldest son will be 7 soon. And while I miss the 3 year old
version of him, I at least don't have any regrets over what might have
been. We had the best time a daddy and a 3 year old ever had.

Relentlessly prune bullshit, don't wait to do things that matter, and
savor the time you have. That's what you do when life is short.


\textbf{Thanks} to Jessica Livingston and Geoff Ralston for reading
drafts of this.
\chapter{How to Make Pittsburgh a Startup Hub}

April 2016

\emph{(This is a talk I gave at an event called Opt412 in Pittsburgh.
Much of it will apply to other towns. But not all, because as I say in
the talk, Pittsburgh has some important advantages over most would-be
startup hubs.)}

What would it take to make Pittsburgh into a startup hub, like Silicon
Valley? I understand Pittsburgh pretty well, because I grew up here, in
Monroeville. And I understand Silicon Valley pretty well because that's
where I live now. Could you get that kind of startup ecosystem going
here?

When I agreed to speak here, I didn't think I'd be able to give a very
optimistic talk. I thought I'd be talking about what Pittsburgh could do
to become a startup hub, very much in the subjunctive. Instead I'm going
to talk about what Pittsburgh can do.

What changed my mind was an article I read in, of all places, the
\emph{New York Times} food section. The title was
``\href{http://www.nytimes.com/2016/03/16/dining/pittsburgh-restaurants.html}{Pittsburgh's
Youth-Driven Food Boom}.'' To most people that might not even sound
interesting, let alone something related to startups. But it was
electrifying to me to read that title. I don't think I could pick a more
promising one if I tried. And when I read the article I got even more
excited. It said ``people ages 25 to 29 now make up 7.6 percent of all
residents, up from 7 percent about a decade ago.'' Wow, I thought,
Pittsburgh could be the next Portland. It could become the cool place
all the people in their twenties want to go live.

When I got here a couple days ago, I could feel the difference. I lived
here from 1968 to 1984. I didn't realize it at the time, but during that
whole period the city was in free fall. On top of the flight to the
suburbs that happened everywhere, the steel and nuclear businesses were
both dying. Boy are things different now. It's not just that downtown
seems a lot more prosperous. There is an energy here that was not here
when I was a kid.

When I was a kid, this was a place young people left. Now it's a place
that attracts them.

What does that have to do with startups? Startups are made of people,
and the average age of the people in a typical startup is right in that
25 to 29 bracket.

I've seen how powerful it is for a city to have those people. Five years
ago they shifted the center of gravity of Silicon Valley from the
peninsula to San Francisco. Google and Facebook are on the peninsula,
but the next generation of big winners are all in SF. The reason the
center of gravity shifted was the talent war, for programmers
especially. Most 25 to 29 year olds want to live in the city, not down
in the boring suburbs. So whether they like it or not, founders know
they have to be in the city. I know multiple founders who would have
preferred to live down in the Valley proper, but who made themselves
move to SF because they knew otherwise they'd lose the talent war.

So being a magnet for people in their twenties is a very promising thing
to be. It's hard to imagine a place becoming a startup hub without also
being that. When I read that statistic about the increasing percentage
of 25 to 29 year olds, I had exactly the same feeling of excitement I
get when I see a startup's graphs start to creep upward off the x axis.

Nationally the percentage of 25 to 29 year olds is 6.8\%. That means
you're .8\% ahead. The population is 306,000, so we're talking about a
surplus of about 2500 people. That's the population of a small town, and
that's just the surplus. So you have a toehold. Now you just have to
expand it.

And though ``youth-driven food boom'' may sound frivolous, it is
anything but. Restaurants and cafes are a big part of the personality of
a city. Imagine walking down a street in Paris. What are you walking
past? Little restaurants and cafes. Imagine driving through some
depressing random exurb. What are you driving past? Starbucks and
McDonalds and Pizza Hut. As Gertrude Stein said, there is no there
there. You could be anywhere.

These independent restaurants and cafes are not just feeding people.
They're making there be a there here.

So here is my first concrete recommendation for turning Pittsburgh into
the next Silicon Valley: do everything you can to encourage this
youth-driven food boom. What could the city do? Treat the people
starting these little restaurants and cafes as your users, and go ask
them what they want. I can guess at least one thing they might want: a
fast permit process. San Francisco has left you a huge amount of room to
beat them in that department.

I know restaurants aren't the prime mover though. The prime mover, as
the Times article said, is cheap housing. That's a big advantage. But
that phrase ``cheap housing'' is a bit misleading. There are plenty of
places that are cheaper. What's special about Pittsburgh is not that
it's cheap, but that it's a cheap place you'd actually want to live.

Part of that is the buildings themselves. I realized a long time ago,
back when I was a poor twenty-something myself, that the best deals were
places that had once been rich, and then became poor. If a place has
always been rich, it's nice but too expensive. If a place has always
been poor, it's cheap but grim. But if a place was once rich and then
got poor, you can find palaces for cheap. And that's what's bringing
people here. When Pittsburgh was rich, a hundred years ago, the people
who lived here built big solid buildings. Not always in the best taste,
but definitely solid. So here is another piece of advice for becoming a
startup hub: don't destroy the buildings that are bringing people here.
When cities are on the way back up, like Pittsburgh is now, developers
race to tear down the old buildings. Don't let that happen. Focus on
historic preservation. Big real estate development projects are not
what's bringing the twenty-somethings here. They're the opposite of the
new restaurants and cafes; they subtract personality from the city.

The empirical evidence suggests you cannot be too strict about historic
preservation. The tougher cities are about it, the better they seem to
do.

But the appeal of Pittsburgh is not just the buildings themselves. It's
the neighborhoods they're in. Like San Francisco and New York,
Pittsburgh is fortunate in being a pre-car city. It's not too spread
out. Because those 25 to 29 year olds do not like driving. They prefer
walking, or bicycling, or taking public transport. If you've been to San
Francisco recently you can't help noticing the huge number of
bicyclists. And this is not just a fad that the twenty-somethings have
adopted. In this respect they have discovered a better way to live. The
beards will go, but not the bikes. Cities where you can get around
without driving are just better period. So I would suggest you do
everything you can to capitalize on this. As with historic preservation,
it seems impossible to go too far.

Why not make Pittsburgh the most bicycle and pedestrian friendly city in
the country? See if you can go so far that you make San Francisco seem
backward by comparison. If you do, it's very unlikely you'll regret it.
The city will seem like a paradise to the young people you want to
attract. If they do leave to get jobs elsewhere, it will be with regret
at leaving behind such a place. And what's the downside? Can you imagine
a headline ``City ruined by becoming too bicycle-friendly?'' It just
doesn't happen.

So suppose cool old neighborhoods and cool little restaurants make this
the next Portland. Will that be enough? It will put you in a way better
position than Portland itself, because Pittsburgh has something Portland
lacks: a first-rate research university. CMU plus little cafes means you
have more than hipsters drinking lattes. It means you have hipsters
drinking lattes while talking about distributed systems. Now you're
getting really close to San Francisco.

In fact you're better off than San Francisco in one way, because CMU is
downtown, but Stanford and Berkeley are out in the suburbs.

What can CMU do to help Pittsburgh become a startup hub? Be an even
better research university. CMU is one of the best universities in the
world, but imagine what things would be like if it were the very best,
and everyone knew it. There are a lot of ambitious people who must go to
the best place, wherever it is. If CMU were it, they would all come
here. There would be kids in Kazakhstan dreaming of one day living in
Pittsburgh.

Being that kind of talent magnet is the most important contribution
universities can make toward making their city a startup hub. In fact it
is practically the only contribution they can make.

But wait, shouldn't universities be setting up programs with words like
``innovation'' and ``entrepreneurship'' in their names? No, they should
not. These kind of things almost always turn out to be disappointments.
They're pursuing the wrong targets. The way to get innovation is not to
aim for innovation but to aim for something more specific, like better
batteries or better 3D printing. And the way to learn about
entrepreneurship is to do it, which you \href{before.html}{can't in
school}.

I know it may disappoint some administrators to hear that the best thing
a university can do to encourage startups is to be a great university.
It's like telling people who want to lose weight that the way to do it
is to eat less.

But if you want to know where startups come from, look at the empirical
evidence. Look at the histories of the most successful startups, and
you'll find they grow organically out of a couple of founders building
something that starts as an interesting side project. Universities are
great at bringing together founders, but beyond that the best thing they
can do is get out of the way. For example, by not claiming ownership of
``intellectual property'' that students and faculty develop, and by
having liberal rules about deferred admission and leaves of absence.

In fact, one of the most effective things a university could do to
encourage startups is an elaborate form of getting out of the way
invented by Harvard. Harvard used to have exams for the fall semester
after Christmas. At the beginning of January they had something called
``Reading Period'' when you were supposed to be studying for exams. And
Microsoft and Facebook have something in common that few people realize:
they were both started during Reading Period. It's the perfect situation
for producing the sort of side projects that turn into startups. The
students are all on campus, but they don't have to do anything because
they're supposed to be studying for exams.

Harvard may have closed this window, because a few years ago they moved
exams before Christmas and shortened reading period from 11 days to 7.
But if a university really wanted to help its students start startups,
the empirical evidence, weighted by market cap, suggests the best thing
they can do is literally nothing.

The culture of Pittsburgh is another of its strengths. It seems like a
city has to be socially liberal to be a startup hub, and it's pretty
clear why. A city has to tolerate strangeness to be a home for startups,
because startups are so strange. And you can't choose to allow just the
forms of strangeness that will turn into big startups, because they're
all intermingled. You have to tolerate all strangeness.

That immediately rules out
\href{http://www.nytimes.com/2016/04/06/us/gay-rights-mississippi-north-carolina.html}{big
chunks of the US}. I'm optimistic it doesn't rule out Pittsburgh. One of
the things I remember from growing up here, though I didn't realize at
the time that there was anything unusual about it, is how well people
got along. I'm still not sure why. Maybe one reason was that everyone
felt like an immigrant. When I was a kid in Monroeville, people didn't
call themselves American. They called themselves Italian or Serbian or
Ukranian. Just imagine what it must have been like here a hundred years
ago, when people were pouring in from twenty different countries.
Tolerance was the only option.

What I remember about the culture of Pittsburgh is that it was both
tolerant and pragmatic. That's how I'd describe the culture of Silicon
Valley too. And it's not a coincidence, because Pittsburgh was the
Silicon Valley of its time. This was a city where people built new
things. And while the things people build have changed, the spirit you
need to do that kind of work is the same.

So although an influx of latte-swilling hipsters may be annoying in some
ways, I would go out of my way to encourage them. And more generally to
tolerate strangeness, even unto the degree wacko Californians do. For
Pittsburgh that is a conservative choice: it's a return to the city's
roots.

Unfortunately I saved the toughest part for last. There is one more
thing you need to be a startup hub, and Pittsburgh hasn't got it:
investors. Silicon Valley has a big investor community because it's had
50 years to grow one. New York has a big investor community because it's
full of people who like money a lot and are quick to notice new ways to
get it. But Pittsburgh has neither of these. And the cheap housing that
draws other people here has no effect on investors.

If an investor community grows up here, it will happen the same way it
did in Silicon Valley: slowly and organically. So I would not bet on
having a big investor community in the short term. But fortunately there
are three trends that make that less necessary than it used to be. One
is that startups are increasingly cheap to start, so you just don't need
as much outside money as you used to. The second is that thanks to
things like Kickstarter, a startup can get to revenue faster. You can
put something on Kickstarter from anywhere. The third is programs like Y
Combinator. A startup from anywhere in the world can go to YC for 3
months, pick up funding, and then return home if they want.

My advice is to make Pittsburgh a great place for startups, and
gradually more of them will stick. Some of those will succeed; some of
their founders will become investors; and still more startups will
stick.

This is not a fast path to becoming a startup hub. But it is at least a
path, which is something few other cities have. And it's not as if you
have to make painful sacrifices in the meantime. Think about what I've
suggested you should do. Encourage local restaurants, save old
buildings, take advantage of density, make CMU the best, promote
tolerance. These are the things that make Pittsburgh good to live in
now. All I'm saying is that you should do even more of them.

And that's an encouraging thought. If Pittsburgh's path to becoming a
startup hub is to be even more itself, then it has a good chance of
succeeding. In fact it probably has the best chance of any city its
size. It will take some effort, and a lot of time, but if any city can
do it, Pittsburgh can.

\textbf{Thanks} to Charlie Cheever and Jessica Livingston for reading
drafts of this, and to Meg Cheever for organizing Opt412 and inviting me
to speak.
\chapter{The Risk of Discovery}

January 2017

Because biographies of famous scientists tend to edit out their
mistakes, we underestimate the degree of risk they were willing to take.
And because anything a famous scientist did that wasn't a mistake has
probably now become the conventional wisdom, those choices don't seem
risky either.

Biographies of Newton, for example, understandably focus more on physics
than alchemy or theology. The impression we get is that his unerring
judgment led him straight to truths no one else had noticed. How to
explain all the time he spent on alchemy and theology? Well, smart
people are often kind of crazy.

But maybe there is a simpler explanation. Maybe the smartness and the
craziness were not as separate as we think. Physics seems to us a
promising thing to work on, and alchemy and theology obvious wastes of
time. But that's because we know how things turned out. In Newton's day
the three problems seemed roughly equally promising. No one knew yet
what the payoff would be for inventing what we now call physics; if they
had, more people would have been working on it. And alchemy and theology
were still then in the category Marc Andreessen would describe as
``huge, if true.''

Newton made three bets. One of them worked. But they were all risky.
\chapter{Charisma / Power}

January 2017

People who are powerful but uncharismatic will tend to be disliked.
Their power makes them a target for criticism that they don't have the
charisma to disarm. That was Hillary Clinton's problem. It also tends to
be a problem for any CEO who is more of a builder than a schmoozer. And
yet the builder-type CEO is (like Hillary) probably the best person for
the job.

I don't think there is any solution to this problem. It's human nature.
The best we can do is to recognize that it's happening, and to
understand that being a magnet for criticism is sometimes a sign not
that someone is the wrong person for a job, but that they're the right
one.
\chapter{General and Surprising}

September 2017

The most valuable insights are both general and surprising. F~=~ma for
example. But general and surprising is a hard combination to achieve.
That territory tends to be picked clean, precisely because those
insights are so valuable.

Ordinarily, the best that people can do is one without the other: either
surprising without being general (e.g.~gossip), or general without being
surprising (e.g.~platitudes).

Where things get interesting is the moderately valuable insights. You
get those from small additions of whichever quality was missing. The
more common case is a small addition of generality: a piece of gossip
that's more than just gossip, because it teaches something interesting
about the world. But another less common approach is to focus on the
most general ideas and see if you can find something new to say about
them. Because these start out so general, you only need a small delta of
novelty to produce a useful insight.

A small delta of novelty is all you'll be able to get most of the time.
Which means if you take this route, your ideas will seem a lot like ones
that already exist. Sometimes you'll find you've merely rediscovered an
idea that did already exist. But don't be discouraged. Remember the huge
multiplier that kicks in when you do manage to think of something even a
little new.

Corollary: the more general the ideas you're talking about, the less you
should worry about repeating yourself. If you write enough, it's
inevitable you will. Your brain is much the same from year to year and
so are the stimuli that hit it. I feel slightly bad when I find I've
said something close to what I've said before, as if I were plagiarizing
myself. But rationally one shouldn't. You won't say something exactly
the same way the second time, and that variation increases the chance
you'll get that tiny but critical delta of novelty.

And of course, ideas beget ideas. (That sounds
\href{ecw.html}{familiar}.) An idea with a small amount of novelty could
lead to one with more. But only if you keep going. So it's doubly
important not to let yourself be discouraged by people who say there's
not much new about something you've discovered. ``Not much new'' is a
real achievement when you're talking about the most general ideas.

It's not true that there's nothing new under the sun. There are some
domains where there's almost nothing new. But there's a big difference
between nothing and almost nothing, when it's multiplied by the area
under the sun.

\textbf{Thanks} to Sam Altman, Patrick Collison, and Jessica Livingston
for reading drafts of this.
\chapter{The Bus Ticket Theory of Genius}

November 2019

Everyone knows that to do great work you need both natural ability and
determination. But there's a third ingredient that's not as well
understood: an obsessive interest in a particular topic.

To explain this point I need to burn my reputation with some group of
people, and I'm going to choose bus ticket collectors. There are people
who collect old bus tickets. Like many collectors, they have an
obsessive interest in the minutiae of what they collect. They can keep
track of distinctions between different types of bus tickets that would
be hard for the rest of us to remember. Because we don't care enough.
What's the point of spending so much time thinking about old bus
tickets?

Which leads us to the second feature of this kind of obsession: there is
no point. A bus ticket collector's love is disinterested. They're not
doing it to impress us or to make themselves rich, but for its own sake.

When you look at the lives of people who've done great work, you see a
consistent pattern. They often begin with a bus ticket collector's
obsessive interest in something that would have seemed pointless to most
of their contemporaries. One of the most striking features of Darwin's
book about his voyage on the Beagle is the sheer depth of his interest
in natural history. His curiosity seems infinite. Ditto for Ramanujan,
sitting by the hour working out on his slate what happens to series.

It's a mistake to think they were ``laying the groundwork'' for the
discoveries they made later. There's too much intention in that
metaphor. Like bus ticket collectors, they were doing it because they
liked it.

But there is a difference between Ramanujan and a bus ticket collector.
Series matter, and bus tickets don't.

If I had to put the recipe for genius into one sentence, that might be
it: to have a disinterested obsession with something that matters.

Aren't I forgetting about the other two ingredients? Less than you might
think. An obsessive interest in a topic is both a proxy for ability and
a substitute for determination. Unless you have sufficient mathematical
aptitude, you won't find series interesting. And when you're obsessively
interested in something, you don't need as much determination: you don't
need to push yourself as hard when curiosity is pulling you.

An obsessive interest will even bring you luck, to the extent anything
can. Chance, as Pasteur said, favors the prepared mind, and if there's
one thing an obsessed mind is, it's prepared.

The disinterestedness of this kind of obsession is its most important
feature. Not just because it's a filter for earnestness, but because it
helps you discover new ideas.

The paths that lead to new ideas tend to look unpromising. If they
looked promising, other people would already have explored them. How do
the people who do great work discover these paths that others overlook?
The popular story is that they simply have better vision: because
they're so talented, they see paths that others miss. But if you look at
the way great discoveries are made, that's not what happens. Darwin
didn't pay closer attention to individual species than other people
because he saw that this would lead to great discoveries, and they
didn't. He was just really, really interested in such things.

Darwin couldn't turn it off. Neither could Ramanujan. They didn't
discover the hidden paths that they did because they seemed promising,
but because they couldn't help it. That's what allowed them to follow
paths that someone who was merely ambitious would have ignored.

What rational person would decide that the way to write great novels was
to begin by spending several years creating an imaginary elvish
language, like Tolkien, or visiting every household in southwestern
Britain, like Trollope? No one, including Tolkien and Trollope.

The bus ticket theory is similar to Carlyle's famous definition of
genius as an infinite capacity for taking pains. But there are two
differences. The bus ticket theory makes it clear that the source of
this infinite capacity for taking pains is not infinite diligence, as
Carlyle seems to have meant, but the sort of infinite interest that
collectors have. It also adds an important qualification: an infinite
capacity for taking pains about something that matters.

So what matters? You can never be sure. It's precisely because no one
can tell in advance which paths are promising that you can discover new
ideas by working on what you're interested in.

But there are some heuristics you can use to guess whether an obsession
might be one that matters. For example, it's more promising if you're
creating something, rather than just consuming something someone else
creates. It's more promising if something you're interested in is
difficult, especially if it's \href{work.html}{more difficult for other
people} than it is for you. And the obsessions of talented people are
more likely to be promising. When talented people become interested in
random things, they're not truly random.

But you can never be sure. In fact, here's an interesting idea that's
also rather alarming if it's true: it may be that to do great work, you
also have to waste a lot of time.

In many different areas, reward is proportionate to risk. If that rule
holds here, then the way to find paths that lead to truly great work is
to be willing to expend a lot of effort on things that turn out to be
every bit as unpromising as they seem.

I'm not sure if this is true. On one hand, it seems surprisingly
difficult to waste your time so long as you're working hard on something
interesting. So much of what you do ends up being useful. But on the
other hand, the rule about the relationship between risk and reward is
so powerful that it seems to hold wherever risk occurs.
\href{disc.html}{Newton's} case, at least, suggests that the risk/reward
rule holds here. He's famous for one particular obsession of his that
turned out to be unprecedentedly fruitful: using math to describe the
world. But he had two other obsessions, alchemy and theology, that seem
to have been complete wastes of time. He ended up net ahead. His bet on
what we now call physics paid off so well that it more than compensated
for the other two. But were the other two necessary, in the sense that
he had to take big risks to make such big discoveries? I don't know.

Here's an even more alarming idea: might one make all bad bets? It
probably happens quite often. But we don't know how often, because these
people don't become famous.

It's not merely that the returns from following a path are hard to
predict. They change dramatically over time. 1830 was a really good time
to be obsessively interested in natural history. If Darwin had been born
in 1709 instead of 1809, we might never have heard of him.

What can one do in the face of such uncertainty? One solution is to
hedge your bets, which in this case means to follow the obviously
promising paths instead of your own private obsessions. But as with any
hedge, you're decreasing reward when you decrease risk. If you forgo
working on what you like in order to follow some more conventionally
ambitious path, you might miss something wonderful that you'd otherwise
have discovered. That too must happen all the time, perhaps even more
often than the genius whose bets all fail.

The other solution is to let yourself be interested in lots of different
things. You don't decrease your upside if you switch between equally
genuine interests based on which seems to be working so far. But there
is a danger here too: if you work on too many different projects, you
might not get deeply enough into any of them.

One interesting thing about the bus ticket theory is that it may help
explain why different types of people excel at different kinds of work.
Interest is much more unevenly distributed than ability. If natural
ability is all you need to do great work, and natural ability is evenly
distributed, you have to invent elaborate theories to explain the skewed
distributions we see among those who actually do great work in various
fields. But it may be that much of the skew has a simpler explanation:
different people are interested in different things.

The bus ticket theory also explains why people are less likely to do
great work after they have children. Here interest has to compete not
just with external obstacles, but with another interest, and one that
for most people is extremely powerful. It's harder to find time for work
after you have kids, but that's the easy part. The real change is that
you don't want to.

But the most exciting implication of the bus ticket theory is that it
suggests ways to encourage great work. If the recipe for genius is
simply natural ability plus hard work, all we can do is hope we have a
lot of ability, and work as hard as we can. But if interest is a
critical ingredient in genius, we may be able, by cultivating interest,
to cultivate genius.

For example, for the very ambitious, the bus ticket theory suggests that
the way to do great work is to relax a little. Instead of gritting your
teeth and diligently pursuing what all your peers agree is the most
promising line of research, maybe you should try doing something just
for fun. And if you're stuck, that may be the vector along which to
break out.

I've always liked \href{hamming.html}{Hamming's} famous double-barrelled
question: what are the most important problems in your field, and why
aren't you working on one of them? It's a great way to shake yourself
up. But it may be overfitting a bit. It might be at least as useful to
ask yourself: if you could take a year off to work on something that
probably wouldn't be important but would be really interesting, what
would it be?

The bus ticket theory also suggests a way to avoid slowing down as you
get older. Perhaps the reason people have fewer new ideas as they get
older is not simply that they're losing their edge. It may also be
because once you become established, you can no longer mess about with
irresponsible side projects the way you could when you were young and no
one cared what you did.

The solution to that is obvious: remain irresponsible. It will be hard,
though, because the apparently random projects you take up to stave off
decline will read to outsiders as evidence of it. And you yourself won't
know for sure that they're wrong. But it will at least be more fun to
work on what you want.

It may even be that we can cultivate a habit of intellectual bus ticket
collecting in kids. The usual plan in education is to start with a
broad, shallow focus, then gradually become more specialized. But I've
done the opposite with my kids. I know I can count on their school to
handle the broad, shallow part, so I take them deep.

When they get interested in something, however random, I encourage them
to go preposterously, bus ticket collectorly, deep. I don't do this
because of the bus ticket theory. I do it because I want them to feel
the joy of learning, and they're never going to feel that about
something I'm making them learn. It has to be something they're
interested in. I'm just following the path of least resistance; depth is
a byproduct. But if in trying to show them the joy of learning I also
end up training them to go deep, so much the better.

Will it have any effect? I have no idea. But that uncertainty may be the
most interesting point of all. There is so much more to learn about how
to do great work. As old as human civilization feels, it's really still
very young if we haven't nailed something so basic. It's exciting to
think there are still discoveries to make about discovery. If that's the
sort of thing you're interested in.


\footnote{There are other types of collecting that illustrate this point
better than bus tickets, but they're also more popular. It seemed just
as well to use an inferior example rather than offend more people by
telling them their hobby doesn't matter.}\footnote{I worried a little about using the word ``disinterested,'' since
some people mistakenly believe it means not interested. But anyone who
expects to be a genius will have to know the meaning of such a basic
word, so I figure they may as well start now.}\footnote{Think how often genius must have been nipped in the bud by
people being told, or telling themselves, to stop messing about and be
responsible. Ramanujan's mother was a huge enabler. Imagine if she
hadn't been. Imagine if his parents had made him go out and get a job
instead of sitting around at home doing math.

On the other hand, anyone quoting the preceding paragraph to justify not
getting a job is probably mistaken.}\footnote{1709 Darwin is to time what the \href{cities.html}{Milanese
Leonardo} is to space.}\footnote{``An infinite capacity for taking pains'' is a paraphrase of
what Carlyle wrote. What he wrote, in his \emph{History of Frederick the
Great}, was ``\ldots{} it is the fruit of `genius' (which means
transcendent capacity of taking trouble, first of all)\ldots.'' Since
the paraphrase seems the name of the idea at this point, I kept it.

Carlyle's \emph{History} was published in 1858. In 1785 Hérault de
Séchelles quoted Buffon as saying ``Le génie n'est qu'une plus grande
aptitude à la patience.'' (Genius is only a greater aptitude for
patience.)}\footnote{Trollope was establishing the system of postal routes. He
himself sensed the obsessiveness with which he pursued this goal.

\begin{quote}
It is amusing to watch how a passion will grow upon a man. During those
two years it was the ambition of my life to cover the country with rural
letter-carriers.
\end{quote}

Even Newton occasionally sensed the degree of his obsessiveness. After
computing pi to 15 digits, he wrote in a letter to a friend:

\begin{quote}
I am ashamed to tell you to how many figures I carried these
computations, having no other business at the time.
\end{quote}

Incidentally, Ramanujan was also a compulsive calculator. As Kanigel
writes in his excellent biography:

\begin{quote}
One Ramanujan scholar, B. M. Wilson, later told how Ramanujan's research
into number theory was often ``preceded by a table of numerical results,
carried usually to a length from which most of us would shrink.''
\end{quote}}\footnote{Working to understand the natural world counts as creating
rather than consuming.

Newton tripped over this distinction when he chose to work on theology.
His beliefs did not allow him to see it, but chasing down paradoxes in
nature is fruitful in a way that chasing down paradoxes in sacred texts
is not.}\footnote{How much of people's propensity to become interested in a topic
is inborn? My experience so far suggests the answer is: most of it.
Different kids get interested in different things, and it's hard to make
a child interested in something they wouldn't otherwise be. Not in a way
that sticks. The most you can do on behalf of a topic is to make sure it
gets a fair showing --- to make it clear to them, for example, that
there's more to math than the dull drills they do in school. After that
it's up to the child.}

\textbf{Thanks} to Marc Andreessen, Trevor Blackwell, Patrick Collison,
Kevin Lacker, Jessica Livingston, Jackie McDonough, Robert Morris, Lisa
Randall, Zak Stone, and
\href{https://twitter.com/paulg/status/1196537802621669376}{my 7 year
old} for reading drafts of this.
\chapter{Novelty and Heresy}

November 2019

If you discover something new, there's a significant chance you'll be
accused of some form of heresy.

To discover new things, you have to work on ideas that are good but
non-obvious; if an idea is obviously good, other people are probably
already working on it. One common way for a good idea to be non-obvious
is for it to be hidden in the shadow of some mistaken assumption that
people are very attached to. But anything you discover from working on
such an idea will tend to contradict the mistaken assumption that was
concealing it. And you will thus get a lot of heat from people attached
to the mistaken assumption. Galileo and Darwin are famous examples of
this phenomenon, but it's probably always an ingredient in the
resistance to new ideas.

So it's particularly dangerous for an organization or society to have a
culture of pouncing on heresy. When you suppress heresies, you don't
just prevent people from contradicting the mistaken assumption you're
trying to protect. You also suppress any idea that implies indirectly
that it's false.

Every cherished mistaken assumption has a dead zone of unexplored ideas
around it. And the more preposterous the assumption, the bigger the dead
zone it creates.

There is a positive side to this phenomenon though. If you're looking
for new ideas, one way to find them is by \href{say.html}{looking for
heresies}. When you look at the question this way, the depressingly
large dead zones around mistaken assumptions become excitingly large
mines of new ideas.
\chapter{The Lesson to Unlearn}

December 2019

The most damaging thing you learned in school wasn't something you
learned in any specific class. It was learning to get good grades.

When I was in college, a particularly earnest philosophy grad student
once told me that he never cared what grade he got in a class, only what
he learned in it. This stuck in my mind because it was the only time I
ever heard anyone say such a thing.

For me, as for most students, the measurement of what I was learning
completely dominated actual learning in college. I was fairly earnest; I
was genuinely interested in most of the classes I took, and I worked
hard. And yet I worked by far the hardest when I was studying for a
test.

In theory, tests are merely what their name implies: tests of what
you've learned in the class. In theory you shouldn't have to prepare for
a test in a class any more than you have to prepare for a blood test. In
theory you learn from taking the class, from going to the lectures and
doing the reading and/or assignments, and the test that comes afterward
merely measures how well you learned.

In practice, as almost everyone reading this will know, things are so
different that hearing this explanation of how classes and tests are
meant to work is like hearing the etymology of a word whose meaning has
changed completely. In practice, the phrase ``studying for a test'' was
almost redundant, because that was when one really studied. The
difference between diligent and slack students was that the former
studied hard for tests and the latter didn't. No one was pulling
all-nighters two weeks into the semester.

Even though I was a diligent student, almost all the work I did in
school was aimed at getting a good grade on something.

To many people, it would seem strange that the preceding sentence has a
``though'' in it. Aren't I merely stating a tautology? Isn't that what a
diligent student is, a straight-A student? That's how deeply the
conflation of learning with grades has infused our culture.

Is it so bad if learning is conflated with grades? Yes, it is bad. And
it wasn't till decades after college, when I was running Y~Combinator,
that I realized how bad it is.

I knew of course when I was a student that studying for a test is far
from identical with actual learning. At the very least, you don't retain
knowledge you cram into your head the night before an exam. But the
problem is worse than that. The real problem is that most tests don't
come close to measuring what they're supposed to.

If tests truly were tests of learning, things wouldn't be so bad.
Getting good grades and learning would converge, just a little late. The
problem is that nearly all tests given to students are terribly
hackable. Most people who've gotten good grades know this, and know it
so well they've ceased even to question it. You'll see when you realize
how naive it sounds to act otherwise.

Suppose you're taking a class on medieval history and the final exam is
coming up. The final exam is supposed to be a test of your knowledge of
medieval history, right? So if you have a couple days between now and
the exam, surely the best way to spend the time, if you want to do well
on the exam, is to read the best books you can find about medieval
history. Then you'll know a lot about it, and do well on the exam.

No, no, no, experienced students are saying to themselves. If you merely
read good books on medieval history, most of the stuff you learned
wouldn't be on the test. It's not good books you want to read, but the
lecture notes and assigned reading in this class. And even most of that
you can ignore, because you only have to worry about the sort of thing
that could turn up as a test question. You're looking for
sharply-defined chunks of information. If one of the assigned readings
has an interesting digression on some subtle point, you can safely
ignore that, because it's not the sort of thing that could be turned
into a test question. But if the professor tells you that there were
three underlying causes of the Schism of 1378, or three main
consequences of the Black Death, you'd better know them. And whether
they were in fact the causes or consequences is beside the point. For
the purposes of this class they are.

At a university there are often copies of old exams floating around, and
these narrow still further what you have to learn. As well as learning
what kind of questions this professor asks, you'll often get actual exam
questions. Many professors re-use them. After teaching a class for 10
years, it would be hard not to, at least inadvertently.

In some classes, your professor will have had some sort of political axe
to grind, and if so you'll have to grind it too. The need for this
varies. In classes in math or the hard sciences or engineering it's
rarely necessary, but at the other end of the spectrum there are classes
where you couldn't get a good grade without it.

Getting a good grade in a class on x is so different from learning a lot
about x that you have to choose one or the other, and you can't blame
students if they choose grades. Everyone judges them by their grades ---
graduate programs, employers, scholarships, even their own parents.

I liked learning, and I really enjoyed some of the papers and programs I
wrote in college. But did I ever, after turning in a paper in some
class, sit down and write another just for fun? Of course not. I had
things due in other classes. If it ever came to a choice of learning or
grades, I chose grades. I hadn't come to college to do badly.

Anyone who cares about getting good grades has to play this game, or
they'll be surpassed by those who do. And at elite universities, that
means nearly everyone, since someone who didn't care about getting good
grades probably wouldn't be there in the first place. The result is that
students compete to maximize the difference between learning and getting
good grades.

Why are tests so bad? More precisely, why are they so hackable? Any
experienced programmer could answer that. How hackable is software whose
author hasn't paid any attention to preventing it from being hacked?
Usually it's as porous as a colander.

Hackable is the default for any test imposed by an authority. The reason
the tests you're given are so consistently bad --- so consistently far
from measuring what they're supposed to measure --- is simply that the
people creating them haven't made much effort to prevent them from being
hacked.

But you can't blame teachers if their tests are hackable. Their job is
to teach, not to create unhackable tests. The real problem is grades, or
more precisely, that grades have been overloaded. If grades were merely
a way for teachers to tell students what they were doing right and
wrong, like a coach giving advice to an athlete, students wouldn't be
tempted to hack tests. But unfortunately after a certain age grades
become more than advice. After a certain age, whenever you're being
taught, you're usually also being judged.

I've used college tests as an example, but those are actually the least
hackable. All the tests most students take their whole lives are at
least as bad, including, most spectacularly of all, the test that gets
them into college. If getting into college were merely a matter of
having the quality of one's mind measured by admissions officers the way
scientists measure the mass of an object, we could tell teenage kids
``learn a lot'' and leave it at that. You can tell how bad college
admissions are, as a test, from how unlike high school that sounds. In
practice, the freakishly specific nature of the stuff ambitious kids
have to do in high school is directly proportionate to the hackability
of college admissions. The classes you don't care about that are mostly
memorization, the random ``extracurricular activities'' you have to
participate in to show you're ``well-rounded,'' the standardized tests
as artificial as chess, the ``essay'' you have to write that's
presumably meant to hit some very specific target, but you're not told
what.

As well as being bad in what it does to kids, this test is also bad in
the sense of being very hackable. So hackable that whole industries have
grown up to hack it. This is the explicit purpose of test-prep companies
and admissions counsellors, but it's also a significant part of the
function of private schools.

Why is this particular test so hackable? I think because of what it's
measuring. Although the popular story is that the way to get into a good
college is to be really smart, admissions officers at elite colleges
neither are, nor claim to be, looking only for that. What are they
looking for? They're looking for people who are not simply smart, but
admirable in some more general sense. And how is this more general
admirableness measured? The admissions officers feel it. In other words,
they accept who they like.

So what college admissions is a test of is whether you suit the taste of
some group of people. Well, of course a test like that is going to be
hackable. And because it's both very hackable and there's (thought to
be) a lot at stake, it's hacked like nothing else. That's why it
distorts your life so much for so long.

It's no wonder high school students often feel alienated. The shape of
their lives is completely artificial.

But wasting your time is not the worst thing the educational system does
to you. The worst thing it does is to train you that the way to win is
by hacking bad tests. This is a much subtler problem that I didn't
recognize until I saw it happening to other people.

When I started advising startup founders at Y Combinator, especially
young ones, I was puzzled by the way they always seemed to make things
overcomplicated. How, they would ask, do you raise money? What's the
trick for making venture capitalists want to invest in you? The best way
to make VCs want to invest in you, I would explain, is to actually be a
good investment. Even if you could trick VCs into investing in a bad
startup, you'd be tricking yourselves too. You're investing time in the
same company you're asking them to invest money in. If it's not a good
investment, why are you even doing it?

Oh, they'd say, and then after a pause to digest this revelation, they'd
ask: What makes a startup a good investment?

So I would explain that what makes a startup promising, not just in the
eyes of investors but in fact, is \href{growth.html}{growth}. Ideally in
revenue, but failing that in usage. What they needed to do was get lots
of users.

How does one get lots of users? They had all kinds of ideas about that.
They needed to do a big launch that would get them ``exposure.'' They
needed influential people to talk about them. They even knew they needed
to launch on a tuesday, because that's when one gets the most attention.

No, I would explain, that is not how to get lots of users. The way you
get lots of users is to make the product really great. Then people will
not only use it but recommend it to their friends, so your growth will
be exponential once you \href{ds.html}{get it started}.

At this point I've told the founders something you'd think would be
completely obvious: that they should make a good company by making a
good product. And yet their reaction would be something like the
reaction many physicists must have had when they first heard about the
theory of relativity: a mixture of astonishment at its apparent genius,
combined with a suspicion that anything so weird couldn't possibly be
right. Ok, they would say, dutifully. And could you introduce us to
such-and-such influential person? And remember, we want to launch on
Tuesday.

It would sometimes take founders years to grasp these simple lessons.
And not because they were lazy or stupid. They just seemed blind to what
was right in front of them.

Why, I would ask myself, do they always make things so complicated? And
then one day I realized this was not a rhetorical question.

Why did founders tie themselves in knots doing the wrong things when the
answer was right in front of them? Because that was what they'd been
trained to do. Their education had taught them that the way to win was
to hack the test. And without even telling them they were being trained
to do this. The younger ones, the recent graduates, had never faced a
non-artificial test. They thought this was just how the world worked:
that the first thing you did, when facing any kind of challenge, was to
figure out what the trick was for hacking the test. That's why the
conversation would always start with how to raise money, because that
read as the test. It came at the end of YC. It had numbers attached to
it, and higher numbers seemed to be better. It must be the test.

There are certainly big chunks of the world where the way to win is to
hack the test. This phenomenon isn't limited to schools. And some
people, either due to ideology or ignorance, claim that this is true of
startups too. But it isn't. In fact, one of the most striking things
about startups is the degree to which you win by simply doing good work.
There are edge cases, as there are in anything, but in general you win
by getting users, and what users care about is whether the product does
what they want.

Why did it take me so long to understand why founders made startups
overcomplicated? Because I hadn't realized explicitly that schools train
us to win by hacking bad tests. And not just them, but me! I'd been
trained to hack bad tests too, and hadn't realized it till decades
later.

I had lived as if I realized it, but without knowing why. For example, I
had avoided working for big companies. But if you'd asked why, I'd have
said it was because they were bogus, or bureaucratic. Or just yuck. I
never understood how much of my dislike of big companies was due to the
fact that you win by hacking bad tests.

Similarly, the fact that the tests were unhackable was a lot of what
attracted me to startups. But again, I hadn't realized that explicitly.

I had in effect achieved by successive approximations something that may
have a closed-form solution. I had gradually undone my training in
hacking bad tests without knowing I was doing it. Could someone coming
out of school banish this demon just by knowing its name, and saying
begone? It seems worth trying.

Merely talking explicitly about this phenomenon is likely to make things
better, because much of its power comes from the fact that we take it
for granted. After you've noticed it, it seems the elephant in the room,
but it's a pretty well camouflaged elephant. The phenomenon is so old,
and so pervasive. And it's simply the result of neglect. No one meant
things to be this way. This is just what happens when you combine
learning with grades, competition, and the naive assumption of
unhackability.

It was mind-blowing to realize that two of the things I'd puzzled about
the most --- the bogusness of high school, and the difficulty of getting
founders to see the obvious --- both had the same cause. It's rare for
such a big block to slide into place so late.

Usually when that happens it has implications in a lot of different
areas, and this case seems no exception. For example, it suggests both
that education could be done better, and how you might fix it. But it
also suggests a potential answer to the question all big companies seem
to have: how can we be more like a startup? I'm not going to chase down
all the implications now. What I want to focus on here is what it means
for individuals.

To start with, it means that most ambitious kids graduating from college
have something they may want to unlearn. But it also changes how you
look at the world. Instead of looking at all the different kinds of work
people do and thinking of them vaguely as more or less appealing, you
can now ask a very specific question that will sort them in an
interesting way: to what extent do you win at this kind of work by
hacking bad tests?

It would help if there was a way to recognize bad tests quickly. Is
there a pattern here? It turns out there is.

Tests can be divided into two kinds: those that are imposed by
authorities, and those that aren't. Tests that aren't imposed by
authorities are inherently unhackable, in the sense that no one is
claiming they're tests of anything more than they actually test. A
football match, for example, is simply a test of who wins, not which
team is better. You can tell that from the fact that commentators
sometimes say afterward that the better team won. Whereas tests imposed
by authorities are usually proxies for something else. A test in a class
is supposed to measure not just how well you did on that particular
test, but how much you learned in the class. While tests that aren't
imposed by authorities are inherently unhackable, those imposed by
authorities have to be made unhackable. Usually they aren't. So as a
first approximation, bad tests are roughly equivalent to tests imposed
by authorities.

You might actually like to win by hacking bad tests. Presumably some
people do. But I bet most people who find themselves doing this kind of
work don't like it. They just take it for granted that this is how the
world works, unless you want to drop out and be some kind of hippie
artisan.

I suspect many people implicitly assume that working in a field with bad
tests is the price of making lots of money. But that, I can tell you, is
false. It used to be true. In the mid-twentieth century, when the
economy was \href{re.html}{composed of oligopolies}, the only way to the
top was by playing their game. But it's not true now. There are now ways
to get rich by doing good work, and that's part of the reason people are
so much more excited about getting rich than they used to be. When I was
a kid, you could either become an engineer and make cool things, or make
lots of money by becoming an ``executive.'' Now you can make lots of
money by making cool things.

Hacking bad tests is becoming less important as the link between work
and authority erodes. The erosion of that link is one of the most
important trends happening now, and we see its effects in almost every
kind of work people do. Startups are one of the most visible examples,
but we see much the same thing in writing. Writers no longer have to
submit to publishers and editors to reach readers; now they can go
direct.

The more I think about this question, the more optimistic I get. This
seems one of those situations where we don't realize how much something
was holding us back until it's eliminated. And I can foresee the whole
bogus edifice crumbling. Imagine what happens as more and more people
start to ask themselves if they want to win by hacking bad tests, and
decide that they don't. The kinds of work where you win by hacking bad
tests will be starved of talent, and the kinds where you win by doing
good work will see an influx of the most ambitious people. And as
hacking bad tests shrinks in importance, education will evolve to stop
training us to do it. Imagine what the world could look like if that
happened.

This is not just a lesson for individuals to unlearn, but one for
society to unlearn, and we'll be amazed at the energy that's liberated
when we do.


\footnote{If using tests only to measure learning sounds impossibly
utopian, that is already the way things work at Lambda School. Lambda
School doesn't have grades. You either graduate or you don't. The only
purpose of tests is to decide at each stage of the curriculum whether
you can continue to the next. So in effect the whole school is
pass/fail.}\footnote{If the final exam consisted of a long conversation with the
professor, you could prepare for it by reading good books on medieval
history. A lot of the hackability of tests in schools is due to the fact
that the same test has to be given to large numbers of students.}\footnote{Learning is the naive algorithm for getting good grades.}\footnote{\href{gba.html}{Hacking} has multiple senses. There's a narrow
sense in which it means to compromise something. That's the sense in
which one hacks a bad test. But there's another, more general sense,
meaning to find a surprising solution to a problem, often by thinking
differently about it. Hacking in this sense is a wonderful thing. And
indeed, some of the hacks people use on bad tests are impressively
ingenious; the problem is not so much the hacking as that, because the
tests are hackable, they don't test what they're meant to.}\footnote{The people who pick startups at Y Combinator are similar to
admissions officers, except that instead of being arbitrary, their
acceptance criteria are trained by a very tight feedback loop. If you
accept a bad startup or reject a good one, you will usually know it
within a year or two at the latest, and often within a month.}\footnote{I'm sure admissions officers are tired of reading applications
from kids who seem to have no personality beyond being willing to seem
however they're supposed to seem to get accepted. What they don't
realize is that they are, in a sense, looking in a mirror. The lack of
authenticity in the applicants is a reflection of the arbitrariness of
the application process. A dictator might just as well complain about
the lack of authenticity in the people around him.}\footnote{By good work, I don't mean morally good, but good in the sense
in which a good craftsman does good work.}\footnote{There are borderline cases where it's hard to say which category
a test falls in. For example, is raising venture capital like college
admissions, or is it like selling to a customer?}\footnote{Note that a good test is merely one that's unhackable. Good here
doesn't mean morally good, but good in the sense of working well. The
difference between fields with bad tests and good ones is not that the
former are bad and the latter are good, but that the former are bogus
and the latter aren't. But those two measures are not unrelated. As Tara
Ploughman said, the path from good to evil goes through bogus.}\footnote{People who think the recent increase in
\href{ineq.html}{economic inequality} is due to changes in tax policy
seem very naive to anyone with experience in startups. Different people
are getting rich now than used to, and they're getting much richer than
mere tax savings could make them.}\footnote{Note to tiger parents: you may think you're training your kids
to win, but if you're training them to win by hacking bad tests, you
are, as parents so often do, training them to fight the last war.}

\textbf{Thanks} to Austen Allred, Trevor Blackwell, Patrick Collison,
Jessica Livingston, Robert Morris, and Harj Taggar for reading drafts of
this.
\chapter{Having Kids}

December 2019

Before I had kids, I was afraid of having kids. Up to that point I felt
about kids the way the young Augustine felt about living virtuously. I'd
have been sad to think I'd never have children. But did I want them now?
No.~

If I had kids, I'd become a parent, and parents, as I'd known since I
was a kid, were uncool. They were dull and responsible and had no fun.
And while it's not surprising that kids would believe that, to be honest
I hadn't seen much as an adult to change my mind. Whenever I'd noticed
parents with kids, the kids seemed to be terrors, and the parents
pathetic harried creatures, even when they prevailed.

When people had babies, I congratulated them enthusiastically, because
that seemed to be what one did. But I didn't feel it at all. ``Better
you than me,'' I was thinking.

Now when people have babies I congratulate them enthusiastically and I
mean it. Especially the first one. I feel like they just got the best
gift in the world.

What changed, of course, is that I had kids. Something I dreaded turned
out to be wonderful.

Partly, and I won't deny it, this is because of serious chemical changes
that happened almost instantly when our first child was born. It was
like someone flipped a switch. I suddenly felt protective not just
toward our child, but toward all children. As I was driving my wife and
new son home from the hospital, I approached a crosswalk full of
pedestrians, and I found myself thinking ``I have to be really careful
of all these people. Every one of them is someone's child!''

So to some extent you can't trust me when I say having kids is great. To
some extent I'm like a religious cultist telling you that you'll be
happy if you join the cult too --- but only because joining the cult
will alter your mind in a way that will make you happy to be a cult
member.

But not entirely. There were some things about having kids that I
clearly got wrong before I had them.

For example, there was a huge amount of selection bias in my
observations of parents and children. Some parents may have noticed that
I wrote ``Whenever I'd noticed parents with kids.'' Of course the times
I noticed kids were when things were going wrong. I only noticed them
when they made noise. And where was I when I noticed them? Ordinarily I
never went to places with kids, so the only times I encountered them
were in shared bottlenecks like airplanes. Which is not exactly a
representative sample. Flying with a toddler is something very few
parents enjoy.

What I didn't notice, because they tend to be much quieter, were all the
great moments parents had with kids. People don't talk about these much
--- the magic is hard to put into words, and all other parents know
about them anyway --- but one of the great things about having kids is
that there are so many times when you feel there is nowhere else you'd
rather be, and nothing else you'd rather be doing. You don't have to be
doing anything special. You could just be going somewhere together, or
putting them to bed, or pushing them on the swings at the park. But you
wouldn't trade these moments for anything. One doesn't tend to associate
kids with peace, but that's what you feel. You don't need to look any
further than where you are right now.

Before I had kids, I had moments of this kind of peace, but they were
rarer. With kids it can happen several times a day.

My other source of data about kids was my own childhood, and that was
similarly misleading. I was pretty bad, and was always in trouble for
something or other. So it seemed to me that parenthood was essentially
law enforcement. I didn't realize there were good times too.

I remember my mother telling me once when I was about 30 that she'd
really enjoyed having me and my sister. My god, I thought, this woman is
a saint. She not only endured all the pain we subjected her to, but
actually enjoyed it? Now I realize she was simply telling the truth.

She said that one reason she liked having us was that we'd been
interesting to talk to. That took me by surprise when I had kids. You
don't just love them. They become your friends too. They're really
interesting. And while I admit small children are disastrously fond of
repetition (anything worth doing once is worth doing fifty times) it's
often genuinely fun to play with them. That surprised me too. Playing
with a 2 year old was fun when I was 2 and definitely not fun when I was
6. Why would it become fun again later? But it does.

There are of course times that are pure drudgery. Or worse still,
terror. Having kids is one of those intense types of experience that are
hard to imagine unless you've had them. But it is not, as I implicitly
believed before having kids, simply your DNA heading for the lifeboats.

Some of my worries about having kids were right, though. They definitely
make you less productive. I know having kids makes some people get their
act together, but if your act was already together, you're going to have
less time to do it in. In particular, you're going to have to work to a
schedule. Kids have schedules. I'm not sure if it's because that's how
kids are, or because it's the only way to integrate their lives with
adults', but once you have kids, you tend to have to work on their
schedule.

You will have chunks of time to work. But you can't let work spill
promiscuously through your whole life, like I used to before I had kids.
You're going to have to work at the same time every day, whether
inspiration is flowing or not, and there are going to be times when you
have to stop, even if it is.

I've been able to adapt to working this way. Work, like love, finds a
way. If there are only certain times it can happen, it happens at those
times. So while I don't get as much done as before I had kids, I get
enough done.

I hate to say this, because being ambitious has always been a part of my
identity, but having kids may make one less ambitious. It hurts to see
that sentence written down. I squirm to avoid it. But if there weren't
something real there, why would I squirm? The fact is, once you have
kids, you're probably going to care more about them than you do about
yourself. And attention is a zero-sum game. Only one idea at a time can
be the \href{top.html}{top idea in your mind}. Once you have kids, it
will often be your kids, and that means it will less often be some
project you're working on.

I have some hacks for sailing close to this wind. For example, when I
write essays, I think about what I'd want my kids to know. That drives
me to get things right. And when I was writing \href{bel.html}{Bel}, I
told my kids that once I finished it I'd take them to Africa. When you
say that sort of thing to a little kid, they treat it as a promise.
Which meant I had to finish or I'd be taking away their trip to Africa.
Maybe if I'm really lucky such tricks could put me net ahead. But the
wind is there, no question.

On the other hand, what kind of wimpy ambition do you have if it won't
survive having kids? Do you have so little to spare?

And while having kids may be warping my present judgement, it hasn't
overwritten my memory. I remember perfectly well what life was like
before. Well enough to miss some things a lot, like the ability to take
off for some other country at a moment's notice. That was so great. Why
did I never do that?

See what I did there? The fact is, most of the freedom I had before
kids, I never used. I paid for it in loneliness, but I never used it.

I had plenty of happy times before I had kids. But if I count up happy
moments, not just potential happiness but actual happy moments, there
are more after kids than before. Now I practically have it on tap,
almost any bedtime.

People's experiences as parents vary a lot, and I know I've been lucky.
But I think the worries I had before having kids must be pretty common,
and judging by other parents' faces when they see their kids, so must
the happiness that kids bring.

\textbf{Note}

\footnote{Adults are sophisticated enough to see 2 year olds for the
fascinatingly complex characters they are, whereas to most 6 year olds,
2 year olds are just defective 6 year olds.}

\textbf{Thanks} to Trevor Blackwell, Jessica Livingston, and Robert
Morris for reading drafts of this.
\chapter{Fashionable Problems}

December 2019

I've seen the same pattern in many different fields: even though lots of
people have worked hard in the field, only a small fraction of the space
of possibilities has been explored, because they've all worked on
similar things.

Even the smartest, most imaginative people are surprisingly conservative
when deciding what to work on. People who would never dream of being
fashionable in any other way get sucked into working on fashionable
problems.

If you want to try working on unfashionable problems, one of the best
places to look is in fields that people think have already been fully
explored: essays, Lisp, venture funding --- you may notice a pattern
here. If you can find a new approach into a big but apparently played
out field, the value of whatever you discover will be
\href{sun.html}{multiplied} by its enormous surface area.

The best protection against getting drawn into working on the same
things as everyone else may be to \href{genius.html}{genuinely love}
what you're doing. Then you'll continue to work on it even if you make
the same mistake as other people and think that it's too marginal to
matter.
\chapter{The Two Kinds of Moderate}

December 2019

There are two distinct ways to be politically moderate: on purpose and
by accident. Intentional moderates are trimmers, deliberately choosing a
position mid-way between the extremes of right and left. Accidental
moderates end up in the middle, on average, because they make up their
own minds about each question, and the far right and far left are
roughly equally wrong.

You can distinguish intentional from accidental moderates by the
distribution of their opinions. If the far left opinion on some matter
is 0 and the far right opinion 100, an intentional moderate's opinion on
every question will be near 50. Whereas an accidental moderate's
opinions will be scattered over a broad range, but will, like those of
the intentional moderate, average to about 50.

Intentional moderates are similar to those on the far left and the far
right in that their opinions are, in a sense, not their own. The
defining quality of an ideologue, whether on the left or the right, is
to acquire one's opinions in bulk. You don't get to pick and choose.
Your opinions about taxation can be predicted from your opinions about
sex. And although intentional moderates might seem to be the opposite of
ideologues, their beliefs (though in their case the word ``positions''
might be more accurate) are also acquired in bulk. If the median opinion
shifts to the right or left, the intentional moderate must shift with
it. Otherwise they stop being moderate.

Accidental moderates, on the other hand, not only choose their own
answers, but choose their own questions. They may not care at all about
questions that the left and right both think are terribly important. So
you can only even measure the politics of an accidental moderate from
the intersection of the questions they care about and those the left and
right care about, and this can sometimes be vanishingly small.

It is not merely a manipulative rhetorical trick to say ``if you're not
with us, you're against us,'' but often simply false.

Moderates are sometimes derided as cowards, particularly by the extreme
left. But while it may be accurate to call intentional moderates
cowards, openly being an accidental moderate requires the most courage
of all, because you get attacked from both right and left, and you don't
have the comfort of being an orthodox member of a large group to sustain
you.

Nearly all the most impressive people I know are accidental moderates.
If I knew a lot of professional athletes, or people in the entertainment
business, that might be different. Being on the far left or far right
doesn't affect how fast you run or how well you sing. But someone who
works with ideas has to be independent-minded to do it well.

Or more precisely, you have to be independent-minded about the ideas you
work with. You could be mindlessly doctrinaire in your politics and
still be a good mathematician. In the 20th century, a lot of very smart
people were Marxists --- just no one who was smart about the subjects
Marxism involves. But if the ideas you use in your work intersect with
the politics of your time, you have two choices: be an accidental
moderate, or be mediocre.


\footnote{It's possible in theory for one side to be entirely right and
the other to be entirely wrong. Indeed, ideologues must always believe
this is the case. But historically it rarely has been.}\footnote{For some reason the far right tend to ignore moderates rather
than despise them as backsliders. I'm not sure why. Perhaps it means
that the far right is less ideological than the far left. Or perhaps
that they are more confident, or more resigned, or simply more
disorganized. I just don't know.}\footnote{Having heretical opinions doesn't mean you have to express them
openly. It may be \href{say.html}{easier to have them} if you don't.}

\textbf{Thanks} to Austen Allred, Trevor Blackwell, Patrick Collison,
Jessica Livingston, Amjad Masad, Ryan Petersen, and Harj Taggar for
reading drafts of this.
\chapter{Haters}

January 2020

\emph{(I originally intended this for startup founders, who are often
surprised by the attention they get as their companies grow, but it
applies equally to anyone who becomes famous.)}

If you become sufficiently famous, you'll acquire some fans who like you
too much. These people are sometimes called ``fanboys,'' and though I
dislike that term, I'm going to have to use it here. We need some word
for them, because this is a distinct phenomenon from someone simply
liking your work.

A fanboy is obsessive and uncritical. Liking you becomes part of their
identity, and they create an image of you in their own head that is much
better than reality. Everything you do is good, because you do it. If
you do something bad, they find a way to see it as good. And their love
for you is not, usually, a quiet, private one. They want everyone to
know how great you are.

Well, you may be thinking, I could do without this kind of obsessive
fan, but I know there are all kinds of people in the world, and if this
is the worst consequence of fame, that's not so bad.

Unfortunately this is not the worst consequence of fame. As well as
fanboys, you'll have haters.

A hater is obsessive and uncritical. Disliking you becomes part of their
identity, and they create an image of you in their own head that is much
worse than reality. Everything you do is bad, because you do it. If you
do something good, they find a way to see it as bad. And their dislike
for you is not, usually, a quiet, private one. They want everyone to
know how awful you are.

If you're thinking of checking, I'll save you the trouble. The second
and fifth paragraphs are identical except for ``good'' being switched to
``bad'' and so on.

I spent years puzzling about haters. What are they, and where do they
come from? Then one day it dawned on me. Haters are just fanboys with
the sign switched.

Note that by haters, I don't simply mean trolls. I'm not talking about
people who say bad things about you and then move on. I'm talking about
the much smaller group of people for whom this becomes a kind of
obsession and who do it repeatedly over a long period.

Like fans, haters seem to be an automatic consequence of fame. Anyone
sufficiently famous will have them. And like fans, haters are energized
by the fame of whoever they hate. They hear a song by some pop singer.
They don't like it much. If the singer were an obscure one, they'd just
forget about it. But instead they keep hearing her name, and this seems
to drive some people crazy. Everyone's always going on about this
singer, but she's no good! She's a fraud!

That word ``fraud'' is an important one. It's the spectral signature of
a hater to regard the object of their hatred as a
\href{https://twitter.com/search?q=Musk\%20fraud&src=typed_query&f=live}{fraud}.
They can't deny their fame. Indeed, their fame is if anything
exaggerated in the hater's mind. They notice every mention of the
singer's name, because every mention makes them angrier. In their own
minds they exaggerate both the singer's fame and her lack of talent, and
the only way to reconcile those two ideas is to conclude that she has
tricked everyone.

What sort of people become haters? Can anyone become one? I'm not sure
about this, but I've noticed some patterns. Haters are generally losers
in a very specific sense: although they are occasionally talented, they
have never achieved much. And indeed, anyone successful enough to have
achieved significant fame would be unlikely to regard another famous
person as a fraud on that account, because anyone famous knows how
random fame is.

But haters are not always complete losers. They are not always the
proverbial guy living in his mom's basement. Many are, but some have
some amount of talent. In fact I suspect that a sense of frustrated
talent is what drives some people to become haters. They're not just
saying ``It's unfair that so-and-so is famous,'' but ``It's unfair that
so-and-so is famous, and not me.''

Could a hater be cured if they achieved something impressive? My guess
is that's a moot point, because they \href{mean.html}{never will}. I've
been able to observe for long enough that I'm fairly confident the
pattern works both ways: not only do people who do great work never
become haters, haters never do great work. Although I dislike the word
``fanboy,'' it's evocative of something important about both haters and
fanboys. It implies that the fanboy is so slavishly predictable in his
admiration that he's diminished as a result, that he's less than a man.

Haters seem even more diminished. I can imagine being a fanboy. I can
think of people whose work I admire so much that I could abase myself
before them out of sheer gratitude. If P. G. Wodehouse were still alive,
I could see myself being a Wodehouse fanboy. But I could not imagine
being a hater.

Knowing that haters are just fanboys with the sign bit flipped makes it
much easier to deal with them. We don't need a separate theory of
haters. We can just use existing techniques for dealing with obsessive
fans.

The most important of which is simply not to think much about them. If
you're like most people who become famous enough to acquire haters, your
initial reaction will be one of mystification. Why does this guy seem to
have it in for me? Where does his obsessive energy come from, and what
makes him so appallingly nasty? What did I do to set him off? Is it
something I can fix?

The mistake here is to think of the hater as someone you have a dispute
with. When you have a dispute with someone, it's usually a good idea to
try to understand why they're upset and then fix things if you can.
Disputes are distracting. But it's a false analogy to think of a hater
as someone you have a dispute with. It's an understandable mistake, if
you've never encountered haters before. But when you realize that you're
dealing with a hater, and what a hater is, it's clear that it's a waste
of time even to think about them. If you have obsessive fans, do you
spend any time wondering what makes them love you so much? No, you just
think ``some people are kind of crazy,'' and that's the end of it.

Since haters are equivalent to fanboys, that's the way to deal with them
too. There may have been something that set them off. But it's not
something that would have set off a normal person, so there's no reason
to spend any time thinking about it. It's not you, it's them.


\footnote{There are of course some people who are genuine frauds. How can
you distinguish between x calling y a fraud because x is a hater, and
because y is a fraud? Look at neutral opinion. Actual frauds are usually
pretty conspicuous. Thoughtful people are rarely taken in by them. So if
there are some thoughtful people who like y, you can usually assume y is
not a fraud.}\footnote{I would make an exception for teenagers, who sometimes act in
such extreme ways that they are literally not themselves. I can imagine
a teenage kid being a hater and then growing out of it. But not anyone
over 25.}\footnote{I have a much worse memory for misdeeds than my wife Jessica,
who is a connoisseur of character, but I don't wish it were better. Most
disputes are a waste of time even if you're in the right, and it's easy
to bury the hatchet with someone if you can't remember why you were mad
at them.}\footnote{A competent hater will not merely attack you individually but
will try to get mobs after you. In some cases you may want to refute
whatever bogus claim they made in order to do so. But err on the side of
not, because ultimately it probably won't matter.}

\textbf{Thanks} to Austen Allred, Trevor Blackwell, Patrick Collison,
Christine Ford, Daniel Gackle, Jessica Livingston, Robert Morris, Elon
Musk, Harj Taggar, and Peter Thiel for reading drafts of this.
\chapter{Being a Noob}

January 2020

When I was young, I thought old people had everything figured out. Now
that I'm old, I know this isn't true.

I constantly feel like a noob. It seems like I'm always talking to some
startup working in a new field I know nothing about, or reading a book
about a topic I don't understand well enough, or visiting some new
country where I don't know how things work.

It's not pleasant to feel like a noob. And the word ``noob'' is
certainly not a compliment. And yet today I realized something
encouraging about being a noob: the more of a noob you are locally, the
less of a noob you are globally.

For example, if you stay in your home country, you'll feel less of a
noob than if you move to Farawavia, where everything works differently.
And yet you'll know more if you move. So the feeling of being a noob is
inversely correlated with actual ignorance.

But if the feeling of being a noob is good for us, why do we dislike it?
What evolutionary purpose could such an aversion serve?

I think the answer is that there are two sources of feeling like a noob:
being stupid, and doing something novel. Our dislike of feeling like a
noob is our brain telling us ``Come on, come on, figure this out.''
Which was the right thing to be thinking for most of human history. The
life of hunter-gatherers was complex, but it didn't change as much as
life does now. They didn't suddenly have to figure out what to do about
cryptocurrency. So it made sense to be biased toward competence at
existing problems over the discovery of new ones. It made sense for
humans to dislike the feeling of being a noob, just as, in a world where
food was scarce, it made sense for them to dislike the feeling of being
hungry.

Now that too much food is more of a problem than too little, our dislike
of feeling hungry leads us astray. And I think our dislike of feeling
like a noob does too.

Though it feels unpleasant, and people will sometimes ridicule you for
it, the more you feel like a noob, the better.
\chapter{How to Write Usefully}

February 2020

What should an essay be? Many people would say persuasive. That's what a
lot of us were taught essays should be. But I think we can aim for
something more ambitious: that an essay should be useful.

To start with, that means it should be correct. But it's not enough
merely to be correct. It's easy to make a statement correct by making it
vague. That's a common flaw in academic writing, for example. If you
know nothing at all about an issue, you can't go wrong by saying that
the issue is a complex one, that there are many factors to be
considered, that it's a mistake to take too simplistic a view of it, and
so on.

Though no doubt correct, such statements tell the reader nothing. Useful
writing makes claims that are as strong as they can be made without
becoming false.

For example, it's more useful to say that Pike's Peak is near the middle
of Colorado than merely somewhere in Colorado. But if I say it's in the
exact middle of Colorado, I've now gone too far, because it's a bit east
of the middle.

Precision and correctness are like opposing forces. It's easy to satisfy
one if you ignore the other. The converse of vaporous academic writing
is the bold, but false, rhetoric of demagogues. Useful writing is bold,
but true.

It's also two other things: it tells people something important, and
that at least some of them didn't already know.

Telling people something they didn't know doesn't always mean surprising
them. Sometimes it means telling them something they knew unconsciously
but had never put into words. In fact those may be the more valuable
insights, because they tend to be more fundamental.

Let's put them all together. Useful writing tells people something true
and important that they didn't already know, and tells them as
unequivocally as possible.

Notice these are all a matter of degree. For example, you can't expect
an idea to be novel to everyone. Any insight that you have will probably
have already been had by at least one of the world's 7 billion people.
But it's sufficient if an idea is novel to a lot of readers.

Ditto for correctness, importance, and strength. In effect the four
components are like numbers you can multiply together to get a score for
usefulness. Which I realize is almost awkwardly reductive, but
nonetheless true.

\_\_\_\_\_

How can you ensure that the things you say are true and novel and
important? Believe it or not, there is a trick for doing this. I learned
it from my friend Robert Morris, who has a horror of saying anything
dumb. His trick is not to say anything unless he's sure it's worth
hearing. This makes it hard to get opinions out of him, but when you do,
they're usually right.

Translated into essay writing, what this means is that if you write a
bad sentence, you don't publish it. You delete it and try again. Often
you abandon whole branches of four or five paragraphs. Sometimes a whole
essay.

You can't ensure that every idea you have is good, but you can ensure
that every one you publish is, by simply not publishing the ones that
aren't.

In the sciences, this is called publication bias, and is considered bad.
When some hypothesis you're exploring gets inconclusive results, you're
supposed to tell people about that too. But with essay writing,
publication bias is the way to go.

My strategy is loose, then tight. I write the first draft of an essay
fast, trying out all kinds of ideas. Then I spend days rewriting it very
carefully.

I've never tried to count how many times I proofread essays, but I'm
sure there are sentences I've read 100 times before publishing them.
When I proofread an essay, there are usually passages that stick out in
an annoying way, sometimes because they're clumsily written, and
sometimes because I'm not sure they're true. The annoyance starts out
unconscious, but after the tenth reading or so I'm saying ``Ugh, that
part'' each time I hit it. They become like briars that catch your
sleeve as you walk past. Usually I won't publish an essay till they're
all gone --- till I can read through the whole thing without the feeling
of anything catching.

I'll sometimes let through a sentence that seems clumsy, if I can't
think of a way to rephrase it, but I will never knowingly let through
one that doesn't seem correct. You never have to. If a sentence doesn't
seem right, all you have to do is ask why it doesn't, and you've usually
got the replacement right there in your head.

This is where essayists have an advantage over journalists. You don't
have a deadline. You can work for as long on an essay as you need to get
it right. You don't have to publish the essay at all, if you can't get
it right. Mistakes seem to lose courage in the face of an enemy with
unlimited resources. Or that's what it feels like. What's really going
on is that you have different expectations for yourself. You're like a
parent saying to a child ``we can sit here all night till you eat your
vegetables.'' Except you're the child too.

I'm not saying no mistake gets through. For example, I added condition
(c) in \href{bias.html}{``A Way to Detect Bias''} after readers pointed
out that I'd omitted it. But in practice you can catch nearly all of
them.

There's a trick for getting importance too. It's like the trick I
suggest to young founders for getting startup ideas: to make something
you yourself want. You can use yourself as a proxy for the reader. The
reader is not completely unlike you, so if you write about topics that
seem important to you, they'll probably seem important to a significant
number of readers as well.

Importance has two factors. It's the number of people something matters
to, times how much it matters to them. Which means of course that it's
not a rectangle, but a sort of ragged comb, like a Riemann sum.

The way to get novelty is to write about topics you've thought about a
lot. Then you can use yourself as a proxy for the reader in this
department too. Anything you notice that surprises you, who've thought
about the topic a lot, will probably also surprise a significant number
of readers. And here, as with correctness and importance, you can use
the Morris technique to ensure that you will. If you don't learn
anything from writing an essay, don't publish it.

You need humility to measure novelty, because acknowledging the novelty
of an idea means acknowledging your previous ignorance of it. Confidence
and humility are often seen as opposites, but in this case, as in many
others, confidence helps you to be humble. If you know you're an expert
on some topic, you can freely admit when you learn something you didn't
know, because you can be confident that most other people wouldn't know
it either.

The fourth component of useful writing, strength, comes from two things:
thinking well, and the skillful use of qualification. These two
counterbalance each other, like the accelerator and clutch in a car with
a manual transmission. As you try to refine the expression of an idea,
you adjust the qualification accordingly. Something you're sure of, you
can state baldly with no qualification at all, as I did the four
components of useful writing. Whereas points that seem dubious have to
be held at arm's length with perhapses.

As you refine an idea, you're pushing in the direction of less
qualification. But you can rarely get it down to zero. Sometimes you
don't even want to, if it's a side point and a fully refined version
would be too long.

Some say that qualifications weaken writing. For example, that you
should never begin a sentence in an essay with ``I think,'' because if
you're saying it, then of course you think it. And it's true that ``I
think x'' is a weaker statement than simply ``x.'' Which is exactly why
you need ``I think.'' You need it to express your degree of certainty.

But qualifications are not scalars. They're not just experimental error.
There must be 50 things they can express: how broadly something applies,
how you know it, how happy you are it's so, even how it could be
falsified. I'm not going to try to explore the structure of
qualification here. It's probably more complex than the whole topic of
writing usefully. Instead I'll just give you a practical tip: Don't
underestimate qualification. It's an important skill in its own right,
not just a sort of tax you have to pay in order to avoid saying things
that are false. So learn and use its full range. It may not be fully
half of having good ideas, but it's part of having them.

There's one other quality I aim for in essays: to say things as simply
as possible. But I don't think this is a component of usefulness. It's
more a matter of consideration for the reader. And it's a practical aid
in getting things right; a mistake is more obvious when expressed in
simple language. But I'll admit that the main reason I write simply is
not for the reader's sake or because it helps get things right, but
because it bothers me to use more or fancier words than I need to. It
seems inelegant, like a program that's too long.

I realize florid writing works for some people. But unless you're sure
you're one of them, the best advice is to write as simply as you can.

\_\_\_\_\_

I believe the formula I've given you, importance + novelty + correctness
+ strength, is the recipe for a good essay. But I should warn you that
it's also a recipe for making people mad.

The root of the problem is novelty. When you tell people something they
didn't know, they don't always thank you for it. Sometimes the reason
people don't know something is because they don't want to know it.
Usually because it contradicts some cherished belief. And indeed, if
you're looking for novel ideas, popular but mistaken beliefs are a good
place to find them. Every popular mistaken belief creates a
\href{nov.html}{dead zone} of ideas around it that are relatively
unexplored because they contradict it.

The strength component just makes things worse. If there's anything that
annoys people more than having their cherished assumptions contradicted,
it's having them flatly contradicted.

Plus if you've used the Morris technique, your writing will seem quite
confident. Perhaps offensively confident, to people who disagree with
you. The reason you'll seem confident is that you are confident: you've
cheated, by only publishing the things you're sure of. It will seem to
people who try to disagree with you that you never admit you're wrong.
In fact you constantly admit you're wrong. You just do it before
publishing instead of after.

And if your writing is as simple as possible, that just makes things
worse. Brevity is the diction of command. If you watch someone
delivering unwelcome news from a position of inferiority, you'll notice
they tend to use lots of words, to soften the blow. Whereas to be short
with someone is more or less to be rude to them.

It can sometimes work to deliberately phrase statements more weakly than
you mean. To put ``perhaps'' in front of something you're actually quite
sure of. But you'll notice that when writers do this, they usually do it
with a wink.

I don't like to do this too much. It's cheesy to adopt an ironic tone
for a whole essay. I think we just have to face the fact that elegance
and curtness are two names for the same thing.

You might think that if you work sufficiently hard to ensure that an
essay is correct, it will be invulnerable to attack. That's sort of
true. It will be invulnerable to valid attacks. But in practice that's
little consolation.

In fact, the strength component of useful writing will make you
particularly vulnerable to misrepresentation. If you've stated an idea
as strongly as you could without making it false, all anyone has to do
is to exaggerate slightly what you said, and now it is false.

Much of the time they're not even doing it deliberately. One of the most
surprising things you'll discover, if you start writing essays, is that
people who disagree with you rarely disagree with what you've actually
written. Instead they make up something you said and disagree with that.

For what it's worth, the countermove is to ask someone who does this to
quote a specific sentence or passage you wrote that they believe is
false, and explain why. I say ``for what it's worth'' because they never
do. So although it might seem that this could get a broken discussion
back on track, the truth is that it was never on track in the first
place.

Should you explicitly forestall likely misinterpretations? Yes, if
they're misinterpretations a reasonably smart and well-intentioned
person might make. In fact it's sometimes better to say something
slightly misleading and then add the correction than to try to get an
idea right in one shot. That can be more efficient, and can also model
the way such an idea would be discovered.

But I don't think you should explicitly forestall intentional
misinterpretations in the body of an essay. An essay is a place to meet
honest readers. You don't want to spoil your house by putting bars on
the windows to protect against dishonest ones. The place to protect
against intentional misinterpretations is in end-notes. But don't think
you can predict them all. People are as ingenious at misrepresenting you
when you say something they don't want to hear as they are at coming up
with rationalizations for things they want to do but know they
shouldn't. I suspect it's the same skill.

\_\_\_\_\_

As with most other things, the way to get better at writing essays is to
practice. But how do you start? Now that we've examined the structure of
useful writing, we can rephrase that question more precisely. Which
constraint do you relax initially? The answer is, the first component of
importance: the number of people who care about what you write.

If you narrow the topic sufficiently, you can probably find something
you're an expert on. Write about that to start with. If you only have
ten readers who care, that's fine. You're helping them, and you're
writing. Later you can expand the breadth of topics you write about.

The other constraint you can relax is a little surprising: publication.
Writing essays doesn't have to mean publishing them. That may seem
strange now that the trend is to publish every random thought, but it
worked for me. I wrote what amounted to essays in notebooks for about 15
years. I never published any of them and never expected to. I wrote them
as a way of figuring things out. But when the web came along I'd had a
lot of practice.

Incidentally,
\href{http://www.foundersatwork.com/steve-wozniak.html}{Steve Wozniak}
did the same thing. In high school he designed computers on paper for
fun. He couldn't build them because he couldn't afford the components.
But when Intel launched 4K DRAMs in 1975, he was ready.

\_\_\_\_\_

How many essays are there left to write though? The answer to that
question is probably the most exciting thing I've learned about essay
writing. Nearly all of them are left to write.

Although \href{essay.html}{the essay} is an old form, it hasn't been
assiduously cultivated. In the print era, publication was expensive, and
there wasn't enough demand for essays to publish that many. You could
publish essays if you were already well known for writing something
else, like novels. Or you could write book reviews that you took over to
express your own ideas. But there was not really a direct path to
becoming an essayist. Which meant few essays got written, and those that
did tended to be about a narrow range of subjects.

Now, thanks to the internet, there's a path. Anyone can publish essays
online. You start in obscurity, perhaps, but at least you can start. You
don't need anyone's permission.

It sometimes happens that an area of knowledge sits quietly for years,
till some change makes it explode. Cryptography did this to number
theory. The internet is doing it to the essay.

The exciting thing is not that there's a lot left to write, but that
there's a lot left to discover. There's a certain kind of idea that's
best discovered by writing essays. If most essays are still unwritten,
most such ideas are still undiscovered.


\footnote{Put railings on the balconies, but don't put bars on the
windows.}\footnote{Even now I sometimes write essays that are not meant for
publication. I wrote several to figure out what Y Combinator should do,
and they were really helpful.}

\textbf{Thanks} to Trevor Blackwell, Daniel Gackle, Jessica Livingston,
and Robert Morris for reading drafts of this.
\chapter{Coronavirus and Credibility}

April 2020

I recently saw a
\href{https://www.youtube.com/watch?v=NAh4uS4f78o}{video} of TV
journalists and politicians confidently saying that the coronavirus
would be no worse than the flu. What struck me about it was not just how
mistaken they seemed, but how daring. How could they feel safe saying
such things?

The answer, I realized, is that they didn't think they could get caught.
They didn't realize there was any danger in making false predictions.
These people constantly make false predictions, and get away with it,
because the things they make predictions about either have mushy enough
outcomes that they can bluster their way out of trouble, or happen so
far in the future that few remember what they said.

An epidemic is different. It falsifies your predictions rapidly and
unequivocally.

But epidemics are rare enough that these people clearly didn't realize
this was even a possibility. Instead they just continued to use their
ordinary m.o., which, as the epidemic has made clear, is to talk
confidently about things they don't understand.

An event like this is thus a uniquely powerful way of taking people's
measure. As Warren Buffett said, ``It's only when the tide goes out that
you learn who's been swimming naked.'' And the tide has just gone out
like never before.

Now that we've seen the results, let's remember what we saw, because
this is the most accurate test of credibility we're ever likely to have.
I hope.
\chapter{Orthodox Privilege}

July 2020

``Few people are capable of expressing with equanimity opinions which
differ from the prejudices of their social environment. Most people are
even incapable of forming such opinions.''

--- Einstein

There has been a lot of talk about privilege lately. Although the
concept is overused, there is something to it, and in particular to the
idea that privilege makes you blind --- that you can't see things that
are visible to someone whose life is very different from yours.

But one of the most pervasive examples of this kind of blindness is one
that I haven't seen mentioned explicitly. I'm going to call it
\emph{orthodox privilege}: The more conventional-minded someone is, the
more it seems to them that it's safe for everyone to express their
opinions.

It's safe for \emph{them} to express their opinions, because the source
of their opinions is whatever it's currently acceptable to believe. So
it seems to them that it must be safe for everyone. They literally can't
imagine a true statement that would get you in trouble.

And yet at every point in history, there \href{say.html}{were} true
things that would get you in trouble to say. Is ours the first where
this isn't so? What an amazing coincidence that would be.

Surely it should at least be the default assumption that our time is not
unique, and that there are true things you can't say now, just as there
have always been. You would think. But even in the face of such
overwhelming historical evidence, most people will go with their gut on
this one.

In the most extreme cases, people suffering from orthodox privilege will
not only deny that there's anything true that you can't say, but will
accuse you of heresy merely for saying there is. Though if there's more
than one heresy current in your time, these accusations will be weirdly
non-deterministic: you must either be an xist or a yist.

Frustrating as it is to deal with these people, it's important to
realize that they're in earnest. They're not pretending they think it's
impossible for an idea to be both unorthodox and true. The world really
looks that way to them.

Indeed, this is a uniquely tenacious form of privilege. People can
overcome the blindness induced by most forms of privilege by learning
more about whatever they're not. But they can't overcome orthodox
privilege just by learning more. They'd have to become more
independent-minded. If that happens at all, it doesn't happen on the
time scale of one conversation.

It may be possible to convince some people that orthodox privilege must
exist even though they can't sense it, just as one can with, say, dark
matter. There may be some who could be convinced, for example, that it's
very unlikely that this is the first point in history at which there's
nothing true you can't say, even if they can't imagine specific
examples.

But in general I don't think it will work to say ``check your
privilege'' about this type of privilege, because those in its
demographic don't realize they're in it. It doesn't seem to
conventional-minded people that they're conventional-minded. It just
seems to them that they're right. Indeed, they tend to be particularly
sure of it.

Perhaps the solution is to appeal to politeness. If someone says they
can hear a high-pitched noise that you can't, it's only polite to take
them at their word, instead of demanding evidence that's impossible to
produce, or simply denying that they hear anything. Imagine how rude
that would seem. Similarly, if someone says they can think of things
that are true but that cannot be said, it's only polite to take them at
their word, even if you can't think of any yourself.

\textbf{Thanks} to Sam Altman, Trevor Blackwell, Patrick Collison,
Antonio Garcia-Martinez, Jessica Livingston, Robert Morris, Michael
Nielsen, Geoff Ralston, Max Roser, and Harj Taggar for reading drafts of
this.
\chapter{The Four Quadrants of Conformism}

July 2020

One of the most revealing ways to classify people is by the degree and
aggressiveness of their conformism. Imagine a Cartesian coordinate
system whose horizontal axis runs from conventional-minded on the left
to independent-minded on the right, and whose vertical axis runs from
passive at the bottom to aggressive at the top. The resulting four
quadrants define four types of people. Starting in the upper left and
going counter-clockwise: aggressively conventional-minded, passively
conventional-minded, passively independent-minded, and aggressively
independent-minded.

I think that you'll find all four types in most societies, and that
which quadrant people fall into depends more on their own personality
than the beliefs prevalent in their society. \footnote{I realize of course that if people's personalities vary in any
two ways, you can use them as axes and call the resulting four quadrants
personality types. So what I'm really claiming is that the axes are
orthogonal and that there's significant variation in both.}

Young children offer some of the best evidence for both points. Anyone
who's been to primary school has seen the four types, and the fact that
school rules are so arbitrary is strong evidence that which quadrant
people fall into depends more on them than the rules.

The kids in the upper left quadrant, the aggressively
conventional-minded ones, are the tattletales. They believe not only
that rules must be obeyed, but that those who disobey them must be
punished.

The kids in the lower left quadrant, the passively conventional-minded,
are the sheep. They're careful to obey the rules, but when other kids
break them, their impulse is to worry that those kids will be punished,
not to ensure that they will.

The kids in the lower right quadrant, the passively independent-minded,
are the dreamy ones. They don't care much about rules and probably
aren't 100\% sure what the rules even are.

And the kids in the upper right quadrant, the aggressively
independent-minded, are the naughty ones. When they see a rule, their
first impulse is to question it. Merely being told what to do makes them
inclined to do the opposite.

When measuring conformism, of course, you have to say with respect to
what, and this changes as kids get older. For younger kids it's the
rules set by adults. But as kids get older, the source of rules becomes
their peers. So a pack of teenagers who all flout school rules in the
same way are not independent-minded; rather the opposite.

In adulthood we can recognize the four types by their distinctive calls,
much as you could recognize four species of birds. The call of the
aggressively conventional-minded is ``Crush !'' (It's rather alarming to
see an exclamation point after a variable, but that's the whole problem
with the aggressively conventional-minded.) The call of the passively
conventional-minded is ``What will the neighbors think?'' The call of
the passively independent-minded is ``To each his own.'' And the call of
the aggressively independent-minded is ``Eppur si muove.''

The four types are not equally common. There are more passive people
than aggressive ones, and far more conventional-minded people than
independent-minded ones. So the passively conventional-minded are the
largest group, and the aggressively independent-minded the smallest.

Since one's quadrant depends more on one's personality than the nature
of the rules, most people would occupy the same quadrant even if they'd
grown up in a quite different society.

Princeton professor Robert George recently wrote:

\begin{quote}
I sometimes ask students what their position on slavery would have been
had they been white and living in the South before abolition. Guess
what? They all would have been abolitionists! They all would have
bravely spoken out against slavery, and worked tirelessly against it.
\end{quote}

He's too polite to say so, but of course they wouldn't. And indeed, our
default assumption should not merely be that his students would, on
average, have behaved the same way people did at the time, but that the
ones who are aggressively conventional-minded today would have been
aggressively conventional-minded then too. In other words, that they'd
not only not have fought against slavery, but that they'd have been
among its staunchest defenders.

I'm biased, I admit, but it seems to me that aggressively
conventional-minded people are responsible for a disproportionate amount
of the trouble in the world, and that a lot of the customs we've evolved
since the Enlightenment have been designed to protect the rest of us
from them. In particular, the retirement of the concept of heresy and
its replacement by the principle of freely debating all sorts of
different ideas, even ones that are currently considered unacceptable,
without any punishment for those who try them out to see if they work.
\footnote{The aggressively conventional-minded aren't responsible for all
the trouble in the world. Another big source of trouble is the sort of
charismatic leader who gains power by appealing to them. They become
much more dangerous when such leaders emerge.}

Why do the independent-minded need to be protected, though? Because they
have all the new ideas. To be a successful scientist, for example, it's
not enough just to be right. You have to be right when everyone else is
wrong. Conventional-minded people can't do that. For similar reasons,
all successful startup CEOs are not merely independent-minded, but
aggressively so. So it's no coincidence that societies prosper only to
the extent that they have customs for keeping the conventional-minded at
bay. \footnote{I never worried about writing things that offended the
conventional-minded when I was running Y Combinator. If YC were a cookie
company, I'd have faced a difficult moral choice. Conventional-minded
people eat cookies too. But they don't start successful startups. So if
I deterred them from applying to YC, the only effect was to save us work
reading applications.}

In the last few years, many of us have noticed that the customs
protecting free inquiry have been weakened. Some say we're overreacting
--- that they haven't been weakened very much, or that they've been
weakened in the service of a greater good. The latter I'll dispose of
immediately. When the conventional-minded get the upper hand, they
always say it's in the service of a greater good. It just happens to be
a different, incompatible greater good each time.

As for the former worry, that the independent-minded are being
oversensitive, and that free inquiry hasn't been shut down that much,
you can't judge that unless you are yourself independent-minded. You
can't know how much of the space of ideas is being lopped off unless you
have them, and only the independent-minded have the ones at the edges.
Precisely because of this, they tend to be very sensitive to changes in
how freely one can explore ideas. They're the canaries in this coalmine.

The conventional-minded say, as they always do, that they don't want to
shut down the discussion of all ideas, just the bad ones.

You'd think it would be obvious just from that sentence what a dangerous
game they're playing. But I'll spell it out. There are two reasons why
we need to be able to discuss even ``bad'' ideas.

The first is that any process for deciding which ideas to ban is bound
to make mistakes. All the more so because no one intelligent wants to
undertake that kind of work, so it ends up being done by the stupid. And
when a process makes a lot of mistakes, you need to leave a margin for
error. Which in this case means you need to ban fewer ideas than you'd
like to. But that's hard for the aggressively conventional-minded to do,
partly because they enjoy seeing people punished, as they have since
they were children, and partly because they compete with one another.
Enforcers of orthodoxy can't allow a borderline idea to exist, because
that gives other enforcers an opportunity to one-up them in the moral
purity department, and perhaps even to turn enforcer upon them. So
instead of getting the margin for error we need, we get the opposite: a
race to the bottom in which any idea that seems at all bannable ends up
being banned. \footnote{There has been progress in one area: the punishments for talking
about banned ideas are less severe than in the past. There's little
danger of being killed, at least in richer countries. The aggressively
conventional-minded are mostly satisfied with getting people fired.}

The second reason it's dangerous to ban the discussion of ideas is that
ideas are more closely related than they look. Which means if you
restrict the discussion of some topics, it doesn't only affect those
topics. The restrictions propagate back into any topic that yields
implications in the forbidden ones. And that is not an edge case. The
best ideas do exactly that: they have consequences in fields far removed
from their origins. Having ideas in a world where some ideas are banned
is like playing soccer on a pitch that has a minefield in one corner.
You don't just play the same game you would have, but on a different
shaped pitch. You play a much more subdued game even on the ground
that's safe.

In the past, the way the independent-minded protected themselves was to
congregate in a handful of places --- first in courts, and later in
universities --- where they could to some extent make their own rules.
Places where people work with ideas tend to have customs protecting free
inquiry, for the same reason wafer fabs have powerful air filters, or
recording studios good sound insulation. For the last couple centuries
at least, when the aggressively conventional-minded were on the rampage
for whatever reason, universities were the safest places to be.

That may not work this time though, due to the unfortunate fact that the
latest wave of intolerance began in universities. It began in the mid
1980s, and by 2000 seemed to have died down, but it has recently flared
up again with the arrival of social media. This seems, unfortunately, to
have been an own goal by Silicon Valley. Though the people who run
Silicon Valley are almost all independent-minded, they've handed the
aggressively conventional-minded a tool such as they could only have
dreamed of.

On the other hand, perhaps the decline in the spirit of free inquiry
within universities is as much the symptom of the departure of the
independent-minded as the cause. People who would have become professors
50 years ago have other options now. Now they can become quants or start
startups. You have to be independent-minded to succeed at either of
those. If these people had been professors, they'd have put up a stiffer
resistance on behalf of academic freedom. So perhaps the picture of the
independent-minded fleeing declining universities is too gloomy. Perhaps
the universities are declining because so many have already left.
\footnote{Many professors are independent-minded --- especially in math,
the hard sciences, and engineering, where you have to be to succeed. But
students are more representative of the general population, and thus
mostly conventional-minded. So when professors and students are in
conflict, it's not just a conflict between generations but also between
different types of people.}

Though I've spent a lot of time thinking about this situation, I can't
predict how it plays out. Could some universities reverse the current
trend and remain places where the independent-minded want to congregate?
Or will the independent-minded gradually abandon them? I worry a lot
about what we might lose if that happened.

But I'm hopeful long term. The independent-minded are good at protecting
themselves. If existing institutions are compromised, they'll create new
ones. That may require some imagination. But imagination is, after all,
their specialty.


\textbf{Thanks} to Sam Altman, Trevor Blackwell, Nicholas Christakis,
Patrick Collison, Sam Gichuru, Jessica Livingston, Patrick McKenzie,
Geoff Ralston, and Harj Taggar for reading drafts of this.
\chapter{Modeling a Wealth Tax}

August 2020

Some politicians are proposing to introduce wealth taxes in addition to
income and capital gains taxes. Let's try modeling the effects of
various levels of wealth tax to see what they would mean in practice for
a startup founder.

Suppose you start a successful startup in your twenties, and then live
for another 60 years. How much of your stock will a wealth tax consume?

If the wealth tax applies to all your assets, it's easy to calculate its
effect. A wealth tax of 1\% means you get to keep 99\% of your stock
each year. After 60 years the proportion of stock you'll have left will
be .99\^{}60, or .547. So a straight 1\% wealth tax means the government
will over the course of your life take 45\% of your stock.

(Losing shares does not, obviously, mean becoming \emph{net} poorer
unless the value per share is increasing by less than the wealth tax
rate.)

Here's how much stock the government would take over 60 years at various
levels of wealth tax:

wealth tax

government takes

0.1\%

6\%

0.5\%

26\%

1.0\%

45\%

2.0\%

70\%

3.0\%

84\%

4.0\%

91\%

5.0\%

95\%

A wealth tax will usually have a threshold at which it starts. How much
difference would a high threshold make? To model that, we need to make
some assumptions about the initial value of your stock and the growth
rate.

Suppose your stock is initially worth \$2 million, and the company's
trajectory is as follows: the value of your stock grows 3x for 2 years,
then 2x for 2 years, then 50\% for 2 years, after which you just get a
typical public company growth rate, which we'll call 8\%.
\footnote{In practice, eventually some of this 8\% would come in the form
of dividends, which are taxed as income at issue, so this model actually
represents the most optimistic case for the founder.} Suppose the wealth tax threshold is \$50
million. How much stock does the government take now?

wealth tax

government takes

0.1\%

5\%

0.5\%

23\%

1.0\%

41\%

2.0\%

65\%

3.0\%

79\%

4.0\%

88\%

5.0\%

93\%

It may at first seem surprising that such apparently small tax rates
produce such dramatic effects. A 2\% wealth tax with a \$50~million
threshold takes about two thirds of a successful founder's stock.

The reason wealth taxes have such dramatic effects is that they're
applied over and over to the same money. Income tax happens every year,
but only to that year's income. Whereas if you live for 60 years after
acquiring some asset, a wealth tax will tax that same asset 60 times. A
wealth tax compounds.

\textbf{Note}
\chapter{Early Work}

October 2020

One of the biggest things holding people back from doing great work is
the fear of making something lame. And this fear is not an irrational
one. Many great projects go through a stage early on where they don't
seem very impressive, even to their creators. You have to push through
this stage to reach the great work that lies beyond. But many people
don't. Most people don't even reach the stage of making something
they're embarrassed by, let alone continue past it. They're too
frightened even to start.

Imagine if we could turn off the fear of making something lame. Imagine
how much more we'd do.

Is there any hope of turning it off? I think so. I think the habits at
work here are not very deeply rooted.

Making new things is itself a new thing for us as a species. It has
always happened, but till the last few centuries it happened so slowly
as to be invisible to individual humans. And since we didn't need
customs for dealing with new ideas, we didn't develop any.

We just don't have enough experience with early versions of ambitious
projects to know how to respond to them. We judge them as we would judge
more finished work, or less ambitious projects. We don't realize they're
a special case.

Or at least, most of us don't. One reason I'm confident we can do better
is that it's already starting to happen. There are already a few places
that are living in the future in this respect. Silicon Valley is one of
them: an unknown person working on a strange-sounding idea won't
automatically be dismissed the way they would back home. In Silicon
Valley, people have learned how dangerous that is.

The right way to deal with new ideas is to treat them as a challenge to
your imagination --- not just to have lower standards, but to
\href{altair.html}{switch polarity} entirely, from listing the reasons
an idea won't work to trying to think of ways it could. That's what I do
when I meet people with new ideas. I've become quite good at it, but
I've had a lot of practice. Being a partner at Y Combinator means being
practically immersed in strange-sounding ideas proposed by unknown
people. Every six months you get thousands of new ones thrown at you and
have to sort through them, knowing that in a world with a power-law
distribution of outcomes, it will be painfully obvious if you miss the
needle in this haystack. Optimism becomes urgent.

But I'm hopeful that, with time, this kind of optimism can become
widespread enough that it becomes a social custom, not just a trick used
by a few specialists. It is after all an extremely lucrative trick, and
those tend to spread quickly.

Of course, inexperience is not the only reason people are too harsh on
early versions of ambitious projects. They also do it to seem clever.
And in a field where the new ideas are risky, like startups, those who
dismiss them are in fact more likely to be right. Just not when their
predictions are \href{swan.html}{weighted by outcome}.

But there is another more sinister reason people dismiss new ideas. If
you try something ambitious, many of those around you will hope,
consciously or unconsciously, that you'll fail. They worry that if you
try something ambitious and succeed, it will put you above them. In some
countries this is not just an individual failing but part of the
national culture.

I wouldn't claim that people in Silicon Valley overcome these impulses
because they're morally better. \footnote{This assumption may be too conservative. There is some evidence
that historically the Bay Area has attracted a
\href{cities.html}{different sort of person} than, say, New York City.} The reason many
hope you'll succeed is that they hope to rise with you. For investors
this incentive is particularly explicit. They want you to succeed
because they hope you'll make them rich in the process. But many other
people you meet can hope to benefit in some way from your success. At
the very least they'll be able to say, when you're famous, that they've
known you since way back.

But even if Silicon Valley's encouraging attitude is rooted in
self-interest, it has over time actually grown into a sort of
benevolence. Encouraging startups has been practiced for so long that it
has become a custom. Now it just seems that that's what one does with
startups.

Maybe Silicon Valley is too optimistic. Maybe it's too easily fooled by
impostors. Many less optimistic journalists want to believe that. But
the lists of impostors they cite are suspiciously short, and plagued
with asterisks. \footnote{One of their great favorites is Theranos. But the most
conspicuous feature of Theranos's cap table is the absence of Silicon
Valley firms. Journalists were fooled by Theranos, but Silicon Valley
investors weren't.} If you use revenue as the test,
Silicon Valley's optimism seems better tuned than the rest of the
world's. And because it works, it will spread.

There's a lot more to new ideas than new startup ideas, of course. The
fear of making something lame holds people back in every field. But
Silicon Valley shows how quickly customs can evolve to support new
ideas. And that in turn proves that dismissing new ideas is not so
deeply rooted in human nature that it can't be unlearnt.

\_\_\_\_\_\_\_\_\_\_\_

Unfortunately, if you want to do new things, you'll face a force more
powerful than other people's skepticism: your own skepticism. You too
will judge your early work too harshly. How do you avoid that?

This is a difficult problem, because you don't want to completely
eliminate your horror of making something lame. That's what steers you
toward doing good work. You just want to turn it off temporarily, the
way a painkiller temporarily turns off pain.

People have already discovered several techniques that work. Hardy
mentions two in \emph{A Mathematician's Apology}:

\begin{quote}
Good work is not done by ``humble'' men. It is one of the first duties
of a professor, for example, in any subject, to exaggerate a little both
the importance of his subject and his importance in it.
\end{quote}

If you overestimate the importance of what you're working on, that will
compensate for your mistakenly harsh judgment of your initial results.
If you look at something that's 20\% of the way to a goal worth 100 and
conclude that it's 10\% of the way to a goal worth 200, your estimate of
its expected value is correct even though both components are wrong.

It also helps, as Hardy suggests, to be slightly overconfident. I've
noticed in many fields that the most successful people are slightly
overconfident. On the face of it this seems implausible. Surely it would
be optimal to have exactly the right estimate of one's abilities. How
could it be an advantage to be mistaken? Because this error compensates
for other sources of error in the opposite direction: being slightly
overconfident armors you against both other people's skepticism and your
own.

Ignorance has a similar effect. It's safe to make the mistake of judging
early work as finished work if you're a sufficiently lax judge of
finished work. I doubt it's possible to cultivate this kind of
ignorance, but empirically it's a real advantage, especially for the
young.

Another way to get through the lame phase of ambitious projects is to
surround yourself with the right people --- to create an eddy in the
social headwind. But it's not enough to collect people who are always
encouraging. You'd learn to discount that. You need colleagues who can
actually tell an ugly duckling from a baby swan. The people best able to
do this are those working on similar projects of their own, which is why
university departments and research labs work so well. You don't need
institutions to collect colleagues. They naturally coalesce, given the
chance. But it's very much worth accelerating this process by seeking
out other people trying to do new things.

Teachers are in effect a special case of colleagues. It's a teacher's
job both to see the promise of early work and to encourage you to
continue. But teachers who are good at this are unfortunately quite
rare, so if you have the opportunity to learn from one, take it.
\footnote{I made two mistakes about teachers when I was younger. I cared
more about professors' research than their reputations as teachers, and
I was also wrong about what it meant to be a good teacher. I thought it
simply meant to be good at explaining things.}

For some it might work to rely on sheer discipline: to tell yourself
that you just have to press on through the initial crap phase and not
get discouraged. But like a lot of ``just tell yourself'' advice, this
is harder than it sounds. And it gets still harder as you get older,
because your standards rise. The old do have one compensating advantage
though: they've been through this before.

It can help if you focus less on where you are and more on the rate of
change. You won't worry so much about doing bad work if you can see it
improving. Obviously the faster it improves, the easier this is. So when
you start something new, it's good if you can spend a lot of time on it.
That's another advantage of being young: you tend to have bigger blocks
of time.

Another common trick is to start by considering new work to be of a
different, less exacting type. To start a painting saying that it's just
a sketch, or a new piece of software saying that it's just a quick hack.
Then you judge your initial results by a lower standard. Once the
project is rolling you can sneakily convert it to something more.
\footnote{Patrick Collison points out that you can go past treating
something as a hack in the sense of a prototype and onward to the sense
of the word that means something closer to a practical joke:

\begin{quote}
I think there may be something related to being a hack that can be
powerful --- the idea of making the tenuousness and implausibility
\emph{a feature}. ``Yes, it's a bit ridiculous, right? I'm just trying
to see how far such a naive approach can get.'' YC seemed to me to have
this characteristic.
\end{quote}}

This will be easier if you use a medium that lets you work fast and
doesn't require too much commitment up front. It's easier to convince
yourself that something is just a sketch when you're drawing in a
notebook than when you're carving stone. Plus you get initial results
faster. \footnote{Much of the advantage of switching from physical to digital
media is not the software per se but that it lets you start something
new with little upfront commitment.} \footnote{John Carmack adds:

\begin{quote}
The value of a medium without a vast gulf between the early work and the
final work is exemplified in game mods. The original Quake game was a
golden age for mods, because everything was very flexible, but so crude
due to technical limitations, that quick hacks to try out a gameplay
idea weren't all \emph{that} far from the official game. Many careers
were born from that, but as the commercial game quality improved over
the years, it became almost a full time job to make a successful mod
that would be appreciated by the community. This was dramatically
reversed with Minecraft and later Roblox, where the entire esthetic of
the experience was so explicitly crude that innovative gameplay concepts
became the overriding value. These ``crude'' game mods by single authors
are now often bigger deals than massive professional teams' work.
\end{quote}}

It will be easier to try out a risky project if you think of it as a way
to learn and not just as a way to make something. Then even if the
project truly is a failure, you'll still have gained by it. If the
problem is sharply enough defined, failure itself is knowledge: if the
theorem you're trying to prove turns out to be false, or you use a
structural member of a certain size and it fails under stress, you've
learned something, even if it isn't what you wanted to learn.
\footnote{Lisa Randall suggests that we

\begin{quote}
treat new things as experiments. That way there's no such thing as
failing, since you learn something no matter what. You treat it like an
experiment in the sense that if it really rules something out, you give
up and move on, but if there's some way to vary it to make it work
better, go ahead and do that
\end{quote}}

One motivation that works particularly well for me is curiosity. I like
to try new things just to see how they'll turn out. We started Y
Combinator in this spirit, and it was one of main things that kept me
going while I was working on \href{bel.html}{Bel}. Having worked for so
long with various dialects of Lisp, I was very curious to see what its
inherent shape was: what you'd end up with if you followed the axiomatic
approach all the way.

But it's a bit strange that you have to play mind games with yourself to
avoid being discouraged by lame-looking early efforts. The thing you're
trying to trick yourself into believing is in fact the truth. A
lame-looking early version of an ambitious project truly is more
valuable than it seems. So the ultimate solution may be to teach
yourself that.

One way to do it is to study the histories of people who've done great
work. What were they thinking early on? What was the very first thing
they did? It can sometimes be hard to get an accurate answer to this
question, because people are often embarrassed by their earliest work
and make little effort to publish it. (They too misjudge it.) But when
you can get an accurate picture of the first steps someone made on the
path to some great work, they're often pretty feeble.
\footnote{Michael Nielsen points out that the internet has made this
easier, because you can see programmers' first commits, musicians' first
videos, and so on.}

Perhaps if you study enough such cases, you can teach yourself to be a
better judge of early work. Then you'll be immune both to other people's
skepticism and your own fear of making something lame. You'll see early
work for what it is.

Curiously enough, the solution to the problem of judging early work too
harshly is to realize that our attitudes toward it are themselves early
work. Holding everything to the same standard is a crude version 1.
We're already evolving better customs, and we can already see signs of
how big the payoff will be.


\textbf{Thanks} to Trevor Blackwell, John Carmack, Patrick Collison,
Jessica Livingston, Michael Nielsen, and Lisa Randall for reading drafts
of this.
\chapter{How to Think for Yourself}

November 2020

There are some kinds of work that you can't do well without thinking
differently from your peers. To be a successful scientist, for example,
it's not enough just to be correct. Your ideas have to be both correct
and novel. You can't publish papers saying things other people already
know. You need to say things no one else has realized yet.

The same is true for investors. It's not enough for a public market
investor to predict correctly how a company will do. If a lot of other
people make the same prediction, the stock price will already reflect
it, and there's no room to make money. The only valuable insights are
the ones most other investors don't share.

You see this pattern with startup founders too. You don't want to start
a startup to do something that everyone agrees is a good idea, or there
will already be other companies doing it. You have to do something that
sounds to most other people like a bad idea, but that you know isn't ---
like writing software for a tiny computer used by a few thousand
hobbyists, or starting a site to let people rent airbeds on strangers'
floors.

Ditto for essayists. An essay that told people things they already knew
would be boring. You have to tell them something
\href{useful.html}{new}.

But this pattern isn't universal. In fact, it doesn't hold for most
kinds of work. In most kinds of work --- to be an administrator, for
example --- all you need is the first half. All you need is to be right.
It's not essential that everyone else be wrong.

There's room for a little novelty in most kinds of work, but in practice
there's a fairly sharp distinction between the kinds of work where it's
essential to be independent-minded, and the kinds where it's not.

I wish someone had told me about this distinction when I was a kid,
because it's one of the most important things to think about when you're
deciding what kind of work you want to do. Do you want to do the kind of
work where you can only win by thinking differently from everyone else?
I suspect most people's unconscious mind will answer that question
before their conscious mind has a chance to. I know mine does.

Independent-mindedness seems to be more a matter of nature than nurture.
Which means if you pick the wrong type of work, you're going to be
unhappy. If you're naturally independent-minded, you're going to find it
frustrating to be a middle manager. And if you're naturally
conventional-minded, you're going to be sailing into a headwind if you
try to do original research.

One difficulty here, though, is that people are often mistaken about
where they fall on the spectrum from conventional- to
independent-minded. Conventional-minded people don't like to think of
themselves as conventional-minded. And in any case, it genuinely feels
to them as if they make up their own minds about everything. It's just a
coincidence that their beliefs are identical to their peers'. And the
independent-minded, meanwhile, are often unaware how different their
ideas are from conventional ones, at least till they state them
publicly. \footnote{One convenient consequence of the fact that no one identifies as
conventional-minded is that you can say what you like about
conventional-minded people without getting in too much trouble. When I
wrote \href{conformism.html}{``The Four Quadrants of Conformism''} I
expected a firestorm of rage from the aggressively conventional-minded,
but in fact it was quite muted. They sensed that there was something
about the essay that they disliked intensely, but they had a hard time
finding a specific passage to pin it on.}

By the time they reach adulthood, most people know roughly how smart
they are (in the narrow sense of ability to solve pre-set problems),
because they're constantly being tested and ranked according to it. But
schools generally ignore independent-mindedness, except to the extent
they try to suppress it. So we don't get anything like the same kind of
feedback about how independent-minded we are.

There may even be a phenomenon like Dunning-Kruger at work, where the
most conventional-minded people are confident that they're
independent-minded, while the genuinely independent-minded worry they
might not be independent-minded enough.

\_\_\_\_\_\_\_\_\_\_\_

Can you make yourself more independent-minded? I think so. This quality
may be largely inborn, but there seem to be ways to magnify it, or at
least not to suppress it.

One of the most effective techniques is one practiced unintentionally by
most nerds: simply to be less aware what conventional beliefs are. It's
hard to be a conformist if you don't know what you're supposed to
conform to. Though again, it may be that such people already are
independent-minded. A conventional-minded person would probably feel
anxious not knowing what other people thought, and make more effort to
find out.

It matters a lot who you surround yourself with. If you're surrounded by
conventional-minded people, it will constrain which ideas you can
express, and that in turn will constrain which ideas you have. But if
you surround yourself with independent-minded people, you'll have the
opposite experience: hearing other people say surprising things will
encourage you to, and to think of more.

Because the independent-minded find it uncomfortable to be surrounded by
conventional-minded people, they tend to self-segregate once they have a
chance to. The problem with high school is that they haven't yet had a
chance to. Plus high school tends to be an inward-looking little world
whose inhabitants lack confidence, both of which magnify the forces of
conformism. So high school is often a \href{nerds.html}{bad time} for
the independent-minded. But there is some advantage even here: it
teaches you what to avoid. If you later find yourself in a situation
that makes you think ``this is like high school,'' you know you should
get out. \footnote{When I ask myself what in my life is like high school, the
answer is Twitter. It's not just full of conventional-minded people, as
anything its size will inevitably be, but subject to violent storms of
conventional-mindedness that remind me of descriptions of Jupiter. But
while it probably is a net loss to spend time there, it has at least
made me think more about the distinction between independent- and
conventional-mindedness, which I probably wouldn't have done otherwise.}

Another place where the independent- and conventional-minded are thrown
together is in successful startups. The founders and early employees are
almost always independent-minded; otherwise the startup wouldn't be
successful. But conventional-minded people greatly outnumber
independent-minded ones, so as the company grows, the original spirit of
independent-mindedness is inevitably diluted. This causes all kinds of
problems besides the obvious one that the company starts to suck. One of
the strangest is that the founders find themselves able to speak more
freely with founders of other companies than with their own employees.
\footnote{The decrease in independent-mindedness in growing startups is
still an open problem, but there may be solutions.

Founders can delay the problem by making a conscious effort only to hire
independent-minded people. Which of course also has the ancillary
benefit that they have better ideas.

Another possible solution is to create policies that somehow disrupt the
force of conformism, much as control rods slow chain reactions, so that
the conventional-minded aren't as dangerous. The physical separation of
Lockheed's Skunk Works may have had this as a side benefit. Recent
examples suggest employee forums like Slack may not be an unmitigated
good.

The most radical solution would be to grow revenues without growing the
company. You think hiring that junior PR person will be cheap, compared
to a programmer, but what will be the effect on the average level of
independent-mindedness in your company? (The growth in staff relative to
faculty seems to have had a similar effect on universities.) Perhaps the
rule about outsourcing work that's not your ``core competency'' should
be augmented by one about outsourcing work done by people who'd ruin
your culture as employees.

Some investment firms already seem to be able to grow revenues without
growing the number of employees. Automation plus the ever increasing
articulation of the ``tech stack'' suggest this may one day be possible
for product companies.}

Fortunately you don't have to spend all your time with
independent-minded people. It's enough to have one or two you can talk
to regularly. And once you find them, they're usually as eager to talk
as you are; they need you too. Although universities no longer have the
kind of monopoly they used to have on education, good universities are
still an excellent way to meet independent-minded people. Most students
will still be conventional-minded, but you'll at least find clumps of
independent-minded ones, rather than the near zero you may have found in
high school.

It also works to go in the other direction: as well as cultivating a
small collection of independent-minded friends, to try to meet as many
different types of people as you can. It will decrease the influence of
your immediate peers if you have several other groups of peers. Plus if
you're part of several different worlds, you can often import ideas from
one to another.

But by different types of people, I don't mean demographically
different. For this technique to work, they have to think differently.
So while it's an excellent idea to go and visit other countries, you can
probably find people who think differently right around the corner. When
I meet someone who knows a lot about something unusual (which includes
practically everyone, if you dig deep enough), I try to learn what they
know that other people don't. There are almost always surprises here.
It's a good way to make conversation when you meet strangers, but I
don't do it to make conversation. I really want to know.

You can expand the source of influences in time as well as space, by
reading history. When I read history I do it not just to learn what
happened, but to try to get inside the heads of people who lived in the
past. How did things look to them? This is hard to do, but worth the
effort for the same reason it's worth travelling far to triangulate a
point.

You can also take more explicit measures to prevent yourself from
automatically adopting conventional opinions. The most general is to
cultivate an attitude of skepticism. When you hear someone say
something, stop and ask yourself ``Is that true?'' Don't say it out
loud. I'm not suggesting that you impose on everyone who talks to you
the burden of proving what they say, but rather that you take upon
yourself the burden of evaluating what they say.

Treat it as a puzzle. You know that some accepted ideas will later turn
out to be wrong. See if you can guess which. The end goal is not to find
flaws in the things you're told, but to find the new ideas that had been
concealed by the broken ones. So this game should be an exciting quest
for novelty, not a boring protocol for intellectual hygiene. And you'll
be surprised, when you start asking ``Is this true?'', how often the
answer is not an immediate yes. If you have any imagination, you're more
likely to have too many leads to follow than too few.

More generally your goal should be not to let anything into your head
unexamined, and things don't always enter your head in the form of
statements. Some of the most powerful influences are implicit. How do
you even notice these? By standing back and watching how other people
get their ideas.

When you stand back at a sufficient distance, you can see ideas
spreading through groups of people like waves. The most obvious are in
fashion: you notice a few people wearing a certain kind of shirt, and
then more and more, until half the people around you are wearing the
same shirt. You may not care much what you wear, but there are
intellectual fashions too, and you definitely don't want to participate
in those. Not just because you want sovereignty over your own thoughts,
but because \href{nov.html}{unfashionable} ideas are disproportionately
likely to lead somewhere interesting. The best place to find
undiscovered ideas is where no one else is looking.
\footnote{There are intellectual fashions in every field, but their
influence varies. One of the reasons politics, for example, tends to be
boring is that it's so extremely subject to them. The threshold for
having opinions about politics is much \href{identity.html}{lower} than
the one for having opinions about set theory. So while there are some
ideas in politics, in practice they tend to be swamped by waves of
intellectual fashion.}

\_\_\_\_\_\_\_\_\_\_\_

To go beyond this general advice, we need to look at the internal
structure of independent-mindedness --- at the individual muscles we
need to exercise, as it were. It seems to me that it has three
components: fastidiousness about truth, resistance to being told what to
think, and curiosity.

Fastidiousness about truth means more than just not believing things
that are false. It means being careful about degree of belief. For most
people, degree of belief rushes unexamined toward the extremes: the
unlikely becomes impossible, and the probable becomes certain.
\footnote{The conventional-minded are often fooled by the strength of
their opinions into believing that they're independent-minded. But
strong convictions are not a sign of independent-mindedness. Rather the
opposite.} To the independent-minded, this seems
unpardonably sloppy. They're willing to have anything in their heads,
from highly speculative hypotheses to (apparent) tautologies, but on
subjects they care about, everything has to be labelled with a carefully
considered degree of belief. \footnote{Fastidiousness about truth doesn't imply that an
independent-minded person won't be dishonest, but that he won't be
deluded. It's sort of like the definition of a gentleman as someone who
is never unintentionally rude.}

The independent-minded thus have a horror of ideologies, which require
one to accept a whole collection of beliefs at once, and to treat them
as articles of faith. To an independent-minded person that would seem
revolting, just as it would seem to someone fastidious about food to
take a bite of a submarine sandwich filled with a large variety of
ingredients of indeterminate age and provenance.

Without this fastidiousness about truth, you can't be truly
independent-minded. It's not enough just to have resistance to being
told what to think. Those kind of people reject conventional ideas only
to replace them with the most random conspiracy theories. And since
these conspiracy theories have often been manufactured to capture them,
they end up being less independent-minded than ordinary people, because
they're subject to a much more exacting master than mere convention.
\footnote{You see this especially among political extremists. They think
themselves nonconformists, but actually they're niche conformists. Their
opinions may be different from the average person's, but they are often
more influenced by their peers' opinions than the average person's are.}

Can you increase your fastidiousness about truth? I would think so. In
my experience, merely thinking about something you're fastidious about
causes that fastidiousness to grow. If so, this is one of those rare
virtues we can have more of merely by wanting it. And if it's like other
forms of fastidiousness, it should also be possible to encourage in
children. I certainly got a strong dose of it from my father.
\footnote{If we broaden the concept of fastidiousness about truth so that
it excludes pandering, bogusness, and pomposity as well as falsehood in
the strict sense, our model of independent-mindedness can expand further
into the arts.}

The second component of independent-mindedness, resistance to being told
what to think, is the most visible of the three. But even this is often
misunderstood. The big mistake people make about it is to think of it as
a merely negative quality. The language we use reinforces that idea.
You're \_un\_conventional. You \emph{don't} care what other people
think. But it's not just a kind of immunity. In the most
independent-minded people, the desire not to be told what to think is a
positive force. It's not mere skepticism, but an active
\href{gba.html}{delight} in ideas that subvert the conventional wisdom,
the more counterintuitive the better.

Some of the most novel ideas seemed at the time almost like practical
jokes. Think how often your reaction to a novel idea is to laugh. I
don't think it's because novel ideas are funny per se, but because
novelty and humor share a certain kind of surprisingness. But while not
identical, the two are close enough that there is a definite correlation
between having a sense of humor and being independent-minded --- just as
there is between being humorless and being conventional-minded.
\footnote{This correlation is far from perfect, though. Gödel and Dirac
don't seem to have been very strong in the humor department. But someone
who is both ``neurotypical'' and humorless is very likely to be
conventional-minded.}

I don't think we can significantly increase our resistance to being told
what to think. It seems the most innate of the three components of
independent-mindedness; people who have this quality as adults usually
showed all too visible signs of it as children. But if we can't increase
our resistance to being told what to think, we can at least shore it up,
by surrounding ourselves with other independent-minded people.

The third component of independent-mindedness, curiosity, may be the
most interesting. To the extent that we can give a brief answer to the
question of where novel ideas come from, it's curiosity. That's what
people are usually feeling before having them.

In my experience, independent-mindedness and curiosity predict one
another perfectly. Everyone I know who's independent-minded is deeply
curious, and everyone I know who's conventional-minded isn't. Except,
curiously, children. All small children are curious. Perhaps the reason
is that even the conventional-minded have to be curious in the
beginning, in order to learn what the conventions are. Whereas the
independent-minded are the gluttons of curiosity, who keep eating even
after they're full. \footnote{Exception: gossip. Almost everyone is curious about gossip.}

The three components of independent-mindedness work in concert:
fastidiousness about truth and resistance to being told what to think
leave space in your brain, and curiosity finds new ideas to fill it.

Interestingly, the three components can substitute for one another in
much the same way muscles can. If you're sufficiently fastidious about
truth, you don't need to be as resistant to being told what to think,
because fastidiousness alone will create sufficient gaps in your
knowledge. And either one can compensate for curiosity, because if you
create enough space in your brain, your discomfort at the resulting
vacuum will add force to your curiosity. Or curiosity can compensate for
them: if you're sufficiently curious, you don't need to clear space in
your brain, because the new ideas you discover will push out the
conventional ones you acquired by default.

Because the components of independent-mindedness are so interchangeable,
you can have them to varying degrees and still get the same result. So
there is not just a single model of independent-mindedness. Some
independent-minded people are openly subversive, and others are quietly
curious. They all know the secret handshake though.

Is there a way to cultivate curiosity? To start with, you want to avoid
situations that suppress it. How much does the work you're currently
doing engage your curiosity? If the answer is ``not much,'' maybe you
should change something.

The most important active step you can take to cultivate your curiosity
is probably to seek out the topics that engage it. Few adults are
equally curious about everything, and it doesn't seem as if you can
choose which topics interest you. So it's up to you to
\href{genius.html}{find} them. Or invent them, if necessary.

Another way to increase your curiosity is to indulge it, by
investigating things you're interested in. Curiosity is unlike most
other appetites in this respect: indulging it tends to increase rather
than to sate it. Questions lead to more questions.

Curiosity seems to be more individual than fastidiousness about truth or
resistance to being told what to think. To the degree people have the
latter two, they're usually pretty general, whereas different people can
be curious about very different things. So perhaps curiosity is the
compass here. Perhaps, if your goal is to discover novel ideas, your
motto should not be ``do what you love'' so much as ``do what you're
curious about.''


\textbf{Thanks} to Trevor Blackwell, Paul Buchheit, Patrick Collison,
Jessica Livingston, Robert Morris, Harj Taggar, and Peter Thiel for
reading drafts of this.
\chapter{The Airbnbs}

December 2020

To celebrate Airbnb's IPO and to help future founders, I thought it
might be useful to explain what was special about Airbnb.

What was special about the Airbnbs was how earnest they were. They did
nothing half-way, and we could sense this even in the interview.
Sometimes after we interviewed a startup we'd be uncertain what to do,
and have to talk it over. Other times we'd just look at one another and
smile. The Airbnbs' interview was that kind. We didn't even like the
idea that much. Nor did users, at that stage; they had no growth. But
the founders seemed so full of energy that it was impossible not to like
them.

That first impression was not misleading. During the batch our nickname
for Brian Chesky was The Tasmanian Devil, because like the
\href{http://www.youtube.com/watch?v=StG2u5qfFRg&t=2m27s}{cartoon
character} he seemed a tornado of energy. All three of them were like
that. No one ever worked harder during YC than the Airbnbs did. When you
talked to the Airbnbs, they took notes. If you suggested an idea to them
in office hours, the next time you talked to them they'd not only have
implemented it, but also implemented two new ideas they had in the
process. ``They probably have the best attitude of any startup we've
funded'' I wrote to Mike Arrington during the batch.

They're still like that. Jessica and I had dinner with Brian in the
summer of 2018, just the three of us. By this point the company is ten
years old. He took a page of notes about ideas for new things Airbnb
could do.

What we didn't realize when we first met Brian and Joe and Nate was that
Airbnb was on its last legs. After working on the company for a year and
getting no growth, they'd agreed to give it one last shot. They'd try
this Y Combinator thing, and if the company still didn't take off,
they'd give up.

Any normal person would have given up already. They'd been funding the
company with credit cards. They had a \emph{binder} full of credit cards
they'd maxed out. Investors didn't think much of the idea. One investor
they met in a cafe walked out in the middle of meeting with them. They
thought he was going to the bathroom, but he never came back. ``He
didn't even finish his smoothie,'' Brian said. And now, in late 2008, it
was the worst recession in decades. The stock market was in free fall
and wouldn't hit bottom for another four months.

Why hadn't they given up? This is a useful question to ask. People, like
matter, reveal their nature under extreme conditions. One thing that's
clear is that they weren't doing this just for the money. As a
money-making scheme, this was pretty lousy: a year's work and all they
had to show for it was a binder full of maxed-out credit cards. So why
were they still working on this startup? Because of the experience
they'd had as the first hosts.

When they first tried renting out airbeds on their floor during a design
convention, all they were hoping for was to make enough money to pay
their rent that month. But something surprising happened: they enjoyed
having those first three guests staying with them. And the guests
enjoyed it too. Both they and the guests had done it because they were
in a sense forced to, and yet they'd all had a great experience. Clearly
there was something new here: for hosts, a new way to make money that
had literally been right under their noses, and for guests, a new way to
travel that was in many ways better than hotels.

That experience was why the Airbnbs didn't give up. They knew they'd
discovered something. They'd seen a glimpse of the future, and they
couldn't let it go.

They knew that once people tried staying in what is now called ``an
airbnb,'' they would also realize that this was the future. But only if
they tried it, and they weren't. That was the problem during Y
Combinator: to get growth started.

Airbnb's goal during YC was to reach what we call
\href{http://paulgraham.com/ramenprofitable.html}{ramen profitability},
which means making enough money that the company can pay the founders'
living expenses, if they live on ramen noodles. Ramen profitability is
not, obviously, the end goal of any startup, but it's the most important
threshold on the way, because this is the point where you're airborne.
This is the point where you no longer need investors' permission to
continue existing. For the Airbnbs, ramen profitability was \$4000 a
month: \$3500 for rent, and \$500 for food. They taped this goal to the
mirror in the bathroom of their apartment.

The way to get growth started in something like Airbnb is to focus on
the hottest subset of the market. If you can get growth started there,
it will spread to the rest. When I asked the Airbnbs where there was
most demand, they knew from searches: New York City. So they focused on
New York. They went there \href{http://paulgraham.com/ds.html}{in
person} to visit their hosts and help them make their listings more
attractive. A big part of that was better pictures. So Joe and Brian
rented a professional camera and took pictures of the hosts' places
themselves.

This didn't just make the listings better. It also taught them about
their hosts. When they came back from their first trip to New York, I
asked what they'd noticed about hosts that surprised them, and they said
the biggest surprise was how many of the hosts were in the same position
they'd been in: they needed this money to pay their rent. This was,
remember, the worst recession in decades, and it had hit New York first.
It definitely added to the Airbnbs' sense of mission to feel that people
needed them.

In late January 2009, about three weeks into Y Combinator, their efforts
started to show results, and their numbers crept upward. But it was hard
to say for sure whether it was growth or just random fluctuation. By
February it was clear that it was real growth. They made \$460 in fees
in the first week of February, \$897 in the second, and \$1428 in the
third. That was it: they were airborne. Brian sent me an email on
February 22 announcing that they were ramen profitable and giving the
last three weeks' numbers.

``I assume you know what you've now set yourself up for next week,'' I
responded.

Brian's reply was seven words: ``We are not going to slow down.''
\chapter{Billionaires Build}

December 2020

As I was deciding what to write about next, I was surprised to find that
two separate essays I'd been planning to write were actually the same.

The first is about how to ace your Y Combinator interview. There has
been so much nonsense written about this topic that I've been meaning
for years to write something telling founders the truth.

The second is about something politicians sometimes say --- that the
only way to become a billionaire is by exploiting people --- and why
this is mistaken.

Keep reading, and you'll learn both simultaneously.

I know the politicians are mistaken because it was my job to predict
which people will become billionaires. I think I can truthfully say that
I know as much about how to do this as anyone. If the key to becoming a
billionaire --- the defining feature of billionaires --- was to exploit
people, then I, as a professional billionaire scout, would surely
realize this and look for people who would be good at it, just as an NFL
scout looks for speed in wide receivers.

But aptitude for exploiting people is not what Y Combinator looks for at
all. In fact, it's the opposite of what they look for. I'll tell you
what they do look for, by explaining how to convince Y~Combinator to
fund you, and you can see for yourself.

What YC looks for, above all, is founders who understand some group of
users and can make what they want. This is so important that it's YC's
motto: ``Make something people want.''

A big company can to some extent force unsuitable products on unwilling
customers, but a startup doesn't have the power to do that. A startup
must sing for its supper, by making things that genuinely delight its
customers. Otherwise it will never get off the ground.

Here's where things get difficult, both for you as a founder and for the
YC partners trying to decide whether to fund you. In a market economy,
it's hard to make something people want that they don't already have.
That's the great thing about market economies. If other people both knew
about this need and were able to satisfy it, they already would be, and
there would be no room for your startup.

Which means the conversation during your YC interview will have to be
about something new: either a new need, or a new way to satisfy one. And
not just new, but uncertain. If it were certain that the need existed
and that you could satisfy it, that certainty would be reflected in
large and rapidly growing revenues, and you wouldn't be seeking seed
funding.

So the YC partners have to guess both whether you've discovered a real
need, and whether you'll be able to satisfy it. That's what they are, at
least in this part of their job: professional guessers. They have 1001
heuristics for doing this, and I'm not going to tell you all of them,
but I'm happy to tell you the most important ones, because these can't
be faked; the only way to ``hack'' them would be to do what you should
be doing anyway as a founder.

The first thing the partners will try to figure out, usually, is whether
what you're making will ever be something a lot of people want. It
doesn't have to be something a lot of people want now. The product and
the market will both evolve, and will influence each other's evolution.
But in the end there has to be something with a huge market. That's what
the partners will be trying to figure out: is there a path to a huge
market? \footnote{The YC partners have so much practice doing this that they
sometimes see paths that the founders themselves haven't seen yet. The
partners don't try to seem skeptical, as buyers in transactions often do
to increase their leverage. Although the founders feel their job is to
convince the partners of the potential of their idea, these roles are
not infrequently reversed, and the founders leave the interview feeling
their idea has more potential than they realized.}

Sometimes it's obvious there will be a huge market. If
\href{https://boomsupersonic.com/}{Boom} manages to ship an airliner at
all, international airlines will have to buy it. But usually it's not
obvious. Usually the path to a huge market is by growing a small market.
This idea is important enough that it's worth coining a phrase for, so
let's call one of these small but growable markets a ``larval market.''

The perfect example of a larval market might be Apple's market when they
were founded in 1976. In 1976, not many people wanted their own
computer. But more and more started to want one, till now every 10 year
old on the planet wants a computer (but calls it a ``phone'').

The ideal combination is the group of founders who are
\href{startupideas.html}{``living in the future''} in the sense of being
at the leading edge of some kind of change, and who are building
something they themselves want. Most super-successful startups are of
this type. Steve Wozniak wanted a computer. Mark Zuckerberg wanted to
engage online with his college friends. Larry and Sergey wanted to find
things on the web. All these founders were building things they and
their peers wanted, and the fact that they were at the leading edge of
change meant that more people would want these things in the future.

But although the ideal larval market is oneself and one's peers, that's
not the only kind. A larval market might also be regional, for example.
You build something to serve one location, and then expand to others.

The crucial feature of the initial market is that it exist. That may
seem like an obvious point, but the lack of it is the biggest flaw in
most startup ideas. There have to be some people who want what you're
building right now, and want it so urgently that they're willing to use
it, bugs and all, even though you're a small company they've never heard
of. There don't have to be many, but there have to be some. As long as
you have some users, there are straightforward ways to get more: build
new features they want, seek out more people like them, get them to
refer you to their friends, and so on. But these techniques all require
some initial seed group of users.

So this is one thing the YC partners will almost certainly dig into
during your interview. Who are your first users going to be, and how do
you know they want this? If I had to decide whether to fund startups
based on a single question, it would be ``How do you know people want
this?''

The most convincing answer is ``Because we and our friends want it.''
It's even better when this is followed by the news that you've already
built a prototype, and even though it's very crude, your friends are
using it, and it's spreading by word of mouth. If you can say that and
you're not lying, the partners will switch from default no to default
yes. Meaning you're in unless there's some other disqualifying flaw.

That is a hard standard to meet, though. Airbnb didn't meet it. They had
the first part. They had made something they themselves wanted. But it
wasn't spreading. So don't feel bad if you don't hit this gold standard
of convincingness. If Airbnb didn't hit it, it must be too high.

In practice, the YC partners will be satisfied if they feel that you
have a deep understanding of your users' needs. And the Airbnbs did have
that. They were able to tell us all about what motivated hosts and
guests. They knew from first-hand experience, because they'd been the
first hosts. We couldn't ask them a question they didn't know the answer
to. We ourselves were not very excited about the idea as users, but we
knew this didn't prove anything, because there were lots of successful
startups we hadn't been excited about as users. We were able to say to
ourselves ``They seem to know what they're talking about. Maybe they're
onto something. It's not growing yet, but maybe they can figure out how
to make it grow during YC.'' Which they did, about three weeks into the
batch.

The best thing you can do in a YC interview is to teach the partners
about your users. So if you want to prepare for your interview, one of
the best ways to do it is to go talk to your users and find out exactly
what they're thinking. Which is what you should be doing anyway.

This may sound strangely credulous, but the YC partners want to rely on
the founders to tell them about the market. Think about how VCs
typically judge the potential market for an idea. They're not ordinarily
domain experts themselves, so they forward the idea to someone who is,
and ask for their opinion. YC doesn't have time to do this, but if the
YC partners can convince themselves that the founders both (a) know what
they're talking about and (b) aren't lying, they don't need outside
domain experts. They can use the founders themselves as domain experts
when evaluating their own idea.

This is why YC interviews aren't pitches. To give as many founders as
possible a chance to get funded, we made interviews as short as we
could: 10 minutes. That is not enough time for the partners to figure
out, through the indirect evidence in a pitch, whether you know what
you're talking about and aren't lying. They need to dig in and ask you
questions. There's not enough time for sequential access. They need
random access. \footnote{In practice, 7 minutes would be enough. You rarely change your
mind at minute 8. But 10 minutes is socially convenient.}

The worst advice I ever heard about how to succeed in a YC interview is
that you should take control of the interview and make sure to deliver
the message you want to. In other words, turn the interview into a
pitch. ⟨elaborate expletive⟩. It is so annoying when people try to do
that. You ask them a question, and instead of answering it, they deliver
some obviously prefabricated blob of pitch. It eats up 10 minutes really
fast.

There is no one who can give you accurate advice about what to do in a
YC interview except a current or former YC partner. People who've merely
been interviewed, even successfully, have no idea of this, but
interviews take all sorts of different forms depending on what the
partners want to know about most. Sometimes they're all about the
founders, other times they're all about the idea. Sometimes some very
narrow aspect of the idea. Founders sometimes walk away from interviews
complaining that they didn't get to explain their idea completely. True,
but they explained enough.

Since a YC interview consists of questions, the way to do it well is to
answer them well. Part of that is answering them candidly. The partners
don't expect you to know everything. But if you don't know the answer to
a question, don't try to bullshit your way out of it. The partners, like
most experienced investors, are professional bullshit detectors, and you
are (hopefully) an amateur bullshitter. And if you try to bullshit them
and fail, they may not even tell you that you failed. So it's better to
be honest than to try to sell them. If you don't know the answer to a
question, say you don't, and tell them how you'd go about finding it, or
tell them the answer to some related question.

If you're asked, for example, what could go wrong, the worst possible
answer is ``nothing.'' Instead of convincing them that your idea is
bullet-proof, this will convince them that you're a fool or a liar. Far
better to go into gruesome detail. That's what experts do when you ask
what could go wrong. The partners know that your idea is risky. That's
what a good bet looks like at this stage: a tiny probability of a huge
outcome.

Ditto if they ask about competitors. Competitors are rarely what kills
startups. Poor execution does. But you should know who your competitors
are, and tell the YC partners candidly what your relative strengths and
weaknesses are. Because the YC partners know that competitors don't kill
startups, they won't hold competitors against you too much. They will,
however, hold it against you if you seem either to be unaware of
competitors, or to be minimizing the threat they pose. They may not be
sure whether you're clueless or lying, but they don't need to be.

The partners don't expect your idea to be perfect. This is seed
investing. At this stage, all they can expect are promising hypotheses.
But they do expect you to be thoughtful and honest. So if trying to make
your idea seem perfect causes you to come off as glib or clueless,
you've sacrificed something you needed for something you didn't.

If the partners are sufficiently convinced that there's a path to a big
market, the next question is whether you'll be able to find it. That in
turn depends on three things: the general qualities of the founders,
their specific expertise in this domain, and the relationship between
them. How determined are the founders? Are they good at building things?
Are they resilient enough to keep going when things go wrong? How strong
is their friendship?

Though the Airbnbs only did ok in the idea department, they did
spectacularly well in this department. The story of how they'd funded
themselves by making Obama- and McCain-themed breakfast cereal was the
single most important factor in our decision to fund them. They didn't
realize it at the time, but what seemed to them an irrelevant story was
in fact fabulously good evidence of their qualities as founders. It
showed they were resourceful and determined, and could work together.

It wasn't just the cereal story that showed that, though. The whole
interview showed that they cared. They weren't doing this just for the
money, or because startups were cool. The reason they were working so
hard on this company was because it was their project. They had
discovered an interesting new idea, and they just couldn't let it go.

Mundane as it sounds, that's the most powerful motivator of all, not
just in startups, but in most ambitious undertakings: to be
\href{genius.html}{genuinely interested} in what you're building. This
is what really drives billionaires, or at least the ones who become
billionaires from starting companies. The company is their project.

One thing few people realize about billionaires is that all of them
could have stopped sooner. They could have gotten acquired, or found
someone else to run the company. Many founders do. The ones who become
really rich are the ones who keep working. And what makes them keep
working is not just money. What keeps them working is the same thing
that keeps anyone else working when they could stop if they wanted to:
that there's nothing else they'd rather do.

That, not exploiting people, is the defining quality of people who
become billionaires from starting companies. So that's what YC looks for
in founders: authenticity. People's motives for starting startups are
usually mixed. They're usually doing it from some combination of the
desire to make money, the desire to seem cool, genuine interest in the
problem, and unwillingness to work for someone else. The last two are
more powerful motivators than the first two. It's ok for founders to
want to make money or to seem cool. Most do. But if the founders seem
like they're doing it \emph{just} to make money or \emph{just} to seem
cool, they're not likely to succeed on a big scale. The founders who are
doing it for the money will take the first sufficiently large
acquisition offer, and the ones who are doing it to seem cool will
rapidly discover that there are much less painful ways of seeming cool.
\footnote{I myself took the first sufficiently large acquisition offer in
my first startup, so I don't blame founders for doing this. There's
nothing wrong with starting a startup to make money. You need to make
money somehow, and for some people startups are the most efficient way
to do it. I'm just saying that these are not the startups that get
really big.}

Y Combinator certainly sees founders whose m.o. is to exploit people. YC
is a magnet for them, because they want the YC brand. But when the YC
partners detect someone like that, they reject them. If bad people made
good founders, the YC partners would face a moral dilemma. Fortunately
they don't, because bad people make bad founders. This exploitative type
of founder is not going to succeed on a large scale, and in fact
probably won't even succeed on a small one, because they're always going
to be taking shortcuts. They see YC itself as a shortcut.

Their exploitation usually begins with their own cofounders, which is
disastrous, since the cofounders' relationship is the foundation of the
company. Then it moves on to the users, which is also disastrous,
because the sort of early adopters a successful startup wants as its
initial users are the hardest to fool. The best this kind of founder can
hope for is to keep the edifice of deception tottering along until some
acquirer can be tricked into buying it. But that kind of acquisition is
never very big. \footnote{Not these days, anyway. There were some big ones during the
Internet Bubble, and indeed some big IPOs.}

If professional billionaire scouts know that exploiting people is not
the skill to look for, why do some politicians think this is the
defining quality of billionaires?

I think they start from the feeling that it's wrong that one person
could have so much more money than another. It's understandable where
that feeling comes from. It's in our DNA, and even in the DNA of other
species.

If they limited themselves to saying that it made them feel bad when one
person had so much more money than other people, who would disagree? It
makes me feel bad too, and I think people who make a lot of money have a
moral obligation to use it for the common good. The mistake they make is
to jump from feeling bad that some people are much richer than others to
the conclusion that there's no legitimate way to make a very large
amount of money. Now we're getting into statements that are not only
falsifiable, but false.

There are certainly some people who become rich by doing bad things. But
there are also plenty of people who behave badly and don't make that
much from it. There is no correlation --- in fact, probably an inverse
correlation --- between how badly you behave and how much money you
make.

The greatest danger of this nonsense may not even be that it sends
policy astray, but that it misleads ambitious people. Can you imagine a
better way to destroy social mobility than by telling poor kids that the
way to get rich is by exploiting people, while the rich kids know, from
having watched the preceding generation do it, how it's really done?

I'll tell you how it's really done, so you can at least tell your own
kids the truth. It's all about users. The most reliable way to become a
billionaire is to start a company that \href{growth.html}{grows fast},
and the way to grow fast is to make what users want. Newly started
startups have no choice but to delight users, or they'll never even get
rolling. But this never stops being the lodestar, and bigger companies
take their eye off it at their peril. Stop delighting users, and
eventually someone else will.

Users are what the partners want to know about in YC interviews, and
what I want to know about when I talk to founders that we funded ten
years ago and who are billionaires now. What do users want? What new
things could you build for them? Founders who've become billionaires are
always eager to talk about that topic. That's how they became
billionaires.


\textbf{Thanks} to Trevor Blackwell, Jessica Livingston, Robert Morris,
Geoff Ralston, and Harj Taggar for reading drafts of this.
\chapter{Earnestness}

December 2020

Jessica and I have certain words that have special significance when
we're talking about startups. The highest compliment we can pay to
founders is to describe them as ``earnest.'' This is not by itself a
guarantee of success. You could be earnest but incapable. But when
founders are both formidable (another of our words) and earnest, they're
as close to unstoppable as you get.

Earnestness sounds like a boring, even Victorian virtue. It seems a bit
of an anachronism that people in Silicon Valley would care about it. Why
does this matter so much?

When you call someone earnest, you're making a statement about their
motives. It means both that they're doing something for the right
reasons, and that they're trying as hard as they can. If we imagine
motives as vectors, it means both the direction and the magnitude are
right. Though these are of course related: when people are doing
something for the right reasons, they try harder.
\footnote{It's interesting how many different ways there are \emph{not} to
be earnest: to be cleverly cynical, to be superficially brilliant, to be
conspicuously virtuous, to be cool, to be sophisticated, to be orthodox,
to be a snob, to bully, to pander, to be on the make. This pattern
suggests that earnestness is not one end of a continuum, but a target
one can fall short of in multiple dimensions.

Another thing I notice about this list is that it sounds like a list of
the ways people behave on Twitter. Whatever else social media is, it's a
vivid catalogue of ways not to be earnest.}

The reason motives matter so much in Silicon Valley is that so many
people there have the wrong ones. Starting a successful startup makes
you rich and famous. So a lot of the people trying to start them are
doing it for those reasons. Instead of what? Instead of interest in the
problem for its own sake. That is the root of earnestness.
\footnote{People's motives are as mixed in Silicon Valley as anywhere
else. Even the founders motivated mostly by money tend to be at least
somewhat interested in the problem they're solving, and even the
founders most interested in the problem they're solving also like the
idea of getting rich. But there's great variation in the relative
proportions of different founders' motivations.

And when I talk about ``wrong'' motives, I don't mean morally wrong.
There's nothing morally wrong with starting a startup to make money. I
just mean that those startups don't do as well.}

It's also the hallmark of a nerd. Indeed, when people describe
themselves as ``x nerds,'' what they mean is that they're interested in
x for its own sake, and not because it's cool to be interested in x, or
because of what they can get from it. They're saying they care so much
about x that they're willing to sacrifice seeming cool for its sake.

A \href{genius.html}{genuine interest} in something is a very powerful
motivator --- for some people, the most powerful motivator of all.
\footnote{The most powerful motivator for most people is probably family.
But there are some for whom intellectual curiosity comes first. In his
(wonderful) autobiography, Paul Halmos says explicitly that for a
mathematician, math must come before anything else, including family.
Which at least implies that it did for him.} Which is why it's what Jessica and I look for in
founders. But as well as being a source of strength, it's also a source
of vulnerability. Caring constrains you. The earnest can't easily reply
in kind to mocking banter, or put on a cool facade of nihil admirari.
They care too much. They are doomed to be the straight man. That's a
real disadvantage in your \href{nerds.html}{teenage years}, when mocking
banter and nihil admirari often have the upper hand. But it becomes an
advantage later.

It's a commonplace now that the kids who were nerds in high school
become the cool kids' bosses later on. But people misunderstand why this
happens. It's not just because the nerds are smarter, but also because
they're more earnest. When the problems get harder than the fake ones
you're given in high school, caring about them starts to matter.

Does it always matter? Do the earnest always win? Not always. It
probably doesn't matter much in politics, or in crime, or in certain
types of business that are similar to crime, like gambling, personal
injury law, patent trolling, and so on. Nor does it matter in academic
fields at the more
\href{https://scholar.google.com/scholar?hl=en&as_sdt=0\%2C5&q=hermeneutic+dialectics+hegemonic+phenomenology+intersectionality}{bogus}
end of the spectrum. And though I don't know enough to say for sure, it
may not matter in some kinds of humor: it may be possible to be
completely cynical and still be very funny. \footnote{Interestingly, just as the word ``nerd'' implies earnestness
even when used as a metaphor, the word ``politics'' implies the
opposite. It's not only in actual politics that earnestness seems to be
a handicap, but also in office politics and academic politics.}

Looking at the list of fields I mentioned, there's an obvious pattern.
Except possibly for humor, these are all types of work I'd avoid like
the plague. So that could be a useful heuristic for deciding which
fields to work in: how much does earnestness matter? Which can in turn
presumably be inferred from the prevalence of nerds at the top.

Along with ``nerd,'' another word that tends to be associated with
earnestness is ``naive.'' The earnest often seem naive. It's not just
that they don't have the motives other people have. They often don't
fully grasp that such motives exist. Or they may know intellectually
that they do, but because they don't feel them, they forget about them.
\footnote{It's a bigger social error to seem naive in most European
countries than it is in America, and this may be one of subtler reasons
startups are less common there. Founder culture is completely at odds
with sophisticated cynicism.

The most earnest part of Europe is Scandinavia, and not surprisingly
this is also the region with the highest number of successful startups
per capita.}

It works to be slightly naive not just about motives but also, believe
it or not, about the problems you're working on. Naive optimism can
compensate for the bit rot that \href{ecw.html}{rapid change} causes in
established beliefs. You plunge into some problem saying ``How hard can
it be?'', and then after solving it you learn that it was till recently
insoluble.

Naivete is an obstacle for anyone who wants to seem sophisticated, and
this is one reason would-be intellectuals find it so difficult to
understand Silicon Valley. It hasn't been safe for such people to use
the word ``earnest'' outside scare quotes since Oscar Wilde wrote ``The
Importance of Being Earnest'' in 1895. And yet when you zoom in on
Silicon Valley, right into \href{jessica.html}{Jessica Livingston's
brain}, that's what her x-ray vision is seeking out in founders.
Earnestness! Who'd have guessed? Reporters literally can't believe it
when founders making piles of money say that they started their
companies to make the world better. The situation seems made for
mockery. How can these founders be so naive as not to realize how
implausible they sound?

Though those asking this question don't realize it, that's not a
rhetorical question.

A lot of founders are faking it, of course, particularly the smaller
fry, and the soon to be smaller fry. But not all of them. There are a
significant number of founders who really are interested in the problem
they're solving mainly for its own sake.

Why shouldn't there be? We have no difficulty believing that people
would be interested in history or math or even old bus tickets for their
own sake. Why can't there be people interested in self-driving cars or
social networks for their own sake? When you look at the question from
this side, it seems obvious there would be. And isn't it likely that
having a deep interest in something would be a source of great energy
and resilience? It is in every other field.

The question really is why we have a blind spot about business. And the
answer to that is obvious if you know enough history. For most of
history, making large amounts of money has not been very intellectually
interesting. In preindustrial times it was never far from robbery, and
some areas of business still retain that character, except using lawyers
instead of soldiers.

But there are other areas of business where the work is genuinely
interesting. Henry Ford got to spend much of his time working on
interesting technical problems, and for the last several decades the
trend in that direction has been accelerating. It's much easier now to
make a lot of money by working on something you're interested in than it
was \href{re.html}{50 years ago}. And that, rather than how fast they
grow, may be the most important change that startups represent. Though
indeed, the fact that the work is genuinely interesting is a big part of
why it gets done so fast. \footnote{Much of business is schleps, and probably always will be. But
even being a professor is largely schleps. It would be interesting to
collect statistics about the schlep ratios of different jobs, but I
suspect they'd rarely be less than 30\%.}

Can you imagine a more important change than one in the relationship
between intellectual curiosity and money? These are two of the most
powerful forces in the world, and in my lifetime they've become
significantly more aligned. How could you not be fascinated to watch
something like this happening in real time?

I meant this essay to be about earnestness generally, and now I've gone
and talked about startups again. But I suppose at least it serves as an
example of an x nerd in the wild.


\textbf{Thanks} to Trevor Blackwell, Patrick Collison, Suhail Doshi,
Jessica Livingston, Mattias Ljungman, Harj Taggar, and Kyle Vogt for
reading drafts of this.
\chapter{What I Worked On}

February 2021

Before college the two main things I worked on, outside of school, were
writing and programming. I didn't write essays. I wrote what beginning
writers were supposed to write then, and probably still are: short
stories. My stories were awful. They had hardly any plot, just
characters with strong feelings, which I imagined made them deep.

The first programs I tried writing were on the IBM 1401 that our school
district used for what was then called ``data processing.'' This was in
9th grade, so I was 13 or 14. The school district's 1401 happened to be
in the basement of our junior high school, and my friend Rich Draves and
I got permission to use it. It was like a mini Bond villain's lair down
there, with all these alien-looking machines --- CPU, disk drives,
printer, card reader --- sitting up on a raised floor under bright
fluorescent lights.

The language we used was an early version of Fortran. You had to type
programs on punch cards, then stack them in the card reader and press a
button to load the program into memory and run it. The result would
ordinarily be to print something on the spectacularly loud printer.

I was puzzled by the 1401. I couldn't figure out what to do with it. And
in retrospect there's not much I could have done with it. The only form
of input to programs was data stored on punched cards, and I didn't have
any data stored on punched cards. The only other option was to do things
that didn't rely on any input, like calculate approximations of pi, but
I didn't know enough math to do anything interesting of that type. So
I'm not surprised I can't remember any programs I wrote, because they
can't have done much. My clearest memory is of the moment I learned it
was possible for programs not to terminate, when one of mine didn't. On
a machine without time-sharing, this was a social as well as a technical
error, as the data center manager's expression made clear.

With microcomputers, everything changed. Now you could have a computer
sitting right in front of you, on a desk, that could respond to your
keystrokes as it was running instead of just churning through a stack of
punch cards and then stopping. \footnote{My experience skipped a step in the evolution of computers:
time-sharing machines with interactive OSes. I went straight from batch
processing to microcomputers, which made microcomputers seem all the
more exciting.}

The first of my friends to get a microcomputer built it himself. It was
sold as a kit by Heathkit. I remember vividly how impressed and envious
I felt watching him sitting in front of it, typing programs right into
the computer.

Computers were expensive in those days and it took me years of nagging
before I convinced my father to buy one, a TRS-80, in about 1980. The
gold standard then was the Apple II, but a TRS-80 was good enough. This
was when I really started programming. I wrote simple games, a program
to predict how high my model rockets would fly, and a word processor
that my father used to write at least one book. There was only room in
memory for about 2 pages of text, so he'd write 2 pages at a time and
then print them out, but it was a lot better than a typewriter.

Though I liked programming, I didn't plan to study it in college. In
college I was going to study philosophy, which sounded much more
powerful. It seemed, to my naive high school self, to be the study of
the ultimate truths, compared to which the things studied in other
fields would be mere domain knowledge. What I discovered when I got to
college was that the other fields took up so much of the space of ideas
that there wasn't much left for these supposed ultimate truths. All that
seemed left for philosophy were edge cases that people in other fields
felt could safely be ignored.

I couldn't have put this into words when I was 18. All I knew at the
time was that I kept taking philosophy courses and they kept being
boring. So I decided to switch to AI.

AI was in the air in the mid 1980s, but there were two things especially
that made me want to work on it: a novel by Heinlein called \emph{The
Moon is a Harsh Mistress}, which featured an intelligent computer called
Mike, and a PBS documentary that showed Terry Winograd using SHRDLU. I
haven't tried rereading \emph{The Moon is a Harsh Mistress}, so I don't
know how well it has aged, but when I read it I was drawn entirely into
its world. It seemed only a matter of time before we'd have Mike, and
when I saw Winograd using SHRDLU, it seemed like that time would be a
few years at most. All you had to do was teach SHRDLU more words.

There weren't any classes in AI at Cornell then, not even graduate
classes, so I started trying to teach myself. Which meant learning Lisp,
since in those days Lisp was regarded as the language of AI. The
commonly used programming languages then were pretty primitive, and
programmers' ideas correspondingly so. The default language at Cornell
was a Pascal-like language called PL/I, and the situation was similar
elsewhere. Learning Lisp expanded my concept of a program so fast that
it was years before I started to have a sense of where the new limits
were. This was more like it; this was what I had expected college to do.
It wasn't happening in a class, like it was supposed to, but that was
ok. For the next couple years I was on a roll. I knew what I was going
to do.

For my undergraduate thesis, I reverse-engineered SHRDLU. My God did I
love working on that program. It was a pleasing bit of code, but what
made it even more exciting was my belief --- hard to imagine now, but
not unique in 1985 --- that it was already climbing the lower slopes of
intelligence.

I had gotten into a program at Cornell that didn't make you choose a
major. You could take whatever classes you liked, and choose whatever
you liked to put on your degree. I of course chose ``Artificial
Intelligence.'' When I got the actual physical diploma, I was dismayed
to find that the quotes had been included, which made them read as
scare-quotes. At the time this bothered me, but now it seems amusingly
accurate, for reasons I was about to discover.

I applied to 3 grad schools: MIT and Yale, which were renowned for AI at
the time, and Harvard, which I'd visited because Rich Draves went there,
and was also home to Bill Woods, who'd invented the type of parser I
used in my SHRDLU clone. Only Harvard accepted me, so that was where I
went.

I don't remember the moment it happened, or if there even was a specific
moment, but during the first year of grad school I realized that AI, as
practiced at the time, was a hoax. By which I mean the sort of AI in
which a program that's told ``the dog is sitting on the chair''
translates this into some formal representation and adds it to the list
of things it knows.

What these programs really showed was that there's a subset of natural
language that's a formal language. But a very proper subset. It was
clear that there was an unbridgeable gap between what they could do and
actually understanding natural language. It was not, in fact, simply a
matter of teaching SHRDLU more words. That whole way of doing AI, with
explicit data structures representing concepts, was not going to work.
Its brokenness did, as so often happens, generate a lot of opportunities
to write papers about various band-aids that could be applied to it, but
it was never going to get us Mike.

So I looked around to see what I could salvage from the wreckage of my
plans, and there was Lisp. I knew from experience that Lisp was
interesting for its own sake and not just for its association with AI,
even though that was the main reason people cared about it at the time.
So I decided to focus on Lisp. In fact, I decided to write a book about
Lisp hacking. It's scary to think how little I knew about Lisp hacking
when I started writing that book. But there's nothing like writing a
book about something to help you learn it. The book, \emph{On Lisp},
wasn't published till 1993, but I wrote much of it in grad school.

Computer Science is an uneasy alliance between two halves, theory and
systems. The theory people prove things, and the systems people build
things. I wanted to build things. I had plenty of respect for theory ---
indeed, a sneaking suspicion that it was the more admirable of the two
halves --- but building things seemed so much more exciting.

The problem with systems work, though, was that it didn't last. Any
program you wrote today, no matter how good, would be obsolete in a
couple decades at best. People might mention your software in footnotes,
but no one would actually use it. And indeed, it would seem very feeble
work. Only people with a sense of the history of the field would even
realize that, in its time, it had been good.

There were some surplus Xerox Dandelions floating around the computer
lab at one point. Anyone who wanted one to play around with could have
one. I was briefly tempted, but they were so slow by present standards;
what was the point? No one else wanted one either, so off they went.
That was what happened to systems work.

I wanted not just to build things, but to build things that would last.

In this dissatisfied state I went in 1988 to visit Rich Draves at CMU,
where he was in grad school. One day I went to visit the Carnegie
Institute, where I'd spent a lot of time as a kid. While looking at a
painting there I realized something that might seem obvious, but was a
big surprise to me. There, right on the wall, was something you could
make that would last. Paintings didn't become obsolete. Some of the best
ones were hundreds of years old.

And moreover this was something you could make a living doing. Not as
easily as you could by writing software, of course, but I thought if you
were really industrious and lived really cheaply, it had to be possible
to make enough to survive. And as an artist you could be truly
independent. You wouldn't have a boss, or even need to get research
funding.

I had always liked looking at paintings. Could I make them? I had no
idea. I'd never imagined it was even possible. I knew intellectually
that people made art --- that it didn't just appear spontaneously ---
but it was as if the people who made it were a different species. They
either lived long ago or were mysterious geniuses doing strange things
in profiles in \emph{Life} magazine. The idea of actually being able to
make art, to put that verb before that noun, seemed almost miraculous.

That fall I started taking art classes at Harvard. Grad students could
take classes in any department, and my advisor, Tom Cheatham, was very
easy going. If he even knew about the strange classes I was taking, he
never said anything.

So now I was in a PhD program in computer science, yet planning to be an
artist, yet also genuinely in love with Lisp hacking and working away at
\emph{On Lisp}. In other words, like many a grad student, I was working
energetically on multiple projects that were not my thesis.

I didn't see a way out of this situation. I didn't want to drop out of
grad school, but how else was I going to get out? I remember when my
friend Robert Morris got kicked out of Cornell for writing the internet
worm of 1988, I was envious that he'd found such a spectacular way to
get out of grad school.

Then one day in April 1990 a crack appeared in the wall. I ran into
professor Cheatham and he asked if I was far enough along to graduate
that June. I didn't have a word of my dissertation written, but in what
must have been the quickest bit of thinking in my life, I decided to
take a shot at writing one in the 5 weeks or so that remained before the
deadline, reusing parts of \emph{On Lisp} where I could, and I was able
to respond, with no perceptible delay ``Yes, I think so. I'll give you
something to read in a few days.''

I picked applications of continuations as the topic. In retrospect I
should have written about macros and embedded languages. There's a whole
world there that's barely been explored. But all I wanted was to get out
of grad school, and my rapidly written dissertation sufficed, just
barely.

Meanwhile I was applying to art schools. I applied to two: RISD in the
US, and the Accademia di Belli Arti in Florence, which, because it was
the oldest art school, I imagined would be good. RISD accepted me, and I
never heard back from the Accademia, so off to Providence I went.

I'd applied for the BFA program at RISD, which meant in effect that I
had to go to college again. This was not as strange as it sounds,
because I was only 25, and art schools are full of people of different
ages. RISD counted me as a transfer sophomore and said I had to do the
foundation that summer. The foundation means the classes that everyone
has to take in fundamental subjects like drawing, color, and design.

Toward the end of the summer I got a big surprise: a letter from the
Accademia, which had been delayed because they'd sent it to Cambridge
England instead of Cambridge Massachusetts, inviting me to take the
entrance exam in Florence that fall. This was now only weeks away. My
nice landlady let me leave my stuff in her attic. I had some money saved
from consulting work I'd done in grad school; there was probably enough
to last a year if I lived cheaply. Now all I had to do was learn
Italian.

Only \emph{stranieri} (foreigners) had to take this entrance exam. In
retrospect it may well have been a way of excluding them, because there
were so many \emph{stranieri} attracted by the idea of studying art in
Florence that the Italian students would otherwise have been
outnumbered. I was in decent shape at painting and drawing from the RISD
foundation that summer, but I still don't know how I managed to pass the
written exam. I remember that I answered the essay question by writing
about Cezanne, and that I cranked up the intellectual level as high as I
could to make the most of my limited vocabulary. \footnote{Italian words for abstract concepts can nearly always be
predicted from their English cognates (except for occasional traps like
\emph{polluzione}). It's the everyday words that differ. So if you
string together a lot of abstract concepts with a few simple verbs, you
can make a little Italian go a long way.}

I'm only up to age 25 and already there are such conspicuous patterns.
Here I was, yet again about to attend some august institution in the
hopes of learning about some prestigious subject, and yet again about to
be disappointed. The students and faculty in the painting department at
the Accademia were the nicest people you could imagine, but they had
long since arrived at an arrangement whereby the students wouldn't
require the faculty to teach anything, and in return the faculty
wouldn't require the students to learn anything. And at the same time
all involved would adhere outwardly to the conventions of a 19th century
atelier. We actually had one of those little stoves, fed with kindling,
that you see in 19th century studio paintings, and a nude model sitting
as close to it as possible without getting burned. Except hardly anyone
else painted her besides me. The rest of the students spent their time
chatting or occasionally trying to imitate things they'd seen in
American art magazines.

Our model turned out to live just down the street from me. She made a
living from a combination of modelling and making fakes for a local
antique dealer. She'd copy an obscure old painting out of a book, and
then he'd take the copy and maltreat it to make it look old.
\footnote{I lived at Piazza San Felice 4, so my walk to the Accademia went
straight down the spine of old Florence: past the Pitti, across the
bridge, past Orsanmichele, between the Duomo and the Baptistery, and
then up Via Ricasoli to Piazza San Marco. I saw Florence at street level
in every possible condition, from empty dark winter evenings to
sweltering summer days when the streets were packed with tourists.}

While I was a student at the Accademia I started painting still lives in
my bedroom at night. These paintings were tiny, because the room was,
and because I painted them on leftover scraps of canvas, which was all I
could afford at the time. Painting still lives is different from
painting people, because the subject, as its name suggests, can't move.
People can't sit for more than about 15 minutes at a time, and when they
do they don't sit very still. So the traditional m.o. for painting
people is to know how to paint a generic person, which you then modify
to match the specific person you're painting. Whereas a still life you
can, if you want, copy pixel by pixel from what you're seeing. You don't
want to stop there, of course, or you get merely photographic accuracy,
and what makes a still life interesting is that it's been through a
head. You want to emphasize the visual cues that tell you, for example,
that the reason the color changes suddenly at a certain point is that
it's the edge of an object. By subtly emphasizing such things you can
make paintings that are more realistic than photographs not just in some
metaphorical sense, but in the strict information-theoretic sense.
\footnote{You can of course paint people like still lives if you want to,
and they're willing. That sort of portrait is arguably the apex of still
life painting, though the long sitting does tend to produce pained
expressions in the sitters.}

I liked painting still lives because I was curious about what I was
seeing. In everyday life, we aren't consciously aware of much we're
seeing. Most visual perception is handled by low-level processes that
merely tell your brain ``that's a water droplet'' without telling you
details like where the lightest and darkest points are, or ``that's a
bush'' without telling you the shape and position of every leaf. This is
a feature of brains, not a bug. In everyday life it would be distracting
to notice every leaf on every bush. But when you have to paint
something, you have to look more closely, and when you do there's a lot
to see. You can still be noticing new things after days of trying to
paint something people usually take for granted, just as you can after
days of trying to write an essay about something people usually take for
granted.

This is not the only way to paint. I'm not 100\% sure it's even a good
way to paint. But it seemed a good enough bet to be worth trying.

Our teacher, professor Ulivi, was a nice guy. He could see I worked
hard, and gave me a good grade, which he wrote down in a sort of
passport each student had. But the Accademia wasn't teaching me anything
except Italian, and my money was running out, so at the end of the first
year I went back to the US.

I wanted to go back to RISD, but I was now broke and RISD was very
expensive, so I decided to get a job for a year and then return to RISD
the next fall. I got one at a company called Interleaf, which made
software for creating documents. You mean like Microsoft Word? Exactly.
That was how I learned that low end software tends to eat high end
software. But Interleaf still had a few years to live yet.
\footnote{Interleaf was one of many companies that had smart people and
built impressive technology, and yet got crushed by Moore's Law. In the
1990s the exponential growth in the power of commodity (i.e.~Intel)
processors rolled up high-end, special-purpose hardware and software
companies like a bulldozer.}

Interleaf had done something pretty bold. Inspired by Emacs, they'd
added a scripting language, and even made the scripting language a
dialect of Lisp. Now they wanted a Lisp hacker to write things in it.
This was the closest thing I've had to a normal job, and I hereby
apologize to my boss and coworkers, because I was a bad employee. Their
Lisp was the thinnest icing on a giant C cake, and since I didn't know C
and didn't want to learn it, I never understood most of the software.
Plus I was terribly irresponsible. This was back when a programming job
meant showing up every day during certain working hours. That seemed
unnatural to me, and on this point the rest of the world is coming
around to my way of thinking, but at the time it caused a lot of
friction. Toward the end of the year I spent much of my time
surreptitiously working on \emph{On Lisp}, which I had by this time
gotten a contract to publish.

The good part was that I got paid huge amounts of money, especially by
art student standards. In Florence, after paying my part of the rent, my
budget for everything else had been \$7 a day. Now I was getting paid
more than 4 times that every hour, even when I was just sitting in a
meeting. By living cheaply I not only managed to save enough to go back
to RISD, but also paid off my college loans.

I learned some useful things at Interleaf, though they were mostly about
what not to do. I learned that it's better for technology companies to
be run by product people than sales people (though sales is a real skill
and people who are good at it are really good at it), that it leads to
bugs when code is edited by too many people, that cheap office space is
no bargain if it's depressing, that planned meetings are inferior to
corridor conversations, that big, bureaucratic customers are a dangerous
source of money, and that there's not much overlap between conventional
office hours and the optimal time for hacking, or conventional offices
and the optimal place for it.

But the most important thing I learned, and which I used in both Viaweb
and Y Combinator, is that the low end eats the high end: that it's good
to be the ``entry level'' option, even though that will be less
prestigious, because if you're not, someone else will be, and will
squash you against the ceiling. Which in turn means that prestige is a
danger sign.

When I left to go back to RISD the next fall, I arranged to do freelance
work for the group that did projects for customers, and this was how I
survived for the next several years. When I came back to visit for a
project later on, someone told me about a new thing called HTML, which
was, as he described it, a derivative of SGML. Markup language
enthusiasts were an occupational hazard at Interleaf and I ignored him,
but this HTML thing later became a big part of my life.

In the fall of 1992 I moved back to Providence to continue at RISD. The
foundation had merely been intro stuff, and the Accademia had been a
(very civilized) joke. Now I was going to see what real art school was
like. But alas it was more like the Accademia than not. Better
organized, certainly, and a lot more expensive, but it was now becoming
clear that art school did not bear the same relationship to art that
medical school bore to medicine. At least not the painting department.
The textile department, which my next door neighbor belonged to, seemed
to be pretty rigorous. No doubt illustration and architecture were too.
But painting was post-rigorous. Painting students were supposed to
express themselves, which to the more worldly ones meant to try to cook
up some sort of distinctive signature style.

A signature style is the visual equivalent of what in show business is
known as a ``schtick'': something that immediately identifies the work
as yours and no one else's. For example, when you see a painting that
looks like a certain kind of cartoon, you know it's by Roy Lichtenstein.
So if you see a big painting of this type hanging in the apartment of a
hedge fund manager, you know he paid millions of dollars for it. That's
not always why artists have a signature style, but it's usually why
buyers pay a lot for such work. \footnote{The signature style seekers at RISD weren't specifically
mercenary. In the art world, money and coolness are tightly coupled.
Anything expensive comes to be seen as cool, and anything seen as cool
will soon become equally expensive.}

There were plenty of earnest students too: kids who ``could draw'' in
high school, and now had come to what was supposed to be the best art
school in the country, to learn to draw even better. They tended to be
confused and demoralized by what they found at RISD, but they kept
going, because painting was what they did. I was not one of the kids who
could draw in high school, but at RISD I was definitely closer to their
tribe than the tribe of signature style seekers.

I learned a lot in the color class I took at RISD, but otherwise I was
basically teaching myself to paint, and I could do that for free. So in
1993 I dropped out. I hung around Providence for a bit, and then my
college friend Nancy Parmet did me a big favor. A rent-controlled
apartment in a building her mother owned in New York was becoming
vacant. Did I want it? It wasn't much more than my current place, and
New York was supposed to be where the artists were. So yes, I wanted it!
\footnote{Technically the apartment wasn't rent-controlled but
rent-stabilized, but this is a refinement only New Yorkers would know or
care about. The point is that it was really cheap, less than half market
price.}

Asterix comics begin by zooming in on a tiny corner of Roman Gaul that
turns out not to be controlled by the Romans. You can do something
similar on a map of New York City: if you zoom in on the Upper East
Side, there's a tiny corner that's not rich, or at least wasn't in 1993.
It's called Yorkville, and that was my new home. Now I was a New York
artist --- in the strictly technical sense of making paintings and
living in New York.

I was nervous about money, because I could sense that Interleaf was on
the way down. Freelance Lisp hacking work was very rare, and I didn't
want to have to program in another language, which in those days would
have meant C++ if I was lucky. So with my unerring nose for financial
opportunity, I decided to write another book on Lisp. This would be a
popular book, the sort of book that could be used as a textbook. I
imagined myself living frugally off the royalties and spending all my
time painting. (The painting on the cover of this book, \emph{ANSI
Common Lisp}, is one that I painted around this time.)

The best thing about New York for me was the presence of Idelle and
Julian Weber. Idelle Weber was a painter, one of the early
photorealists, and I'd taken her painting class at Harvard. I've never
known a teacher more beloved by her students. Large numbers of former
students kept in touch with her, including me. After I moved to New York
I became her de facto studio assistant.

She liked to paint on big, square canvases, 4 to 5 feet on a side. One
day in late 1994 as I was stretching one of these monsters there was
something on the radio about a famous fund manager. He wasn't that much
older than me, and was super rich. The thought suddenly occurred to me:
why don't I become rich? Then I'll be able to work on whatever I want.

Meanwhile I'd been hearing more and more about this new thing called the
World Wide Web. Robert Morris showed it to me when I visited him in
Cambridge, where he was now in grad school at Harvard. It seemed to me
that the web would be a big deal. I'd seen what graphical user
interfaces had done for the popularity of microcomputers. It seemed like
the web would do the same for the internet.

If I wanted to get rich, here was the next train leaving the station. I
was right about that part. What I got wrong was the idea. I decided we
should start a company to put art galleries online. I can't honestly
say, after reading so many Y Combinator applications, that this was the
worst startup idea ever, but it was up there. Art galleries didn't want
to be online, and still don't, not the fancy ones. That's not how they
sell. I wrote some software to generate web sites for galleries, and
Robert wrote some to resize images and set up an http server to serve
the pages. Then we tried to sign up galleries. To call this a difficult
sale would be an understatement. It was difficult to give away. A few
galleries let us make sites for them for free, but none paid us.

Then some online stores started to appear, and I realized that except
for the order buttons they were identical to the sites we'd been
generating for galleries. This impressive-sounding thing called an
``internet storefront'' was something we already knew how to build.

So in the summer of 1995, after I submitted the camera-ready copy of
\emph{ANSI Common Lisp} to the publishers, we started trying to write
software to build online stores. At first this was going to be normal
desktop software, which in those days meant Windows software. That was
an alarming prospect, because neither of us knew how to write Windows
software or wanted to learn. We lived in the Unix world. But we decided
we'd at least try writing a prototype store builder on Unix. Robert
wrote a shopping cart, and I wrote a new site generator for stores ---
in Lisp, of course.

We were working out of Robert's apartment in Cambridge. His roommate was
away for big chunks of time, during which I got to sleep in his room.
For some reason there was no bed frame or sheets, just a mattress on the
floor. One morning as I was lying on this mattress I had an idea that
made me sit up like a capital L. What if we ran the software on the
server, and let users control it by clicking on links? Then we'd never
have to write anything to run on users' computers. We could generate the
sites on the same server we'd serve them from. Users wouldn't need
anything more than a browser.

This kind of software, known as a web app, is common now, but at the
time it wasn't clear that it was even possible. To find out, we decided
to try making a version of our store builder that you could control
through the browser. A couple days later, on August 12, we had one that
worked. The UI was horrible, but it proved you could build a whole store
through the browser, without any client software or typing anything into
the command line on the server.

Now we felt like we were really onto something. I had visions of a whole
new generation of software working this way. You wouldn't need versions,
or ports, or any of that crap. At Interleaf there had been a whole group
called Release Engineering that seemed to be at least as big as the
group that actually wrote the software. Now you could just update the
software right on the server.

We started a new company we called Viaweb, after the fact that our
software worked via the web, and we got \$10,000 in seed funding from
Idelle's husband Julian. In return for that and doing the initial legal
work and giving us business advice, we gave him 10\% of the company. Ten
years later this deal became the model for Y Combinator's. We knew
founders needed something like this, because we'd needed it ourselves.

At this stage I had a negative net worth, because the thousand dollars
or so I had in the bank was more than counterbalanced by what I owed the
government in taxes. (Had I diligently set aside the proper proportion
of the money I'd made consulting for Interleaf? No, I had not.) So
although Robert had his graduate student stipend, I needed that seed
funding to live on.

We originally hoped to launch in September, but we got more ambitious
about the software as we worked on it. Eventually we managed to build a
WYSIWYG site builder, in the sense that as you were creating pages, they
looked exactly like the static ones that would be generated later,
except that instead of leading to static pages, the links all referred
to closures stored in a hash table on the server.

It helped to have studied art, because the main goal of an online store
builder is to make users look legit, and the key to looking legit is
high production values. If you get page layouts and fonts and colors
right, you can make a guy running a store out of his bedroom look more
legit than a big company.

(If you're curious why my site looks so old-fashioned, it's because it's
still made with this software. It may look clunky today, but in 1996 it
was the last word in slick.)

In September, Robert rebelled. ``We've been working on this for a
month,'' he said, ``and it's still not done.'' This is funny in
retrospect, because he would still be working on it almost 3 years
later. But I decided it might be prudent to recruit more programmers,
and I asked Robert who else in grad school with him was really good. He
recommended Trevor Blackwell, which surprised me at first, because at
that point I knew Trevor mainly for his plan to reduce everything in his
life to a stack of notecards, which he carried around with him. But Rtm
was right, as usual. Trevor turned out to be a frighteningly effective
hacker.

It was a lot of fun working with Robert and Trevor. They're the two most
\href{think.html}{independent-minded} people I know, and in completely
different ways. If you could see inside Rtm's brain it would look like a
colonial New England church, and if you could see inside Trevor's it
would look like the worst excesses of Austrian Rococo.

We opened for business, with 6 stores, in January 1996. It was just as
well we waited a few months, because although we worried we were late,
we were actually almost fatally early. There was a lot of talk in the
press then about ecommerce, but not many people actually wanted online
stores. \footnote{Most software you can launch as soon as it's done. But when the
software is an online store builder and you're hosting the stores, if
you don't have any users yet, that fact will be painfully obvious. So
before we could launch publicly we had to launch privately, in the sense
of recruiting an initial set of users and making sure they had
decent-looking stores.}

There were three main parts to the software: the editor, which people
used to build sites and which I wrote, the shopping cart, which Robert
wrote, and the manager, which kept track of orders and statistics, and
which Trevor wrote. In its time, the editor was one of the best
general-purpose site builders. I kept the code tight and didn't have to
integrate with any other software except Robert's and Trevor's, so it
was quite fun to work on. If all I'd had to do was work on this
software, the next 3 years would have been the easiest of my life.
Unfortunately I had to do a lot more, all of it stuff I was worse at
than programming, and the next 3 years were instead the most stressful.

There were a lot of startups making ecommerce software in the second
half of the 90s. We were determined to be the Microsoft Word, not the
Interleaf. Which meant being easy to use and inexpensive. It was lucky
for us that we were poor, because that caused us to make Viaweb even
more inexpensive than we realized. We charged \$100 a month for a small
store and \$300 a month for a big one. This low price was a big
attraction, and a constant thorn in the sides of competitors, but it
wasn't because of some clever insight that we set the price low. We had
no idea what businesses paid for things. \$300 a month seemed like a lot
of money to us.

We did a lot of things right by accident like that. For example, we did
what's now called ``doing things that \href{ds.html}{don't scale},''
although at the time we would have described it as ``being so lame that
we're driven to the most desperate measures to get users.'' The most
common of which was building stores for them. This seemed particularly
humiliating, since the whole raison d'etre of our software was that
people could use it to make their own stores. But anything to get users.

We learned a lot more about retail than we wanted to know. For example,
that if you could only have a small image of a man's shirt (and all
images were small then by present standards), it was better to have a
closeup of the collar than a picture of the whole shirt. The reason I
remember learning this was that it meant I had to rescan about 30 images
of men's shirts. My first set of scans were so beautiful too.

Though this felt wrong, it was exactly the right thing to be doing.
Building stores for users taught us about retail, and about how it felt
to use our software. I was initially both mystified and repelled by
``business'' and thought we needed a ``business person'' to be in charge
of it, but once we started to get users, I was converted, in much the
same way I was converted to \href{kids.html}{fatherhood} once I had
kids. Whatever users wanted, I was all theirs. Maybe one day we'd have
so many users that I couldn't scan their images for them, but in the
meantime there was nothing more important to do.

Another thing I didn't get at the time is that \href{growth.html}{growth
rate} is the ultimate test of a startup. Our growth rate was fine. We
had about 70 stores at the end of 1996 and about 500 at the end of 1997.
I mistakenly thought the thing that mattered was the absolute number of
users. And that is the thing that matters in the sense that that's how
much money you're making, and if you're not making enough, you might go
out of business. But in the long term the growth rate takes care of the
absolute number. If we'd been a startup I was advising at Y Combinator,
I would have said: Stop being so stressed out, because you're doing
fine. You're growing 7x a year. Just don't hire too many more people and
you'll soon be profitable, and then you'll control your own destiny.

Alas I hired lots more people, partly because our investors wanted me
to, and partly because that's what startups did during the Internet
Bubble. A company with just a handful of employees would have seemed
amateurish. So we didn't reach breakeven until about when Yahoo bought
us in the summer of 1998. Which in turn meant we were at the mercy of
investors for the entire life of the company. And since both we and our
investors were noobs at startups, the result was a mess even by startup
standards.

It was a huge relief when Yahoo bought us. In principle our Viaweb stock
was valuable. It was a share in a business that was profitable and
growing rapidly. But it didn't feel very valuable to me; I had no idea
how to value a business, but I was all too keenly aware of the
near-death experiences we seemed to have every few months. Nor had I
changed my grad student lifestyle significantly since we started. So
when Yahoo bought us it felt like going from rags to riches. Since we
were going to California, I bought a car, a yellow 1998 VW GTI. I
remember thinking that its leather seats alone were by far the most
luxurious thing I owned.

The next year, from the summer of 1998 to the summer of 1999, must have
been the least productive of my life. I didn't realize it at the time,
but I was worn out from the effort and stress of running Viaweb. For a
while after I got to California I tried to continue my usual m.o. of
programming till 3 in the morning, but fatigue combined with Yahoo's
prematurely aged \href{yahoo.html}{culture} and grim cube farm in Santa
Clara gradually dragged me down. After a few months it felt
disconcertingly like working at Interleaf.

Yahoo had given us a lot of options when they bought us. At the time I
thought Yahoo was so overvalued that they'd never be worth anything, but
to my astonishment the stock went up 5x in the next year. I hung on till
the first chunk of options vested, then in the summer of 1999 I left. It
had been so long since I'd painted anything that I'd half forgotten why
I was doing this. My brain had been entirely full of software and men's
shirts for 4 years. But I had done this to get rich so I could paint, I
reminded myself, and now I was rich, so I should go paint.

When I said I was leaving, my boss at Yahoo had a long conversation with
me about my plans. I told him all about the kinds of pictures I wanted
to paint. At the time I was touched that he took such an interest in me.
Now I realize it was because he thought I was lying. My options at that
point were worth about \$2 million a month. If I was leaving that kind
of money on the table, it could only be to go and start some new
startup, and if I did, I might take people with me. This was the height
of the Internet Bubble, and Yahoo was ground zero of it. My boss was at
that moment a billionaire. Leaving then to start a new startup must have
seemed to him an insanely, and yet also plausibly, ambitious plan.

But I really was quitting to paint, and I started immediately. There was
no time to lose. I'd already burned 4 years getting rich. Now when I
talk to founders who are leaving after selling their companies, my
advice is always the same: take a vacation. That's what I should have
done, just gone off somewhere and done nothing for a month or two, but
the idea never occurred to me.

So I tried to paint, but I just didn't seem to have any energy or
ambition. Part of the problem was that I didn't know many people in
California. I'd compounded this problem by buying a house up in the
Santa Cruz Mountains, with a beautiful view but miles from anywhere. I
stuck it out for a few more months, then in desperation I went back to
New York, where unless you understand about rent control you'll be
surprised to hear I still had my apartment, sealed up like a tomb of my
old life. Idelle was in New York at least, and there were other people
trying to paint there, even though I didn't know any of them.

When I got back to New York I resumed my old life, except now I was
rich. It was as weird as it sounds. I resumed all my old patterns,
except now there were doors where there hadn't been. Now when I was
tired of walking, all I had to do was raise my hand, and (unless it was
raining) a taxi would stop to pick me up. Now when I walked past
charming little restaurants I could go in and order lunch. It was
exciting for a while. Painting started to go better. I experimented with
a new kind of still life where I'd paint one painting in the old way,
then photograph it and print it, blown up, on canvas, and then use that
as the underpainting for a second still life, painted from the same
objects (which hopefully hadn't rotted yet).

Meanwhile I looked for an apartment to buy. Now I could actually choose
what neighborhood to live in. Where, I asked myself and various real
estate agents, is the Cambridge of New York? Aided by occasional visits
to actual Cambridge, I gradually realized there wasn't one. Huh.

Around this time, in the spring of 2000, I had an idea. It was clear
from our experience with Viaweb that web apps were the future. Why not
build a web app for making web apps? Why not let people edit code on our
server through the browser, and then host the resulting applications for
them? \footnote{We'd had a code editor in Viaweb for users to define their own
page styles. They didn't know it, but they were editing Lisp expressions
underneath. But this wasn't an app editor, because the code ran when the
merchants' sites were generated, not when shoppers visited them.} You could run all sorts of services on the
servers that these applications could use just by making an API call:
making and receiving phone calls, manipulating images, taking credit
card payments, etc.

I got so excited about this idea that I couldn't think about anything
else. It seemed obvious that this was the future. I didn't particularly
want to start another company, but it was clear that this idea would
have to be embodied as one, so I decided to move to Cambridge and start
it. I hoped to lure Robert into working on it with me, but there I ran
into a hitch. Robert was now a postdoc at MIT, and though he'd made a
lot of money the last time I'd lured him into working on one of my
schemes, it had also been a huge time sink. So while he agreed that it
sounded like a plausible idea, he firmly refused to work on it.

Hmph. Well, I'd do it myself then. I recruited Dan Giffin, who had
worked for Viaweb, and two undergrads who wanted summer jobs, and we got
to work trying to build what it's now clear is about twenty companies
and several open source projects worth of software. The language for
defining applications would of course be a dialect of Lisp. But I wasn't
so naive as to assume I could spring an overt Lisp on a general
audience; we'd hide the parentheses, like Dylan did.

By then there was a name for the kind of company Viaweb was, an
``application service provider,'' or ASP. This name didn't last long
before it was replaced by ``software as a service,'' but it was current
for long enough that I named this new company after it: it was going to
be called Aspra.

I started working on the application builder, Dan worked on network
infrastructure, and the two undergrads worked on the first two services
(images and phone calls). But about halfway through the summer I
realized I really didn't want to run a company --- especially not a big
one, which it was looking like this would have to be. I'd only started
Viaweb because I needed the money. Now that I didn't need money anymore,
why was I doing this? If this vision had to be realized as a company,
then screw the vision. I'd build a subset that could be done as an open
source project.

Much to my surprise, the time I spent working on this stuff was not
wasted after all. After we started Y Combinator, I would often encounter
startups working on parts of this new architecture, and it was very
useful to have spent so much time thinking about it and even trying to
write some of it.

The subset I would build as an open source project was the new Lisp,
whose parentheses I now wouldn't even have to hide. A lot of Lisp
hackers dream of building a new Lisp, partly because one of the
distinctive features of the language is that it has dialects, and
partly, I think, because we have in our minds a Platonic form of Lisp
that all existing dialects fall short of. I certainly did. So at the end
of the summer Dan and I switched to working on this new dialect of Lisp,
which I called Arc, in a house I bought in Cambridge.

The following spring, lightning struck. I was invited to give a talk at
a Lisp conference, so I gave one about how we'd used Lisp at Viaweb.
Afterward I put a postscript file of this talk online, on
paulgraham.com, which I'd created years before using Viaweb but had
never used for anything. In one day it got 30,000 page views. What on
earth had happened? The referring urls showed that someone had posted it
on Slashdot. \footnote{This was the first instance of what is now a familiar
experience, and so was what happened next, when I read the comments and
found they were full of angry people. How could I claim that Lisp was
better than other languages? Weren't they all Turing complete? People
who see the responses to essays I write sometimes tell me how sorry they
feel for me, but I'm not exaggerating when I reply that it has always
been like this, since the very beginning. It comes with the territory.
An essay must tell readers things they \href{useful.html}{don't already
know}, and some people dislike being told such things.}

Wow, I thought, there's an audience. If I write something and put it on
the web, anyone can read it. That may seem obvious now, but it was
surprising then. In the print era there was a narrow channel to readers,
guarded by fierce monsters known as editors. The only way to get an
audience for anything you wrote was to get it published as a book, or in
a newspaper or magazine. Now anyone could publish anything.

This had been possible in principle since 1993, but not many people had
realized it yet. I had been intimately involved with building the
infrastructure of the web for most of that time, and a writer as well,
and it had taken me 8 years to realize it. Even then it took me several
years to understand the implications. It meant there would be a whole
new generation of \href{essay.html}{essays}. \footnote{People put plenty of stuff on the internet in the 90s of
course, but putting something online is not the same as publishing it
online. Publishing online means you treat the online version as the (or
at least a) primary version.}

In the print era, the channel for publishing essays had been vanishingly
small. Except for a few officially anointed thinkers who went to the
right parties in New York, the only people allowed to publish essays
were specialists writing about their specialties. There were so many
essays that had never been written, because there had been no way to
publish them. Now they could be, and I was going to write them.
\footnote{There is a general lesson here that our experience with Y
Combinator also teaches: Customs continue to constrain you long after
the restrictions that caused them have disappeared. Customary VC
practice had once, like the customs about publishing essays, been based
on real constraints. Startups had once been much more expensive to
start, and proportionally rare. Now they could be cheap and common, but
the VCs' customs still reflected the old world, just as customs about
writing essays still reflected the constraints of the print era.

Which in turn implies that people who are independent-minded (i.e.~less
influenced by custom) will have an advantage in fields affected by rapid
change (where customs are more likely to be obsolete).

Here's an interesting point, though: you can't always predict which
fields will be affected by rapid change. Obviously software and venture
capital will be, but who would have predicted that essay writing would
be?}

I've worked on several different things, but to the extent there was a
turning point where I figured out what to work on, it was when I started
publishing essays online. From then on I knew that whatever else I did,
I'd always write essays too.

I knew that online essays would be a \href{marginal.html}{marginal}
medium at first. Socially they'd seem more like rants posted by nutjobs
on their GeoCities sites than the genteel and beautifully typeset
compositions published in \emph{The New Yorker}. But by this point I
knew enough to find that encouraging instead of discouraging.

One of the most conspicuous patterns I've noticed in my life is how well
it has worked, for me at least, to work on things that weren't
prestigious. Still life has always been the least prestigious form of
painting. Viaweb and Y Combinator both seemed lame when we started them.
I still get the glassy eye from strangers when they ask what I'm
writing, and I explain that it's an essay I'm going to publish on my web
site. Even Lisp, though prestigious intellectually in something like the
way Latin is, also seems about as hip.

It's not that unprestigious types of work are good per se. But when you
find yourself drawn to some kind of work despite its current lack of
prestige, it's a sign both that there's something real to be discovered
there, and that you have the right kind of motives. Impure motives are a
big danger for the ambitious. If anything is going to lead you astray,
it will be the desire to impress people. So while working on things that
aren't prestigious doesn't guarantee you're on the right track, it at
least guarantees you're not on the most common type of wrong one.

Over the next several years I wrote lots of essays about all kinds of
different topics. O'Reilly reprinted a collection of them as a book,
called \emph{Hackers \& Painters} after one of the essays in it. I also
worked on spam filters, and did some more painting. I used to have
dinners for a group of friends every thursday night, which taught me how
to cook for groups. And I bought another building in Cambridge, a former
candy factory (and later, twas said, porn studio), to use as an office.

One night in October 2003 there was a big party at my house. It was a
clever idea of my friend Maria Daniels, who was one of the thursday
diners. Three separate hosts would all invite their friends to one
party. So for every guest, two thirds of the other guests would be
people they didn't know but would probably like. One of the guests was
someone I didn't know but would turn out to like a lot: a woman called
Jessica Livingston. A couple days later I asked her out.

Jessica was in charge of marketing at a Boston investment bank. This
bank thought it understood startups, but over the next year, as she met
friends of mine from the startup world, she was surprised how different
reality was. And how colorful their stories were. So she decided to
compile a book of
\href{https://www.amazon.com/Founders-Work-Stories-Startups-Early/dp/1430210788}{interviews}
with startup founders.

When the bank had financial problems and she had to fire half her staff,
she started looking for a new job. In early 2005 she interviewed for a
marketing job at a Boston VC firm. It took them weeks to make up their
minds, and during this time I started telling her about all the things
that needed to be fixed about venture capital. They should make a larger
number of smaller investments instead of a handful of giant ones, they
should be funding younger, more technical founders instead of MBAs, they
should let the founders remain as CEO, and so on.

One of my tricks for writing essays had always been to give talks. The
prospect of having to stand up in front of a group of people and tell
them something that won't waste their time is a great spur to the
imagination. When the Harvard Computer Society, the undergrad computer
club, asked me to give a talk, I decided I would tell them how to start
a startup. Maybe they'd be able to avoid the worst of the mistakes we'd
made.

So I gave this talk, in the course of which I told them that the best
sources of seed funding were successful startup founders, because then
they'd be sources of advice too. Whereupon it seemed they were all
looking expectantly at me. Horrified at the prospect of having my inbox
flooded by business plans (if I'd only known), I blurted out ``But not
me!'' and went on with the talk. But afterward it occurred to me that I
should really stop procrastinating about angel investing. I'd been
meaning to since Yahoo bought us, and now it was 7 years later and I
still hadn't done one angel investment.

Meanwhile I had been scheming with Robert and Trevor about projects we
could work on together. I missed working with them, and it seemed like
there had to be something we could collaborate on.

As Jessica and I were walking home from dinner on March 11, at the
corner of Garden and Walker streets, these three threads converged.
Screw the VCs who were taking so long to make up their minds. We'd start
our own investment firm and actually implement the ideas we'd been
talking about. I'd fund it, and Jessica could quit her job and work for
it, and we'd get Robert and Trevor as partners too.
\footnote{Y Combinator was not the original name. At first we were called
Cambridge Seed. But we didn't want a regional name, in case someone
copied us in Silicon Valley, so we renamed ourselves after one of the
coolest tricks in the lambda calculus, the Y combinator.

I picked orange as our color partly because it's the warmest, and partly
because no VC used it. In 2005 all the VCs used staid colors like
maroon, navy blue, and forest green, because they were trying to appeal
to LPs, not founders. The YC logo itself is an inside joke: the Viaweb
logo had been a white V on a red circle, so I made the YC logo a white Y
on an orange square.}

Once again, ignorance worked in our favor. We had no idea how to be
angel investors, and in Boston in 2005 there were no Ron Conways to
learn from. So we just made what seemed like the obvious choices, and
some of the things we did turned out to be novel.

There are multiple components to Y Combinator, and we didn't figure them
all out at once. The part we got first was to be an angel firm. In those
days, those two words didn't go together. There were VC firms, which
were organized companies with people whose job it was to make
investments, but they only did big, million dollar investments. And
there were angels, who did smaller investments, but these were
individuals who were usually focused on other things and made
investments on the side. And neither of them helped founders enough in
the beginning. We knew how helpless founders were in some respects,
because we remembered how helpless we'd been. For example, one thing
Julian had done for us that seemed to us like magic was to get us set up
as a company. We were fine writing fairly difficult software, but
actually getting incorporated, with bylaws and stock and all that stuff,
how on earth did you do that? Our plan was not only to make seed
investments, but to do for startups everything Julian had done for us.

YC was not organized as a fund. It was cheap enough to run that we
funded it with our own money. That went right by 99\% of readers, but
professional investors are thinking ``Wow, that means they got all the
returns.'' But once again, this was not due to any particular insight on
our part. We didn't know how VC firms were organized. It never occurred
to us to try to raise a fund, and if it had, we wouldn't have known
where to start. \footnote{YC did become a fund for a couple years starting in 2009,
because it was getting so big I could no longer afford to fund it
personally. But after Heroku got bought we had enough money to go back
to being self-funded.}

The most distinctive thing about YC is the batch model: to fund a bunch
of startups all at once, twice a year, and then to spend three months
focusing intensively on trying to help them. That part we discovered by
accident, not merely implicitly but explicitly due to our ignorance
about investing. We needed to get experience as investors. What better
way, we thought, than to fund a whole bunch of startups at once? We knew
undergrads got temporary jobs at tech companies during the summer. Why
not organize a summer program where they'd start startups instead? We
wouldn't feel guilty for being in a sense fake investors, because they
would in a similar sense be fake founders. So while we probably wouldn't
make much money out of it, we'd at least get to practice being investors
on them, and they for their part would probably have a more interesting
summer than they would working at Microsoft.

We'd use the building I owned in Cambridge as our headquarters. We'd all
have dinner there once a week --- on tuesdays, since I was already
cooking for the thursday diners on thursdays --- and after dinner we'd
bring in experts on startups to give talks.

We knew undergrads were deciding then about summer jobs, so in a matter
of days we cooked up something we called the Summer Founders Program,
and I posted an \href{summerfounder.html}{announcement} on my site,
inviting undergrads to apply. I had never imagined that writing essays
would be a way to get ``deal flow,'' as investors call it, but it turned
out to be the perfect source. \footnote{I've never liked the term ``deal flow,'' because it implies
that the number of new startups at any given time is fixed. This is not
only false, but it's the purpose of YC to falsify it, by causing
startups to be founded that would not otherwise have existed.} We got 225
applications for the Summer Founders Program, and we were surprised to
find that a lot of them were from people who'd already graduated, or
were about to that spring. Already this SFP thing was starting to feel
more serious than we'd intended.

We invited about 20 of the 225 groups to interview in person, and from
those we picked 8 to fund. They were an impressive group. That first
batch included reddit, Justin Kan and Emmett Shear, who went on to found
Twitch, Aaron Swartz, who had already helped write the RSS spec and
would a few years later become a martyr for open access, and Sam Altman,
who would later become the second president of YC. I don't think it was
entirely luck that the first batch was so good. You had to be pretty
bold to sign up for a weird thing like the Summer Founders Program
instead of a summer job at a legit place like Microsoft or Goldman
Sachs.

The deal for startups was based on a combination of the deal we did with
Julian (\$10k for 10\%) and what Robert said MIT grad students got for
the summer (\$6k). We invested \$6k per founder, which in the typical
two-founder case was \$12k, in return for 6\%. That had to be fair,
because it was twice as good as the deal we ourselves had taken. Plus
that first summer, which was really hot, Jessica brought the founders
free air conditioners. \footnote{She reports that they were all different shapes and sizes,
because there was a run on air conditioners and she had to get whatever
she could, but that they were all heavier than she could carry now.}

Fairly quickly I realized that we had stumbled upon the way to scale
startup funding. Funding startups in batches was more convenient for us,
because it meant we could do things for a lot of startups at once, but
being part of a batch was better for the startups too. It solved one of
the biggest problems faced by founders: the isolation. Now you not only
had colleagues, but colleagues who understood the problems you were
facing and could tell you how they were solving them.

As YC grew, we started to notice other advantages of scale. The alumni
became a tight community, dedicated to helping one another, and
especially the current batch, whose shoes they remembered being in. We
also noticed that the startups were becoming one another's customers. We
used to refer jokingly to the ``YC GDP,'' but as YC grows this becomes
less and less of a joke. Now lots of startups get their initial set of
customers almost entirely from among their batchmates.

I had not originally intended YC to be a full-time job. I was going to
do three things: hack, write essays, and work on YC. As YC grew, and I
grew more excited about it, it started to take up a lot more than a
third of my attention. But for the first few years I was still able to
work on other things.

In the summer of 2006, Robert and I started working on a new version of
Arc. This one was reasonably fast, because it was compiled into Scheme.
To test this new Arc, I wrote Hacker News in it. It was originally meant
to be a news aggregator for startup founders and was called Startup
News, but after a few months I got tired of reading about nothing but
startups. Plus it wasn't startup founders we wanted to reach. It was
future startup founders. So I changed the name to Hacker News and the
topic to whatever engaged one's intellectual curiosity.

HN was no doubt good for YC, but it was also by far the biggest source
of stress for me. If all I'd had to do was select and help founders,
life would have been so easy. And that implies that HN was a mistake.
Surely the biggest source of stress in one's work should at least be
something close to the core of the work. Whereas I was like someone who
was in pain while running a marathon not from the exertion of running,
but because I had a blister from an ill-fitting shoe. When I was dealing
with some urgent problem during YC, there was about a 60\% chance it had
to do with HN, and a 40\% chance it had do with everything else
combined. \footnote{Another problem with HN was a bizarre edge case that occurs
when you both write essays and run a forum. When you run a forum, you're
assumed to see if not every conversation, at least every conversation
involving you. And when you write essays, people post highly imaginative
misinterpretations of them on forums. Individually these two phenomena
are tedious but bearable, but the combination is disastrous. You
actually have to respond to the misinterpretations, because the
assumption that you're present in the conversation means that not
responding to any sufficiently upvoted misinterpretation reads as a
tacit admission that it's correct. But that in turn encourages more;
anyone who wants to pick a fight with you senses that now is their
chance.}

As well as HN, I wrote all of YC's internal software in Arc. But while I
continued to work a good deal \emph{in} Arc, I gradually stopped working
\emph{on} Arc, partly because I didn't have time to, and partly because
it was a lot less attractive to mess around with the language now that
we had all this infrastructure depending on it. So now my three projects
were reduced to two: writing essays and working on YC.

YC was different from other kinds of work I've done. Instead of deciding
for myself what to work on, the problems came to me. Every 6 months
there was a new batch of startups, and their problems, whatever they
were, became our problems. It was very engaging work, because their
problems were quite varied, and the good founders were very effective.
If you were trying to learn the most you could about startups in the
shortest possible time, you couldn't have picked a better way to do it.

There were parts of the job I didn't like. Disputes between cofounders,
figuring out when people were lying to us, fighting with people who
maltreated the startups, and so on. But I worked hard even at the parts
I didn't like. I was haunted by something Kevin Hale once said about
companies: ``No one works harder than the boss.'' He meant it both
descriptively and prescriptively, and it was the second part that scared
me. I wanted YC to be good, so if how hard I worked set the upper bound
on how hard everyone else worked, I'd better work very hard.

One day in 2010, when he was visiting California for interviews, Robert
Morris did something astonishing: he offered me unsolicited advice. I
can only remember him doing that once before. One day at Viaweb, when I
was bent over double from a kidney stone, he suggested that it would be
a good idea for him to take me to the hospital. That was what it took
for Rtm to offer unsolicited advice. So I remember his exact words very
clearly. ``You know,'' he said, ``you should make sure Y Combinator
isn't the last cool thing you do.''

At the time I didn't understand what he meant, but gradually it dawned
on me that he was saying I should quit. This seemed strange advice,
because YC was doing great. But if there was one thing rarer than Rtm
offering advice, it was Rtm being wrong. So this set me thinking. It was
true that on my current trajectory, YC would be the last thing I did,
because it was only taking up more of my attention. It had already eaten
Arc, and was in the process of eating essays too. Either YC was my
life's work or I'd have to leave eventually. And it wasn't, so I would.

In the summer of 2012 my mother had a stroke, and the cause turned out
to be a blood clot caused by colon cancer. The stroke destroyed her
balance, and she was put in a nursing home, but she really wanted to get
out of it and back to her house, and my sister and I were determined to
help her do it. I used to fly up to Oregon to visit her regularly, and I
had a lot of time to think on those flights. On one of them I realized I
was ready to hand YC over to someone else.

I asked Jessica if she wanted to be president, but she didn't, so we
decided we'd try to recruit Sam Altman. We talked to Robert and Trevor
and we agreed to make it a complete changing of the guard. Up till that
point YC had been controlled by the original LLC we four had started.
But we wanted YC to last for a long time, and to do that it couldn't be
controlled by the founders. So if Sam said yes, we'd let him reorganize
YC. Robert and I would retire, and Jessica and Trevor would become
ordinary partners.

When we asked Sam if he wanted to be president of YC, initially he said
no. He wanted to start a startup to make nuclear reactors. But I kept at
it, and in October 2013 he finally agreed. We decided he'd take over
starting with the winter 2014 batch. For the rest of 2013 I left running
YC more and more to Sam, partly so he could learn the job, and partly
because I was focused on my mother, whose cancer had returned.

She died on January 15, 2014. We knew this was coming, but it was still
hard when it did.

I kept working on YC till March, to help get that batch of startups
through Demo Day, then I checked out pretty completely. (I still talk to
alumni and to new startups working on things I'm interested in, but that
only takes a few hours a week.)

What should I do next? Rtm's advice hadn't included anything about that.
I wanted to do something completely different, so I decided I'd paint. I
wanted to see how good I could get if I really focused on it. So the day
after I stopped working on YC, I started painting. I was rusty and it
took a while to get back into shape, but it was at least completely
engaging. \footnote{The worst thing about leaving YC was not working with Jessica
anymore. We'd been working on YC almost the whole time we'd known each
other, and we'd neither tried nor wanted to separate it from our
personal lives, so leaving was like pulling up a deeply rooted tree.}

I spent most of the rest of 2014 painting. I'd never been able to work
so uninterruptedly before, and I got to be better than I had been. Not
good enough, but better. Then in November, right in the middle of a
painting, I ran out of steam. Up till that point I'd always been curious
to see how the painting I was working on would turn out, but suddenly
finishing this one seemed like a chore. So I stopped working on it and
cleaned my brushes and haven't painted since. So far anyway.

I realize that sounds rather wimpy. But attention is a zero sum game. If
you can choose what to work on, and you choose a project that's not the
best one (or at least a good one) for you, then it's getting in the way
of another project that is. And at 50 there was some opportunity cost to
screwing around.

I started writing essays again, and wrote a bunch of new ones over the
next few months. I even wrote a couple that \href{know.html}{weren't}
about startups. Then in March 2015 I started working on Lisp again.

The distinctive thing about Lisp is that its core is a language defined
by writing an interpreter in itself. It wasn't originally intended as a
programming language in the ordinary sense. It was meant to be a formal
model of computation, an alternative to the Turing machine. If you want
to write an interpreter for a language in itself, what's the minimum set
of predefined operators you need? The Lisp that John McCarthy invented,
or more accurately discovered, is an answer to that question.
\footnote{One way to get more precise about the concept of invented vs
discovered is to talk about space aliens. Any sufficiently advanced
alien civilization would certainly know about the Pythagorean theorem,
for example. I believe, though with less certainty, that they would also
know about the Lisp in McCarthy's 1960 paper.

But if so there's no reason to suppose that this is the limit of the
language that might be known to them. Presumably aliens need numbers and
errors and I/O too. So it seems likely there exists at least one path
out of McCarthy's Lisp along which discoveredness is preserved.}

McCarthy didn't realize this Lisp could even be used to program
computers till his grad student Steve Russell suggested it. Russell
translated McCarthy's interpreter into IBM 704 machine language, and
from that point Lisp started also to be a programming language in the
ordinary sense. But its origins as a model of computation gave it a
power and elegance that other languages couldn't match. It was this that
attracted me in college, though I didn't understand why at the time.

McCarthy's 1960 Lisp did nothing more than interpret Lisp expressions.
It was missing a lot of things you'd want in a programming language. So
these had to be added, and when they were, they weren't defined using
McCarthy's original axiomatic approach. That wouldn't have been feasible
at the time. McCarthy tested his interpreter by hand-simulating the
execution of programs. But it was already getting close to the limit of
interpreters you could test that way --- indeed, there was a bug in it
that McCarthy had overlooked. To test a more complicated interpreter,
you'd have had to run it, and computers then weren't powerful enough.

Now they are, though. Now you could continue using McCarthy's axiomatic
approach till you'd defined a complete programming language. And as long
as every change you made to McCarthy's Lisp was a
discoveredness-preserving transformation, you could, in principle, end
up with a complete language that had this quality. Harder to do than to
talk about, of course, but if it was possible in principle, why not try?
So I decided to take a shot at it. It took 4 years, from March 26, 2015
to October 12, 2019. It was fortunate that I had a precisely defined
goal, or it would have been hard to keep at it for so long.

I wrote this new Lisp, called \href{bel.html}{Bel}, in itself in Arc.
That may sound like a contradiction, but it's an indication of the sort
of trickery I had to engage in to make this work. By means of an
egregious collection of hacks I managed to make something close enough
to an interpreter written in itself that could actually run. Not fast,
but fast enough to test.

I had to ban myself from writing essays during most of this time, or I'd
never have finished. In late 2015 I spent 3 months writing essays, and
when I went back to working on Bel I could barely understand the code.
Not so much because it was badly written as because the problem is so
convoluted. When you're working on an interpreter written in itself,
it's hard to keep track of what's happening at what level, and errors
can be practically encrypted by the time you get them.

So I said no more essays till Bel was done. But I told few people about
Bel while I was working on it. So for years it must have seemed that I
was doing nothing, when in fact I was working harder than I'd ever
worked on anything. Occasionally after wrestling for hours with some
gruesome bug I'd check Twitter or HN and see someone asking ``Does Paul
Graham still code?''

Working on Bel was hard but satisfying. I worked on it so intensively
that at any given time I had a decent chunk of the code in my head and
could write more there. I remember taking the boys to the coast on a
sunny day in 2015 and figuring out how to deal with some problem
involving continuations while I watched them play in the tide pools. It
felt like I was doing life right. I remember that because I was slightly
dismayed at how novel it felt. The good news is that I had more moments
like this over the next few years.

In the summer of 2016 we moved to England. We wanted our kids to see
what it was like living in another country, and since I was a British
citizen by birth, that seemed the obvious choice. We only meant to stay
for a year, but we liked it so much that we still live there. So most of
Bel was written in England.

In the fall of 2019, Bel was finally finished. Like McCarthy's original
Lisp, it's a spec rather than an implementation, although like
McCarthy's Lisp it's a spec expressed as code.

Now that I could write essays again, I wrote a bunch about topics I'd
had stacked up. I kept writing essays through 2020, but I also started
to think about other things I could work on. How should I choose what to
do? Well, how had I chosen what to work on in the past? I wrote an essay
for myself to answer that question, and I was surprised how long and
messy the answer turned out to be. If this surprised me, who'd lived it,
then I thought perhaps it would be interesting to other people, and
encouraging to those with similarly messy lives. So I wrote a more
detailed version for others to read, and this is the last sentence of
it.


\textbf{Thanks} to Trevor Blackwell, John Collison, Patrick Collison,
Daniel Gackle, Ralph Hazell, Jessica Livingston, Robert Morris, and Harj
Taggar for reading drafts of this.
\chapter{Donate Unrestricted}

March 2021

The secret curse of the nonprofit world is restricted donations. If you
haven't been involved with nonprofits, you may never have heard this
phrase before. But if you have been, it probably made you wince.

Restricted donations mean donations where the donor limits what can be
done with the money. This is common with big donations, perhaps the
default. And yet it's usually a bad idea. Usually the way the donor
wants the money spent is not the way the nonprofit would have chosen.
Otherwise there would have been no need to restrict the donation. But
who has a better understanding of where money needs to be spent, the
nonprofit or the donor?

If a nonprofit doesn't understand better than its donors where money
needs to be spent, then it's incompetent and you shouldn't be donating
to it at all.

Which means a restricted donation is inherently suboptimal. It's either
a donation to a bad nonprofit, or a donation for the wrong things.

There are a couple exceptions to this principle. One is when the
nonprofit is an umbrella organization. It's reasonable to make a
restricted donation to a university, for example, because a university
is only nominally a single nonprofit. Another exception is when the
donor actually does know as much as the nonprofit about where money
needs to be spent. The Gates Foundation, for example, has specific goals
and often makes restricted donations to individual nonprofits to
accomplish them. But unless you're a domain expert yourself or donating
to an umbrella organization, your donation would do more good if it were
unrestricted.

If restricted donations do less good than unrestricted ones, why do
donors so often make them? Partly because doing good isn't donors' only
motive. They often have other motives as well --- to make a mark, or to
generate good publicity \footnote{Unfortunately restricted donations tend to generate more
publicity than unrestricted ones. ``X donates money to build a school in
Africa'' is not only more interesting than ``X donates money to Y
nonprofit to spend as Y chooses,'' but also focuses more attention on X.}, or to comply with
regulations or corporate policies. Many donors may simply never have
considered the distinction between restricted and unrestricted
donations. They may believe that donating money for some specific
purpose is just how donation works. And to be fair, nonprofits don't try
very hard to discourage such illusions. They can't afford to. People
running nonprofits are almost always anxious about money. They can't
afford to talk back to big donors.

You can't expect candor in a relationship so asymmetric. So I'll tell
you what nonprofits wish they could tell you. If you want to donate to a
nonprofit, donate unrestricted. If you trust them to spend your money,
trust them to decide how.

\textbf{Note}

\textbf{Thanks} to Chase Adam, Ingrid Bassett, Trevor Blackwell, and
Edith Elliot for reading drafts of this.
\chapter{Write Simply}

March 2021

I try to write using ordinary words and simple sentences.

That kind of writing is easier to read, and the easier something is to
read, the more deeply readers will engage with it. The less energy they
expend on your prose, the more they'll have left for your ideas.

And the further they'll read. Most readers' energy tends to flag part
way through an article or essay. If the friction of reading is low
enough, more keep going till the end.

There's an Italian dish called \emph{saltimbocca}, which means ``leap
into the mouth.'' My goal when writing might be called
\emph{saltintesta}: the ideas leap into your head and you barely notice
the words that got them there.

It's too much to hope that writing could ever be pure ideas. You might
not even want it to be. But for most writers, most of the time, that's
the goal to aim for. The gap between most writing and pure ideas is not
filled with poetry.

Plus it's more considerate to write simply. When you write in a fancy
way to impress people, you're making them do extra work just so you can
seem cool. It's like trailing a long train behind you that readers have
to carry.

And remember, if you're writing in English, that a lot of your readers
won't be native English speakers. Their understanding of ideas may be
way ahead of their understanding of English. So you can't assume that
writing about a difficult topic means you can use difficult words.

Of course, fancy writing doesn't just conceal ideas. It can also conceal
the lack of them. That's why some people write that way, to conceal the
fact that they have
\href{https://scholar.google.com/scholar?hl=en&as_sdt=0\%2C5&q=hermeneutic+dialectics+hegemonic+modalities}{}nothing
to say. Whereas writing simply keeps you honest. If you say nothing
simply, it will be obvious to everyone, including you.

Simple writing also lasts better. People reading your stuff in the
future will be in much the same position as people from other countries
reading it today. The culture and the language will have changed. It's
not vain to care about that, any more than it's vain for a woodworker to
build a chair to last.

Indeed, lasting is not merely an accidental quality of chairs, or
writing. It's a sign you did a good job.

But although these are all real advantages of writing simply, none of
them are why I do it. The main reason I write simply is that it offends
me not to. When I write a sentence that seems too complicated, or that
uses unnecessarily intellectual words, it doesn't seem fancy to me. It
seems clumsy.

There are of course times when you want to use a complicated sentence or
fancy word for effect. But you should never do it by accident.

The other reason my writing ends up being simple is the way I do it. I
write the first draft fast, then spend days editing it, trying to get
everything just right. Much of this editing is cutting, and that makes
simple writing even simpler.
\chapter{How People Get Rich Now}

April 2021

Every year since 1982, \emph{Forbes} magazine has published a list of
the richest Americans. If we compare the 100 richest people in 1982 to
the 100 richest in 2020, we notice some big differences.

In 1982 the most common source of wealth was inheritance. Of the 100
richest people, 60 inherited from an ancestor. There were 10 du Pont
heirs alone. By 2020 the number of heirs had been cut in half,
accounting for only 27 of the biggest 100 fortunes.

Why would the percentage of heirs decrease? Not because inheritance
taxes increased. In fact, they decreased significantly during this
period. The reason the percentage of heirs has decreased is not that
fewer people are inheriting great fortunes, but that more people are
making them.

How are people making these new fortunes? Roughly 3/4 by starting
companies and 1/4 by investing. Of the 73 new fortunes in 2020, 56
derive from founders' or early employees' equity (52 founders, 2 early
employees, and 2 wives of founders), and 17 from managing investment
funds.

There were no fund managers among the 100 richest Americans in 1982.
Hedge funds and private equity firms existed in 1982, but none of their
founders were rich enough yet to make it into the top 100. Two things
changed: fund managers discovered new ways to generate high returns, and
more investors were willing to trust them with their money.
\footnote{Investment firms grew rapidly after a regulatory change by the
Labor Department in 1978 allowed pension funds to invest in them, but
the effects of this growth were not yet visible in the top 100 fortunes
in 1982.}

But the main source of new fortunes now is starting companies, and when
you look at the data, you see big changes there too. People get richer
from starting companies now than they did in 1982, because the companies
do different things.

In 1982, there were two dominant sources of new wealth: oil and real
estate. Of the 40 new fortunes in 1982, at least 24 were due primarily
to oil or real estate. Now only a small number are: of the 73 new
fortunes in 2020, 4 were due to real estate and only 2 to oil.

By 2020 the biggest source of new wealth was what are sometimes called
``tech'' companies. Of the 73 new fortunes, about 30 derive from such
companies. These are particularly common among the richest of the rich:
8 of the top 10 fortunes in 2020 were new fortunes of this type.

Arguably it's slightly misleading to treat tech as a category. Isn't
Amazon really a retailer, and Tesla a car maker? Yes and no. Maybe in 50
years, when what we call tech is taken for granted, it won't seem right
to put these two businesses in the same category. But at the moment at
least, there is definitely something they share in common that
distinguishes them. What retailer starts AWS? What car maker is run by
someone who also has a rocket company?

The tech companies behind the top 100 fortunes also form a
well-differentiated group in the sense that they're all companies that
venture capitalists would readily invest in, and the others mostly not.
And there's a reason why: these are mostly companies that win by having
better technology, rather than just a CEO who's really driven and good
at making deals.

To that extent, the rise of the tech companies represents a qualitative
change. The oil and real estate magnates of the 1982 Forbes 400 didn't
win by making better technology. They won by being really driven and
good at making deals. \footnote{George Mitchell deserves mention as an exception. Though really
driven and good at making deals, he was also the first to figure out how
to use fracking to get natural gas out of shale.} And indeed, that way of
getting rich is so old that it predates the Industrial Revolution. The
courtiers who got rich in the (nominal) service of European royal houses
in the 16th and 17th centuries were also, as a rule, really driven and
good at making deals.

People who don't look any deeper than the Gini coefficient look back on
the world of 1982 as the good old days, because those who got rich then
didn't get as rich. But if you dig into \emph{how} they got rich, the
old days don't look so good. In 1982, 84\% of the richest 100 people got
rich by inheritance, extracting natural resources, or doing real estate
deals. Is that really better than a world in which the richest people
get rich by starting tech companies?

Why are people starting so many more new companies than they used to,
and why are they getting so rich from it? The answer to the first
question, curiously enough, is that it's misphrased. We shouldn't be
asking why people are starting companies, but why they're starting
companies \emph{again}. \footnote{When I say people are starting more companies, I mean the type
of company meant to \href{growth.html}{grow} very big. There has
actually been a decrease in the last couple decades in the overall
number of new companies. But the vast majority of companies are small
retail and service businesses. So what the statistics about the
decreasing number of new businesses mean is that people are starting
fewer shoe stores and barber shops.

People sometimes get
\href{https://www.inc.com/magazine/201505/leigh-buchanan/the-vanishing-startups-in-decline.html}{confused}
when they see a graph labelled ``startups'' that's going down, because
there are two senses of the word ``startup'': (1) the founding of a
company, and (2) a particular type of company designed to grow big fast.
The statistics mean startup in sense (1), not sense (2).}

In 1892, the \emph{New York Herald Tribune} compiled a list of all the
millionaires in America. They found 4047 of them. How many had inherited
their wealth then? Only about 20\%, which is less than the proportion of
heirs today. And when you investigate the sources of the new fortunes,
1892 looks even more like today. Hugh Rockoff found that ``many of the
richest \ldots{} gained their initial edge from the new technology of
mass production.'' \footnote{Rockoff, Hugh. ``Great Fortunes of the Gilded Age.'' NBER
Working Paper 14555, 2008.}

So it's not 2020 that's the anomaly here, but 1982. The real question is
why so few people had gotten rich from starting companies in 1982. And
the answer is that even as the \emph{Herald Tribune}'s list was being
compiled, a wave of \href{re.html}{consolidation} was sweeping through
the American economy. In the late 19th and early 20th centuries,
financiers like J. P. Morgan combined thousands of smaller companies
into a few hundred giant ones with commanding economies of scale. By the
end of World War II, as Michael Lind writes, ``the major sectors of the
economy were either organized as government-backed cartels or dominated
by a few oligopolistic corporations.'' \footnote{Lind, Michael. \emph{Land of Promise.} HarperCollins, 2012.

It's also likely that the high tax rates in the mid 20th century
deterred people from starting their own companies. Starting one's own
company is risky, and when risk isn't rewarded, people opt for
\href{inequality.html}{safety} instead.

But it wasn't simply cause and effect. The oligopolies and high tax
rates of the mid 20th century were all of a piece. Lower taxes are not
just a cause of entrepreneurship, but an effect as well: the people
getting rich in the mid 20th century from real estate and oil
exploration lobbied for and got huge tax loopholes that made their
effective tax rate much lower, and presumably if it had been more common
to grow big companies by building new technology, the people doing that
would have lobbied for their own loopholes as well.}

In 1960, most of the people who start startups today would have gone to
work for one of them. You could get rich from starting your own company
in 1890 and in 2020, but in 1960 it was not really a viable option. You
couldn't break through the oligopolies to get at the markets. So the
prestigious route in 1960 was not to start your own company, but to work
your way up the corporate ladder at an existing one.
\footnote{That's why the people who did get rich in the mid 20th century
so often got rich from oil exploration or real estate. Those were the
two big areas of the economy that weren't susceptible to consolidation.}

Making everyone a corporate employee decreased economic inequality (and
every other kind of variation), but if your model of normal is the mid
20th century, you have a very misleading model in that respect. J. P.
Morgan's economy turned out to be just a phase, and starting in the
1970s, it began to break up.

Why did it break up? Partly senescence. The big companies that seemed
models of scale and efficiency in 1930 had by 1970 become slack and
bloated. By 1970 the rigid structure of the economy was full of cosy
nests that various groups had built to insulate themselves from market
forces. During the Carter administration the federal government realized
something was amiss and began, in a process they called
``deregulation,'' to roll back the policies that propped up the
oligopolies.

But it wasn't just decay from within that broke up J. P. Morgan's
economy. There was also pressure from without, in the form of new
technology, and particularly microelectronics. The best way to envision
what happened is to imagine a pond with a crust of ice on top. Initially
the only way from the bottom to the surface is around the edges. But as
the ice crust weakens, you start to be able to punch right through the
middle.

The edges of the pond were pure tech: companies that actually described
themselves as being in the electronics or software business. When you
used the word ``startup'' in 1990, that was what you meant. But now
startups are punching right through the middle of the ice crust and
displacing incumbents like retailers and TV networks and car companies.
\footnote{The pure tech companies used to be called ``high technology''
startups. But now that startups can punch through the middle of the ice
crust, we don't need a separate name for the edges, and the term
``high-tech'' has a decidedly
\href{https://books.google.com/ngrams/graph?content=high+tech&year_start=1900&year_end=2019&corpus=en-2019&smoothing=3}{retro}
sound.}

But though the breakup of J. P. Morgan's economy created a new world in
the technological sense, it was a reversion to the norm in the social
sense. If you only look back as far as the mid 20th century, it seems
like people getting rich by starting their own companies is a recent
phenomenon. But if you look back further, you realize it's actually the
default. So what we should expect in the future is more of the same.
Indeed, we should expect both the number and wealth of founders to grow,
because every decade it gets easier to start a startup.

Part of the reason it's getting easier to start a startup is social.
Society is (re)assimilating the concept. If you start one now, your
parents won't freak out the way they would have a generation ago, and
knowledge about how to do it is much more widespread. But the main
reason it's easier to start a startup now is that it's cheaper.
Technology has driven down the cost of both building products and
acquiring customers.

The decreasing cost of starting a startup has in turn changed the
balance of power between founders and investors. Back when starting a
startup meant building a factory, you needed investors' permission to do
it at all. But now investors need founders more than founders need
investors, and that, combined with the increasing amount of venture
capital available, has driven up valuations. \footnote{Higher valuations mean you either sell less stock to get a given
amount of money, or get more money for a given amount of stock. The
typical startup does some of each. Obviously you end up richer if you
keep more stock, but you should also end up richer if you raise more
money, because (a) it should make the company more successful, and (b)
you should be able to last longer before the next round, or not even
need one. Notice all those shoulds though. In practice a lot of money
slips through them.

It might seem that the huge rounds raised by startups nowadays
contradict the claim that it has become cheaper to start one. But
there's no contradiction here; the startups that raise the most are the
ones doing it by choice, in order to grow faster, not the ones doing it
because they need the money to survive. There's nothing like not needing
money to make people offer it to you.

You would think, after having been on the side of labor in its fight
with capital for almost two centuries, that the far left would be happy
that labor has finally prevailed. But none of them seem to be. You can
almost hear them saying ``No, no, not \emph{that} way.''}

So the decreasing cost of starting a startup increases the number of
rich people in two ways: it means that more people start them, and that
those who do can raise money on better terms.

But there's also a third factor at work: the companies themselves are
more valuable, because newly founded companies grow faster than they
used to. Technology hasn't just made it cheaper to build and distribute
things, but faster too.

This trend has been running for a long time. IBM, founded in 1896, took
45 years to reach a billion 2020 dollars in revenue. Hewlett-Packard,
founded in 1939, took 25 years. Microsoft, founded in 1975, took 13
years. Now the norm for fast-growing companies is 7 or 8 years.
\footnote{IBM was created in 1911 by merging three companies, the most
important of which was Herman Hollerith's Tabulating Machine Company,
founded in 1896. In 1941 its revenues were \$60 million.

Hewlett-Packard's revenues in 1964 were \$125 million.

Microsoft's revenues in 1988 were \$590 million.}

Fast growth has a double effect on the value of founders' stock. The
value of a company is a function of its revenue and its growth rate. So
if a company grows faster, you not only get to a billion dollars in
revenue sooner, but the company is more valuable when it reaches that
point than it would be if it were growing slower.

That's why founders sometimes get so rich so young now. The low initial
cost of starting a startup means founders can start young, and the fast
growth of companies today means that if they succeed they could be
surprisingly rich just a few years later.

It's easier now to start and grow a company than it has ever been. That
means more people start them, that those who do get better terms from
investors, and that the resulting companies become more valuable. Once
you understand how these mechanisms work, and that startups were
suppressed for most of the 20th century, you don't have to resort to
some vague right turn the country took under Reagan to explain why
America's Gini coefficient is increasing. Of course the Gini coefficient
is increasing. With more people starting more valuable companies, how
could it not be?


\textbf{Thanks} to Trevor Blackwell, Jessica Livingston, Bob Lesko,
Robert Morris, Russ Roberts, and Alex Tabarrok for reading drafts of
this, and to Jon Erlichman for growth data.
\chapter{The Real Reason to End the Death Penalty}

April 2021

When intellectuals talk about the death penalty, they talk about things
like whether it's permissible for the state to take someone's life,
whether the death penalty acts as a deterrent, and whether more death
sentences are given to some groups than others. But in practice the
debate about the death penalty is not about whether it's ok to kill
murderers. It's about whether it's ok to kill innocent people, because
at least 4\% of people on death row are
\href{https://www.pnas.org/content/111/20/7230}{innocent}.

When I was a kid I imagined that it was unusual for people to be
convicted of crimes they hadn't committed, and that in murder cases
especially this must be very rare. Far from it. Now, thanks to
organizations like the
\href{https://innocenceproject.org/all-cases}{Innocence Project}, we see
a constant stream of stories about murder convictions being overturned
after new evidence emerges. Sometimes the police and prosecutors were
just very sloppy. Sometimes they were crooked, and knew full well they
were convicting an innocent person.

Kenneth Adams and three other men spent 18 years in prison on a murder
conviction. They were exonerated after DNA testing implicated three
different men, two of whom later confessed. The police had been told
about the other men early in the investigation, but never followed up
the lead.

Keith Harward spent 33 years in prison on a murder conviction. He was
convicted because ``experts'' said his teeth matched photos of bite
marks on one victim. He was exonerated after DNA testing showed the
murder had been committed by another man, Jerry Crotty.

Ricky Jackson and two other men spent 39 years in prison after being
convicted of murder on the testimony of a 12 year old boy, who later
recanted and said he'd been coerced by police. Multiple people have
confirmed the boy was elsewhere at the time. The three men were
exonerated after the county prosecutor dropped the charges, saying ``The
state is conceding the obvious.''

Alfred Brown spent 12 years in prison on a murder conviction, including
10 years on death row. He was exonerated after it was discovered that
the assistant district attorney had concealed phone records proving he
could not have committed the crimes.

Glenn Ford spent 29 years on death row after having been convicted of
murder. He was exonerated after new evidence proved he was not even at
the scene when the murder occurred. The attorneys assigned to represent
him had never tried a jury case before.

Cameron Willingham was actually executed in 2004 by lethal injection.
The ``expert'' who testified that he deliberately set fire to his house
has since been discredited. A re-examination of the case ordered by the
state of Texas in 2009 concluded that ``a finding of arson could not be
sustained.''

\href{https://saverichardglossip.com/facts}{Rich Glossip} has spent 20
years on death row after being convicted of murder on the testimony of
the actual killer, who escaped with a life sentence in return for
implicating him. In 2015 he came within minutes of execution before it
emerged that Oklahoma had been planning to kill him with an illegal
combination of drugs. They still plan to go ahead with the execution,
perhaps as soon as this summer, despite
\href{https://www.usnews.com/news/best-states/oklahoma/articles/2020-10-14/attorney-for-oklahoma-death-row-inmate-claims-new-evidence}{new
evidence} exonerating him.

I could go on. There are hundreds of similar cases. In Florida alone, 29
death row prisoners have been exonerated so far.

Far from being rare, wrongful murder convictions are
\href{https://deathpenaltyinfo.org/policy-issues/innocence/description-of-innocence-cases}{very
common}. Police are under pressure to solve a crime that has gotten a
lot of attention. When they find a suspect, they want to believe he's
guilty, and ignore or even destroy evidence suggesting otherwise.
District attorneys want to be seen as effective and tough on crime, and
in order to win convictions are willing to manipulate witnesses and
withhold evidence. Court-appointed defense attorneys are overworked and
often incompetent. There's a ready supply of criminals willing to give
false testimony in return for a lighter sentence, suggestible witnesses
who can be made to say whatever police want, and bogus ``experts'' eager
to claim that science proves the defendant is guilty. And juries want to
believe them, since otherwise some terrible crime remains unsolved.

This circus of incompetence and dishonesty is the real issue with the
death penalty. We don't even reach the point where theoretical questions
about the moral justification or effectiveness of capital punishment
start to matter, because so many of the people sentenced to death are
actually innocent. Whatever it means in theory, in practice capital
punishment means killing innocent people.

\textbf{Thanks} to Trevor Blackwell, Jessica Livingston, and Don Knight
for reading drafts of this.

\textbf{Related:}
\chapter{An NFT That Saves Lives}

May 2021

\href{https://www.noorahealth.org/}{Noora Health}, a nonprofit I've
supported for years, just launched a new NFT. It has a dramatic name,
\href{http://bit.ly/NooraNFT}{Save Thousands of Lives}, because that's
what the proceeds will do.

Noora has been saving lives for 7 years. They run programs in hospitals
in South Asia to teach new mothers how to take care of their babies once
they get home. They're in 165 hospitals now. And because they know the
numbers before and after they start at a new hospital, they can measure
the impact they have. It is massive. For every 1000 live births, they
save 9 babies.

This number comes from a \href{http://bit.ly/NFT-research}{study} of
133,733 families at 28 different hospitals that Noora conducted in
collaboration with the Better Birth team at Ariadne Labs, a joint center
for health systems innovation at Brigham and Women's Hospital and
Harvard T.H. Chan School of Public Health.

Noora is so effective that even if you measure their costs in the most
conservative way, by dividing their entire budget by the number of lives
saved, the cost of saving a life is the lowest I've seen. \$1,235.

For this NFT, they're going to issue a public report tracking how this
specific tranche of money is spent, and estimating the number of lives
saved as a result.

NFTs are a new territory, and this way of using them is especially new,
but I'm excited about its potential. And I'm excited to see what happens
with this particular auction, because unlike an NFT representing
something that has already happened, this NFT gets better as the price
gets higher.

The reserve price was about \$2.5 million, because that's what it takes
for the name to be accurate: that's what it costs to save 2000 lives.
But the higher the price of this NFT goes, the more lives will be saved.
What a sentence to be able to write.
\chapter{Crazy New Ideas}

May 2021

There's one kind of opinion I'd be very afraid to express publicly. If
someone I knew to be both a domain expert and a reasonable person
proposed an idea that sounded preposterous, I'd be very reluctant to say
``That will never work.''

Anyone who has studied the history of ideas, and especially the history
of science, knows that's how big things start. Someone proposes an idea
that sounds crazy, most people dismiss it, then it gradually takes over
the world.

Most implausible-sounding ideas are in fact bad and could be safely
dismissed. But not when they're proposed by reasonable domain experts.
If the person proposing the idea is reasonable, then they know how
implausible it sounds. And yet they're proposing it anyway. That
suggests they know something you don't. And if they have deep domain
expertise, that's probably the source of it. \footnote{This domain expertise could be in another field. Indeed, such
crossovers tend to be particularly promising.}

Such ideas are not merely unsafe to dismiss, but disproportionately
likely to be interesting. When the average person proposes an
implausible-sounding idea, its implausibility is evidence of their
incompetence. But when a reasonable domain expert does it, the situation
is reversed. There's something like an efficient market here: on average
the ideas that seem craziest will, if correct, have the biggest effect.
So if you can eliminate the theory that the person proposing an
implausible-sounding idea is incompetent, its implausibility switches
from evidence that it's boring to evidence that it's exciting.
\footnote{I'm not claiming this principle extends much beyond math,
engineering, and the hard sciences. In politics, for example,
crazy-sounding ideas generally are as bad as they sound. Though arguably
this is not an exception, because the people who propose them are not in
fact domain experts; politicians are domain experts in political
tactics, like how to get elected and how to get legislation passed, but
not in the world that policy acts upon. Perhaps no one could be.}

Such ideas are not guaranteed to work. But they don't have to be. They
just have to be sufficiently good bets --- to have sufficiently high
expected value. And I think on average they do. I think if you bet on
the entire set of implausible-sounding ideas proposed by reasonable
domain experts, you'd end up net ahead.

The reason is that everyone is too conservative. The word ``paradigm''
is overused, but this is a case where it's warranted. Everyone is too
much in the grip of the current paradigm. Even the people who have the
new ideas undervalue them initially. Which means that before they reach
the stage of proposing them publicly, they've already subjected them to
an excessively strict filter. \footnote{This sense of ``paradigm'' was defined by Thomas Kuhn in his
\emph{Structure of Scientific Revolutions}, but I also recommend his
\emph{Copernican Revolution}, where you can see him at work developing
the idea.}

The wise response to such an idea is not to make statements, but to ask
questions, because there's a real mystery here. Why has this smart and
reasonable person proposed an idea that seems so wrong? Are they
mistaken, or are you? One of you has to be. If you're the one who's
mistaken, that would be good to know, because it means there's a hole in
your model of the world. But even if they're mistaken, it should be
interesting to learn why. A trap that an expert falls into is one you
have to worry about too.

This all seems pretty obvious. And yet there are clearly a lot of people
who don't share my fear of dismissing new ideas. Why do they do it? Why
risk looking like a jerk now and a fool later, instead of just reserving
judgement?

One reason they do it is envy. If you propose a radical new idea and it
succeeds, your reputation (and perhaps also your wealth) will increase
proportionally. Some people would be envious if that happened, and this
potential envy propagates back into a conviction that you must be wrong.

Another reason people dismiss new ideas is that it's an easy way to seem
sophisticated. When a new idea first emerges, it usually seems pretty
feeble. It's a mere hatchling. Received wisdom is a full-grown eagle by
comparison. So it's easy to launch a devastating attack on a new idea,
and anyone who does will seem clever to those who don't understand this
asymmetry.

This phenomenon is exacerbated by the difference between how those
working on new ideas and those attacking them are rewarded. The rewards
for working on new ideas are weighted by the value of the outcome. So
it's worth working on something that only has a 10\% chance of
succeeding if it would make things more than 10x better. Whereas the
rewards for attacking new ideas are roughly constant; such attacks seem
roughly equally clever regardless of the target.

People will also attack new ideas when they have a vested interest in
the old ones. It's not surprising, for example, that some of Darwin's
harshest critics were churchmen. People build whole careers on some
ideas. When someone claims they're false or obsolete, they feel
threatened.

The lowest form of dismissal is mere factionalism: to automatically
dismiss any idea associated with the opposing faction. The lowest form
of all is to dismiss an idea because of who proposed it.

But the main thing that leads reasonable people to dismiss new ideas is
the same thing that holds people back from proposing them: the sheer
pervasiveness of the current paradigm. It doesn't just affect the way we
think; it is the Lego blocks we build thoughts out of. Popping out of
the current paradigm is something only a few people can do. And even
they usually have to suppress their intuitions at first, like a pilot
flying through cloud who has to trust his instruments over his sense of
balance. \footnote{This is one reason people with a touch of Asperger's may have an
advantage in discovering new ideas. They're always flying on
instruments.}

Paradigms don't just define our present thinking. They also vacuum up
the trail of crumbs that led to them, making our standards for new ideas
impossibly high. The current paradigm seems so perfect to us, its
offspring, that we imagine it must have been accepted completely as soon
as it was discovered --- that whatever the church thought of the
heliocentric model, astronomers must have been convinced as soon as
Copernicus proposed it. Far, in fact, from it. Copernicus published the
heliocentric model in 1532, but it wasn't till the mid seventeenth
century that the balance of scientific opinion shifted in its favor.
\footnote{Hall, Rupert. \emph{From Galileo to Newton.} Collins, 1963. This
book is particularly good at getting into contemporaries' heads.}

Few understand how feeble new ideas look when they first appear. So if
you want to have new ideas yourself, one of the most valuable things you
can do is to learn what they look like when they're born. Read about how
new ideas happened, and try to get yourself into the heads of people at
the time. How did things look to them, when the new idea was only
half-finished, and even the person who had it was only half-convinced it
was right?

But you don't have to stop at history. You can observe big new ideas
being born all around you right now. Just look for a reasonable domain
expert proposing something that sounds wrong.

If you're nice, as well as wise, you won't merely resist attacking such
people, but encourage them. Having new ideas is a lonely business. Only
those who've tried it know how lonely. These people need your help. And
if you help them, you'll probably learn something in the process.


\textbf{Thanks} to Trevor Blackwell, Patrick Collison, Suhail Doshi,
Daniel Gackle, Jessica Livingston, and Robert Morris for reading drafts
of this.
\chapter{Fierce Nerds}

May 2021

Most people think of nerds as quiet, diffident people. In ordinary
social situations they are --- as quiet and diffident as the star
quarterback would be if he found himself in the middle of a physics
symposium. And for the same reason: they are fish out of water. But the
apparent diffidence of nerds is an illusion due to the fact that when
non-nerds observe them, it's usually in ordinary social situations. In
fact some nerds are quite fierce.

The fierce nerds are a small but interesting group. They are as a rule
extremely competitive --- more competitive, I'd say, than highly
competitive non-nerds. Competition is more personal for them. Partly
perhaps because they're not emotionally mature enough to distance
themselves from it, but also because there's less randomness in the
kinds of competition they engage in, and they are thus more justified in
taking the results personally.

Fierce nerds also tend to be somewhat overconfident, especially when
young. It might seem like it would be a disadvantage to be mistaken
about one's abilities, but empirically it isn't. Up to a point,
confidence is a self-fullfilling prophecy.

Another quality you find in most fierce nerds is intelligence. Not all
nerds are smart, but the fierce ones are always at least moderately so.
If they weren't, they wouldn't have the confidence to be fierce.
\footnote{To be a nerd is to be socially awkward, and there are two
distinct ways to do that: to be playing the same game as everyone else,
but badly, and to be playing a different game. The smart nerds are the
latter type.}

There's also a natural connection between nerdiness and
\href{think.html}{independent-mindedness}. It's hard to be
independent-minded without being somewhat socially awkward, because
conventional beliefs are so often mistaken, or at least arbitrary. No
one who was both independent-minded and ambitious would want to waste
the effort it takes to fit in. And the independent-mindedness of the
fierce nerds will obviously be of the \href{conformism.html}{aggressive}
rather than the passive type: they'll be annoyed by rules, rather than
dreamily unaware of them.

I'm less sure why fierce nerds are impatient, but most seem to be. You
notice it first in conversation, where they tend to interrupt you. This
is merely annoying, but in the more promising fierce nerds it's
connected to a deeper impatience about solving problems. Perhaps the
competitiveness and impatience of fierce nerds are not separate
qualities, but two manifestations of a single underlying drivenness.

When you combine all these qualities in sufficient quantities, the
result is quite formidable. The most vivid example of fierce nerds in
action may be James Watson's \emph{The Double Helix}. The first sentence
of the book is ``I have never seen Francis Crick in a modest mood,'' and
the portrait he goes on to paint of Crick is the quintessential fierce
nerd: brilliant, socially awkward, competitive, independent-minded,
overconfident. But so is the implicit portrait he paints of himself.
Indeed, his lack of social awareness makes both portraits that much more
realistic, because he baldly states all sorts of opinions and
motivations that a smoother person would conceal. And moreover it's
clear from the story that Crick and Watson's fierce nerdiness was
integral to their success. Their independent-mindedness caused them to
consider approaches that most others ignored, their overconfidence
allowed them to work on problems they only half understood (they were
literally described as ``clowns'' by one eminent insider), and their
impatience and competitiveness got them to the answer ahead of two other
groups that would otherwise have found it within the next year, if not
the next several months. \footnote{The same qualities that make fierce nerds so effective can also
make them very annoying. Fierce nerds would do well to remember this,
and (a) try to keep a lid on it, and (b) seek out organizations and
types of work where getting the right answer matters more than
preserving social harmony. In practice that means small groups working
on hard problems. Which fortunately is the most fun kind of environment
anyway.}

The idea that there could be fierce nerds is an unfamiliar one not just
to many normal people but even to some young nerds. Especially early on,
nerds spend so much of their time in ordinary social situations and so
little doing real work that they get a lot more evidence of their
awkwardness than their power. So there will be some who read this
description of the fierce nerd and realize ``Hmm, that's me.'' And it is
to you, young fierce nerd, that I now turn.

I have some good news, and some bad news. The good news is that your
fierceness will be a great help in solving difficult problems. And not
just the kind of scientific and technical problems that nerds have
traditionally solved. As the world progresses, the number of things you
can win at by getting the right answer increases. Recently
\href{richnow.html}{getting rich} became one of them: 7 of the 8 richest
people in America are now fierce nerds.

Indeed, being a fierce nerd is probably even more helpful in business
than in nerds' original territory of scholarship. Fierceness seems
optional there. Darwin for example doesn't seem to have been especially
fierce. Whereas it's impossible to be the CEO of a company over a
certain size without being fierce, so now that nerds can win at
business, fierce nerds will increasingly monopolize the really big
successes.

The bad news is that if it's not exercised, your fierceness will turn to
bitterness, and you will become an intellectual playground bully: the
grumpy sysadmin, the forum troll, the \href{fh.html}{hater}, the shooter
down of \href{newideas.html}{new ideas}.

How do you avoid this fate? Work on ambitious projects. If you succeed,
it will bring you a kind of satisfaction that neutralizes bitterness.
But you don't need to have succeeded to feel this; merely working on
hard projects gives most fierce nerds some feeling of satisfaction. And
those it doesn't, it at least keeps busy. \footnote{If success neutralizes bitterness, why are there some people who
are at least moderately successful and yet still quite bitter? Because
people's potential bitterness varies depending on how naturally bitter
their personality is, and how ambitious they are: someone who's
naturally very bitter will still have a lot left after success
neutralizes some of it, and someone who's very ambitious will need
proportionally more success to satisfy that ambition.

So the worst-case scenario is someone who's both naturally bitter and
extremely ambitious, and yet only moderately successful.}

Another solution may be to somehow turn off your fierceness, by devoting
yourself to meditation or psychotherapy or something like that. Maybe
that's the right answer for some people. I have no idea. But it doesn't
seem the optimal solution to me. If you're given a sharp knife, it seems
to me better to use it than to blunt its edge to avoid cutting yourself.

If you do choose the ambitious route, you'll have a tailwind behind you.
There has never been a better time to be a nerd. In the past century
we've seen a continuous transfer of power from dealmakers to technicians
--- from the charismatic to the competent --- and I don't see anything
on the horizon that will end it. At least not till the nerds end it
themselves by bringing about the singularity.


\textbf{Thanks} to Trevor Blackwell, Steve Blank, Patrick Collison,
Jessica Livingston, Amjad Masad, and Robert Morris for reading drafts of
this.
\chapter{A Project of One's Own}

June 2021

A few days ago, on the way home from school, my nine year old son told
me he couldn't wait to get home to write more of the story he was
working on. This made me as happy as anything I've heard him say --- not
just because he was excited about his story, but because he'd discovered
this way of working. Working on a project of your own is as different
from ordinary work as skating is from walking. It's more fun, but also
much more productive.

What proportion of great work has been done by people who were skating
in this sense? If not all of it, certainly a lot.

There is something special about working on a project of your own. I
wouldn't say exactly that you're happier. A better word would be
excited, or engaged. You're happy when things are going well, but often
they aren't. When I'm writing an essay, most of the time I'm worried and
puzzled: worried that the essay will turn out badly, and puzzled because
I'm groping for some idea that I can't see clearly enough. Will I be
able to pin it down with words? In the end I usually can, if I take long
enough, but I'm never sure; the first few attempts often fail.

You have moments of happiness when things work out, but they don't last
long, because then you're on to the next problem. So why do it at all?
Because to the kind of people who like working this way, nothing else
feels as right. You feel as if you're an animal in its natural habitat,
doing what you were meant to do --- not always happy, maybe, but awake
and alive.

Many kids experience the excitement of working on projects of their own.
The hard part is making this converge with the work you do as an adult.
And our customs make it harder. We treat ``playing'' and ``hobbies'' as
qualitatively different from ``work''. It's not clear to a kid building
a treehouse that there's a direct (though long) route from that to
architecture or engineering. And instead of pointing out the route, we
conceal it, by implicitly treating the stuff kids do as different from
real work. \footnote{``Hobby'' is a curious word. Now it means work that isn't
\emph{real} work --- work that one is not to be judged by --- but
originally it just meant an obsession in a fairly general sense (even a
political opinion, for example) that one metaphorically rode as a child
rides a hobby-horse. It's hard to say if its recent, narrower meaning is
a change for the better or the worse. For sure there are lots of false
positives --- lots of projects that end up being important but are
dismissed initially as mere hobbies. But on the other hand, the concept
provides valuable cover for projects in the early, ugly duckling phase.}

Instead of telling kids that their treehouses could be on the path to
the work they do as adults, we tell them the path goes through school.
And unfortunately schoolwork tends to be very different from working on
projects of one's own. It's usually neither a project, nor one's own. So
as school gets more serious, working on projects of one's own is
something that survives, if at all, as a thin thread off to the side.

It's a bit sad to think of all the high school kids turning their backs
on building treehouses and sitting in class dutifully learning about
Darwin or Newton to pass some exam, when the work that made Darwin and
Newton famous was actually closer in spirit to building treehouses than
studying for exams.

If I had to choose between my kids getting good grades and working on
ambitious projects of their own, I'd pick the projects. And not because
I'm an indulgent parent, but because I've been on the other end and I
know which has more predictive value. When I was picking startups for Y
Combinator, I didn't care about applicants' grades. But if they'd worked
on projects of their own, I wanted to hear all about those.
\footnote{Tiger parents, as parents so often do, are fighting the last
war. Grades mattered more in the old days when the route to success was
to acquire \href{credentials.html}{credentials} while ascending some
predefined ladder. But it's just as well that their tactics are focused
on grades. How awful it would be if they invaded the territory of
projects, and thereby gave their kids a distaste for this kind of work
by forcing them to do it. Grades are already a grim, fake world, and
aren't harmed much by parental interference, but working on one's own
projects is a more delicate, private thing that could be damaged very
easily.}

It may be inevitable that school is the way it is. I'm not saying we
have to redesign it (though I'm not saying we don't), just that we
should understand what it does to our attitudes to work --- that it
steers us toward the dutiful plodding kind of work, often using
competition as bait, and away from skating.

There are occasionally times when schoolwork becomes a project of one's
own. Whenever I had to write a paper, that would become a project of my
own --- except in English classes, ironically, because the things one
has to write in English classes are so \href{essay.html}{bogus}. And
when I got to college and started taking CS classes, the programs I had
to write became projects of my own. Whenever I was writing or
programming, I was usually skating, and that has been true ever since.

So where exactly is the edge of projects of one's own? That's an
interesting question, partly because the answer is so complicated, and
partly because there's so much at stake. There turn out to be two senses
in which work can be one's own: 1) that you're doing it voluntarily,
rather than merely because someone told you to, and 2) that you're doing
it by yourself.

The edge of the former is quite sharp. People who care a lot about their
work are usually very sensitive to the difference between pulling, and
being pushed, and work tends to fall into one category or the other. But
the test isn't simply whether you're told to do something. You can
choose to do something you're told to do. Indeed, you can own it far
more thoroughly than the person who told you to do it.

For example, math homework is for most people something they're told to
do. But for my father, who was a mathematician, it wasn't. Most of us
think of the problems in a math book as a way to test or develop our
knowledge of the material explained in each section. But to my father
the problems were the part that mattered, and the text was merely a sort
of annotation. Whenever he got a new math book it was to him like being
given a puzzle: here was a new set of problems to solve, and he'd
immediately set about solving all of them.

The other sense of a project being one's own --- working on it by
oneself --- has a much softer edge. It shades gradually into
collaboration. And interestingly, it shades into collaboration in two
different ways. One way to collaborate is to share a single project. For
example, when two mathematicians collaborate on a proof that takes shape
in the course of a conversation between them. The other way is when
multiple people work on separate projects of their own that fit together
like a jigsaw puzzle. For example, when one person writes the text of a
book and another does the graphic design. \footnote{The complicated, gradual edge between working on one's own
projects and collaborating with others is one reason there is so much
disagreement about the idea of the ``lone genius.'' In practice people
collaborate (or not) in all kinds of different ways, but the idea of the
lone genius is definitely not a myth. There's a core of truth to it that
goes with a certain way of working.}

These two paths into collaboration can of course be combined. But under
the right conditions, the excitement of working on a project of one's
own can be preserved for quite a while before disintegrating into the
turbulent flow of work in a large organization. Indeed, the history of
successful organizations is partly the history of techniques for
preserving that excitement. \footnote{Collaboration is powerful too. The optimal organization would
combine collaboration and ownership in such a way as to do the least
damage to each. Interestingly, companies and university departments
approach this ideal from opposite directions: companies insist on
collaboration, and occasionally also manage both to recruit skaters and
allow them to skate, and university departments insist on the ability to
do independent research (which is by custom treated as skating, whether
it is or not), and the people they hire collaborate as much as they
choose.}

The team that made the original Macintosh were a great example of this
phenomenon. People like Burrell Smith and Andy Hertzfeld and Bill
Atkinson and Susan Kare were not just following orders. They were not
tennis balls hit by Steve Jobs, but rockets let loose by Steve Jobs.
There was a lot of collaboration between them, but they all seem to have
individually felt the excitement of working on a project of one's own.

In Andy Hertzfeld's book on the Macintosh, he describes how they'd come
back into the office after dinner and work late into the night. People
who've never experienced the thrill of working on a project they're
excited about can't distinguish this kind of working long hours from the
kind that happens in sweatshops and boiler rooms, but they're at
opposite ends of the spectrum. That's why it's a mistake to insist
dogmatically on ``work/life balance.'' Indeed, the mere expression
``work/life'' embodies a mistake: it assumes work and life are distinct.
For those to whom the word ``work'' automatically implies the dutiful
plodding kind, they are. But for the skaters, the relationship between
work and life would be better represented by a dash than a slash. I
wouldn't want to work on anything that I didn't want to take over my
life.

Of course, it's easier to achieve this level of motivation when you're
making something like the Macintosh. It's easy for something new to feel
like a project of your own. That's one of the reasons for the tendency
programmers have to rewrite things that don't need rewriting, and to
write their own versions of things that already exist. This sometimes
alarms managers, and measured by total number of characters typed, it's
rarely the optimal solution. But it's not always driven simply by
arrogance or cluelessness. Writing code from scratch is also much more
rewarding --- so much more rewarding that a good programmer can end up
net ahead, despite the shocking waste of characters. Indeed, it may be
one of the advantages of capitalism that it encourages such rewriting. A
company that needs software to do something can't use the software
already written to do it at another company, and thus has to write their
own, which often turns out better. \footnote{If a company could design its software in such a way that the
best newly arrived programmers always got a clean sheet, it could have a
kind of eternal youth. That might not be impossible. If you had a
software backbone defining a game with sufficiently clear rules,
individual programmers could write their own players.}

The natural alignment between skating and solving new problems is one of
the reasons the payoffs from startups are so high. Not only is the
market price of unsolved problems higher, you also get a discount on
productivity when you work on them. In fact, you get a double increase
in productivity: when you're doing a clean-sheet design, it's easier to
recruit skaters, and they get to spend all their time skating.

Steve Jobs knew a thing or two about skaters from having watched Steve
Wozniak. If you can find the right people, you only have to tell them
what to do at the highest level. They'll handle the details. Indeed,
they insist on it. For a project to feel like your own, you must have
sufficient autonomy. You can't be working to order, or
\href{artistsship.html}{slowed down} by bureaucracy.

One way to ensure autonomy is not to have a boss at all. There are two
ways to do that: to be the boss yourself, and to work on projects
outside of work. Though they're at opposite ends of the scale
financially, startups and open source projects have a lot in common,
including the fact that they're often run by skaters. And indeed,
there's a wormhole from one end of the scale to the other: one of the
best ways to discover \href{startupideas.html}{startup ideas} is to work
on a project just for fun.

If your projects are the kind that make money, it's easy to work on
them. It's harder when they're not. And the hardest part, usually, is
morale. That's where adults have it harder than kids. Kids just plunge
in and build their treehouse without worrying about whether they're
wasting their time, or how it compares to other treehouses. And frankly
we could learn a lot from kids here. The high standards most grownups
have for ``real'' work do not always serve us well.

The most important phase in a project of one's own is at the beginning:
when you go from thinking it might be cool to do x to actually doing x.
And at that point high standards are not merely useless but positively
harmful. There are a few people who start too many new projects, but far
more, I suspect, who are deterred by fear of failure from starting
projects that would have succeeded if they had.

But if we couldn't benefit as kids from the knowledge that our
treehouses were on the path to grownup projects, we can at least benefit
as grownups from knowing that our projects are on a path that stretches
back to treehouses. Remember that careless confidence you had as a kid
when starting something new? That would be a powerful thing to
recapture.

If it's harder as adults to retain that kind of confidence, we at least
tend to be more aware of what we're doing. Kids bounce, or are herded,
from one kind of work to the next, barely realizing what's happening to
them. Whereas we know more about different types of work and have more
control over which we do. Ideally we can have the best of both worlds:
to be deliberate in choosing to work on projects of our own, and
carelessly confident in starting new ones.


\textbf{Thanks} to Trevor Blackwell, Paul Buchheit, Andy Hertzfeld,
Jessica Livingston, and Peter Norvig for reading drafts of this.
\chapter{How to Work Hard}

June 2021

It might not seem there's much to learn about how to work hard. Anyone
who's been to school knows what it entails, even if they chose not to do
it. There are 12 year olds who work amazingly hard. And yet when I ask
if I know more about working hard now than when I was in school, the
answer is definitely yes.

One thing I know is that if you want to do great things, you'll have to
work very hard. I wasn't sure of that as a kid. Schoolwork varied in
difficulty; one didn't always have to work super hard to do well. And
some of the things famous adults did, they seemed to do almost
effortlessly. Was there, perhaps, some way to evade hard work through
sheer brilliance? Now I know the answer to that question. There isn't.

The reason some subjects seemed easy was that my school had low
standards. And the reason famous adults seemed to do things effortlessly
was years of practice; they made it look easy.

Of course, those famous adults usually had a lot of natural ability too.
There are three ingredients in great work: natural ability, practice,
and effort. You can do pretty well with just two, but to do the best
work you need all three: you need great natural ability \emph{and} to
have practiced a lot \emph{and} to be trying very hard.
\footnote{In ``The Bus Ticket Theory of Genius'' I said the three
ingredients in great work were natural ability, determination, and
interest. That's the formula in the preceding stage; determination and
interest yield practice and effort.}

Bill Gates, for example, was among the smartest people in business in
his era, but he was also among the hardest working. ``I never took a day
off in my twenties,'' he said. ``Not one.'' It was similar with Lionel
Messi. He had great natural ability, but when his youth coaches talk
about him, what they remember is not his talent but his dedication and
his desire to win. P. G. Wodehouse would probably get my vote for best
English writer of the 20th century, if I had to choose. Certainly no one
ever made it look easier. But no one ever worked harder. At 74, he wrote

\begin{quote}
with each new book of mine I have, as I say, the feeling that this time
I have picked a lemon in the garden of literature. A good thing, really,
I suppose. Keeps one up on one's toes and makes one rewrite every
sentence ten times. Or in many cases twenty times.
\end{quote}

Sounds a bit extreme, you think. And yet Bill Gates sounds even more
extreme. Not one day off in ten years? These two had about as much
natural ability as anyone could have, and yet they also worked about as
hard as anyone could work. You need both.

That seems so obvious, and yet in practice we find it slightly hard to
grasp. There's a faint xor between talent and hard work. It comes partly
from popular culture, where it seems to run very deep, and partly from
the fact that the outliers are so rare. If great talent and great drive
are both rare, then people with both are rare squared. Most people you
meet who have a lot of one will have less of the other. But you'll need
both if you want to be an outlier yourself. And since you can't really
change how much natural talent you have, in practice doing great work,
insofar as you can, reduces to working very hard.

It's straightforward to work hard if you have clearly defined,
externally imposed goals, as you do in school. There is some technique
to it: you have to learn not to lie to yourself, not to procrastinate
(which is a form of lying to yourself), not to get distracted, and not
to give up when things go wrong. But this level of discipline seems to
be within the reach of quite young children, if they want it.

What I've learned since I was a kid is how to work toward goals that are
neither clearly defined nor externally imposed. You'll probably have to
learn both if you want to do really great things.

The most basic level of which is simply to feel you should be working
without anyone telling you to. Now, when I'm not working hard, alarm
bells go off. I can't be sure I'm getting anywhere when I'm working
hard, but I can be sure I'm getting nowhere when I'm not, and it feels
awful. \footnote{I mean this at a resolution of days, not hours. You'll often get
somewhere while not working in the sense that the solution to a problem
comes to you while taking a \href{top.html}{shower}, or even in your
sleep, but only because you were working hard on it the day before.

It's good to go on vacation occasionally, but when I go on vacation, I
like to learn new things. I wouldn't like just sitting on a beach.}

There wasn't a single point when I learned this. Like most little kids,
I enjoyed the feeling of achievement when I learned or did something
new. As I grew older, this morphed into a feeling of disgust when I
wasn't achieving anything. The one precisely dateable landmark I have is
when I stopped watching TV, at age~13.

Several people I've talked to remember getting serious about work around
this age. When I asked Patrick Collison when he started to find idleness
distasteful, he said

\begin{quote}
I think around age 13 or 14. I have a clear memory from around then of
sitting in the sitting room, staring outside, and wondering why I was
wasting my summer holiday.
\end{quote}

Perhaps something changes at adolescence. That would make sense.

Strangely enough, the biggest obstacle to getting serious about work was
probably school, which made work (what they called work) seem boring and
pointless. I had to learn what real work was before I could
wholeheartedly desire to do it. That took a while, because even in
college a lot of the work is pointless; there are entire departments
that are pointless. But as I learned the shape of real work, I found
that my desire to do it slotted into it as if they'd been made for each
other.

I suspect most people have to learn what work is before they can love
it. Hardy wrote eloquently about this in \emph{A Mathematician's
Apology}:

\begin{quote}
I do not remember having felt, as a boy, any \emph{passion} for
mathematics, and such notions as I may have had of the career of a
mathematician were far from noble. I thought of mathematics in terms of
examinations and scholarships: I wanted to beat other boys, and this
seemed to be the way in which I could do so most decisively.
\end{quote}

He didn't learn what math was really about till part way through
college, when he read Jordan's \emph{Cours d'analyse}.

\begin{quote}
I shall never forget the astonishment with which I read that remarkable
work, the first inspiration for so many mathematicians of my generation,
and learnt for the first time as I read it what mathematics really
meant.
\end{quote}

There are two separate kinds of fakeness you need to learn to discount
in order to understand what real work is. One is the kind Hardy
encountered in school. Subjects get distorted when they're adapted to be
taught to kids --- often so distorted that they're nothing like the work
done by actual practitioners. \footnote{The thing kids do in school that's most like the real version is
sports. Admittedly because many sports originated as games played in
schools. But in this one area, at least, kids are doing exactly what
adults do.

In the average American high school, you have a choice of pretending to
do something serious, or seriously doing something pretend. Arguably the
latter is no worse.} The other kind of
fakeness is intrinsic to certain types of work. Some types of work are
inherently bogus, or at best mere busywork.

There's a kind of solidity to real work. It's not all writing the
\emph{Principia}, but it all feels necessary. That's a vague criterion,
but it's deliberately vague, because it has to cover a lot of different
types. \footnote{Knowing what you want to work on doesn't mean you'll be able to.
Most people have to spend a lot of their time working on things they
don't want to, especially early on. But if you know what you want to do,
you at least know what direction to nudge your life in.}

Once you know the shape of real work, you have to learn how many hours a
day to spend on it. You can't solve this problem by simply working every
waking hour, because in many kinds of work there's a point beyond which
the quality of the result will start to decline.

That limit varies depending on the type of work and the person. I've
done several different kinds of work, and the limits were different for
each. My limit for the harder types of writing or programming is about
five hours a day. Whereas when I was running a startup, I could work all
the time. At least for the three years I did it; if I'd kept going much
longer, I'd probably have needed to take occasional vacations.
\footnote{The lower time limits for intense work suggest a solution to the
problem of having less time to work after you have kids: switch to
harder problems. In effect I did that, though not deliberately.}

The only way to find the limit is by crossing it. Cultivate a
sensitivity to the quality of the work you're doing, and then you'll
notice if it decreases because you're working too hard. Honesty is
critical here, in both directions: you have to notice when you're being
lazy, but also when you're working too hard. And if you think there's
something admirable about working too hard, get that idea out of your
head. You're not merely getting worse results, but getting them because
you're showing off --- if not to other people, then to yourself.
\footnote{Some cultures have a tradition of performative hard work. I
don't love this idea, because (a) it makes a parody of something
important and (b) it causes people to wear themselves out doing things
that don't matter. I don't know enough to say for sure whether it's net
good or bad, but my guess is bad.}

Finding the limit of working hard is a constant, ongoing process, not
something you do just once. Both the difficulty of the work and your
ability to do it can vary hour to hour, so you need to be constantly
judging both how hard you're trying and how well you're doing.

Trying hard doesn't mean constantly pushing yourself to work, though.
There may be some people who do, but I think my experience is fairly
typical, and I only have to push myself occasionally when I'm starting a
project or when I encounter some sort of check. That's when I'm in
danger of procrastinating. But once I get rolling, I tend to keep going.

What keeps me going depends on the type of work. When I was working on
Viaweb, I was driven by fear of failure. I barely procrastinated at all
then, because there was always something that needed doing, and if I
could put more distance between me and the pursuing beast by doing it,
why wait? \footnote{One of the reasons people work so hard on startups is that
startups can fail, and when they do, that failure tends to be both
decisive and conspicuous.} Whereas what drives me now, writing
essays, is the flaws in them. Between essays I fuss for a few days, like
a dog circling while it decides exactly where to lie down. But once I
get started on one, I don't have to push myself to work, because there's
always some error or omission already pushing me.

I do make some amount of effort to focus on important topics. Many
problems have a hard core at the center, surrounded by easier stuff at
the edges. Working hard means aiming toward the center to the extent you
can. Some days you may not be able to; some days you'll only be able to
work on the easier, peripheral stuff. But you should always be aiming as
close to the center as you can without stalling.

The bigger question of what to do with your life is one of these
problems with a hard core. There are important problems at the center,
which tend to be hard, and less important, easier ones at the edges. So
as well as the small, daily adjustments involved in working on a
specific problem, you'll occasionally have to make big, lifetime-scale
adjustments about which type of work to do. And the rule is the same:
working hard means aiming toward the center --- toward the most
ambitious problems.

By center, though, I mean the actual center, not merely the current
consensus about the center. The consensus about which problems are most
important is often mistaken, both in general and within specific fields.
If you disagree with it, and you're right, that could represent a
valuable opportunity to do something new.

The more ambitious types of work will usually be harder, but although
you should not be in denial about this, neither should you treat
difficulty as an infallible guide in deciding what to do. If you
discover some ambitious type of work that's a bargain in the sense of
being easier for you than other people, either because of the abilities
you happen to have, or because of some new way you've found to approach
it, or simply because you're more excited about it, by all means work on
that. Some of the best work is done by people who find an easy way to do
something hard.

As well as learning the shape of real work, you need to figure out which
kind you're suited for. And that doesn't just mean figuring out which
kind your natural abilities match the best; it doesn't mean that if
you're 7 feet tall, you have to play basketball. What you're suited for
depends not just on your talents but perhaps even more on your
interests. A \href{genius.html}{deep interest} in a topic makes people
work harder than any amount of discipline can.

It can be harder to discover your interests than your talents. There are
fewer types of talent than interest, and they start to be judged early
in childhood, whereas interest in a topic is a subtle thing that may not
mature till your twenties, or even later. The topic may not even exist
earlier. Plus there are some powerful sources of error you need to learn
to discount. Are you really interested in x, or do you want to work on
it because you'll make a lot of money, or because other people will be
impressed with you, or because your parents want you to?
\footnote{It's ok to work on something to make a lot of money. You need to
solve the money problem somehow, and there's nothing wrong with doing
that efficiently by trying to make a lot at once. I suppose it would
even be ok to be interested in money for its own sake; whatever floats
your boat. Just so long as you're conscious of your motivations. The
thing to avoid is \emph{unconsciously} letting the need for money warp
your ideas about what kind of work you find most interesting.}

The difficulty of figuring out what to work on varies enormously from
one person to another. That's one of the most important things I've
learned about work since I was a kid. As a kid, you get the impression
that everyone has a calling, and all they have to do is figure out what
it is. That's how it works in movies, and in the streamlined biographies
fed to kids. Sometimes it works that way in real life. Some people
figure out what to do as children and just do it, like Mozart. But
others, like Newton, turn restlessly from one kind of work to another.
Maybe in retrospect we can identify one as their calling --- we can wish
Newton spent more time on math and physics and less on alchemy and
theology --- but this is an \href{disc.html}{illusion} induced by
hindsight bias. There was no voice calling to him that he could have
heard.

So while some people's lives converge fast, there will be others whose
lives never converge. And for these people, figuring out what to work on
is not so much a prelude to working hard as an ongoing part of it, like
one of a set of simultaneous equations. For these people, the process I
described earlier has a third component: along with measuring both how
hard you're working and how well you're doing, you have to think about
whether you should keep working in this field or switch to another. If
you're working hard but not getting good enough results, you should
switch. It sounds simple expressed that way, but in practice it's very
difficult. You shouldn't give up on the first day just because you work
hard and don't get anywhere. You need to give yourself time to get
going. But how much time? And what should you do if work that was going
well stops going well? How much time do you give yourself then?
\footnote{Many people face this question on a smaller scale with
individual projects. But it's easier both to recognize and to accept a
dead end in a single project than to abandon some type of work entirely.
The more determined you are, the harder it gets. Like a Spanish Flu
victim, you're fighting your own immune system: Instead of giving up,
you tell yourself, I should just try harder. And who can say you're not
right?}

What even counts as good results? That can be really hard to decide. If
you're exploring an area few others have worked in, you may not even
know what good results look like. History is full of examples of people
who misjudged the importance of what they were working on.

The best test of whether it's worthwhile to work on something is whether
you find it interesting. That may sound like a dangerously subjective
measure, but it's probably the most accurate one you're going to get.
You're the one working on the stuff. Who's in a better position than you
to judge whether it's important, and what's a better predictor of its
importance than whether it's interesting?

For this test to work, though, you have to be honest with yourself.
Indeed, that's the most striking thing about the whole question of
working hard: how at each point it depends on being honest with
yourself.

Working hard is not just a dial you turn up to 11. It's a complicated,
dynamic system that has to be tuned just right at each point. You have
to understand the shape of real work, see clearly what kind you're best
suited for, aim as close to the true core of it as you can, accurately
judge at each moment both what you're capable of and how you're doing,
and put in as many hours each day as you can without harming the quality
of the result. This network is too complicated to trick. But if you're
consistently honest and clear-sighted, it will automatically assume an
optimal shape, and you'll be productive in a way few people are.


\textbf{Thanks} to Trevor Blackwell, John Carmack, John Collison,
Patrick Collison, Robert Morris, Geoff Ralston, and Harj Taggar for
reading drafts of this.
\chapter{Weird Languages}

August 2021

When people say that in their experience all programming languages are
basically equivalent, they're making a statement not about languages but
about the kind of programming they've done.

99.5\% of programming consists of gluing together calls to library
functions. All popular languages are equally good at this. So one can
easily spend one's whole career operating in the intersection of popular
programming languages.

But the other .5\% of programming is disproportionately interesting. If
you want to learn what it consists of, the weirdness of weird languages
is a good clue to follow.

Weird languages aren't weird by accident. Not the good ones, at least.
The weirdness of the good ones usually implies the existence of some
form of programming that's not just the usual gluing together of library
calls.

A concrete example: Lisp macros. Lisp macros seem weird even to many
Lisp programmers. They're not only not in the intersection of popular
languages, but by their nature would be hard to implement properly in a
language without turning it into a dialect of Lisp. And macros are
definitely evidence of techniques that go beyond glue programming. For
example, solving problems by first writing a language for problems of
that type, and then writing your specific application in it. Nor is this
all you can do with macros; it's just one region in a space of
program-manipulating techniques that even now is far from fully
explored.

So if you want to expand your concept of what programming can be, one
way to do it is by learning weird languages. Pick a language that most
programmers consider weird but whose median user is smart, and then
focus on the differences between this language and the intersection of
popular languages. What can you say in this language that would be
impossibly inconvenient to say in others? In the process of learning how
to say things you couldn't previously say, you'll probably be learning
how to think things you couldn't previously think.

\textbf{Thanks} to Trevor Blackwell, Patrick Collison, Daniel Gackle,
Amjad Masad, and Robert Morris for reading drafts of this.
\chapter{Beyond Smart}

October 2021

If you asked people what was special about Einstein, most would say that
he was really smart. Even the ones who tried to give you a more
sophisticated-sounding answer would probably think this first. Till a
few years ago I would have given the same answer myself. But that wasn't
what was special about Einstein. What was special about him was that he
had important new ideas. Being very smart was a necessary precondition
for having those ideas, but the two are not identical.

It may seem a hair-splitting distinction to point out that intelligence
and its consequences are not identical, but it isn't. There's a big gap
between them. Anyone who's spent time around universities and research
labs knows how big. There are a lot of genuinely smart people who don't
achieve very much.

I grew up thinking that being smart was the thing most to be desired.
Perhaps you did too. But I bet it's not what you really want. Imagine
you had a choice between being really smart but discovering nothing new,
and being less smart but discovering lots of new ideas. Surely you'd
take the latter. I would. The choice makes me uncomfortable, but when
you see the two options laid out explicitly like that, it's obvious
which is better.

The reason the choice makes me uncomfortable is that being smart still
feels like the thing that matters, even though I know intellectually
that it isn't. I spent so many years thinking it was. The circumstances
of childhood are a perfect storm for fostering this illusion.
Intelligence is much easier to measure than the value of new ideas, and
you're constantly being judged by it. Whereas even the kids who will
ultimately discover new things aren't usually discovering them yet. For
kids that way inclined, intelligence is the only game in town.

There are more subtle reasons too, which persist long into adulthood.
Intelligence wins in conversation, and thus becomes the basis of the
dominance hierarchy. \footnote{What wins in conversation depends on who with. It ranges from
mere aggressiveness at the bottom, through quick-wittedness in the
middle, to something closer to actual intelligence at the top, though
probably always with some component of quick-wittedness.} Plus having new ideas is
such a new thing historically, and even now done by so few people, that
society hasn't yet assimilated the fact that this is the actual
destination, and intelligence merely a means to an end.
\footnote{Just as intelligence isn't the only ingredient in having new
ideas, having new ideas isn't the only thing intelligence is useful for.
It's also useful, for example, in diagnosing problems and figuring out
how to fix them. Both overlap with having new ideas, but both have an
end that doesn't.

Those ways of using intelligence are much more common than having new
ideas. And in such cases intelligence is even harder to distinguish from
its consequences.}

Why do so many smart people fail to discover anything new? Viewed from
that direction, the question seems a rather depressing one. But there's
another way to look at it that's not just more optimistic, but more
interesting as well. Clearly intelligence is not the only ingredient in
having new ideas. What are the other ingredients? Are they things we
could cultivate?

Because the trouble with intelligence, they say, is that it's mostly
inborn. The evidence for this seems fairly convincing, especially
considering that most of us don't want it to be true, and the evidence
thus has to face a stiff headwind. But I'm not going to get into that
question here, because it's the other ingredients in new ideas that I
care about, and it's clear that many of them can be cultivated.

That means the truth is excitingly different from the story I got as a
kid. If intelligence is what matters, and also mostly inborn, the
natural consequence is a sort of \emph{Brave New World} fatalism. The
best you can do is figure out what sort of work you have an ``aptitude''
for, so that whatever intelligence you were born with will at least be
put to the best use, and then work as hard as you can at it. Whereas if
intelligence isn't what matters, but only one of several ingredients in
what does, and many of those aren't inborn, things get more interesting.
You have a lot more control, but the problem of how to arrange your life
becomes that much more complicated.

So what are the other ingredients in having new ideas? The fact that I
can even ask this question proves the point I raised earlier --- that
society hasn't assimilated the fact that it's this and not intelligence
that matters. Otherwise we'd all know the answers to such a fundamental
question. \footnote{Some would attribute the difference between intelligence and
having new ideas to ``creativity,'' but this doesn't seem a very useful
term. As well as being pretty vague, it's shifted half a frame sideways
from what we care about: it's neither separable from intelligence, nor
responsible for all the difference between intelligence and having new
ideas.}

I'm not going to try to provide a complete catalogue of the other
ingredients here. This is the first time I've posed the question to
myself this way, and I think it may take a while to answer. But I wrote
recently about one of the most important: an obsessive
\href{genius.html}{interest} in a particular topic. And this can
definitely be cultivated.

Another quality you need in order to discover new ideas is
\href{think.html}{independent-mindedness}. I wouldn't want to claim that
this is distinct from intelligence --- I'd be reluctant to call someone
smart who wasn't independent-minded --- but though largely inborn, this
quality seems to be something that can be cultivated to some extent.

There are general techniques for having new ideas --- for example, for
working on your own \href{own.html}{projects} and for overcoming the
obstacles you face with \href{early.html}{early} work --- and these can
all be learned. Some of them can be learned by societies. And there are
also collections of techniques for generating specific types of new
ideas, like \href{startupideas.html}{startup ideas} and
\href{essay.html}{essay topics}.

And of course there are a lot of fairly mundane ingredients in
discovering new ideas, like \href{hwh.html}{working hard}, getting
enough sleep, avoiding certain kinds of stress, having the right
colleagues, and finding tricks for working on what you want even when
it's not what you're supposed to be working on. Anything that prevents
people from doing great work has an inverse that helps them to. And this
class of ingredients is not as boring as it might seem at first. For
example, having new ideas is generally associated with youth. But
perhaps it's not youth per se that yields new ideas, but specific things
that come with youth, like good health and lack of responsibilities.
Investigating this might lead to strategies that will help people of any
age to have better ideas.

One of the most surprising ingredients in having new ideas is writing
ability. There's a class of new ideas that are best discovered by
writing essays and books. And that ``by'' is deliberate: you don't think
of the ideas first, and then merely write them down. There is a kind of
thinking that one does by writing, and if you're clumsy at writing, or
don't enjoy doing it, that will get in your way if you try to do this
kind of thinking. \footnote{Curiously enough, this essay is an example. It started out as an
essay about writing ability. But when I came to the distinction between
intelligence and having new ideas, that seemed so much more important
that I turned the original essay inside out, making that the topic and
my original topic one of the points in it. As in many other fields, that
level of reworking is easier to contemplate once you've had a lot of
practice.}

I predict the gap between intelligence and new ideas will turn out to be
an interesting place. If we think of this gap merely as a measure of
unrealized potential, it becomes a sort of wasteland that we try to
hurry through with our eyes averted. But if we flip the question, and
start inquiring into the other ingredients in new ideas that it implies
must exist, we can mine this gap for discoveries about discovery.


\textbf{Thanks} to Trevor Blackwell, Patrick Collison, Jessica
Livingston, Robert Morris, Michael Nielsen, and Lisa Randall for reading
drafts of this.
\chapter{Is There Such a Thing as Good Taste?}

November 2021

\emph{(This essay is derived from a talk at the Cambridge Union.)}

When I was a kid, I'd have said there wasn't. My father told me so. Some
people like some things, and other people like other things, and who's
to say who's right?

It seemed so obvious that there was no such thing as good taste that it
was only through indirect evidence that I realized my father was wrong.
And that's what I'm going to give you here: a proof by reductio ad
absurdum. If we start from the premise that there's no such thing as
good taste, we end up with conclusions that are obviously false, and
therefore the premise must be wrong.

We'd better start by saying what good taste is. There's a narrow sense
in which it refers to aesthetic judgements and a broader one in which it
refers to preferences of any kind. The strongest proof would be to show
that taste exists in the narrowest sense, so I'm going to talk about
taste in art. You have better taste than me if the art you like is
better than the art I like.

If there's no such thing as good taste, then there's no such thing as
\href{goodart.html}{good art}. Because if there is such a thing as good
art, it's easy to tell which of two people has better taste. Show them a
lot of works by artists they've never seen before and ask them to choose
the best, and whoever chooses the better art has better taste.

So if you want to discard the concept of good taste, you also have to
discard the concept of good art. And that means you have to discard the
possibility of people being good at making it. Which means there's no
way for artists to be good at their jobs. And not just visual artists,
but anyone who is in any sense an artist. You can't have good actors, or
novelists, or composers, or dancers either. You can have popular
novelists, but not good ones.

We don't realize how far we'd have to go if we discarded the concept of
good taste, because we don't even debate the most obvious cases. But it
doesn't just mean we can't say which of two famous painters is better.
It means we can't say that any painter is better than a randomly chosen
eight year old.

That was how I realized my father was wrong. I started studying
painting. And it was just like other kinds of work I'd done: you could
do it well, or badly, and if you tried hard, you could get better at it.
And it was obvious that Leonardo and Bellini were much better at it than
me. That gap between us was not imaginary. They were so good. And if
they could be good, then art could be good, and there was such a thing
as good taste after all.

Now that I've explained how to show there is such a thing as good taste,
I should also explain why people think there isn't. There are two
reasons. One is that there's always so much disagreement about taste.
Most people's response to art is a tangle of unexamined impulses. Is the
artist famous? Is the subject attractive? Is this the sort of art
they're supposed to like? Is it hanging in a famous museum, or
reproduced in a big, expensive book? In practice most people's response
to art is dominated by such extraneous factors.

And the people who do claim to have good taste are so often mistaken.
The paintings admired by the so-called experts in one generation are
often so different from those admired a few generations later. It's easy
to conclude there's nothing real there at all. It's only when you
isolate this force, for example by trying to paint and comparing your
work to Bellini's, that you can see that it does in fact exist.

The other reason people doubt that art can be good is that there doesn't
seem to be any room in the art for this goodness. The argument goes like
this. Imagine several people looking at a work of art and judging how
good it is. If being good art really is a property of objects, it should
be in the object somehow. But it doesn't seem to be; it seems to be
something happening in the heads of each of the observers. And if they
disagree, how do you choose between them?

The solution to this puzzle is to realize that the purpose of art is to
work on its human audience, and humans have a lot in common. And to the
extent the things an object acts upon respond in the same way, that's
arguably what it means for the object to have the corresponding
property. If everything a particle interacts with behaves as if the
particle had a mass of \emph{m}, then it has a mass of \emph{m}. So the
distinction between ``objective'' and ``subjective'' is not binary, but
a matter of degree, depending on how much the subjects have in common.
Particles interacting with one another are at one pole, but people
interacting with art are not all the way at the other; their reactions
aren't \emph{random}.

Because people's responses to art aren't random, art can be designed to
operate on people, and be good or bad depending on how effectively it
does so. Much as a vaccine can be. If someone were talking about the
ability of a vaccine to confer immunity, it would seem very frivolous to
object that conferring immunity wasn't really a property of vaccines,
because acquiring immunity is something that happens in the immune
system of each individual person. Sure, people's immune systems vary,
and a vaccine that worked on one might not work on another, but that
doesn't make it meaningless to talk about the effectiveness of a
vaccine.

The situation with art is messier, of course. You can't measure
effectiveness by simply taking a vote, as you do with vaccines. You have
to imagine the responses of subjects with a deep knowledge of art, and
enough clarity of mind to be able to ignore extraneous influences like
the fame of the artist. And even then you'd still see some disagreement.
People do vary, and judging art is hard, especially recent art. There is
definitely not a total order either of works or of people's ability to
judge them. But there is equally definitely a partial order of both. So
while it's not possible to have perfect taste, it is possible to have
good taste.

\textbf{Thanks} to the Cambridge Union for inviting me, and to Trevor
Blackwell, Jessica Livingston, and Robert Morris for reading drafts of
this.
\chapter{Putting Ideas into Words}

February 2022

Writing about something, even something you know well, usually shows you
that you didn't know it as well as you thought. Putting ideas into words
is a severe test. The first words you choose are usually wrong; you have
to rewrite sentences over and over to get them exactly right. And your
ideas won't just be imprecise, but incomplete too. Half the ideas that
end up in an essay will be ones you thought of while you were writing
it. Indeed, that's why I write them.

Once you publish something, the convention is that whatever you wrote
was what you thought before you wrote it. These were your ideas, and now
you've expressed them. But you know this isn't true. You know that
putting your ideas into words changed them. And not just the ideas you
published. Presumably there were others that turned out to be too broken
to fix, and those you discarded instead.

It's not just having to commit your ideas to specific words that makes
writing so exacting. The real test is reading what you've written. You
have to pretend to be a neutral reader who knows nothing of what's in
your head, only what you wrote. When he reads what you wrote, does it
seem correct? Does it seem complete? If you make an effort, you can read
your writing as if you were a complete stranger, and when you do the
news is usually bad. It takes me many cycles before I can get an essay
past the stranger. But the stranger is rational, so you always can, if
you ask him what he needs. If he's not satisfied because you failed to
mention x or didn't qualify some sentence sufficiently, then you mention
x or add more qualifications. Happy now? It may cost you some nice
sentences, but you have to resign yourself to that. You just have to
make them as good as you can and still satisfy the stranger.

This much, I assume, won't be that controversial. I think it will accord
with the experience of anyone who has tried to write about anything
nontrivial. There may exist people whose thoughts are so perfectly
formed that they just flow straight into words. But I've never known
anyone who could do this, and if I met someone who said they could, it
would seem evidence of their limitations rather than their ability.
Indeed, this is a trope in movies: the guy who claims to have a plan for
doing some difficult thing, and who when questioned further, taps his
head and says ``It's all up here.'' Everyone watching the movie knows
what that means. At best the plan is vague and incomplete. Very likely
there's some undiscovered flaw that invalidates it completely. At best
it's a plan for a plan.

In precisely defined domains it's possible to form complete ideas in
your head. People can play chess in their heads, for example. And
mathematicians can do some amount of math in their heads, though they
don't seem to feel sure of a proof over a certain length till they write
it down. But this only seems possible with ideas you can express in a
formal language. \footnote{Machinery and circuits are formal languages.} Arguably what such people are
doing is putting ideas into words in their heads. I can to some extent
write essays in my head. I'll sometimes think of a paragraph while
walking or lying in bed that survives nearly unchanged in the final
version. But really I'm writing when I do this. I'm doing the mental
part of writing; my fingers just aren't moving as I do it.
\footnote{I thought of this sentence as I was walking down the street in
Palo Alto.}

You can know a great deal about something without writing about it. Can
you ever know so much that you wouldn't learn more from trying to
explain what you know? I don't think so. I've written about at least two
subjects I know well --- Lisp hacking and startups --- and in both cases
I learned a lot from writing about them. In both cases there were things
I didn't consciously realize till I had to explain them. And I don't
think my experience was anomalous. A great deal of knowledge is
unconscious, and experts have if anything a higher proportion of
unconscious knowledge than beginners.

I'm not saying that writing is the best way to explore all ideas. If you
have ideas about architecture, presumably the best way to explore them
is to build actual buildings. What I'm saying is that however much you
learn from exploring ideas in other ways, you'll still learn new things
from writing about them.

Putting ideas into words doesn't have to mean writing, of course. You
can also do it the old way, by talking. But in my experience, writing is
the stricter test. You have to commit to a single, optimal sequence of
words. Less can go unsaid when you don't have tone of voice to carry
meaning. And you can focus in a way that would seem excessive in
conversation. I'll often spend 2 weeks on an essay and reread drafts 50
times. If you did that in conversation it would seem evidence of some
kind of mental disorder. If you're lazy, of course, writing and talking
are equally useless. But if you want to push yourself to get things
right, writing is the steeper hill. \footnote{There are two senses of talking to someone: a strict sense in
which the conversation is verbal, and a more general sense in which it
can take any form, including writing. In the limit case (e.g.~Seneca's
letters), conversation in the latter sense becomes essay writing.

It can be very useful to talk (in either sense) with other people as
you're writing something. But a verbal conversation will never be more
exacting than when you're talking about something you're writing.}

The reason I've spent so long establishing this rather obvious point is
that it leads to another that many people will find shocking. If writing
down your ideas always makes them more precise and more complete, then
no one who hasn't written about a topic has fully formed ideas about it.
And someone who never writes has no fully formed ideas about anything
nontrivial.

It feels to them as if they do, especially if they're not in the habit
of critically examining their own thinking. Ideas can feel complete.
It's only when you try to put them into words that you discover they're
not. So if you never subject your ideas to that test, you'll not only
never have fully formed ideas, but also never realize it.

Putting ideas into words is certainly no guarantee that they'll be
right. Far from it. But though it's not a sufficient condition, it is a
necessary one.


\textbf{Thanks} to Trevor Blackwell, Patrick Collison, and Robert Morris
for reading drafts of this.
\chapter{Heresy}

April 2022

One of the most surprising things I've witnessed in my lifetime is the
rebirth of the concept of heresy.

In his excellent biography of Newton, Richard Westfall writes about the
moment when he was elected a fellow of Trinity College:

\begin{quote}
Supported comfortably, Newton was free to devote himself wholly to
whatever he chose. To remain on, he had only to avoid the three
unforgivable sins: crime, heresy, and marriage. \footnote{Or more accurately, biographies of Newton, since Westfall wrote
two: a long version called \emph{Never at Rest}, and a shorter one
called \emph{The Life of Isaac Newton}. Both are great. The short
version moves faster, but the long one is full of interesting and often
very funny details. This passage is the same in both.}
\end{quote}

The first time I read that, in the 1990s, it sounded amusingly medieval.
How strange, to have to avoid committing heresy. But when I reread it 20
years later it sounded like a description of contemporary employment.

There are an ever-increasing number of opinions you can be fired for.
Those doing the firing don't use the word ``heresy'' to describe them,
but structurally they're equivalent. Structurally there are two
distinctive things about heresy: (1) that it takes priority over the
question of truth or falsity, and (2) that it outweighs everything else
the speaker has done.

For example, when someone calls a statement ``x-ist,'' they're also
implicitly saying that this is the end of the discussion. They do not,
having said this, go on to consider whether the statement is true or
not. Using such labels is the conversational equivalent of signalling an
exception. That's one of the reasons they're used: to end a discussion.

If you find yourself talking to someone who uses these labels a lot, it
might be worthwhile to ask them explicitly if they believe any babies
are being thrown out with the bathwater. Can a statement be x-ist, for
whatever value of x, and also true? If the answer is yes, then they're
admitting to banning the truth. That's obvious enough that I'd guess
most would answer no. But if they answer no, it's easy to show that
they're mistaken, and that in practice such labels are applied to
statements regardless of their truth or falsity.

The clearest evidence of this is that whether a statement is considered
x-ist often depends on who said it. Truth doesn't work that way. The
same statement can't be true when one person says it, but x-ist, and
therefore false, when another person does. \footnote{Another more subtle but equally damning bit of evidence is that
claims of x-ism are never qualified. You never hear anyone say that a
statement is ``probably x-ist'' or ``almost certainly y-ist.'' If claims
of x-ism were actually claims about truth, you'd expect to see
``probably'' in front of ``x-ist'' as often as you see it in front of
``fallacious.''}

The other distinctive thing about heresies, compared to ordinary
opinions, is that the public expression of them outweighs everything
else the speaker has done. In ordinary matters, like knowledge of
history, or taste in music, you're judged by the average of your
opinions. A heresy is qualitatively different. It's like dropping a
chunk of uranium onto the scale.

Back in the day (and still, in some places) the punishment for heresy
was death. You could have led a life of exemplary goodness, but if you
publicly doubted, say, the divinity of Christ, you were going to burn.
Nowadays, in civilized countries, heretics only get fired in the
metaphorical sense, by losing their jobs. But the structure of the
situation is the same: the heresy outweighs everything else. You could
have spent the last ten years saving children's lives, but if you
express certain opinions, you're automatically fired.

It's much the same as if you committed a crime. No matter how virtuously
you've lived, if you commit a crime, you must still suffer the penalty
of the law. Having lived a previously blameless life might mitigate the
punishment, but it doesn't affect whether you're guilty or not.

A heresy is an opinion whose expression is treated like a crime --- one
that makes some people feel not merely that you're mistaken, but that
you should be punished. Indeed, their desire to see you punished is
often stronger than it would be if you'd committed an actual crime.
There are many on the far left who believe strongly in the reintegration
of felons (as I do myself), and yet seem to feel that anyone guilty of
certain heresies should never work again.

There are always some heresies --- some opinions you'd be punished for
expressing. But there are a lot more now than there were a few decades
ago, and even those who are happy about this would have to agree that
it's so.

Why? Why has this antiquated-sounding religious concept come back in a
secular form? And why now?

You need two ingredients for a wave of intolerance: intolerant people,
and an ideology to guide them. The intolerant people are always there.
They exist in every sufficiently large society. That's why waves of
intolerance can arise so suddenly; all they need is something to set
them off.

I've already written an \href{conformism.html}{essay} describing the
aggressively conventional-minded. The short version is that people can
be classified in two dimensions according to (1) how independent- or
conventional-minded they are, and (2) how aggressive they are about it.
The aggressively conventional-minded are the enforcers of orthodoxy.

Normally they're only locally visible. They're the grumpy, censorious
people in a group --- the ones who are always first to complain when
something violates the current rules of propriety. But occasionally,
like a vector field whose elements become aligned, a large number of
aggressively conventional-minded people unite behind some ideology all
at once. Then they become much more of a problem, because a mob dynamic
takes over, where the enthusiasm of each participant is increased by the
enthusiasm of the others.

The most notorious 20th century case may have been the Cultural
Revolution. Though initiated by Mao to undermine his rivals, the
Cultural Revolution was otherwise mostly a grass-roots phenomenon. Mao
said in essence: There are heretics among us. Seek them out and punish
them. And that's all the aggressively conventional-minded ever need to
hear. They went at it with the delight of dogs chasing squirrels.

To unite the conventional-minded, an ideology must have many of the
features of a religion. In particular it must have strict and arbitrary
rules that adherents can demonstrate their
\href{https://www.youtube.com/watch?v=qaHLd8de6nM}{purity} by obeying,
and its adherents must believe that anyone who obeys these rules is ipso
facto morally superior to anyone who doesn't. \footnote{The rules must be strict, but they need not be demanding. So the
most effective type of rules are those about superficial matters, like
doctrinal minutiae, or the precise words adherents must use. Such rules
can be made extremely complicated, and yet don't repel potential
converts by requiring significant sacrifice.

The superficial demands of orthodoxy make it an inexpensive substitute
for virtue. And that in turn is one of the reasons orthodoxy is so
attractive to bad people. You could be a horrible person, and yet as
long as you're orthodox, you're better than everyone who isn't.}

In the late 1980s a new ideology of this type appeared in US
universities. It had a very strong component of moral purity, and the
aggressively conventional-minded seized upon it with their usual
eagerness --- all the more because the relaxation of social norms in the
preceding decades meant there had been less and less to forbid. The
resulting wave of intolerance has been eerily similar in form to the
Cultural Revolution, though fortunately much smaller in magnitude.
\footnote{Arguably there were two. The first had died down somewhat by
2000, but was followed by a second in the 2010s, probably caused by
social media.}

I've deliberately avoided mentioning any specific heresies here. Partly
because one of the universal tactics of heretic hunters, now as in the
past, is to accuse those who disapprove of the way in which they
suppress ideas of being heretics themselves. Indeed, this tactic is so
consistent that you could use it as a way of detecting witch hunts in
any era.

And that's the second reason I've avoided mentioning any specific
heresies. I want this essay to work in the future, not just now. And
unfortunately it probably will. The aggressively conventional-minded
will always be among us, looking for things to forbid. All they need is
an ideology to tell them what. And it's unlikely the current one will be
the last.

There are aggressively conventional-minded people on both the right and
the left. The reason the current wave of intolerance comes from the left
is simply because the new unifying ideology happened to come from the
left. The next one might come from the right. Imagine what that would be
like.

Fortunately in western countries the suppression of heresies is nothing
like as bad as it used to be. Though the window of opinions you can
express publicly has narrowed in the last decade, it's still much wider
than it was a few hundred years ago. The problem is the derivative. Up
till about 1985 the window had been growing ever wider. Anyone looking
into the future in 1985 would have expected freedom of expression to
continue to increase. Instead it has decreased. \footnote{Fortunately most of those trying to suppress ideas today still
respect Enlightenment principles enough to pay lip service to them. They
know they're not supposed to ban ideas per se, so they have to recast
the ideas as causing ``harm,'' which sounds like something that can be
banned. The more extreme try to claim speech itself is violence, or even
that silence is. But strange as it may sound, such gymnastics are a good
sign. We'll know we're really in trouble when they stop bothering to
invent pretenses for banning ideas --- when, like the medieval church,
they say ``Damn right we're banning ideas, and in fact here's a list of
them.''}

The situation is similar to what's happened with infectious diseases
like measles. Anyone looking into the future in 2010 would have expected
the number of measles cases in the US to continue to decrease. Instead,
thanks to anti-vaxxers, it has increased. The absolute number is still
not that high. The problem is the derivative. \footnote{People only have the luxury of ignoring the medical consensus
about vaccines because vaccines have worked so well. If we didn't have
any vaccines at all, the mortality rate would be so high that most
current anti-vaxxers would be begging for them. And the situation with
freedom of expression is similar. It's only because they live in a world
created by the Enlightenment that kids from the suburbs can play at
banning ideas.}

In both cases it's hard to know how much to worry. Is it really
dangerous to society as a whole if a handful of extremists refuse to get
their kids vaccinated, or shout down speakers at universities? The point
to start worrying is presumably when their efforts start to spill over
into everyone else's lives. And in both cases that does seem to be
happening.

So it's probably worth spending some amount of effort on pushing back to
keep open the window of free expression. My hope is that this essay will
help form social antibodies not just against current efforts to suppress
ideas, but against the concept of heresy in general. That's the real
prize. How do you disable the concept of heresy? Since the
Enlightenment, western societies have discovered many techniques for
doing that, but there are surely more to be discovered.

Overall I'm optimistic. Though the trend in freedom of expression has
been bad over the last decade, it's been good over the longer term. And
there are signs that the current wave of intolerance is peaking.
Independent-minded people I talk to seem more confident than they did a
few years ago. On the other side, even some of the
\href{https://www.nytimes.com/2022/03/18/opinion/cancel-culture-free-speech-poll.html}{leaders}
are starting to wonder if things have gone too far. And popular culture
among the young has already moved on. All we have to do is keep pushing
back, and the wave collapses. And then we'll be net ahead, because as
well as having defeated this wave, we'll also have developed new tactics
for resisting the next one.


\textbf{Thanks} to Marc Andreessen, Chris Best, Trevor Blackwell,
Nicholas Christakis, Daniel Gackle, Jonathan Haidt, Claire Lehmann,
Jessica Livingston, Greg Lukianoff, Robert Morris, and Garry Tan for
reading drafts of this.
\chapter{What I've Learned from Users}

September 2022

I recently told applicants to Y Combinator that the best advice I could
give for getting in, per word, was

\begin{quote}
Explain what you've learned from users.
\end{quote}

That tests a lot of things: whether you're paying attention to users,
how well you understand them, and even how much they need what you're
making.

Afterward I asked myself the same question. What have I learned from
YC's users, the startups we've funded?

The first thing that came to mind was that most startups have the same
problems. No two have exactly the same problems, but it's surprising how
much the problems remain the same, regardless of what they're making.
Once you've advised 100 startups all doing different things, you rarely
encounter problems you haven't seen before.

This fact is one of the things that makes YC work. But I didn't know it
when we started YC. I only had a few data points: our own startup, and
those started by friends. It was a surprise to me how often the same
problems recur in different forms. Many later stage investors might
never realize this, because later stage investors might not advise 100
startups in their whole career, but a YC partner will get this much
experience in the first year or two.

That's one advantage of funding large numbers of early stage companies
rather than smaller numbers of later-stage ones. You get a lot of data.
Not just because you're looking at more companies, but also because more
goes wrong.

But knowing (nearly) all the problems startups can encounter doesn't
mean that advising them can be automated, or reduced to a formula.
There's no substitute for individual office hours with a YC partner.
Each startup is unique, which means they have to be advised by specific
partners who know them well. \footnote{This is why I've never liked it when people refer to YC as a
``bootcamp.'' It's intense like a bootcamp, but the opposite in
structure. Instead of everyone doing the same thing, they're each
talking to YC partners to figure out what their specific startup needs.}

We learned that the hard way, in the notorious ``batch that broke YC''
in the summer of 2012. Up till that point we treated the partners as a
pool. When a startup requested office hours, they got the next available
slot posted by any partner. That meant every partner had to know every
startup. This worked fine up to 60 startups, but when the batch grew to
80, everything broke. The founders probably didn't realize anything was
wrong, but the partners were confused and unhappy because halfway
through the batch they still didn't know all the companies yet.
\footnote{When I say the summer 2012 batch was broken, I mean it felt to
the partners that something was wrong. Things weren't yet so broken that
the startups had a worse experience. In fact that batch did unusually
well.}

At first I was puzzled. How could things be fine at 60 startups and
broken at 80? It was only a third more. Then I realized what had
happened. We were using an \emph{O(n2)} algorithm. So of course it blew
up.

The solution we adopted was the classic one in these situations. We
sharded the batch into smaller groups of startups, each overseen by a
dedicated group of partners. That fixed the problem, and has worked fine
ever since. But the batch that broke YC was a powerful demonstration of
how individualized the process of advising startups has to be.

Another related surprise is how bad founders can be at realizing what
their problems are. Founders will sometimes come in to talk about some
problem, and we'll discover another much bigger one in the course of the
conversation. For example (and this case is all too common), founders
will come in to talk about the difficulties they're having raising
money, and after digging into their situation, it turns out the reason
is that the company is doing badly, and investors can tell. Or founders
will come in worried that they still haven't cracked the problem of user
acquisition, and the reason turns out to be that their product isn't
good enough. There have been times when I've asked ``Would you use this
yourself, if you hadn't built it?'' and the founders, on thinking about
it, said ``No.'' Well, there's the reason you're having trouble getting
users.

Often founders know what their problems are, but not their relative
importance. \footnote{This situation reminds me of the research showing that people
are much better at answering questions than they are at judging how
accurate their answers are. The two phenomena feel very similar.} They'll come in to talk about three
problems they're worrying about. One is of moderate importance, one
doesn't matter at all, and one will kill the company if it isn't
addressed immediately. It's like watching one of those horror movies
where the heroine is deeply upset that her boyfriend cheated on her, and
only mildly curious about the door that's mysteriously ajar. You want to
say: never mind about your boyfriend, think about that door! Fortunately
in office hours you can. So while startups still die with some
regularity, it's rarely because they wandered into a room containing a
murderer. The YC partners can warn them where the murderers are.

Not that founders listen. That was another big surprise: how often
founders don't listen to us. A couple weeks ago I talked to a partner
who had been working for YC for a couple batches and was starting to see
the pattern. ``They come back a year later,'' she said, ``and say `We
wish we'd listened to you.'\,''

It took me a long time to figure out why founders don't listen. At first
I thought it was mere stubbornness. That's part of the reason, but
another and probably more important reason is that so much about
startups is \href{before.html}{counterintuitive}. And when you tell
someone something counterintuitive, what it sounds to them is wrong. So
the reason founders don't listen to us is that they don't \emph{believe}
us. At least not till experience teaches them otherwise.
\footnote{The \href{airbnbs.html}{Airbnbs} were particularly good at
listening --- partly because they were flexible and disciplined, but
also because they'd had such a rough time during the preceding year.
They were ready to listen.}

The reason startups are so counterintuitive is that they're so different
from most people's other experiences. No one knows what it's like except
those who've done it. Which is why YC partners should usually have been
founders themselves. But strangely enough, the counterintuitiveness of
startups turns out to be another of the things that make YC work. If it
weren't counterintuitive, founders wouldn't need our advice about how to
do it.

Focus is doubly important for early stage startups, because not only do
they have a hundred different problems, they don't have anyone to work
on them except the founders. If the founders focus on things that don't
matter, there's no one focusing on the things that do. So the essence of
what happens at YC is to figure out which problems matter most, then
cook up ideas for solving them --- ideally at a resolution of a week or
less --- and then try those ideas and measure how well they worked. The
focus is on action, with measurable, near-term results.

This doesn't imply that founders should rush forward regardless of the
consequences. If you correct course at a high enough frequency, you can
be simultaneously decisive at a micro scale and tentative at a macro
scale. The result is a somewhat winding path, but executed very rapidly,
like the path a running back takes downfield. And in practice there's
less backtracking than you might expect. Founders usually guess right
about which direction to run in, especially if they have someone
experienced like a YC partner to bounce their hypotheses off. And when
they guess wrong, they notice fast, because they'll talk about the
results at office hours the next week. \footnote{The optimal unit of decisiveness depends on how long it takes to
get results, and that depends on the type of problem you're solving.
When you're negotiating with investors, it could be a couple days,
whereas if you're building hardware it could be months.}

A small improvement in navigational ability can make you a lot faster,
because it has a double effect: the path is shorter, and you can travel
faster along it when you're more certain it's the right one. That's
where a lot of YC's value lies, in helping founders get an extra
increment of focus that lets them move faster. And since moving fast is
the essence of a startup, YC in effect makes startups more startup-like.

Speed defines startups. Focus enables speed. YC improves focus.

Why are founders uncertain about what to do? Partly because startups
almost by definition are doing something new, which means no one knows
how to do it yet, or in most cases even what ``it'' is. Partly because
startups are so counterintuitive generally. And partly because many
founders, especially young and ambitious ones, have been trained to win
the wrong way. That took me years to figure out. The educational system
in most countries trains you to win by \href{lesson.html}{hacking the
test} instead of actually doing whatever it's supposed to measure. But
that stops working when you start a startup. So part of what YC does is
to retrain founders to stop trying to hack the test. (It takes a
surprisingly long time. A year in, you still see them reverting to their
old habits.)

YC is not simply more experienced founders passing on their knowledge.
It's more like specialization than apprenticeship. The knowledge of the
YC partners and the founders have different shapes: It wouldn't be
worthwhile for a founder to acquire the encyclopedic knowledge of
startup problems that a YC partner has, just as it wouldn't be
worthwhile for a YC partner to acquire the depth of domain knowledge
that a founder has. That's why it can still be valuable for an
experienced founder to do YC, just as it can still be valuable for an
experienced athlete to have a coach.

The other big thing YC gives founders is colleagues, and this may be
even more important than the advice of partners. If you look at history,
great work clusters around certain places and institutions: Florence in
the late 15th century, the University of Göttingen in the late 19th,
\emph{The New Yorker} under Ross, Bell Labs, Xerox PARC. However good
you are, good colleagues make you better. Indeed, very ambitious people
probably need colleagues more than anyone else, because they're so
starved for them in everyday life.

Whether or not YC manages one day to be listed alongside those famous
clusters, it won't be for lack of trying. We were very aware of this
historical phenomenon and deliberately designed YC to be one. By this
point it's not bragging to say that it's the biggest cluster of great
startup founders. Even people trying to attack YC concede that.

Colleagues and startup founders are two of the most powerful forces in
the world, so you'd expect it to have a big effect to combine them.
Before YC, to the extent people thought about the question at all, most
assumed they couldn't be combined --- that loneliness was the price of
independence. That was how it felt to us when we started our own startup
in Boston in the 1990s. We had a handful of older people we could go to
for advice (of varying quality), but no peers. There was no one we could
commiserate with about the misbehavior of investors, or speculate with
about the future of technology. I often tell founders to make something
they themselves want, and YC is certainly that: it was designed to be
exactly what we wanted when we were starting a startup.

One thing we wanted was to be able to get seed funding without having to
make the rounds of random rich people. That has become a commodity now,
at least in the US. But great colleagues can never become a commodity,
because the fact that they cluster in some places means they're
proportionally absent from the rest.

Something magical happens where they do cluster though. The energy in
the room at a YC dinner is like nothing else I've experienced. We would
have been happy just to have one or two other startups to talk to. When
you have a whole roomful it's another thing entirely.

YC founders aren't just inspired by one another. They also help one
another. That's the happiest thing I've learned about startup founders:
how generous they can be in helping one another. We noticed this in the
first batch and consciously designed YC to magnify it. The result is
something far more intense than, say, a university. Between the
partners, the alumni, and their batchmates, founders are surrounded by
people who want to help them, and can.


\textbf{Thanks} to Trevor Blackwell, Jessica Livingston, Harj Taggar,
and Garry Tan for reading drafts of this.
\chapter{Alien Truth}

October 2022

If there were intelligent beings elsewhere in the universe, they'd share
certain truths in common with us. The truths of mathematics would be the
same, because they're true by definition. Ditto for the truths of
physics; the mass of a carbon atom would be the same on their planet.
But I think we'd share other truths with aliens besides the truths of
math and physics, and that it would be worthwhile to think about what
these might be.

For example, I think we'd share the principle that a controlled
experiment testing some hypothesis entitles us to have proportionally
increased belief in it. It seems fairly likely, too, that it would be
true for aliens that one can get better at something by practicing. We'd
probably share Occam's razor. There doesn't seem anything specifically
human about any of these ideas.

We can only guess, of course. We can't say for sure what forms
intelligent life might take. Nor is it my goal here to explore that
question, interesting though it is. The point of the idea of alien truth
is not that it gives us a way to speculate about what forms intelligent
life might take, but that it gives us a threshold, or more precisely a
target, for truth. If you're trying to find the most general truths
short of those of math or physics, then presumably they'll be those we'd
share in common with other forms of intelligent life.

Alien truth will work best as a heuristic if we err on the side of
generosity. If an idea might plausibly be relevant to aliens, that's
enough. Justice, for example. I wouldn't want to bet that all
intelligent beings would understand the concept of justice, but I
wouldn't want to bet against it either.

The idea of alien truth is related to Erdos's idea of God's book. He
used to describe a particularly good proof as being in God's book, the
implication being (a) that a sufficiently good proof was more discovered
than invented, and (b) that its goodness would be universally
recognized. If there's such a thing as alien truth, then there's more in
God's book than math.

What should we call the search for alien truth? The obvious choice is
``philosophy.'' Whatever else philosophy includes, it should probably
include this. I'm fairly sure Aristotle would have thought so. One could
even make the case that the search for alien truth is, if not an
accurate description \emph{of} philosophy, a good definition \emph{for}
it. I.e. that it's what people who call themselves philosophers should
be doing, whether or not they currently are. But I'm not wedded to that;
doing it is what matters, not what we call it.

We may one day have something like alien life among us in the form of
AIs. And that may in turn allow us to be precise about what truths an
intelligent being would have to share with us. We might find, for
example, that it's impossible to create something we'd consider
intelligent that doesn't use Occam's razor. We might one day even be
able to prove that. But though this sort of research would be very
interesting, it's not necessary for our purposes, or even the same
field; the goal of philosophy, if we're going to call it that, would be
to see what ideas we come up with using alien truth as a target, not to
say precisely where the threshold of it is. Those two questions might
one day converge, but they'll converge from quite different directions,
and till they do, it would be too constraining to restrict ourselves to
thinking only about things we're certain would be alien truths.
Especially since this will probably be one of those areas where the best
guesses turn out to be surprisingly close to optimal. (Let's see if that
one does.)

Whatever we call it, the attempt to discover alien truths would be a
worthwhile undertaking. And curiously enough, that is itself probably an
alien truth.

\textbf{Thanks} to Trevor Blackwell, Greg Brockman, Patrick Collison,
Robert Morris, and Michael Nielsen for reading drafts of this.
\chapter{What You (Want to)* Want}

November 2022

Since I was about 9 I've been puzzled by the apparent contradiction
between being made of matter that behaves in a predictable way, and the
feeling that I could choose to do whatever I wanted. At the time I had a
self-interested motive for exploring the question. At that age (like
most succeeding ages) I was always in trouble with the authorities, and
it seemed to me that there might possibly be some way to get out of
trouble by arguing that I wasn't responsible for my actions. I gradually
lost hope of that, but the puzzle remained: How do you reconcile being a
machine made of matter with the feeling that you're free to choose what
you do? \footnote{I didn't know when I was 9 that matter might behave randomly,
but I don't think it affects the problem much. Randomness destroys the
ghost in the machine as effectively as determinism.}

The best way to explain the answer may be to start with a slightly wrong
version, and then fix it. The wrong version is: You can do what you
want, but you can't want what you want. Yes, you can control what you
do, but you'll do what you want, and you can't control that.

The reason this is mistaken is that people do sometimes change what they
want. People who don't want to want something --- drug addicts, for
example --- can sometimes make themselves stop wanting it. And people
who want to want something --- who want to like classical music, or
broccoli --- sometimes succeed.

So we modify our initial statement: You can do what you want, but you
can't want to want what you want.

That's still not quite true. It's possible to change what you want to
want. I can imagine someone saying ``I decided to stop wanting to like
classical music.'' But we're getting closer to the truth. It's rare for
people to change what they want to want, and the more ``want to''s we
add, the rarer it gets.

We can get arbitrarily close to a true statement by adding more ``want
to''s in much the same way we can get arbitrarily close to 1 by adding
more 9s to a string of 9s following a decimal point. In practice three
or four ``want to''s must surely be enough. It's hard even to envision
what it would mean to change what you want to want to want to want, let
alone actually do it.

So one way to express the correct answer is to use a regular expression.
You can do what you want, but there's some statement of the form ``you
can't (want to)* want what you want'' that's true. Ultimately you get
back to a want that you don't control. \footnote{If you don't like using an expression, you can make the same
point using higher-order desires: There is some n such that you don't
control your nth-order desires.}


\textbf{Thanks} to Trevor Blackwell, Jessica Livingston, Robert Morris,
and Michael Nielsen for reading drafts of this.
\chapter{The Need to Read}

November 2022

In the science fiction books I read as a kid, reading had often been
replaced by some more efficient way of acquiring knowledge. Mysterious
``tapes'' would load it into one's brain like a program being loaded
into a computer.

That sort of thing is unlikely to happen anytime soon. Not just because
it would be hard to build a replacement for reading, but because even if
one existed, it would be insufficient. Reading about x doesn't just
teach you about x; it also teaches you how to write.
\footnote{Audiobooks can give you examples of good writing, but having
them read to you doesn't teach you as much about writing as reading them
yourself.}

Would that matter? If we replaced reading, would anyone need to be good
at writing?

The reason it would matter is that writing is not just a way to convey
ideas, but also a way to have them.

A good writer doesn't just think, and then write down what he thought,
as a sort of transcript. A good writer will almost always discover new
things in the process of writing. And there is, as far as I know, no
substitute for this kind of discovery. Talking about your ideas with
other people is a good way to develop them. But even after doing this,
you'll find you still discover new things when you sit down to write.
There is a kind of thinking that can only be done by
\href{words.html}{writing}.

There are of course kinds of thinking that can be done without writing.
If you don't need to go too deeply into a problem, you can solve it
without writing. If you're thinking about how two pieces of machinery
should fit together, writing about it probably won't help much. And when
a problem can be described formally, you can sometimes solve it in your
head. But if you need to solve a complicated, ill-defined problem, it
will almost always help to write about it. Which in turn means that
someone who's not good at writing will almost always be at a
disadvantage in solving such problems.

You can't think well without writing well, and you can't write well
without reading well. And I mean that last ``well'' in both senses. You
have to be good at reading, and read good things.
\footnote{By ``good at reading'' I don't mean good at the mechanics of
reading. You don't have to be good at extracting words from the page so
much as extracting meaning from the words.}

People who just want information may find other ways to get it. But
people who want to have ideas can't afford to.
\chapter{How to Get New Ideas}

January 2023

\emph{(\href{https://twitter.com/stef/status/1617222428727586816}{Someone}
fed my essays into GPT to make something that could answer questions
based on them, then asked it where good ideas come from. The answer was
ok, but not what I would have said. This is what I would have said.)}

The way to get new ideas is to notice anomalies: what seems strange, or
missing, or broken? You can see anomalies in everyday life (much of
standup comedy is based on this), but the best place to look for them is
at the frontiers of knowledge.

Knowledge grows fractally. From a distance its edges look smooth, but
when you learn enough to get close to one, you'll notice it's full of
gaps. These gaps will seem obvious; it will seem inexplicable that no
one has tried x or wondered about y. In the best case, exploring such
gaps yields whole new fractal buds.
\chapter{How to Do Great Work}

July 2023

If you collected lists of techniques for doing great work in a lot of
different fields, what would the intersection look like? I decided to
find out by making it.

Partly my goal was to create a guide that could be used by someone
working in any field. But I was also curious about the shape of the
intersection. And one thing this exercise shows is that it does have a
definite shape; it's not just a point labelled ``work hard.''

The following recipe assumes you're very ambitious.

The first step is to decide what to work on. The work you choose needs
to have three qualities: it has to be something you have a natural
aptitude for, that you have a deep interest in, and that offers scope to
do great work.

In practice you don't have to worry much about the third criterion.
Ambitious people are if anything already too conservative about it. So
all you need to do is find something you have an aptitude for and great
interest in. \footnote{I don't think you could give a precise definition of what counts
as great work. Doing great work means doing something important so well
that you expand people's ideas of what's possible. But there's no
threshold for importance. It's a matter of degree, and often hard to
judge at the time anyway. So I'd rather people focused on developing
their interests rather than worrying about whether they're important or
not. Just try to do something amazing, and leave it to future
generations to say if you succeeded.}

That sounds straightforward, but it's often quite difficult. When you're
young you don't know what you're good at or what different kinds of work
are like. Some kinds of work you end up doing may not even exist yet. So
while some people know what they want to do at 14, most have to figure
it out.

The way to figure out what to work on is by working. If you're not sure
what to work on, guess. But pick something and get going. You'll
probably guess wrong some of the time, but that's fine. It's good to
know about multiple things; some of the biggest discoveries come from
noticing connections between different fields.

Develop a habit of working on your own projects. Don't let ``work'' mean
something other people tell you to do. If you do manage to do great work
one day, it will probably be on a project of your own. It may be within
some bigger project, but you'll be driving your part of it.

What should your projects be? Whatever seems to you excitingly
ambitious. As you grow older and your taste in projects evolves,
exciting and important will converge. At 7 it may seem excitingly
ambitious to build huge things out of Lego, then at 14 to teach yourself
calculus, till at 21 you're starting to explore unanswered questions in
physics. But always preserve excitingness.

There's a kind of excited curiosity that's both the engine and the
rudder of great work. It will not only drive you, but if you let it have
its way, will also show you what to work on.

What are you excessively curious about --- curious to a degree that
would bore most other people? That's what you're looking for.

Once you've found something you're excessively interested in, the next
step is to learn enough about it to get you to one of the frontiers of
knowledge. Knowledge expands fractally, and from a distance its edges
look smooth, but once you learn enough to get close to one, they turn
out to be full of gaps.

The next step is to notice them. This takes some skill, because your
brain wants to ignore such gaps in order to make a simpler model of the
world. Many discoveries have come from asking questions about things
that everyone else took for granted. \footnote{A lot of standup comedy is based on noticing anomalies in
everyday life. ``Did you ever notice\ldots?'' New ideas come from doing
this about nontrivial things. Which may help explain why people's
reaction to a new idea is often the first half of laughing: Ha!}

If the answers seem strange, so much the better. Great work often has a
tincture of strangeness. You see this from painting to math. It would be
affected to try to manufacture it, but if it appears, embrace it.

Boldly chase outlier ideas, even if other people aren't interested in
them --- in fact, especially if they aren't. If you're excited about
some possibility that everyone else ignores, and you have enough
expertise to say precisely what they're all overlooking, that's as good
a bet as you'll find. \footnote{That second qualifier is critical. If you're excited about
something most authorities discount, but you can't give a more precise
explanation than ``they don't get it,'' then you're starting to drift
into the territory of cranks.}

Four steps: choose a field, learn enough to get to the frontier, notice
gaps, explore promising ones. This is how practically everyone who's
done great work has done it, from painters to physicists.

Steps two and four will require hard work. It may not be possible to
prove that you have to work hard to do great things, but the empirical
evidence is on the scale of the evidence for mortality. That's why it's
essential to work on something you're deeply interested in. Interest
will drive you to work harder than mere diligence ever could.

The three most powerful motives are curiosity, delight, and the desire
to do something impressive. Sometimes they converge, and that
combination is the most powerful of all.

The big prize is to discover a new fractal bud. You notice a crack in
the surface of knowledge, pry it open, and there's a whole world inside.

Let's talk a little more about the complicated business of figuring out
what to work on. The main reason it's hard is that you can't tell what
most kinds of work are like except by doing them. Which means the four
steps overlap: you may have to work at something for years before you
know how much you like it or how good you are at it. And in the meantime
you're not doing, and thus not learning about, most other kinds of work.
So in the worst case you choose late based on very incomplete
information. \footnote{Finding something to work on is not simply a matter of finding a
match between the current version of you and a list of known problems.
You'll often have to coevolve with the problem. That's why it can
sometimes be so hard to figure out what to work on. The search space is
huge. It's the cartesian product of all possible types of work, both
known and yet to be discovered, and all possible future versions of you.

There's no way you could search this whole space, so you have to rely on
heuristics to generate promising paths through it and hope the best
matches will be clustered. Which they will not always be; different
types of work have been collected together as much by accidents of
history as by the intrinsic similarities between them.}

The nature of ambition exacerbates this problem. Ambition comes in two
forms, one that precedes interest in the subject and one that grows out
of it. Most people who do great work have a mix, and the more you have
of the former, the harder it will be to decide what to do.

The educational systems in most countries pretend it's easy. They expect
you to commit to a field long before you could know what it's really
like. And as a result an ambitious person on an optimal trajectory will
often read to the system as an instance of breakage.

It would be better if they at least admitted it --- if they admitted
that the system not only can't do much to help you figure out what to
work on, but is designed on the assumption that you'll somehow magically
guess as a teenager. They don't tell you, but I will: when it comes to
figuring out what to work on, you're on your own. Some people get lucky
and do guess correctly, but the rest will find themselves scrambling
diagonally across tracks laid down on the assumption that everyone does.

What should you do if you're young and ambitious but don't know what to
work on? What you should \emph{not} do is drift along passively,
assuming the problem will solve itself. You need to take action. But
there is no systematic procedure you can follow. When you read
biographies of people who've done great work, it's remarkable how much
luck is involved. They discover what to work on as a result of a chance
meeting, or by reading a book they happen to pick up. So you need to
make yourself a big target for luck, and the way to do that is to be
curious. Try lots of things, meet lots of people, read lots of books,
ask lots of questions. \footnote{There are many reasons curious people are more likely to do
great work, but one of the more subtle is that, by casting a wide net,
they're more likely to find the right thing to work on in the first
place.}

When in doubt, optimize for interestingness. Fields change as you learn
more about them. What mathematicians do, for example, is very different
from what you do in high school math classes. So you need to give
different types of work a chance to show you what they're like. But a
field should become \emph{increasingly} interesting as you learn more
about it. If it doesn't, it's probably not for you.

Don't worry if you find you're interested in different things than other
people. The stranger your tastes in interestingness, the better. Strange
tastes are often strong ones, and a strong taste for work means you'll
be productive. And you're more likely to find new things if you're
looking where few have looked before.

One sign that you're suited for some kind of work is when you like even
the parts that other people find tedious or frightening.

But fields aren't people; you don't owe them any loyalty. If in the
course of working on one thing you discover another that's more
exciting, don't be afraid to switch.

If you're making something for people, make sure it's something they
actually want. The best way to do this is to make something you yourself
want. Write the story you want to read; build the tool you want to use.
Since your friends probably have similar interests, this will also get
you your initial audience.

This \emph{should} follow from the excitingness rule. Obviously the most
exciting story to write will be the one you want to read. The reason I
mention this case explicitly is that so many people get it wrong.
Instead of making what they want, they try to make what some imaginary,
more sophisticated audience wants. And once you go down that route,
you're lost. \footnote{It can also be dangerous to make things for an audience you feel
is less sophisticated than you, if that causes you to talk down to them.
You can make a lot of money doing that, if you do it in a sufficiently
cynical way, but it's not the route to great work. Not that anyone using
this m.o. would care.}

There are a lot of forces that will lead you astray when you're trying
to figure out what to work on. Pretentiousness, fashion, fear, money,
politics, other people's wishes, eminent frauds. But if you stick to
what you find genuinely interesting, you'll be proof against all of
them. If you're interested, you're not astray.

Following your interests may sound like a rather passive strategy, but
in practice it usually means following them past all sorts of obstacles.
You usually have to risk rejection and failure. So it does take a good
deal of boldness.

But while you need boldness, you don't usually need much planning. In
most cases the recipe for doing great work is simply: work hard on
excitingly ambitious projects, and something good will come of it.
Instead of making a plan and then executing it, you just try to preserve
certain invariants.

The trouble with planning is that it only works for achievements you can
describe in advance. You can win a gold medal or get rich by deciding to
as a child and then tenaciously pursuing that goal, but you can't
discover natural selection that way.

I think for most people who want to do great work, the right strategy is
not to plan too much. At each stage do whatever seems most interesting
and gives you the best options for the future. I call this approach
``staying upwind.'' This is how most people who've done great work seem
to have done it.

Even when you've found something exciting to work on, working on it is
not always straightforward. There will be times when some new idea makes
you leap out of bed in the morning and get straight to work. But there
will also be plenty of times when things aren't like that.

You don't just put out your sail and get blown forward by inspiration.
There are headwinds and currents and hidden shoals. So there's a
technique to working, just as there is to sailing.

For example, while you must work hard, it's possible to work too hard,
and if you do that you'll find you get diminishing returns: fatigue will
make you stupid, and eventually even damage your health. The point at
which work yields diminishing returns depends on the type. Some of the
hardest types you might only be able to do for four or five hours a day.

Ideally those hours will be contiguous. To the extent you can, try to
arrange your life so you have big blocks of time to work in. You'll shy
away from hard tasks if you know you might be interrupted.

It will probably be harder to start working than to keep working. You'll
often have to trick yourself to get over that initial threshold. Don't
worry about this; it's the nature of work, not a flaw in your character.
Work has a sort of activation energy, both per day and per project. And
since this threshold is fake in the sense that it's higher than the
energy required to keep going, it's ok to tell yourself a lie of
corresponding magnitude to get over it.

It's usually a mistake to lie to yourself if you want to do great work,
but this is one of the rare cases where it isn't. When I'm reluctant to
start work in the morning, I often trick myself by saying ``I'll just
read over what I've got so far.'' Five minutes later I've found
something that seems mistaken or incomplete, and I'm off.

Similar techniques work for starting new projects. It's ok to lie to
yourself about how much work a project will entail, for example. Lots of
great things began with someone saying ``How hard could it be?''

This is one case where the young have an advantage. They're more
optimistic, and even though one of the sources of their optimism is
ignorance, in this case ignorance can sometimes beat knowledge.

Try to finish what you start, though, even if it turns out to be more
work than you expected. Finishing things is not just an exercise in
tidiness or self-discipline. In many projects a lot of the best work
happens in what was meant to be the final stage.

Another permissible lie is to exaggerate the importance of what you're
working on, at least in your own mind. If that helps you discover
something new, it may turn out not to have been a lie after all.
\footnote{This idea I learned from Hardy's \emph{A Mathematician's
Apology}, which I recommend to anyone ambitious to do great work, in any
field.}

Since there are two senses of starting work --- per day and per project
--- there are also two forms of procrastination. Per-project
procrastination is far the more dangerous. You put off starting that
ambitious project from year to year because the time isn't quite right.
When you're procrastinating in units of years, you can get a lot not
done. \footnote{Just as we overestimate what we can do in a day and
underestimate what we can do over several years, we overestimate the
damage done by procrastinating for a day and underestimate the damage
done by procrastinating for several years.}

One reason per-project procrastination is so dangerous is that it
usually camouflages itself as work. You're not just sitting around doing
nothing; you're working industriously on something else. So per-project
procrastination doesn't set off the alarms that per-day procrastination
does. You're too busy to notice it.

The way to beat it is to stop occasionally and ask yourself: Am I
working on what I most want to work on? When you're young it's ok if the
answer is sometimes no, but this gets increasingly dangerous as you get
older. \footnote{You can't usually get paid for doing exactly what you want,
especially early on. There are two options: get paid for doing work
close to what you want and hope to push it closer, or get paid for doing
something else entirely and do your own projects on the side. Both can
work, but both have drawbacks: in the first approach your work is
compromised by default, and in the second you have to fight to get time
to do it.}

Great work usually entails spending what would seem to most people an
unreasonable amount of time on a problem. You can't think of this time
as a cost, or it will seem too high. You have to find the work
sufficiently engaging as it's happening.

There may be some jobs where you have to work diligently for years at
things you hate before you get to the good part, but this is not how
great work happens. Great work happens by focusing consistently on
something you're genuinely interested in. When you pause to take stock,
you're surprised how far you've come.

The reason we're surprised is that we underestimate the cumulative
effect of work. Writing a page a day doesn't sound like much, but if you
do it every day you'll write a book a year. That's the key: consistency.
People who do great things don't get a lot done every day. They get
something done, rather than nothing.

If you do work that compounds, you'll get exponential growth. Most
people who do this do it unconsciously, but it's worth stopping to think
about. Learning, for example, is an instance of this phenomenon: the
more you learn about something, the easier it is to learn more. Growing
an audience is another: the more fans you have, the more new fans
they'll bring you.

The trouble with exponential growth is that the curve feels flat in the
beginning. It isn't; it's still a wonderful exponential curve. But we
can't grasp that intuitively, so we underrate exponential growth in its
early stages.

Something that grows exponentially can become so valuable that it's
worth making an extraordinary effort to get it started. But since we
underrate exponential growth early on, this too is mostly done
unconsciously: people push through the initial, unrewarding phase of
learning something new because they know from experience that learning
new things always takes an initial push, or they grow their audience one
fan at a time because they have nothing better to do. If people
consciously realized they could invest in exponential growth, many more
would do it.

Work doesn't just happen when you're trying to. There's a kind of
undirected thinking you do when walking or taking a shower or lying in
bed that can be very powerful. By letting your mind wander a little,
you'll often solve problems you were unable to solve by frontal attack.

You have to be working hard in the normal way to benefit from this
phenomenon, though. You can't just walk around daydreaming. The
daydreaming has to be interleaved with deliberate work that feeds it
questions. \footnote{If you set your life up right, it will deliver the focus-relax
cycle automatically. The perfect setup is an office you work in and that
you walk to and from.}

Everyone knows to avoid distractions at work, but it's also important to
avoid them in the other half of the cycle. When you let your mind
wander, it wanders to whatever you care about most at that moment. So
avoid the kind of distraction that pushes your work out of the top spot,
or you'll waste this valuable type of thinking on the distraction
instead. (Exception: Don't avoid love.)

Consciously cultivate your taste in the work done in your field. Until
you know which is the best and what makes it so, you don't know what
you're aiming for.

And that \emph{is} what you're aiming for, because if you don't try to
be the best, you won't even be good. This observation has been made by
so many people in so many different fields that it might be worth
thinking about why it's true. It could be because ambition is a
phenomenon where almost all the error is in one direction --- where
almost all the shells that miss the target miss by falling short. Or it
could be because ambition to be the best is a qualitatively different
thing from ambition to be good. Or maybe being good is simply too vague
a standard. Probably all three are true. \footnote{There may be some very unworldly people who do great work
without consciously trying to. If you want to expand this rule to cover
that case, it becomes: Don't try to be anything except the best.}

Fortunately there's a kind of economy of scale here. Though it might
seem like you'd be taking on a heavy burden by trying to be the best, in
practice you often end up net ahead. It's exciting, and also strangely
liberating. It simplifies things. In some ways it's easier to try to be
the best than to try merely to be good.

One way to aim high is to try to make something that people will care
about in a hundred years. Not because their opinions matter more than
your contemporaries', but because something that still seems good in a
hundred years is more likely to be genuinely good.

Don't try to work in a distinctive style. Just try to do the best job
you can; you won't be able to help doing it in a distinctive way.

Style is doing things in a distinctive way without trying to. Trying to
is affectation.

Affectation is in effect to pretend that someone other than you is doing
the work. You adopt an impressive but fake persona, and while you're
pleased with the impressiveness, the fakeness is what shows in the work.
\footnote{This gets more complicated in work like acting, where the goal
is to adopt a fake persona. But even here it's possible to be affected.
Perhaps the rule in such fields should be to avoid \emph{unintentional}
affectation.}

The temptation to be someone else is greatest for the young. They often
feel like nobodies. But you never need to worry about that problem,
because it's self-solving if you work on sufficiently ambitious
projects. If you succeed at an ambitious project, you're not a nobody;
you're the person who did it. So just do the work and your identity will
take care of itself.

``Avoid affectation'' is a useful rule so far as it goes, but how would
you express this idea positively? How would you say what to be, instead
of what not to be? The best answer is earnest. If you're earnest you
avoid not just affectation but a whole set of similar vices.

The core of being earnest is being intellectually honest. We're taught
as children to be honest as an unselfish virtue --- as a kind of
sacrifice. But in fact it's a source of power too. To see new ideas, you
need an exceptionally sharp eye for the truth. You're trying to see more
truth than others have seen so far. And how can you have a sharp eye for
the truth if you're intellectually dishonest?

One way to avoid intellectual dishonesty is to maintain a slight
positive pressure in the opposite direction. Be aggressively willing to
admit that you're mistaken. Once you've admitted you were mistaken about
something, you're free. Till then you have to carry it.
\footnote{It's safe to have beliefs that you treat as unquestionable if
and only if they're also unfalsifiable. For example, it's safe to have
the principle that everyone should be treated equally under the law,
because a sentence with a ``should'' in it isn't really a statement
about the world and is therefore hard to disprove. And if there's no
evidence that could disprove one of your principles, there can't be any
facts you'd need to ignore in order to preserve it.}

Another more subtle component of earnestness is informality. Informality
is much more important than its grammatically negative name implies.
It's not merely the absence of something. It means focusing on what
matters instead of what doesn't.

What formality and affectation have in common is that as well as doing
the work, you're trying to seem a certain way as you're doing it. But
any energy that goes into how you seem comes out of being good. That's
one reason nerds have an advantage in doing great work: they expend
little effort on seeming anything. In fact that's basically the
definition of a nerd.

Nerds have a kind of innocent boldness that's exactly what you need in
doing great work. It's not learned; it's preserved from childhood. So
hold onto it. Be the one who puts things out there rather than the one
who sits back and offers sophisticated-sounding criticisms of them.
``It's easy to criticize'' is true in the most literal sense, and the
route to great work is never easy.

There may be some jobs where it's an advantage to be cynical and
pessimistic, but if you want to do great work it's an advantage to be
optimistic, even though that means you'll risk looking like a fool
sometimes. There's an old tradition of doing the opposite. The Old
Testament says it's better to keep quiet lest you look like a fool. But
that's advice for \emph{seeming} smart. If you actually want to discover
new things, it's better to take the risk of telling people your ideas.

Some people are naturally earnest, and with others it takes a conscious
effort. Either kind of earnestness will suffice. But I doubt it would be
possible to do great work without being earnest. It's so hard to do even
if you are. You don't have enough margin for error to accommodate the
distortions introduced by being affected, intellectually dishonest,
orthodox, fashionable, or cool. \footnote{Affectation is easier to cure than intellectual dishonesty.
Affectation is often a shortcoming of the young that burns off in time,
while intellectual dishonesty is more of a character flaw.}

Great work is consistent not only with who did it, but with itself. It's
usually all of a piece. So if you face a decision in the middle of
working on something, ask which choice is more consistent.

You may have to throw things away and redo them. You won't necessarily
have to, but you have to be willing to. And that can take some effort;
when there's something you need to redo, status quo bias and laziness
will combine to keep you in denial about it. To beat this ask: If I'd
already made the change, would I want to revert to what I have now?

Have the confidence to cut. Don't keep something that doesn't fit just
because you're proud of it, or because it cost you a lot of effort.

Indeed, in some kinds of work it's good to strip whatever you're doing
to its essence. The result will be more concentrated; you'll understand
it better; and you won't be able to lie to yourself about whether
there's anything real there.

Mathematical elegance may sound like a mere metaphor, drawn from the
arts. That's what I thought when I first heard the term ``elegant''
applied to a proof. But now I suspect it's conceptually prior --- that
the main ingredient in artistic elegance is mathematical elegance. At
any rate it's a useful standard well beyond math.

Elegance can be a long-term bet, though. Laborious solutions will often
have more prestige in the short term. They cost a lot of effort and
they're hard to understand, both of which impress people, at least
temporarily.

Whereas some of the very best work will seem like it took comparatively
little effort, because it was in a sense already there. It didn't have
to be built, just seen. It's a very good sign when it's hard to say
whether you're creating something or discovering it.

When you're doing work that could be seen as either creation or
discovery, err on the side of discovery. Try thinking of yourself as a
mere conduit through which the ideas take their natural shape.

(Strangely enough, one exception is the problem of choosing a problem to
work on. This is usually seen as search, but in the best case it's more
like creating something. In the best case you create the field in the
process of exploring it.)

Similarly, if you're trying to build a powerful tool, make it
gratuitously unrestrictive. A powerful tool almost by definition will be
used in ways you didn't expect, so err on the side of eliminating
restrictions, even if you don't know what the benefit will be.

Great work will often be tool-like in the sense of being something
others build on. So it's a good sign if you're creating ideas that
others could use, or exposing questions that others could answer. The
best ideas have implications in many different areas.

If you express your ideas in the most general form, they'll be truer
than you intended.

True by itself is not enough, of course. Great ideas have to be true and
new. And it takes a certain amount of ability to see new ideas even once
you've learned enough to get to one of the frontiers of knowledge.

In English we give this ability names like originality, creativity, and
imagination. And it seems reasonable to give it a separate name, because
it does seem to some extent a separate skill. It's possible to have a
great deal of ability in other respects --- to have a great deal of
what's often called \emph{technical} ability --- and yet not have much
of this.

I've never liked the term ``creative process.'' It seems misleading.
Originality isn't a process, but a habit of mind. Original thinkers
throw off new ideas about whatever they focus on, like an angle grinder
throwing off sparks. They can't help it.

If the thing they're focused on is something they don't understand very
well, these new ideas might not be good. One of the most original
thinkers I know decided to focus on dating after he got divorced. He
knew roughly as much about dating as the average 15 year old, and the
results were spectacularly colorful. But to see originality separated
from expertise like that made its nature all the more clear.

I don't know if it's possible to cultivate originality, but there are
definitely ways to make the most of however much you have. For example,
you're much more likely to have original ideas when you're working on
something. Original ideas don't come from trying to have original ideas.
They come from trying to build or understand something slightly too
difficult. \footnote{Obviously you don't have to be working at the exact moment you
have the idea, but you'll probably have been working fairly recently.}

Talking or writing about the things you're interested in is a good way
to generate new ideas. When you try to put ideas into words, a missing
idea creates a sort of vacuum that draws it out of you. Indeed, there's
a kind of thinking that can only be done by writing.

Changing your context can help. If you visit a new place, you'll often
find you have new ideas there. The journey itself often dislodges them.
But you may not have to go far to get this benefit. Sometimes it's
enough just to go for a walk. \footnote{Some say psychoactive drugs have a similar effect. I'm
skeptical, but also almost totally ignorant of their effects.}

It also helps to travel in topic space. You'll have more new ideas if
you explore lots of different topics, partly because it gives the angle
grinder more surface area to work on, and partly because analogies are
an especially fruitful source of new ideas.

Don't divide your attention \emph{evenly} between many topics though, or
you'll spread yourself too thin. You want to distribute it according to
something more like a power law. \footnote{For example you might give the nth most important topic
(m-1)/m\^{}n of your attention, for some m \textgreater{} 1. You
couldn't allocate your attention so precisely, of course, but this at
least gives an idea of a reasonable distribution.} Be
professionally curious about a few topics and idly curious about many
more.

Curiosity and originality are closely related. Curiosity feeds
originality by giving it new things to work on. But the relationship is
closer than that. Curiosity is itself a kind of originality; it's
roughly to questions what originality is to answers. And since questions
at their best are a big component of answers, curiosity at its best is a
creative force.

Having new ideas is a strange game, because it usually consists of
seeing things that were right under your nose. Once you've seen a new
idea, it tends to seem obvious. Why did no one think of this before?

When an idea seems simultaneously novel and obvious, it's probably a
good one.

Seeing something obvious sounds easy. And yet empirically having new
ideas is hard. What's the source of this apparent contradiction? It's
that seeing the new idea usually requires you to change the way you look
at the world. We see the world through models that both help and
constrain us. When you fix a broken model, new ideas become obvious. But
noticing and fixing a broken model is hard. That's how new ideas can be
both obvious and yet hard to discover: they're easy to see after you do
something hard.

One way to discover broken models is to be stricter than other people.
Broken models of the world leave a trail of clues where they bash
against reality. Most people don't want to see these clues. It would be
an understatement to say that they're attached to their current model;
it's what they think in; so they'll tend to ignore the trail of clues
left by its breakage, however conspicuous it may seem in retrospect.

To find new ideas you have to seize on signs of breakage instead of
looking away. That's what Einstein did. He was able to see the wild
implications of Maxwell's equations not so much because he was looking
for new ideas as because he was stricter.

The other thing you need is a willingness to break rules. Paradoxical as
it sounds, if you want to fix your model of the world, it helps to be
the sort of person who's comfortable breaking rules. From the point of
view of the old model, which everyone including you initially shares,
the new model usually breaks at least implicit rules.

Few understand the degree of rule-breaking required, because new ideas
seem much more conservative once they succeed. They seem perfectly
reasonable once you're using the new model of the world they brought
with them. But they didn't at the time; it took the greater part of a
century for the heliocentric model to be generally accepted, even among
astronomers, because it felt so wrong.

Indeed, if you think about it, a good new idea has to seem bad to most
people, or someone would have already explored it. So what you're
looking for is ideas that seem crazy, but the right kind of crazy. How
do you recognize these? You can't with certainty. Often ideas that seem
bad are bad. But ideas that are the right kind of crazy tend to be
exciting; they're rich in implications; whereas ideas that are merely
bad tend to be depressing.

There are two ways to be comfortable breaking rules: to enjoy breaking
them, and to be indifferent to them. I call these two cases being
aggressively and passively independent-minded.

The aggressively independent-minded are the naughty ones. Rules don't
merely fail to stop them; breaking rules gives them additional energy.
For this sort of person, delight at the sheer audacity of a project
sometimes supplies enough activation energy to get it started.

The other way to break rules is not to care about them, or perhaps even
to know they exist. This is why novices and outsiders often make new
discoveries; their ignorance of a field's assumptions acts as a source
of temporary passive independent-mindedness. Aspies also seem to have a
kind of immunity to conventional beliefs. Several I know say that this
helps them to have new ideas.

Strictness plus rule-breaking sounds like a strange combination. In
popular culture they're opposed. But popular culture has a broken model
in this respect. It implicitly assumes that issues are trivial ones, and
in trivial matters strictness and rule-breaking \emph{are} opposed. But
in questions that really matter, only rule-breakers can be truly strict.

An overlooked idea often doesn't lose till the semifinals. You do see
it, subconsciously, but then another part of your subconscious shoots it
down because it would be too weird, too risky, too much work, too
controversial. This suggests an exciting possibility: if you could turn
off such filters, you could see more new ideas.

One way to do that is to ask what would be good ideas for \emph{someone
else} to explore. Then your subconscious won't shoot them down to
protect you.

You could also discover overlooked ideas by working in the other
direction: by starting from what's obscuring them. Every cherished but
mistaken principle is surrounded by a dead zone of valuable ideas that
are unexplored because they contradict it.

Religions are collections of cherished but mistaken principles. So
anything that can be described either literally or metaphorically as a
religion will have valuable unexplored ideas in its shadow. Copernicus
and Darwin both made discoveries of this type. \footnote{The principles defining a religion have to be mistaken.
Otherwise anyone might adopt them, and there would be nothing to
distinguish the adherents of the religion from everyone else.}

What are people in your field religious about, in the sense of being too
attached to some principle that might not be as self-evident as they
think? What becomes possible if you discard it?

People show much more originality in solving problems than in deciding
which problems to solve. Even the smartest can be surprisingly
conservative when deciding what to work on. People who'd never dream of
being fashionable in any other way get sucked into working on
fashionable problems.

One reason people are more conservative when choosing problems than
solutions is that problems are bigger bets. A problem could occupy you
for years, while exploring a solution might only take days. But even so
I think most people are too conservative. They're not merely responding
to risk, but to fashion as well. Unfashionable problems are undervalued.

One of the most interesting kinds of unfashionable problem is the
problem that people think has been fully explored, but hasn't. Great
work often takes something that already exists and shows its latent
potential. Durer and Watt both did this. So if you're interested in a
field that others think is tapped out, don't let their skepticism deter
you. People are often wrong about this.

Working on an unfashionable problem can be very pleasing. There's no
hype or hurry. Opportunists and critics are both occupied elsewhere. The
existing work often has an old-school solidity. And there's a satisfying
sense of economy in cultivating ideas that would otherwise be wasted.

But the most common type of overlooked problem is not explicitly
unfashionable in the sense of being out of fashion. It just doesn't seem
to matter as much as it actually does. How do you find these? By being
self-indulgent --- by letting your curiosity have its way, and tuning
out, at least temporarily, the little voice in your head that says you
should only be working on ``important'' problems.

You do need to work on important problems, but almost everyone is too
conservative about what counts as one. And if there's an important but
overlooked problem in your neighborhood, it's probably already on your
subconscious radar screen. So try asking yourself: if you were going to
take a break from ``serious'' work to work on something just because it
would be really interesting, what would you do? The answer is probably
more important than it seems.

Originality in choosing problems seems to matter even more than
originality in solving them. That's what distinguishes the people who
discover whole new fields. So what might seem to be merely the initial
step --- deciding what to work on --- is in a sense the key to the whole
game.

Few grasp this. One of the biggest misconceptions about new ideas is
about the ratio of question to answer in their composition. People think
big ideas are answers, but often the real insight was in the question.

Part of the reason we underrate questions is the way they're used in
schools. In schools they tend to exist only briefly before being
answered, like unstable particles. But a really good question can be
much more than that. A really good question is a partial discovery. How
do new species arise? Is the force that makes objects fall to earth the
same as the one that keeps planets in their orbits? By even asking such
questions you were already in excitingly novel territory.

Unanswered questions can be uncomfortable things to carry around with
you. But the more you're carrying, the greater the chance of noticing a
solution --- or perhaps even more excitingly, noticing that two
unanswered questions are the same.

Sometimes you carry a question for a long time. Great work often comes
from returning to a question you first noticed years before --- in your
childhood, even --- and couldn't stop thinking about. People talk a lot
about the importance of keeping your youthful dreams alive, but it's
just as important to keep your youthful questions alive.
\footnote{It might be a good exercise to try writing down a list of
questions you wondered about in your youth. You might find you're now in
a position to do something about some of them.}

This is one of the places where actual expertise differs most from the
popular picture of it. In the popular picture, experts are certain. But
actually the more puzzled you are, the better, so long as (a) the things
you're puzzled about matter, and (b) no one else understands them
either.

Think about what's happening at the moment just before a new idea is
discovered. Often someone with sufficient expertise is puzzled about
something. Which means that originality consists partly of puzzlement
--- of confusion! You have to be comfortable enough with the world being
full of puzzles that you're willing to see them, but not so comfortable
that you don't want to solve them. \footnote{The connection between originality and uncertainty causes a
strange phenomenon: because the conventional-minded are more certain
than the independent-minded, this tends to give them the upper hand in
disputes, even though they're generally stupider.

\begin{quote}
The best lack all conviction, while the worst\\
Are full of passionate intensity.
\end{quote}}

It's a great thing to be rich in unanswered questions. And this is one
of those situations where the rich get richer, because the best way to
acquire new questions is to try answering existing ones. Questions don't
just lead to answers, but also to more questions.

The best questions grow in the answering. You notice a thread protruding
from the current paradigm and try pulling on it, and it just gets longer
and longer. So don't require a question to be obviously big before you
try answering it. You can rarely predict that. It's hard enough even to
notice the thread, let alone to predict how much will unravel if you
pull on it.

It's better to be promiscuously curious --- to pull a little bit on a
lot of threads, and see what happens. Big things start small. The
initial versions of big things were often just experiments, or side
projects, or talks, which then grew into something bigger. So start lots
of small things.

Being prolific is underrated. The more different things you try, the
greater the chance of discovering something new. Understand, though,
that trying lots of things will mean trying lots of things that don't
work. You can't have a lot of good ideas without also having a lot of
bad ones. \footnote{Derived from Linus Pauling's ``If you want to have good ideas,
you must have many ideas.''}

Though it sounds more responsible to begin by studying everything that's
been done before, you'll learn faster and have more fun by trying stuff.
And you'll understand previous work better when you do look at it. So
err on the side of starting. Which is easier when starting means
starting small; those two ideas fit together like two puzzle pieces.

How do you get from starting small to doing something great? By making
successive versions. Great things are almost always made in successive
versions. You start with something small and evolve it, and the final
version is both cleverer and more ambitious than anything you could have
planned.

It's particularly useful to make successive versions when you're making
something for people --- to get an initial version in front of them
quickly, and then evolve it based on their response.

Begin by trying the simplest thing that could possibly work.
Surprisingly often, it does. If it doesn't, this will at least get you
started.

Don't try to cram too much new stuff into any one version. There are
names for doing this with the first version (taking too long to ship)
and the second (the second system effect), but these are both merely
instances of a more general principle.

An early version of a new project will sometimes be dismissed as a toy.
It's a good sign when people do this. That means it has everything a new
idea needs except scale, and that tends to follow.
\footnote{Attacking a project as a ``toy'' is similar to attacking a
statement as ``inappropriate.'' It means that no more substantial
criticism can be made to stick.}

The alternative to starting with something small and evolving it is to
plan in advance what you're going to do. And planning does usually seem
the more responsible choice. It sounds more organized to say ``we're
going to do x and then y and then z'' than ``we're going to try x and
see what happens.'' And it is more \emph{organized}; it just doesn't
work as well.

Planning per se isn't good. It's sometimes necessary, but it's a
necessary evil --- a response to unforgiving conditions. It's something
you have to do because you're working with inflexible media, or because
you need to coordinate the efforts of a lot of people. If you keep
projects small and use flexible media, you don't have to plan as much,
and your designs can evolve instead.

Take as much risk as you can afford. In an efficient market, risk is
proportionate to reward, so don't look for certainty, but for a bet with
high expected value. If you're not failing occasionally, you're probably
being too conservative.

Though conservatism is usually associated with the old, it's the young
who tend to make this mistake. Inexperience makes them fear risk, but
it's when you're young that you can afford the most.

Even a project that fails can be valuable. In the process of working on
it, you'll have crossed territory few others have seen, and encountered
questions few others have asked. And there's probably no better source
of questions than the ones you encounter in trying to do something
slightly too hard.

Use the advantages of youth when you have them, and the advantages of
age once you have those. The advantages of youth are energy, time,
optimism, and freedom. The advantages of age are knowledge, efficiency,
money, and power. With effort you can acquire some of the latter when
young and keep some of the former when old.

The old also have the advantage of knowing which advantages they have.
The young often have them without realizing it. The biggest is probably
time. The young have no idea how rich they are in time. The best way to
turn this time to advantage is to use it in slightly frivolous ways: to
learn about something you don't need to know about, just out of
curiosity, or to try building something just because it would be cool,
or to become freakishly good at something.

That ``slightly'' is an important qualification. Spend time lavishly
when you're young, but don't simply waste it. There's a big difference
between doing something you worry might be a waste of time and doing
something you know for sure will be. The former is at least a bet, and
possibly a better one than you think. \footnote{One way to tell whether you're wasting time is to ask if you're
producing or consuming. Writing computer games is less likely to be a
waste of time than playing them, and playing games where you create
something is less likely to be a waste of time than playing games where
you don't.}

The most subtle advantage of youth, or more precisely of inexperience,
is that you're seeing everything with fresh eyes. When your brain
embraces an idea for the first time, sometimes the two don't fit
together perfectly. Usually the problem is with your brain, but
occasionally it's with the idea. A piece of it sticks out awkwardly and
jabs you when you think about it. People who are used to the idea have
learned to ignore it, but you have the opportunity not to.
\footnote{Another related advantage is that if you haven't said anything
publicly yet, you won't be biased toward evidence that supports your
earlier conclusions. With sufficient integrity you could achieve eternal
youth in this respect, but few manage to. For most people, having
previously published opinions has an effect similar to ideology, just in
quantity 1.}

So when you're learning about something for the first time, pay
attention to things that seem wrong or missing. You'll be tempted to
ignore them, since there's a 99\% chance the problem is with you. And
you may have to set aside your misgivings temporarily to keep
progressing. But don't forget about them. When you've gotten further
into the subject, come back and check if they're still there. If they're
still viable in the light of your present knowledge, they probably
represent an undiscovered idea.

One of the most valuable kinds of knowledge you get from experience is
to know what you \emph{don't} have to worry about. The young know all
the things that could matter, but not their relative importance. So they
worry equally about everything, when they should worry much more about a
few things and hardly at all about the rest.

But what you don't know is only half the problem with inexperience. The
other half is what you do know that ain't so. You arrive at adulthood
with your head full of nonsense --- bad habits you've acquired and false
things you've been taught --- and you won't be able to do great work
till you clear away at least the nonsense in the way of whatever type of
work you want to do.

Much of the nonsense left in your head is left there by schools. We're
so used to schools that we unconsciously treat going to school as
identical with learning, but in fact schools have all sorts of strange
qualities that warp our ideas about learning and thinking.

For example, schools induce passivity. Since you were a small child,
there was an authority at the front of the class telling all of you what
you had to learn and then measuring whether you did. But neither classes
nor tests are intrinsic to learning; they're just artifacts of the way
schools are usually designed.

The sooner you overcome this passivity, the better. If you're still in
school, try thinking of your education as your project, and your
teachers as working for you rather than vice versa. That may seem a
stretch, but it's not merely some weird thought experiment. It's the
truth economically, and in the best case it's the truth intellectually
as well. The best teachers don't want to be your bosses. They'd prefer
it if you pushed ahead, using them as a source of advice, rather than
being pulled by them through the material.

Schools also give you a misleading impression of what work is like. In
school they tell you what the problems are, and they're almost always
soluble using no more than you've been taught so far. In real life you
have to figure out what the problems are, and you often don't know if
they're soluble at all.

But perhaps the worst thing schools do to you is train you to win by
hacking the test. You can't do great work by doing that. You can't trick
God. So stop looking for that kind of shortcut. The way to beat the
system is to focus on problems and solutions that others have
overlooked, not to skimp on the work itself.

Don't think of yourself as dependent on some gatekeeper giving you a
``big break.'' Even if this were true, the best way to get it would be
to focus on doing good work rather than chasing influential people.

And don't take rejection by committees to heart. The qualities that
impress admissions officers and prize committees are quite different
from those required to do great work. The decisions of selection
committees are only meaningful to the extent that they're part of a
feedback loop, and very few are.

People new to a field will often copy existing work. There's nothing
inherently bad about that. There's no better way to learn how something
works than by trying to reproduce it. Nor does copying necessarily make
your work unoriginal. Originality is the presence of new ideas, not the
absence of old ones.

There's a good way to copy and a bad way. If you're going to copy
something, do it openly instead of furtively, or worse still,
unconsciously. This is what's meant by the famously misattributed phrase
``Great artists steal.'' The really dangerous kind of copying, the kind
that gives copying a bad name, is the kind that's done without realizing
it, because you're nothing more than a train running on tracks laid down
by someone else. But at the other extreme, copying can be a sign of
superiority rather than subordination. \footnote{In the early 1630s Daniel Mytens made a painting of Henrietta
Maria handing a laurel wreath to Charles I. Van Dyck then painted his
own version to show how much better he was.}

In many fields it's almost inevitable that your early work will be in
some sense based on other people's. Projects rarely arise in a vacuum.
They're usually a reaction to previous work. When you're first starting
out, you don't have any previous work; if you're going to react to
something, it has to be someone else's. Once you're established, you can
react to your own. But while the former gets called derivative and the
latter doesn't, structurally the two cases are more similar than they
seem.

Oddly enough, the very novelty of the most novel ideas sometimes makes
them seem at first to be more derivative than they are. New discoveries
often have to be conceived initially as variations of existing things,
\emph{even by their discoverers}, because there isn't yet the conceptual
vocabulary to express them.

There are definitely some dangers to copying, though. One is that you'll
tend to copy old things --- things that were in their day at the
frontier of knowledge, but no longer are.

And when you do copy something, don't copy every feature of it. Some
will make you ridiculous if you do. Don't copy the manner of an eminent
50 year old professor if you're 18, for example, or the idiom of a
Renaissance poem hundreds of years later.

Some of the features of things you admire are flaws they succeeded
despite. Indeed, the features that are easiest to imitate are the most
likely to be the flaws.

This is particularly true for behavior. Some talented people are jerks,
and this sometimes makes it seem to the inexperienced that being a jerk
is part of being talented. It isn't; being talented is merely how they
get away with it.

One of the most powerful kinds of copying is to copy something from one
field into another. History is so full of chance discoveries of this
type that it's probably worth giving chance a hand by deliberately
learning about other kinds of work. You can take ideas from quite
distant fields if you let them be metaphors.

Negative examples can be as inspiring as positive ones. In fact you can
sometimes learn more from things done badly than from things done well;
sometimes it only becomes clear what's needed when it's missing.

If a lot of the best people in your field are collected in one place,
it's usually a good idea to visit for a while. It will increase your
ambition, and also, by showing you that these people are human, increase
your self-confidence. \footnote{I'm being deliberately vague about what a place is. As of this
writing, being in the same physical place has advantages that are hard
to duplicate, but that could change.}

If you're earnest you'll probably get a warmer welcome than you might
expect. Most people who are very good at something are happy to talk
about it with anyone who's genuinely interested. If they're really good
at their work, then they probably have a hobbyist's interest in it, and
hobbyists always want to talk about their hobbies.

It may take some effort to find the people who are really good, though.
Doing great work has such prestige that in some places, particularly
universities, there's a polite fiction that everyone is engaged in it.
And that is far from true. People within universities can't say so
openly, but the quality of the work being done in different departments
varies immensely. Some departments have people doing great work; others
have in the past; others never have.

Seek out the best colleagues. There are a lot of projects that can't be
done alone, and even if you're working on one that can be, it's good to
have other people to encourage you and to bounce ideas off.

Colleagues don't just affect your work, though; they also affect you. So
work with people you want to become like, because you will.

Quality is more important than quantity in colleagues. It's better to
have one or two great ones than a building full of pretty good ones. In
fact it's not merely better, but necessary, judging from history: the
degree to which great work happens in clusters suggests that one's
colleagues often make the difference between doing great work and not.

How do you know when you have sufficiently good colleagues? In my
experience, when you do, you know. Which means if you're unsure, you
probably don't. But it may be possible to give a more concrete answer
than that. Here's an attempt: sufficiently good colleagues offer
\emph{surprising} insights. They can see and do things that you can't.
So if you have a handful of colleagues good enough to keep you on your
toes in this sense, you're probably over the threshold.

Most of us can benefit from collaborating with colleagues, but some
projects require people on a larger scale, and starting one of those is
not for everyone. If you want to run a project like that, you'll have to
become a manager, and managing well takes aptitude and interest like any
other kind of work. If you don't have them, there is no middle path: you
must either force yourself to learn management as a second language, or
avoid such projects. \footnote{This is false when the work the other people have to do is very
constrained, as with SETI@home or Bitcoin. It may be possible to expand
the area in which it's false by defining similarly restricted protocols
with more freedom of action in the nodes.}

Husband your morale. It's the basis of everything when you're working on
ambitious projects. You have to nurture and protect it like a living
organism.

Morale starts with your view of life. You're more likely to do great
work if you're an optimist, and more likely to if you think of yourself
as lucky than if you think of yourself as a victim.

Indeed, work can to some extent protect you from your problems. If you
choose work that's pure, its very difficulties will serve as a refuge
from the difficulties of everyday life. If this is escapism, it's a very
productive form of it, and one that has been used by some of the
greatest minds in history.

Morale compounds via work: high morale helps you do good work, which
increases your morale and helps you do even better work. But this cycle
also operates in the other direction: if you're not doing good work,
that can demoralize you and make it even harder to. Since it matters so
much for this cycle to be running in the right direction, it can be a
good idea to switch to easier work when you're stuck, just so you start
to get something done.

One of the biggest mistakes ambitious people make is to allow setbacks
to destroy their morale all at once, like a balloon bursting. You can
inoculate yourself against this by explicitly considering setbacks a
part of your process. Solving hard problems always involves some
backtracking.

Doing great work is a depth-first search whose root node is the desire
to. So ``If at first you don't succeed, try, try again'' isn't quite
right. It should be: If at first you don't succeed, either try again, or
backtrack and then try again.

``Never give up'' is also not quite right. Obviously there are times
when it's the right choice to eject. A more precise version would be:
Never let setbacks panic you into backtracking more than you need to.
Corollary: Never abandon the root node.

It's not necessarily a bad sign if work is a struggle, any more than
it's a bad sign to be out of breath while running. It depends how fast
you're running. So learn to distinguish good pain from bad. Good pain is
a sign of effort; bad pain is a sign of damage.

An audience is a critical component of morale. If you're a scholar, your
audience may be your peers; in the arts, it may be an audience in the
traditional sense. Either way it doesn't need to be big. The value of an
audience doesn't grow anything like linearly with its size. Which is bad
news if you're famous, but good news if you're just starting out,
because it means a small but dedicated audience can be enough to sustain
you. If a handful of people genuinely love what you're doing, that's
enough.

To the extent you can, avoid letting intermediaries come between you and
your audience. In some types of work this is inevitable, but it's so
liberating to escape it that you might be better off switching to an
adjacent type if that will let you go direct. \footnote{Corollary: Building something that enables people to go around
intermediaries and engage directly with their audience is probably a
good idea.}

The people you spend time with will also have a big effect on your
morale. You'll find there are some who increase your energy and others
who decrease it, and the effect someone has is not always what you'd
expect. Seek out the people who increase your energy and avoid those who
decrease it. Though of course if there's someone you need to take care
of, that takes precedence.

Don't marry someone who doesn't understand that you need to work, or
sees your work as competition for your attention. If you're ambitious,
you need to work; it's almost like a medical condition; so someone who
won't let you work either doesn't understand you, or does and doesn't
care.

Ultimately morale is physical. You think with your body, so it's
important to take care of it. That means exercising regularly, eating
and sleeping well, and avoiding the more dangerous kinds of drugs.
Running and walking are particularly good forms of exercise because
they're good for thinking. \footnote{It may be helpful always to walk or run the same route, because
that frees attention for thinking. It feels that way to me, and there is
some historical evidence for it.}

People who do great work are not necessarily happier than everyone else,
but they're happier than they'd be if they didn't. In fact, if you're
smart and ambitious, it's dangerous \emph{not} to be productive. People
who are smart and ambitious but don't achieve much tend to become
bitter.

It's ok to want to impress other people, but choose the right people.
The opinion of people you respect is signal. Fame, which is the opinion
of a much larger group you might or might not respect, just adds noise.

The prestige of a type of work is at best a trailing indicator and
sometimes completely mistaken. If you do anything well enough, you'll
make it prestigious. So the question to ask about a type of work is not
how much prestige it has, but how well it could be done.

Competition can be an effective motivator, but don't let it choose the
problem for you; don't let yourself get drawn into chasing something
just because others are. In fact, don't let competitors make you do
anything much more specific than work harder.

Curiosity is the best guide. Your curiosity never lies, and it knows
more than you do about what's worth paying attention to.

Notice how often that word has come up. If you asked an oracle the
secret to doing great work and the oracle replied with a single word, my
bet would be on ``curiosity.''

That doesn't translate directly to advice. It's not enough just to be
curious, and you can't command curiosity anyway. But you can nurture it
and let it drive you.

Curiosity is the key to all four steps in doing great work: it will
choose the field for you, get you to the frontier, cause you to notice
the gaps in it, and drive you to explore them. The whole process is a
kind of dance with curiosity.

Believe it or not, I tried to make this essay as short as I could. But
its length at least means it acts as a filter. If you made it this far,
you must be interested in doing great work. And if so you're already
further along than you might realize, because the set of people willing
to want to is small.

The factors in doing great work are factors in the literal, mathematical
sense, and they are: ability, interest, effort, and luck. Luck by
definition you can't do anything about, so we can ignore that. And we
can assume effort, if you do in fact want to do great work. So the
problem boils down to ability and interest. Can you find a kind of work
where your ability and interest will combine to yield an explosion of
new ideas?

Here there are grounds for optimism. There are so many different ways to
do great work, and even more that are still undiscovered. Out of all
those different types of work, the one you're most suited for is
probably a pretty close match. Probably a comically close match. It's
just a question of finding it, and how far into it your ability and
interest can take you. And you can only answer that by trying.

Many more people could try to do great work than do. What holds them
back is a combination of modesty and fear. It seems presumptuous to try
to be Newton or Shakespeare. It also seems hard; surely if you tried
something like that, you'd fail. Presumably the calculation is rarely
explicit. Few people consciously decide not to try to do great work. But
that's what's going on subconsciously; they shy away from the question.

So I'm going to pull a sneaky trick on you. Do you want to do great
work, or not? Now you have to decide consciously. Sorry about that. I
wouldn't have done it to a general audience. But we already know you're
interested.

Don't worry about being presumptuous. You don't have to tell anyone. And
if it's too hard and you fail, so what? Lots of people have worse
problems than that. In fact you'll be lucky if it's the worst problem
you have.

Yes, you'll have to work hard. But again, lots of people have to work
hard. And if you're working on something you find very interesting,
which you necessarily will if you're on the right path, the work will
probably feel less burdensome than a lot of your peers'.

The discoveries are out there, waiting to be made. Why not by you?


\textbf{Thanks} to Trevor Blackwell, Daniel Gackle, Pam Graham, Tom
Howard, Patrick Hsu, Steve Huffman, Jessica Livingston, Henry
Lloyd-Baker, Bob Metcalfe, Ben Miller, Robert Morris, Michael Nielsen,
Courtenay Pipkin, Joris Poort, Mieke Roos, Rajat Suri, Harj Taggar,
Garry Tan, and my younger son for suggestions and for reading drafts.
\chapter{Superlinear Returns}

October 2023

One of the most important things I didn't understand about the world
when I was a child is the degree to which the returns for performance
are superlinear.

Teachers and coaches implicitly told us the returns were linear. ``You
get out,'' I heard a thousand times, ``what you put in.'' They meant
well, but this is rarely true. If your product is only half as good as
your competitor's, you don't get half as many customers. You get no
customers, and you go out of business.

It's obviously true that the returns for performance are superlinear in
business. Some think this is a flaw of capitalism, and that if we
changed the rules it would stop being true. But superlinear returns for
performance are a feature of the world, not an artifact of rules we've
invented. We see the same pattern in fame, power, military victories,
knowledge, and even benefit to humanity. In all of these, the rich get
richer. \footnote{Evolution itself is probably the most pervasive example of
superlinear returns for performance. But this is hard for us to
empathize with because we're not the recipients; we're the returns.}

You can't understand the world without understanding the concept of
superlinear returns. And if you're ambitious you definitely should,
because this will be the wave you surf on.

It may seem as if there are a lot of different situations with
superlinear returns, but as far as I can tell they reduce to two
fundamental causes: exponential growth and thresholds.

The most obvious case of superlinear returns is when you're working on
something that grows exponentially. For example, growing bacterial
cultures. When they grow at all, they grow exponentially. But they're
tricky to grow. Which means the difference in outcome between someone
who's adept at it and someone who's not is very great.

Startups can also grow exponentially, and we see the same pattern there.
Some manage to achieve high growth rates. Most don't. And as a result
you get qualitatively different outcomes: the companies with high growth
rates tend to become immensely valuable, while the ones with lower
growth rates may not even survive.

Y Combinator encourages founders to focus on growth rate rather than
absolute numbers. It prevents them from being discouraged early on, when
the absolute numbers are still low. It also helps them decide what to
focus on: you can use growth rate as a compass to tell you how to evolve
the company. But the main advantage is that by focusing on growth rate
you tend to get something that grows exponentially.

YC doesn't explicitly tell founders that with growth rate ``you get out
what you put in,'' but it's not far from the truth. And if growth rate
were proportional to performance, then the reward for performance
\emph{p} over time \emph{t} would be proportional to \emph{pt}.

Even after decades of thinking about this, I find that sentence
startling.

Whenever how well you do depends on how well you've done, you'll get
exponential growth. But neither our DNA nor our customs prepare us for
it. No one finds exponential growth natural; every child is surprised,
the first time they hear it, by the story of the man who asks the king
for a single grain of rice the first day and double the amount each
successive day.

What we don't understand naturally we develop customs to deal with, but
we don't have many customs about exponential growth either, because
there have been so few instances of it in human history. In principle
herding should have been one: the more animals you had, the more
offspring they'd have. But in practice grazing land was the limiting
factor, and there was no plan for growing that exponentially.

Or more precisely, no generally applicable plan. There \emph{was} a way
to grow one's territory exponentially: by conquest. The more territory
you control, the more powerful your army becomes, and the easier it is
to conquer new territory. This is why history is full of empires. But so
few people created or ran empires that their experiences didn't affect
customs very much. The emperor was a remote and terrifying figure, not a
source of lessons one could use in one's own life.

The most common case of exponential growth in preindustrial times was
probably scholarship. The more you know, the easier it is to learn new
things. The result, then as now, was that some people were startlingly
more knowledgeable than the rest about certain topics. But this didn't
affect customs much either. Although empires of ideas can overlap and
there can thus be far more emperors, in preindustrial times this type of
empire had little practical effect. \footnote{Knowledge did of course have a practical effect before the
Industrial Revolution. The development of agriculture changed human life
completely. But this kind of change was the result of broad, gradual
improvements in technique, not the discoveries of a few exceptionally
learned people.}

That has changed in the last few centuries. Now the emperors of ideas
can design bombs that defeat the emperors of territory. But this
phenomenon is still so new that we haven't fully assimilated it. Few
even of the participants realize they're benefitting from exponential
growth or ask what they can learn from other instances of it.

The other source of superlinear returns is embodied in the expression
``winner take all.'' In a sports match the relationship between
performance and return is a step function: the winning team gets one win
whether they do much better or just slightly better.
\footnote{It's not mathematically correct to describe a step function as
superlinear, but a step function starting from zero works like a
superlinear function when it describes the reward curve for effort by a
rational actor. If it starts at zero then the part before the step is
below any linearly increasing return, and the part after the step must
be above the necessary return at that point or no one would bother.}

The source of the step function is not competition per se, however. It's
that there are thresholds in the outcome. You don't need competition to
get those. There can be thresholds in situations where you're the only
participant, like proving a theorem or hitting a target.

It's remarkable how often a situation with one source of superlinear
returns also has the other. Crossing thresholds leads to exponential
growth: the winning side in a battle usually suffers less damage, which
makes them more likely to win in the future. And exponential growth
helps you cross thresholds: in a market with network effects, a company
that grows fast enough can shut out potential competitors.

Fame is an interesting example of a phenomenon that combines both
sources of superlinear returns. Fame grows exponentially because
existing fans bring you new ones. But the fundamental reason it's so
concentrated is thresholds: there's only so much room on the A-list in
the average person's head.

The most important case combining both sources of superlinear returns
may be learning. Knowledge grows exponentially, but there are also
thresholds in it. Learning to ride a bicycle, for example. Some of these
thresholds are akin to machine tools: once you learn to read, you're
able to learn anything else much faster. But the most important
thresholds of all are those representing new discoveries. Knowledge
seems to be fractal in the sense that if you push hard at the boundary
of one area of knowledge, you sometimes discover a whole new field. And
if you do, you get first crack at all the new discoveries to be made in
it. Newton did this, and so did Durer and Darwin.

Are there general rules for finding situations with superlinear returns?
The most obvious one is to seek work that compounds.

There are two ways work can compound. It can compound directly, in the
sense that doing well in one cycle causes you to do better in the next.
That happens for example when you're building infrastructure, or growing
an audience or brand. Or work can compound by teaching you, since
learning compounds. This second case is an interesting one because you
may feel you're doing badly as it's happening. You may be failing to
achieve your immediate goal. But if you're learning a lot, then you're
getting exponential growth nonetheless.

This is one reason Silicon Valley is so tolerant of failure. People in
Silicon Valley aren't blindly tolerant of failure. They'll only continue
to bet on you if you're learning from your failures. But if you are, you
are in fact a good bet: maybe your company didn't grow the way you
wanted, but you yourself have, and that should yield results eventually.

Indeed, the forms of exponential growth that don't consist of learning
are so often intermixed with it that we should probably treat this as
the rule rather than the exception. Which yields another heuristic:
always be learning. If you're not learning, you're probably not on a
path that leads to superlinear returns.

But don't overoptimize \emph{what} you're learning. Don't limit yourself
to learning things that are already known to be valuable. You're
learning; you don't know for sure yet what's going to be valuable, and
if you're too strict you'll lop off the outliers.

What about step functions? Are there also useful heuristics of the form
``seek thresholds'' or ``seek competition?'' Here the situation is
trickier. The existence of a threshold doesn't guarantee the game will
be worth playing. If you play a round of Russian roulette, you'll be in
a situation with a threshold, certainly, but in the best case you're no
better off. ``Seek competition'' is similarly useless; what if the prize
isn't worth competing for? Sufficiently fast exponential growth
guarantees both the shape and magnitude of the return curve --- because
something that grows fast enough will grow big even if it's trivially
small at first --- but thresholds only guarantee the shape.
\footnote{Seeking competition could be a good heuristic in the sense that
some people find it motivating. It's also somewhat of a guide to
promising problems, because it's a sign that other people find them
promising. But it's a very imperfect sign: often there's a clamoring
crowd chasing some problem, and they all end up being trumped by someone
quietly working on another one.}

A principle for taking advantage of thresholds has to include a test to
ensure the game is worth playing. Here's one that does: if you come
across something that's mediocre yet still popular, it could be a good
idea to replace it. For example, if a company makes a product that
people dislike yet still buy, then presumably they'd buy a better
alternative if you made one. \footnote{Not always, though. You have to be careful with this rule. When
something is popular despite being mediocre, there's often a hidden
reason why. Perhaps monopoly or regulation make it hard to compete.
Perhaps customers have bad taste or have broken procedures for deciding
what to buy. There are huge swathes of mediocre things that exist for
such reasons.}

It would be great if there were a way to find promising intellectual
thresholds. Is there a way to tell which questions have whole new fields
beyond them? I doubt we could ever predict this with certainty, but the
prize is so valuable that it would be useful to have predictors that
were even a little better than random, and there's hope of finding
those. We can to some degree predict when a research problem
\emph{isn't} likely to lead to new discoveries: when it seems legit but
boring. Whereas the kind that do lead to new discoveries tend to seem
very mystifying, but perhaps unimportant. (If they were mystifying and
obviously important, they'd be famous open questions with lots of people
already working on them.) So one heuristic here is to be driven by
curiosity rather than careerism --- to give free rein to your curiosity
instead of working on what you're supposed to.

The prospect of superlinear returns for performance is an exciting one
for the ambitious. And there's good news in this department: this
territory is expanding in both directions. There are more types of work
in which you can get superlinear returns, and the returns themselves are
growing.

There are two reasons for this, though they're so closely intertwined
that they're more like one and a half: progress in technology, and the
decreasing importance of organizations.

Fifty years ago it used to be much more necessary to be part of an
organization to work on ambitious projects. It was the only way to get
the resources you needed, the only way to have colleagues, and the only
way to get distribution. So in 1970 your prestige was in most cases the
prestige of the organization you belonged to. And prestige was an
accurate predictor, because if you weren't part of an organization, you
weren't likely to achieve much. There were a handful of exceptions, most
notably artists and writers, who worked alone using inexpensive tools
and had their own brands. But even they were at the mercy of
organizations for reaching audiences. \footnote{In my twenties I wanted to be an \href{worked.html}{artist} and
even went to art school to study painting. Mostly because I liked art,
but a nontrivial part of my motivation came from the fact that artists
seemed least at the mercy of organizations.}

A world dominated by organizations damped variation in the returns for
performance. But this world has eroded significantly just in my
lifetime. Now a lot more people can have the freedom that artists and
writers had in the 20th century. There are lots of ambitious projects
that don't require much initial funding, and lots of new ways to learn,
make money, find colleagues, and reach audiences.

There's still plenty of the old world left, but the rate of change has
been dramatic by historical standards. Especially considering what's at
stake. It's hard to imagine a more fundamental change than one in the
returns for performance.

Without the damping effect of institutions, there will be more variation
in outcomes. Which doesn't imply everyone will be better off: people who
do well will do even better, but those who do badly will do worse.
That's an important point to bear in mind. Exposing oneself to
superlinear returns is not for everyone. Most people will be better off
as part of the pool. So who should shoot for superlinear returns?
Ambitious people of two types: those who know they're so good that
they'll be net ahead in a world with higher variation, and those,
particularly the young, who can afford to risk trying it to find out.
\footnote{In principle everyone is getting superlinear returns. Learning
compounds, and everyone learns in the course of their life. But in
practice few push this kind of everyday learning to the point where the
return curve gets really steep.}

The switch away from institutions won't simply be an exodus of their
current inhabitants. Many of the new winners will be people they'd never
have let in. So the resulting democratization of opportunity will be
both greater and more authentic than any tame intramural version the
institutions themselves might have cooked up.

Not everyone is happy about this great unlocking of ambition. It
threatens some vested interests and contradicts some
ideologies.~\footnote{It's unclear exactly what advocates of ``equity'' mean by it.
They seem to disagree among themselves. But whatever they mean is
probably at odds with a world in which institutions have less power to
control outcomes, and a handful of outliers do much better than everyone
else.

It may seem like bad luck for this concept that it arose at just the
moment when the world was shifting in the opposite direction, but I
don't think this was a coincidence. I think one reason it arose now is
because its adherents feel threatened by rapidly increasing variation in
performance.} But if you're an ambitious
individual it's good news for you. How should you take advantage of it?

The most obvious way to take advantage of superlinear returns for
performance is by doing exceptionally good work. At the far end of the
curve, incremental effort is a bargain. All the more so because there's
less competition at the far end --- and not just for the obvious reason
that it's hard to do something exceptionally well, but also because
people find the prospect so intimidating that few even try. Which means
it's not just a bargain to do exceptional work, but a bargain even to
try to.

There are many variables that affect how good your work is, and if you
want to be an outlier you need to get nearly all of them right. For
example, to do something exceptionally well, you have to be interested
in it. Mere diligence is not enough. So in a world with superlinear
returns, it's even more valuable to know what you're interested in, and
to find ways to work on it. \footnote{Corollary: Parents who pressure their kids to work on something
prestigious, like medicine, even though they have no interest in it,
will be hosing them even more than they have in the past.} It will also be
important to choose work that suits your circumstances. For example, if
there's a kind of work that inherently requires a huge expenditure of
time and energy, it will be increasingly valuable to do it when you're
young and don't yet have children.

There's a surprising amount of technique to doing great work. It's not
just a matter of trying hard. I'm going to take a shot giving a recipe
in one paragraph.

Choose work you have a natural aptitude for and a deep interest in.
Develop a habit of working on your own projects; it doesn't matter what
they are so long as you find them excitingly ambitious. Work as hard as
you can without burning out, and this will eventually bring you to one
of the frontiers of knowledge. These look smooth from a distance, but up
close they're full of gaps. Notice and explore such gaps, and if you're
lucky one will expand into a whole new field. Take as much risk as you
can afford; if you're not failing occasionally you're probably being too
conservative. Seek out the best colleagues. Develop good taste and learn
from the best examples. Be honest, especially with yourself. Exercise
and eat and sleep well and avoid the more dangerous drugs. When in
doubt, follow your curiosity. It never lies, and it knows more than you
do about what's worth paying attention to. \footnote{The original version of this paragraph was the first draft of
``\href{greatwork.html}{How to Do Great Work}.'' As soon as I wrote it I
realized it was a more important topic than superlinear returns, so I
paused the present essay to expand this paragraph into its own.
Practically nothing remains of the original version, because after I
finished ``How to Do Great Work'' I rewrote it based on that.}

And there is of course one other thing you need: to be lucky. Luck is
always a factor, but it's even more of a factor when you're working on
your own rather than as part of an organization. And though there are
some valid aphorisms about luck being where preparedness meets
opportunity and so on, there's also a component of true chance that you
can't do anything about. The solution is to take multiple shots. Which
is another reason to start taking risks early.

The best example of a field with superlinear returns is probably
science. It has exponential growth, in the form of learning, combined
with thresholds at the extreme edge of performance --- literally at the
limits of knowledge.

The result has been a level of inequality in scientific discovery that
makes the wealth inequality of even the most stratified societies seem
mild by comparison. Newton's discoveries were arguably greater than all
his contemporaries' combined. \footnote{Before the Industrial Revolution, people who got rich usually
did it like emperors: capturing some resource made them more powerful
and enabled them to capture more. Now it can be done like a scientist,
by discovering or building something uniquely valuable. Most people who
get rich use a mix of the old and the new ways, but in the most advanced
economies the ratio has \href{richnow.html}{shifted dramatically} toward
discovery just in the last half century.}

This point may seem obvious, but it might be just as well to spell it
out. Superlinear returns imply inequality. The steeper the return curve,
the greater the variation in outcomes.

In fact, the correlation between superlinear returns and inequality is
so strong that it yields another heuristic for finding work of this
type: look for fields where a few big winners outperform everyone else.
A kind of work where everyone does about the same is unlikely to be one
with superlinear returns.

What are fields where a few big winners outperform everyone else? Here
are some obvious ones: sports, politics, art, music, acting, directing,
writing, math, science, starting companies, and investing. In sports the
phenomenon is due to externally imposed thresholds; you only need to be
a few percent faster to win every race. In politics, power grows much as
it did in the days of emperors. And in some of the other fields
(including politics) success is driven largely by fame, which has its
own source of superlinear growth. But when we exclude sports and
politics and the effects of fame, a remarkable pattern emerges: the
remaining list is exactly the same as the list of fields where you have
to be \href{think.html}{independent-minded} to succeed --- where your
ideas have to be not just correct, but novel as well.
\footnote{It's not surprising that conventional-minded people would
dislike inequality if independent-mindedness is one of the biggest
drivers of it. But it's not simply that they don't want anyone to have
what they can't. The conventional-minded literally can't imagine what
it's like to have novel ideas. So the whole phenomenon of great
variation in performance seems unnatural to them, and when they
encounter it they assume it must be due to cheating or to some malign
external influence.}

This is obviously the case in science. You can't publish papers saying
things that other people have already said. But it's just as true in
investing, for example. It's only useful to believe that a company will
do well if most other investors don't; if everyone else thinks the
company will do well, then its stock price will already reflect that,
and there's no room to make money.

What else can we learn from these fields? In all of them you have to put
in the initial effort. Superlinear returns seem small at first. \emph{At
this rate,} you find yourself thinking, \emph{I'll never get anywhere.}
But because the reward curve rises so steeply at the far end, it's worth
taking extraordinary measures to get there.

In the startup world, the name for this principle is ``do things that
don't scale.'' If you pay a ridiculous amount of attention to your tiny
initial set of customers, ideally you'll kick off exponential growth by
word of mouth. But this same principle applies to anything that grows
exponentially. Learning, for example. When you first start learning
something, you feel lost. But it's worth making the initial effort to
get a toehold, because the more you learn, the easier it will get.

There's another more subtle lesson in the list of fields with
superlinear returns: not to equate work with a job. For most of the 20th
century the two were identical for nearly everyone, and as a result
we've inherited a custom that equates productivity with having a job.
Even now to most people the phrase ``your work'' means their job. But to
a writer or artist or scientist it means whatever they're currently
studying or creating. For someone like that, their work is something
they carry with them from job to job, if they have jobs at all. It may
be done for an employer, but it's part of their portfolio.

It's an intimidating prospect to enter a field where a few big winners
outperform everyone else. Some people do this deliberately, but you
don't need to. If you have sufficient natural ability and you follow
your curiosity sufficiently far, you'll end up in one. Your curiosity
won't let you be interested in boring questions, and interesting
questions tend to create fields with superlinear returns if they're not
already part of one.

The territory of superlinear returns is by no means static. Indeed, the
most extreme returns come from expanding it. So while both ambition and
curiosity can get you into this territory, curiosity may be the more
powerful of the two. Ambition tends to make you climb existing peaks,
but if you stick close enough to an interesting enough question, it may
grow into a mountain beneath you.


There's a limit to how sharply you can distinguish between effort,
performance, and return, because they're not sharply distinguished in
fact. What counts as return to one person might be performance to
another. But though the borders of these concepts are blurry, they're
not meaningless. I've tried to write about them as precisely as I could
without crossing into error.

\textbf{Thanks} to Trevor Blackwell, Patrick Collison, Tyler Cowen,
Jessica Livingston, Harj Taggar, and Garry Tan for reading drafts of
this.
\chapter{The Best Essay}

March 2024

Despite its title this isn't meant to be the best essay. My goal here is
to figure out what the best essay would be like.

It would be well-written, but you can write well about any topic. What
made it special would be what it was about.

Obviously some topics would be better than others. It probably wouldn't
be about this year's lipstick colors. But it wouldn't be vaporous talk
about elevated themes either. A good essay has to be surprising. It has
to tell people something they don't already know.

The best essay would be on the most important topic you could tell
people something surprising about.

That may sound obvious, but it has some unexpected consequences. One is
that science enters the picture like an elephant stepping into a
rowboat. For example, Darwin first described the idea of natural
selection in an essay written in 1844. Talk about an important topic you
could tell people something surprising about. If that's the test of a
great essay, this was surely the best one written in 1844. And indeed,
the best possible essay at any given time would usually be one
describing the most important scientific or technological discovery it
was possible to make. \footnote{There might be some resistance to this conclusion on the grounds
that some of these discoveries could only be understood by a small
number of readers. But you get into all sorts of difficulties if you
want to disqualify essays on this account. How do you decide where the
cutoff should be? If a virus kills off everyone except a handful of
people sequestered at Los Alamos, could an essay that had been
disqualified now be eligible? Etc.

Darwin's 1844 essay was derived from an earlier version written in 1839.
Extracts from it were published in 1858.}

Another unexpected consequence: I imagined when I started writing this
that the best essay would be fairly timeless --- that the best essay you
could write in 1844 would be much the same as the best one you could
write now. But in fact the opposite seems to be true. It might be true
that the best painting would be timeless in this sense. But it wouldn't
be impressive to write an essay introducing natural selection now. The
best essay \emph{now} would be one describing a great discovery we
didn't yet know about.

If the question of how to write the best possible essay reduces to the
question of how to make great discoveries, then I started with the wrong
question. Perhaps what this exercise shows is that we shouldn't waste
our time writing essays but instead focus on making discoveries in some
specific domain. But I'm interested in essays and what can be done with
them, so I want to see if there's some other question I could have
asked.

There is, and on the face of it, it seems almost identical to the one I
started with. Instead of asking \emph{what would the best essay be?} I
should have asked \emph{how do you write essays well?} Though these seem
only phrasing apart, their answers diverge. The answer to the first
question, as we've seen, isn't really about essay writing. The second
question forces it to be.

Writing essays, at its best, is a way of discovering ideas. How do you
do that well? How do you discover by writing?

An essay should ordinarily start with what I'm going to call a question,
though I mean this in a very general sense: it doesn't have to be a
question grammatically, just something that acts like one in the sense
that it spurs some response.

How do you get this initial question? It probably won't work to choose
some important-sounding topic at random and go at it. Professional
traders won't even trade unless they have what they call an \emph{edge}
--- a convincing story about why in some class of trades they'll win
more than they lose. Similarly, you shouldn't attack a topic unless you
have a way in --- some new insight about it or way of approaching it.

You don't need to have a complete thesis; you just need some kind of gap
you can explore. In fact, merely having questions about something other
people take for granted can be edge enough.

If you come across a question that's sufficiently puzzling, it could be
worth exploring even if it doesn't seem very momentous. Many an
important discovery has been made by pulling on a thread that seemed
insignificant at first. How can they \emph{all} be finches?
\footnote{When you find yourself very curious about an apparently minor
question, that's an exciting sign. Evolution has designed you to pay
attention to things that matter. So when you're very curious about
something random, that could mean you've unconsciously noticed it's less
random than it seems.}

Once you've got a question, then what? You start thinking out loud about
it. Not literally out loud, but you commit to a specific string of words
in response, as you would if you were talking. This initial response is
usually mistaken or incomplete. Writing converts your ideas from vague
to bad. But that's a step forward, because once you can see the
brokenness, you can fix it.

Perhaps beginning writers are alarmed at the thought of starting with
something mistaken or incomplete, but you shouldn't be, because this is
why essay writing works. Forcing yourself to commit to some specific
string of words gives you a starting point, and if it's wrong, you'll
see that when you reread it. At least half of essay writing is rereading
what you've written and asking \emph{is this correct and complete?} You
have to be very strict when rereading, not just because you want to keep
yourself honest, but because a gap between your response and the truth
is often a sign of new ideas to be discovered.

The prize for being strict with what you've written is not just
refinement. When you take a roughly correct answer and try to make it
exactly right, sometimes you find that you can't, and that the reason is
that you were depending on a false assumption. And when you discard it,
the answer turns out to be completely different. \footnote{Corollary: If you're not intellectually honest, your writing
won't just be biased, but also boring, because you'll miss all the ideas
you'd have discovered if you pushed for the truth.}

Ideally the response to a question is two things: the first step in a
process that converges on the truth, and a source of additional
questions (in my very general sense of the word). So the process
continues recursively, as response spurs response.
\footnote{Sometimes this process begins before you start writing.
Sometimes you've already figured out the first few things you want to
say. Schoolchildren are often taught they should decide
\emph{everything} they want to say, and write this down as an outline
before they start writing the essay itself. Maybe that's a good way to
get them started --- or not, I don't know --- but it's antithetical to
the spirit of essay writing. The more detailed your outline, the less
your ideas can benefit from the sort of discovery that essays are for.}

Usually there are several possible responses to a question, which means
you're traversing a tree. But essays are linear, not tree-shaped, which
means you have to choose one branch to follow at each point. How do you
choose? Usually you should follow whichever offers the greatest
combination of generality and novelty. I don't consciously rank branches
this way; I just follow whichever seems most exciting; but generality
and novelty are what make a branch exciting. \footnote{The problem with this type of ``greedy'' algorithm is that you
can end up on a local maximum. If the most valuable question is preceded
by a boring one, you'll overlook it. But I can't imagine a better
strategy. There's no lookahead except by writing. So use a greedy
algorithm and a lot of time.}

If you're willing to do a lot of rewriting, you don't have to guess
right. You can follow a branch and see how it turns out, and if it isn't
good enough, cut it and backtrack. I do this all the time. In this essay
I've already cut a 17-paragraph subtree, in addition to countless
shorter ones. Maybe I'll reattach it at the end, or boil it down to a
footnote, or spin it off as its own essay; we'll see.
\footnote{I ended up reattaching the first 5 of the 17 paragraphs, and
discarding the rest.}

In general you want to be quick to cut. One of the most dangerous
temptations in writing (and in software and painting) is to keep
something that isn't right, just because it contains a few good bits or
cost you a lot of effort.

The most surprising new question being thrown off at this point is
\emph{does it really matter what the initial question is?} If the space
of ideas is highly connected, it shouldn't, because you should be able
to get from any question to the most valuable ones in a few hops. And we
see evidence that it's highly connected in the way, for example, that
people who are obsessed with some topic can turn any conversation toward
it. But that only works if you know where you want to go, and you don't
in an essay. That's the whole point. You don't want to be the obsessive
conversationalist, or all your essays will be about the same thing.
\footnote{Stephen Fry confessed to making use of this phenomenon when
taking exams at Oxford. He had in his head a standard essay about some
general literary topic, and he would find a way to turn the exam
question toward it and then just reproduce it again.

Strictly speaking it's the graph of ideas that would be highly
connected, not the space, but that usage would confuse people who don't
know graph theory, whereas people who do know it will get what I mean if
I say ``space''.}

The other reason the initial question matters is that you usually feel
somewhat obliged to stick to it. I don't think about this when I decide
which branch to follow. I just follow novelty and generality. Sticking
to the question is enforced later, when I notice I've wandered too far
and have to backtrack. But I think this is the optimal solution. You
don't want the hunt for novelty and generality to be constrained in the
moment. Go with it and see what you get. \footnote{Too far doesn't depend just on the distance from the original
topic. It's more like that distance divided by the value of whatever
I've discovered in the subtree.}

Since the initial question does constrain you, in the best case it sets
an upper bound on the quality of essay you'll write. If you do as well
as you possibly can on the chain of thoughts that follow from the
initial question, the initial question itself is the only place where
there's room for variation.

It would be a mistake to let this make you too conservative though,
because you can't predict where a question will lead. Not if you're
doing things right, because doing things right means making discoveries,
and by definition you can't predict those. So the way to respond to this
situation is not to be cautious about which initial question you choose,
but to write a lot of essays. Essays are for taking risks.

Almost any question can get you a good essay. Indeed, it took some
effort to think of a sufficiently unpromising topic in the third
paragraph, because any essayist's first impulse on hearing that the best
essay couldn't be about x would be to try to write it. But if most
questions yield good essays, only some yield great ones.

Can we predict which questions will yield great essays? Considering how
long I've been writing essays, it's alarming how novel that question
feels.

One thing I like in an initial question is outrageousness. I love
questions that seem naughty in some way --- for example, by seeming
counterintuitive or overambitious or heterodox. Ideally all three. This
essay is an example. Writing about the best essay implies there is such
a thing, which pseudo-intellectuals will dismiss as reductive, though it
follows necessarily from the possibility of one essay being better than
another. And thinking about how to do something so ambitious is close
enough to doing it that it holds your attention.

I like to start an essay with a gleam in my eye. This could be just a
taste of mine, but there's one aspect of it that probably isn't: to
write a really good essay on some topic, you have to be interested in
it. A good writer can write well about anything, but to stretch for the
novel insights that are the raison d'etre of the essay, you have to
care.

If caring about it is one of the criteria for a good initial question,
then the optimal question varies from person to person. It also means
you're more likely to write great essays if you care about a lot of
different things. The more curious you are, the greater the probable
overlap between the set of things you're curious about and the set of
topics that yield great essays.

What other qualities would a great initial question have? It's probably
good if it has implications in a lot of different areas. And I find it's
a good sign if it's one that people think has already been thoroughly
explored. But the truth is that I've barely thought about how to choose
initial questions, because I rarely do it. I rarely \emph{choose} what
to write about; I just start thinking about something, and sometimes it
turns into an essay.

Am I going to stop writing essays about whatever I happen to be thinking
about and instead start working my way through some systematically
generated list of topics? That doesn't sound like much fun. And yet I
want to write good essays, and if the initial question matters, I should
care about it.

Perhaps the answer is to go one step earlier: to write about whatever
pops into your head, but try to ensure that what pops into your head is
good. Indeed, now that I think about it, this has to be the answer,
because a mere list of topics wouldn't be any use if you didn't have
edge with any of them. To start writing an essay, you need a topic plus
some initial insight about it, and you can't generate those
systematically. If only. \footnote{Or can you? I should try writing about this. Even if the chance
of succeeding is small, the expected value is huge.}

You can probably cause yourself to have more of them, though. The
quality of the ideas that come out of your head depends on what goes in,
and you can improve that in two dimensions, breadth and depth.

You can't learn everything, so getting breadth implies learning about
topics that are very different from one another. When I tell people
about my book-buying trips to Hay and they ask what I buy books about, I
usually feel a bit sheepish answering, because the topics seem like a
laundry list of unrelated subjects. But perhaps that's actually optimal
in this business.

You can also get ideas by talking to people, by doing and building
things, and by going places and seeing things. I don't think it's
important to talk to new people so much as the sort of people who make
you have new ideas. I get more new ideas after talking for an afternoon
with Robert Morris than from talking to 20 new smart people. I know
because that's what a block of office hours at Y Combinator consists of.

While breadth comes from reading and talking and seeing, depth comes
from doing. The way to really learn about some domain is to have to
solve problems in it. Though this could take the form of writing, I
suspect that to be a good essayist you also have to do, or have done,
some other kind of work. That may not be true for most other fields, but
essay writing is different. You could spend half your time working on
something else and be net ahead, so long as it was hard.

I'm not proposing that as a recipe so much as an encouragement to those
already doing it. If you've spent all your life so far working on other
things, you're already halfway there. Though of course to be good at
writing you have to like it, and if you like writing you'd probably have
spent at least some time doing it.

Everything I've said about initial questions applies also to the
questions you encounter in writing the essay. They're the same thing;
every subtree of an essay is usually a shorter essay, just as every
subtree of a Calder mobile is a smaller mobile. So any technique that
gets you good initial questions also gets you good whole essays.

At some point the cycle of question and response reaches what feels like
a natural end. Which is a little suspicious; shouldn't every answer
suggest more questions? I think what happens is that you start to feel
sated. Once you've covered enough interesting ground, you start to lose
your appetite for new questions. Which is just as well, because the
reader is probably feeling sated too. And it's not lazy to stop asking
questions, because you could instead be asking the initial question of a
new essay.

That's the ultimate source of drag on the connectedness of ideas: the
discoveries you make along the way. If you discover enough starting from
question A, you'll never make it to question B. Though if you keep
writing essays you'll gradually fix this problem by burning off such
discoveries. So bizarrely enough, writing lots of essays makes it as if
the space of ideas were more highly connected.

When a subtree comes to an end, you can do one of two things. You can
either stop, or pull the Cubist trick of laying separate subtrees end to
end by returning to a question you skipped earlier. Usually it requires
some sleight of hand to make the essay flow continuously at this point,
but not this time. This time I actually need an example of the
phenomenon. For example, we discovered earlier that the best possible
essay wouldn't usually be timeless in the way the best painting would.
This seems surprising enough to be worth investigating further.

There are two senses in which an essay can be timeless: to be about a
matter of permanent importance, and always to have the same effect on
readers. With art these two senses blend together. Art that looked
beautiful to the ancient Greeks still looks beautiful to us. But with
essays the two senses diverge, because essays teach, and you can't teach
people something they already know. Natural selection is certainly a
matter of permanent importance, but an essay explaining it couldn't have
the same effect on us that it would have had on Darwin's contemporaries,
precisely because his ideas were so successful that everyone already
knows about them. \footnote{There was a vogue in the 20th century for saying that the
purpose of art was also to teach. Some artists tried to justify their
work by explaining that their goal was not to produce something good,
but to challenge our preconceptions about art. And to be fair, art can
teach somewhat. The ancient Greeks' naturalistic sculptures represented
a new idea, and must have been extra exciting to contemporaries on that
account. But they still look good to us.}

I imagined when I started writing this that the best possible essay
would be timeless in the stricter, evergreen sense: that it would
contain some deep, timeless wisdom that would appeal equally to
Aristotle and Feynman. That doesn't seem to be true. But if the best
possible essay wouldn't usually be timeless in this stricter sense, what
would it take to write essays that were?

The answer to that turns out to be very strange: to be the evergreen
kind of timeless, an essay has to be ineffective, in the sense that its
discoveries aren't assimilated into our shared culture. Otherwise there
will be nothing new in it for the second generation of readers. If you
want to surprise readers not just now but in the future as well, you
have to write essays that won't stick --- essays that, no matter how
good they are, won't become part of what people in the future learn
before they read them. \footnote{Bertrand Russell caused huge controversy in the early 20th
century with his ideas about ``trial marriage.'' But they make boring
reading now, because they prevailed. ``Trial marriage'' is what we call
``dating.''}

I can imagine several ways to do that. One would be to write about
things people never learn. For example, it's a long-established pattern
for ambitious people to chase after various types of prizes, and only
later, perhaps too late, to realize that some of them weren't worth as
much as they thought. If you write about that, you can be confident of a
conveyor belt of future readers to be surprised by it.

Ditto if you write about the tendency of the inexperienced to overdo
things --- of young engineers to produce overcomplicated solutions, for
example. There are some kinds of mistakes people never learn to avoid
except by making them. Any of those should be a timeless topic.

Sometimes when we're slow to grasp things it's not just because we're
obtuse or in denial but because we've been deliberately lied to. There
are a lot of things adults \href{lies.html}{lie} to kids about, and when
you reach adulthood, they don't take you aside and hand you a list of
them. They don't remember which lies they told you, and most were
implicit anyway. So contradicting such lies will be a source of
surprises for as long as adults keep telling them.

Sometimes it's systems that lie to you. For example, the educational
systems in most countries train you to win by \href{lesson.html}{hacking
the test}. But that's not how you win at the most important real-world
tests, and after decades of training, this is hard for new arrivals in
the real world to grasp. Helping them overcome such institutional lies
will work as long as the institutions remain broken.
\footnote{If you'd asked me 10 years ago, I'd have predicted that schools
would continue to teach hacking the test for centuries. But now it seems
plausible that students will soon be taught individually by AIs, and
that exams will be replaced by ongoing, invisible micro-assessments.}

Another recipe for timelessness is to write about things readers already
know, but in much more detail than can be transmitted culturally.
``Everyone knows,'' for example, that it can be rewarding to have
\href{kids.html}{kids}. But till you have them you don't know precisely
what forms that takes, and even then much of what you know you may never
have put into words.

I've written about all these kinds of topics. But I didn't do it in a
deliberate attempt to write essays that were timeless in the stricter
sense. And indeed, the fact that this depends on one's ideas not
sticking suggests that it's not worth making a deliberate attempt to.
You should write about topics of timeless importance, yes, but if you do
such a good job that your conclusions stick and future generations find
your essay obvious instead of novel, so much the better. You've crossed
into Darwin territory.

Writing about topics of timeless importance is an instance of something
even more general, though: breadth of applicability. And there are more
kinds of breadth than chronological --- applying to lots of different
fields, for example. So breadth is the ultimate aim.

I already aim for it. Breadth and novelty are the two things I'm always
chasing. But I'm glad I understand where timelessness fits.

I understand better where a lot of things fit now. This essay has been a
kind of tour of essay writing. I started out hoping to get advice about
topics; if you assume good writing, the only thing left to differentiate
the best essay is its topic. And I did get advice about topics: discover
natural selection. Yeah, that would be nice. But when you step back and
ask what's the best you can do short of making some great discovery like
that, the answer turns out to be about procedure. Ultimately the quality
of an essay is a function of the ideas discovered in it, and the way you
get them is by casting a wide net for questions and then being very
exacting with the answers.

The most striking feature of this map of essay writing are the
alternating stripes of inspiration and effort required. The questions
depend on inspiration, but the answers can be got by sheer persistence.
You don't have to get an answer right the first time, but there's no
excuse for not getting it right eventually, because you can keep
rewriting till you do. And this is not just a theoretical possibility.
It's a pretty accurate description of the way I work. I'm rewriting as
we speak.

But although I wish I could say that writing great essays depends mostly
on effort, in the limit case it's inspiration that makes the difference.
In the limit case, the questions are the harder thing to get. That pool
has no bottom.

How to get more questions? That is the most important question of all.


\textbf{Thanks} to Sam Altman, Trevor Blackwell, Jessica Livingston,
Robert Morris, Courtenay Pipkin, and Harj Taggar for reading drafts of
this.
\chapter{How to Start Google}

March 2024

\emph{(This is a talk I gave to 14 and 15 year olds about what to do now
if they might want to start a startup later. Lots of schools think they
should tell students something about startups. This is what I think they
should tell them.)}

Most of you probably think that when you're released into the so-called
real world you'll eventually have to get some kind of job. That's not
true, and today I'm going to talk about a trick you can use to avoid
ever having to get a job.

The trick is to start your own company. So it's not a trick for avoiding
\emph{work}, because if you start your own company you'll work harder
than you would if you had an ordinary job. But you will avoid many of
the annoying things that come with a job, including a boss telling you
what to do.

It's more exciting to work on your own project than someone else's. And
you can also get a lot richer. In fact, this is the standard way to get
\href{richnow.html}{really rich}. If you look at the lists of the
richest people that occasionally get published in the press, nearly all
of them did it by starting their own companies.

Starting your own company can mean anything from starting a barber shop
to starting Google. I'm here to talk about one extreme end of that
continuum. I'm going to tell you how to start Google.

The companies at the Google end of the continuum are called startups
when they're young. The reason I know about them is that my wife Jessica
and I started something called Y Combinator that is basically a startup
factory. Since 2005, Y Combinator has funded over 4000 startups. So we
know exactly what you need to start a startup, because we've helped
people do it for the last 19 years.

You might have thought I was joking when I said I was going to tell you
how to start Google. You might be thinking ``How could \emph{we} start
Google?'' But that's effectively what the people who did start Google
were thinking before they started it. If you'd told Larry Page and
Sergey Brin, the founders of Google, that the company they were about to
start would one day be worth over a trillion dollars, their heads would
have exploded.

All you can know when you start working on a startup is that it seems
worth pursuing. You can't know whether it will turn into a company worth
billions or one that goes out of business. So when I say I'm going to
tell you how to start Google, I mean I'm going to tell you how to get to
the point where you can start a company that has as much chance of being
Google as Google had of being Google. \footnote{The rhetorical trick in this sentence is that the ``Google''s
refer to different things. What I mean is: a company that has as much
chance of growing as big as Google ultimately did as Larry and Sergey
could have reasonably expected Google itself would at the time they
started it. But I think the original version is zippier.}

How do you get from where you are now to the point where you can start a
successful startup? You need three things. You need to be good at some
kind of technology, you need an idea for what you're going to build, and
you need cofounders to start the company with.

How do you get good at technology? And how do you choose which
technology to get good at? Both of those questions turn out to have the
same answer: work on your own projects. Don't try to guess whether gene
editing or LLMs or rockets will turn out to be the most valuable
technology to know about. No one can predict that. Just work on whatever
interests you the most. You'll work much harder on something you're
interested in than something you're doing because you think you're
supposed to.

If you're not sure what technology to get good at, get good at
programming. That has been the source of the median startup for the last
30 years, and this is probably not going to change in the next 10.

Those of you who are taking computer science classes in school may at
this point be thinking, ok, we've got this sorted. We're already being
taught all about programming. But sorry, this is not enough. You have to
be working on your own projects, not just learning stuff in classes. You
can do well in computer science classes without ever really learning to
program. In fact you can graduate with a degree in computer science from
a top university and still not be any good at programming. That's why
tech companies all make you take a coding test before they'll hire you,
regardless of where you went to university or how well you did there.
They know grades and exam results prove nothing.

If you really want to learn to program, you have to work on your own
projects. You learn so much faster that way. Imagine you're writing a
game and there's something you want to do in it, and you don't know how.
You're going to figure out how a lot faster than you'd learn anything in
a class.

You don't have to learn programming, though. If you're wondering what
counts as technology, it includes practically everything you could
describe using the words ``make'' or ``build.'' So welding would count,
or making clothes, or making videos. Whatever you're most interested in.
The critical distinction is whether you're producing or just consuming.
Are you writing computer games, or just playing them? That's the cutoff.

Steve Jobs, the founder of Apple, spent time when he was a teenager
studying calligraphy --- the sort of beautiful writing that you see in
medieval manuscripts. No one, including him, thought that this would
help him in his career. He was just doing it because he was interested
in it. But it turned out to help him a lot. The computer that made Apple
really big, the Macintosh, came out at just the moment when computers
got powerful enough to make letters like the ones in printed books
instead of the computery-looking letters you see in 8 bit games. Apple
destroyed everyone else at this, and one reason was that Steve was one
of the few people in the computer business who really got graphic
design.

Don't feel like your projects have to be \emph{serious}. They can be as
frivolous as you like, so long as you're building things you're excited
about. Probably 90\% of programmers start out building games. They and
their friends like to play games. So they build the kind of things they
and their friends want. And that's exactly what you should be doing at
15 if you want to start a startup one day.

You don't have to do just one project. In fact it's good to learn about
multiple things. Steve Jobs didn't just learn calligraphy. He also
learned about electronics, which was even more valuable. Whatever you're
interested in. (Do you notice a theme here?)

So that's the first of the three things you need, to get good at some
kind or kinds of technology. You do it the same way you get good at the
violin or football: practice. If you start a startup at 22, and you
start writing your own programs now, then by the time you start the
company you'll have spent at least 7 years practicing writing code, and
you can get pretty good at anything after practicing it for 7 years.

Let's suppose you're 22 and you've succeeded: You're now really good at
some technology. How do you get \href{startupideas.html}{startup ideas}?
It might seem like that's the hard part. Even if you are a good
programmer, how do you get the idea to start Google?

Actually it's easy to get startup ideas once you're good at technology.
Once you're good at some technology, when you look at the world you see
dotted outlines around the things that are missing. You start to be able
to see both the things that are missing from the technology itself, and
all the broken things that could be fixed using it, and each one of
these is a potential startup.

In the town near our house there's a shop with a sign warning that the
door is hard to close. The sign has been there for several years. To the
people in the shop it must seem like this mysterious natural phenomenon
that the door sticks, and all they can do is put up a sign warning
customers about it. But any carpenter looking at this situation would
think ``why don't you just plane off the part that sticks?''

Once you're good at programming, all the missing software in the world
starts to become as obvious as a sticking door to a carpenter. I'll give
you a real world example. Back in the 20th century, American
universities used to publish printed directories with all the students'
names and contact info. When I tell you what these directories were
called, you'll know which startup I'm talking about. They were called
facebooks, because they usually had a picture of each student next to
their name.

So Mark Zuckerberg shows up at Harvard in 2002, and the university still
hasn't gotten the facebook online. Each individual house has an online
facebook, but there isn't one for the whole university. The university
administration has been diligently having meetings about this, and will
probably have solved the problem in another decade or so. Most of the
students don't consciously notice that anything is wrong. But Mark is a
programmer. He looks at this situation and thinks ``Well, this is
stupid. I could write a program to fix this in one night. Just let
people upload their own photos and then combine the data into a new site
for the whole university.'' So he does. And almost literally overnight
he has thousands of users.

Of course Facebook was not a startup yet. It was just a\ldots{} project.
There's that word again. Projects aren't just the best way to learn
about technology. They're also the best source of startup ideas.

Facebook was not unusual in this respect. Apple and Google also began as
projects. Apple wasn't meant to be a company. Steve Wozniak just wanted
to build his own computer. It only turned into a company when Steve Jobs
said ``Hey, I wonder if we could sell plans for this computer to other
people.'' That's how Apple started. They weren't even selling computers,
just plans for computers. Can you imagine how lame this company seemed?

Ditto for Google. Larry and Sergey weren't trying to start a company at
first. They were just trying to make search better. Before Google, most
search engines didn't try to sort the results they gave you in order of
importance. If you searched for ``rugby'' they just gave you every web
page that contained the word ``rugby.'' And the web was so small in 1997
that this actually worked! Kind of. There might only be 20 or 30 pages
with the word ``rugby,'' but the web was growing exponentially, which
meant this way of doing search was becoming exponentially more broken.
Most users just thought, ``Wow, I sure have to look through a lot of
search results to find what I want.'' Door sticks. But like Mark, Larry
and Sergey were programmers. Like Mark, they looked at this situation
and thought ``Well, this is stupid. Some pages about rugby matter more
than others. Let's figure out which those are and show them first.''

It's obvious in retrospect that this was a great idea for a startup. It
wasn't obvious at the time. It's never obvious. If it was obviously a
good idea to start Apple or Google or Facebook, someone else would have
already done it. That's why the best startups grow out of projects that
aren't meant to be startups. You're not trying to start a company.
You're just following your instincts about what's interesting. And if
you're young and good at technology, then your unconscious instincts
about what's interesting are better than your conscious ideas about what
would be a good company.

So it's critical, if you're a young founder, to build things for
yourself and your friends to use. The biggest mistake young founders
make is to build something for some mysterious group of other people.
But if you can make something that you and your friends truly want to
use --- something your friends aren't just using out of loyalty to you,
but would be really sad to lose if you shut it down --- then you almost
certainly have the germ of a good startup idea. It may not seem like a
startup to you. It may not be obvious how to make money from it. But
trust me, there's a way.

What you need in a startup idea, and all you need, is something your
friends actually want. And those ideas aren't hard to see once you're
good at technology. There are sticking doors everywhere.
\footnote{Making something for your friends isn't the only source of
startup ideas. It's just the best source for young founders, who have
the least knowledge of what other people want, and whose own wants are
most predictive of future demand.}

Now for the third and final thing you need: a cofounder, or cofounders.
The optimal startup has two or three founders, so you need one or two
cofounders. How do you find them? Can you predict what I'm going to say
next? It's the same thing: projects. You find cofounders by working on
projects with them. What you need in a cofounder is someone who's good
at what they do and that you work well with, and the only way to judge
this is to work with them on things.

At this point I'm going to tell you something you might not want to
hear. It really matters to do well in your classes, even the ones that
are just memorization or blathering about literature, because you need
to do well in your classes to get into a good university. And if you
want to start a startup you should try to get into the best university
you can, because that's where the best cofounders are. It's also where
the best employees are. When Larry and Sergey started Google, they began
by just hiring all the smartest people they knew out of Stanford, and
this was a real advantage for them.

The empirical evidence is clear on this. If you look at where the
largest numbers of successful startups come from, it's pretty much the
same as the list of the most selective universities.

I don't think it's the prestigious names of these universities that
cause more good startups to come out of them. Nor do I think it's
because the quality of the teaching is better. What's driving this is
simply the difficulty of getting in. You have to be pretty smart and
determined to get into MIT or Cambridge, so if you do manage to get in,
you'll find the other students include a lot of smart and determined
people. \footnote{Strangely enough this is particularly true in countries like the
US where undergraduate admissions are done badly. US admissions
departments make applicants jump through a lot of arbitrary hoops that
have little to do with their intellectual ability. But the more
arbitrary a test, the more it becomes a test of mere determination and
resourcefulness. And those are the two most important qualities in
startup founders. So US admissions departments are better at selecting
founders than they would be if they were better at selecting students.}

You don't have to start a startup with someone you meet at university.
The founders of Twitch met when they were seven. The founders of Stripe,
Patrick and John Collison, met when John was born. But universities are
the main source of cofounders. And because they're where the cofounders
are, they're also where the ideas are, because the best ideas grow out
of projects you do with the people who become your cofounders.

So the list of what you need to do to get from here to starting a
startup is quite short. You need to get good at technology, and the way
to do that is to work on your own projects. And you need to do as well
in school as you can, so you can get into a good university, because
that's where the cofounders and the ideas are.

That's it, just two things, build stuff and do well in school.


\textbf{Thanks} to Jared Friedman, Carolynn Levy, Jessica Livingston,
Harj Taggar, and Garry Tan for reading drafts of this.
\chapter{The Reddits}

March 2024

I met the Reddits before we even started Y Combinator. In fact they were
one of the reasons we started it.

YC grew out of a talk I gave to the Harvard Computer Society (the
undergrad computer club) about how to start a startup. Everyone else in
the audience was probably local, but Steve and Alexis came up on the
train from the University of Virginia, where they were seniors. Since
they'd come so far I agreed to meet them for coffee. They told me about
the startup idea we'd later fund them to drop: a way to order fast food
on your cellphone.

This was before smartphones. They'd have had to make deals with cell
carriers and fast food chains just to get it launched. So it was not
going to happen. It still doesn't exist, 19 years later. But I was
impressed with their brains and their energy. In fact I was so impressed
with them and some of the other people I met at that talk that I decided
to start something to fund them. A few days later I told Steve and
Alexis that we were starting Y~Combinator, and encouraged them to apply.

That first batch we didn't have any way to identify applicants, so we
made up nicknames for them. The Reddits were the ``Cell food muffins.''
``Muffin'' is a term of endearment Jessica uses for things like small
dogs and two year olds. So that gives you some idea what kind of
impression Steve and Alexis made in those days. They had the look of
slightly ruffled surprise that baby birds have.

Their idea was bad though. And since we thought then that we were
funding ideas rather than founders, we rejected them. But we felt bad
about it. Jessica was sad that we'd rejected the muffins. And it seemed
wrong to me to turn down the people we'd been inspired to start YC to
fund.

I don't think the startup sense of the word ``pivot'' had been invented
yet, but we wanted to fund Steve and Alexis, so if their idea was bad,
they'd have to work on something else. And I knew what else. In those
days there was a site called Delicious where you could save links. It
had a page called del.icio.us/popular that listed the most-saved links,
and people were using this page as a de facto Reddit. I knew because a
lot of the traffic to my site was coming from it. There needed to be
something like del.icio.us/popular, but designed for sharing links
instead of being a byproduct of saving them.

So I called Steve and Alexis and said that we liked them, just not their
idea, so we'd fund them if they'd work on something else. They were on
the train home to Virginia at that point. They got off at the next
station and got on the next train north, and by the end of the day were
committed to working on what's now called Reddit.

They would have liked to call it Snoo, as in ``What snoo?'' But snoo.com
was too expensive, so they settled for calling the mascot Snoo and
picked a name for the site that wasn't registered. Early on Reddit was
just a provisional name, or so they told me at least, but it's probably
too late to change it now.

As with all the really great startups, there's an uncannily close match
between the company and the founders. Steve in particular. Reddit has a
certain personality --- curious, skeptical, ready to be amused --- and
that personality is Steve's.

Steve will roll his eyes at this, but he's an intellectual; he's
interested in ideas for their own sake. That was how he came to be in
that audience in Cambridge in the first place. He knew me because he was
interested in a programming language I've written about called Lisp, and
Lisp is one of those languages few people learn except out of
intellectual curiosity. Steve's kind of vacuum-cleaner curiosity is
exactly what you want when you're starting a site that's a list of links
to literally anything interesting.

Steve was not a big fan of authority, so he also liked the idea of a
site without editors. In those days the top forum for programmers was a
site called Slashdot. It was a lot like Reddit, except the stories on
the frontpage were chosen by human moderators. And though they did a
good job, that one small difference turned out to be a big difference.
Being driven by user submissions meant Reddit was fresher than Slashdot.
News there was newer, and users will always go where the newest news is.

I pushed the Reddits to launch fast. A version one didn't need to be
more than a couple hundred lines of code. How could that take more than
a week or two to build? And they did launch comparatively fast, about
three weeks into the first YC batch. The first users were Steve, Alexis,
me, and some of their YC batchmates and college friends. It turns out
you don't need that many users to collect a decent list of interesting
links, especially if you have multiple accounts per user.

Reddit got two more people from their YC batch: Chris Slowe and Aaron
Swartz, and they too were unusually smart. Chris was just finishing his
PhD in physics at Harvard. Aaron was younger, a college freshman, and
even more anti-authority than Steve. It's not exaggerating to describe
him as a martyr for what authority later did to him.

Slowly but inexorably Reddit's traffic grew. At first the numbers were
so small they were hard to distinguish from background noise. But within
a few weeks it was clear that there was a core of real users returning
regularly to the site. And although all kinds of things have happened to
Reddit the company in the years since, Reddit the \emph{site} never
looked back.

Reddit the site (and now app) is such a fundamentally useful thing that
it's almost unkillable. Which is why, despite a long stretch after Steve
left when the management strategy ranged from benign neglect to
spectacular blunders, traffic just kept growing. You can't do that with
most companies. Most companies you take your eye off the ball for six
months and you're in deep trouble. But Reddit was special, and when
Steve came back in 2015, I knew the world was in for a surprise.

People thought they had Reddit's number: one of the players in Silicon
Valley, but not one of the big ones. But those who knew what had been
going on behind the scenes knew there was more to the story than this.
If Reddit could grow to the size it had with management that was
harmless at best, what could it do if Steve came back? We now know the
answer to that question. Or at least a lower bound on the answer. Steve
is not out of ideas yet.
\chapter{The Right Kind of Stubborn}

July 2024

Successful people tend to be persistent. New ideas often don't work at
first, but they're not deterred. They keep trying and eventually find
something that does.

Mere obstinacy, on the other hand, is a recipe for failure. Obstinate
people are so annoying. They won't listen. They beat their heads against
a wall and get nowhere.

But is there any real difference between these two cases? Are persistent
and obstinate people actually behaving differently? Or are they doing
the same thing, and we just label them later as persistent or obstinate
depending on whether they turned out to be right or not?

If that's the only difference then there's nothing to be learned from
the distinction. Telling someone to be persistent rather than obstinate
would just be telling them to be right rather than wrong, and they
already know that. Whereas if persistence and obstinacy are actually
different kinds of behavior, it would be worthwhile to tease them apart.
\footnote{I'm going to use ``persistent'' for the good kind of stubborn
and ``obstinate'' for the bad kind, but I can't claim I'm simply
following current usage. Conventional opinion barely distinguishes
between good and bad kinds of stubbornness, and usage is correspondingly
promiscuous. I could have invented a new word for the good kind, but it
seemed better just to stretch ``persistent.''}

I've talked to a lot of determined people, and it seems to me that
they're different kinds of behavior. I've often walked away from a
conversation thinking either ``Wow, that guy is determined'' or ``Damn,
that guy is stubborn,'' and I don't think I'm just talking about whether
they seemed right or not. That's part of it, but not all of it.

There's something annoying about the obstinate that's not simply due to
being mistaken. They won't listen. And that's not true of all determined
people. I can't think of anyone more determined than the Collison
brothers, and when you point out a problem to them, they not only
listen, but listen with an almost predatory intensity. Is there a hole
in the bottom of their boat? Probably not, but if there is, they want to
know about it.

It's the same with most successful people. They're never \emph{more}
engaged than when you disagree with them. Whereas the obstinate don't
want to hear you. When you point out problems, their eyes glaze over,
and their replies sound like ideologues talking about matters of
doctrine. \footnote{There are some domains where one can succeed by being obstinate.
Some political leaders have been notorious for it. But it won't work in
situations where you have to pass external tests. And indeed the
political leaders who are famous for being obstinate are famous for
getting power, not for using it well.}

The reason the persistent and the obstinate seem similar is that they're
both hard to stop. But they're hard to stop in different senses. The
persistent are like boats whose engines can't be throttled back. The
obstinate are like boats whose rudders can't be turned.
\footnote{There will be some resistance to turning the rudder of a
persistent person, because there's some cost to changing direction.}

In the degenerate case they're indistinguishable: when there's only one
way to solve a problem, your only choice is whether to give up or not,
and persistence and obstinacy both say no. This is presumably why the
two are so often conflated in popular culture. It assumes simple
problems. But as problems get more complicated, we can see the
difference between them. The persistent are much more attached to points
high in the decision tree than to minor ones lower down, while the
obstinate spray ``don't give up'' indiscriminately over the whole tree.

The persistent are attached to the goal. The obstinate are attached to
their ideas about how to reach it.

Worse still, that means they'll tend to be attached to their
\emph{first} ideas about how to solve a problem, even though these are
the least informed by the experience of working on it. So the obstinate
aren't merely attached to details, but disproportionately likely to be
attached to wrong ones.

Why are they like this? Why are the obstinate obstinate? One possibility
is that they're overwhelmed. They're not very capable. They take on a
hard problem. They're immediately in over their head. So they grab onto
ideas the way someone on the deck of a rolling ship might grab onto the
nearest handhold.

That was my initial theory, but on examination it doesn't hold up. If
being obstinate were simply a consequence of being in over one's head,
you could make persistent people become obstinate by making them solve
harder problems. But that's not what happens. If you handed the
Collisons an extremely hard problem to solve, they wouldn't become
obstinate. If anything they'd become less obstinate. They'd know they
had to be open to anything.

Similarly, if obstinacy were caused by the situation, the obstinate
would stop being obstinate when solving easier problems. But they don't.
And if obstinacy isn't caused by the situation, it must come from
within. It must be a feature of one's personality.

Obstinacy is a reflexive resistance to changing one's ideas. This is not
identical with stupidity, but they're closely related. A reflexive
resistance to changing one's ideas becomes a sort of induced stupidity
as contrary evidence mounts. And obstinacy is a form of not giving up
that's easily practiced by the stupid. You don't have to consider
complicated tradeoffs; you just dig in your heels. It even works, up to
a point.

The fact that obstinacy works for simple problems is an important clue.
Persistence and obstinacy aren't opposites. The relationship between
them is more like the relationship between the two kinds of respiration
we can do: aerobic respiration, and the anaerobic respiration we
inherited from our most distant ancestors. Anaerobic respiration is a
more primitive process, but it has its uses. When you leap suddenly away
from a threat, that's what you're using.

The optimal amount of obstinacy is not zero. It can be good if your
initial reaction to a setback is an unthinking ``I won't give up,''
because this helps prevent panic. But unthinking only gets you so far.
The further someone is toward the obstinate end of the continuum, the
less likely they are to succeed in solving hard problems.
\footnote{The obstinate do sometimes succeed in solving hard problems. One
way is through luck: like the stopped clock that's right twice a day,
they seize onto some arbitrary idea, and it turns out to be right.
Another is when their obstinacy cancels out some other form of error.
For example, if a leader has overcautious subordinates, their estimates
of the probability of success will always be off in the same direction.
So if he mindlessly says ``push ahead regardless'' in every borderline
case, he'll usually turn out to be right.}

Obstinacy is a simple thing. Animals have it. But persistence turns out
to have a fairly complicated internal structure.

One thing that distinguishes the persistent is their energy. At the risk
of putting too much weight on words, they persist rather than merely
resisting. They keep trying things. Which means the persistent must also
be imaginative. To keep trying things, you have to keep thinking of
things to try.

Energy and imagination make a wonderful combination. Each gets the best
out of the other. Energy creates demand for the ideas produced by
imagination, which thus produces more, and imagination gives energy
somewhere to go. \footnote{If you stop there, at just energy and imagination, you get the
conventional caricature of an artist or poet.}

Merely having energy and imagination is quite rare. But to solve hard
problems you need three more qualities: resilience, good judgement, and
a focus on some kind of goal.

Resilience means not having one's morale destroyed by setbacks. Setbacks
are inevitable once problems reach a certain size, so if you can't
bounce back from them, you can only do good work on a small scale. But
resilience is not the same as obstinacy. Resilience means setbacks can't
change your morale, not that they can't change your mind.

Indeed, persistence often requires that one change one's mind. That's
where good judgement comes in. The persistent are quite rational. They
focus on expected value. It's this, not recklessness, that lets them
work on things that are unlikely to succeed.

There is one point at which the persistent are often irrational though:
at the very top of the decision tree. When they choose between two
problems of roughly equal expected value, the choice usually comes down
to personal preference. Indeed, they'll often classify projects into
deliberately wide bands of expected value in order to ensure that the
one they want to work on still qualifies.

Empirically this doesn't seem to be a problem. It's ok to be irrational
near the top of the decision tree. One reason is that we humans will
work harder on a problem we love. But there's another more subtle factor
involved as well: our preferences among problems aren't random. When we
love a problem that other people don't, it's often because we've
unconsciously noticed that it's more important than they realize.

Which leads to our fifth quality: there needs to be some overall goal.
If you're like me you began, as a kid, merely with the desire to do
something great. In theory that should be the most powerful motivator of
all, since it includes everything that could possibly be done. But in
practice it's not much use, precisely because it includes too much. It
doesn't tell you what to do at this moment.

So in practice your energy and imagination and resilience and good
judgement have to be directed toward some fairly specific goal. Not too
specific, or you might miss a great discovery adjacent to what you're
searching for, but not too general, or it won't work to motivate you.
\footnote{Start by erring on the small side. If you're inexperienced
you'll inevitably err on one side or the other, and if you err on the
side of making the goal too broad, you won't get anywhere. Whereas if
you err on the small side you'll at least be moving forward. Then, once
you're moving, you expand the goal.}

When you look at the internal structure of persistence, it doesn't
resemble obstinacy at all. It's so much more complex. Five distinct
qualities --- energy, imagination, resilience, good judgement, and focus
on a goal --- combine to produce a phenomenon that seems a bit like
obstinacy in the sense that it causes you not to give up. But the way
you don't give up is completely different. Instead of merely resisting
change, you're driven toward a goal by energy and resilience, through
paths discovered by imagination and optimized by judgement. You'll give
way on any point low down in the decision tree, if its expected value
drops sufficiently, but energy and resilience keep pushing you toward
whatever you chose higher up.

Considering what it's made of, it's not surprising that the right kind
of stubbornness is so much rarer than the wrong kind, or that it gets so
much better results. Anyone can do obstinacy. Indeed, kids and drunks
and fools are best at it. Whereas very few people have enough of all
five of the qualities that produce the right kind of stubbornness, but
when they do the results are magical.


\textbf{Thanks} to Trevor Blackwell, Jessica Livingston, Jackie
McDonough, Courtenay Pipkin, Harj Taggar, and Garry Tan for reading
drafts of this.
\chapter{Founder Mode}

September 2024

At a YC event last week Brian Chesky gave a talk that everyone who was
there will remember. Most founders I talked to afterward said it was the
best they'd ever heard. Ron Conway, for the first time in his life,
forgot to take notes. I'm not going to try to reproduce it here. Instead
I want to talk about a question it raised.

The theme of Brian's talk was that the conventional wisdom about how to
run larger companies is mistaken. As Airbnb grew, well-meaning people
advised him that he had to run the company in a certain way for it to
scale. Their advice could be optimistically summarized as ``hire good
people and give them room to do their jobs.'' He followed this advice
and the results were disastrous. So he had to figure out a better way on
his own, which he did partly by studying how Steve Jobs ran Apple. So
far it seems to be working. Airbnb's free cash flow margin is now among
the best in Silicon Valley.

The audience at this event included a lot of the most successful
founders we've funded, and one after another said that the same thing
had happened to them. They'd been given the same advice about how to run
their companies as they grew, but instead of helping their companies, it
had damaged them.

Why was everyone telling these founders the wrong thing? That was the
big mystery to me. And after mulling it over for a bit I figured out the
answer: what they were being told was how to run a company you hadn't
founded --- how to run a company if you're merely a professional
manager. But this m.o. is so much less effective that to founders it
feels broken. There are things founders can do that managers can't, and
not doing them feels wrong to founders, because it is.

In effect there are two different ways to run a company: founder mode
and manager mode. Till now most people even in Silicon Valley have
implicitly assumed that scaling a startup meant switching to manager
mode. But we can infer the existence of another mode from the dismay of
founders who've tried it, and the success of their attempts to escape
from it.

There are as far as I know no books specifically about founder mode.
Business schools don't know it exists. All we have so far are the
experiments of individual founders who've been figuring it out for
themselves. But now that we know what we're looking for, we can search
for it. I hope in a few years founder mode will be as well understood as
manager mode. We can already guess at some of the ways it will differ.

The way managers are taught to run companies seems to be like modular
design in the sense that you treat subtrees of the org chart as black
boxes. You tell your direct reports what to do, and it's up to them to
figure out how. But you don't get involved in the details of what they
do. That would be micromanaging them, which is bad.

Hire good people and give them room to do their jobs. Sounds great when
it's described that way, doesn't it? Except in practice, judging from
the report of founder after founder, what this often turns out to mean
is: hire professional fakers and let them drive the company into the
ground.

One theme I noticed both in Brian's talk and when talking to founders
afterward was the idea of being gaslit. Founders feel like they're being
gaslit from both sides --- by the people telling them they have to run
their companies like managers, and by the people working for them when
they do. Usually when everyone around you disagrees with you, your
default assumption should be that you're mistaken. But this is one of
the rare exceptions. VCs who haven't been founders themselves don't know
how founders should run companies, and C-level execs, as a class,
include some of the most skillful liars in the world.
\footnote{The more diplomatic way of phrasing this statement would be to
say that experienced C-level execs are often very skilled at managing
up. And I don't think anyone with knowledge of this world would dispute
that.}

Whatever founder mode consists of, it's pretty clear that it's going to
break the principle that the CEO should engage with the company only via
his or her direct reports. ``Skip-level'' meetings will become the norm
instead of a practice so unusual that there's a name for it. And once
you abandon that constraint there are a huge number of permutations to
choose from.

For example, Steve Jobs used to run an annual retreat for what he
considered the 100 most important people at Apple, and these were not
the 100 people highest on the org chart. Can you imagine the force of
will it would take to do this at the average company? And yet imagine
how useful such a thing could be. It could make a big company feel like
a startup. Steve presumably wouldn't have kept having these retreats if
they didn't work. But I've never heard of another company doing this. So
is it a good idea, or a bad one? We still don't know. That's how little
we know about founder mode. \footnote{If the practice of having such retreats became so widespread
that even mature companies dominated by politics started to do it, we
could quantify the senescence of companies by the average depth on the
org chart of those invited.}

Obviously founders can't keep running a 2000 person company the way they
ran it when it had 20. There's going to have to be some amount of
delegation. Where the borders of autonomy end up, and how sharp they
are, will probably vary from company to company. They'll even vary from
time to time within the same company, as managers earn trust. So founder
mode will be more complicated than manager mode. But it will also work
better. We already know that from the examples of individual founders
groping their way toward it.

Indeed, another prediction I'll make about founder mode is that once we
figure out what it is, we'll find that a number of individual founders
were already most of the way there --- except that in doing what they
did they were regarded by many as eccentric or worse.
\footnote{I also have another less optimistic prediction: as soon as the
concept of founder mode becomes established, people will start misusing
it. Founders who are unable to delegate even things they should will use
founder mode as the excuse. Or managers who aren't founders will decide
they should try to act like founders. That may even work, to some
extent, but the results will be messy when it doesn't; the modular
approach does at least limit the damage a bad CEO can do.}

Curiously enough it's an encouraging thought that we still know so
little about founder mode. Look at what founders have achieved already,
and yet they've achieved this against a headwind of bad advice. Imagine
what they'll do once we can tell them how to run their companies like
Steve Jobs instead of John Sculley.


\textbf{Thanks} to Brian Chesky, Patrick Collison, Ron Conway, Jessica
Livingston, Elon Musk, Ryan Petersen, Harj Taggar, and Garry Tan for
reading drafts of this.
\chapter{When To Do What You Love}

September 2024

There's some debate about whether it's a good idea to ``follow your
passion.'' In fact the question is impossible to answer with a simple
yes or no. Sometimes you should and sometimes you shouldn't, but the
border between should and shouldn't is very complicated. The only way to
give a general answer is to trace it.

When people talk about this question, there's always an implicit
``instead of.'' All other things being equal, why wouldn't you work on
what interests you the most? So even raising the question implies that
all other things aren't equal, and that you have to choose between
working on what interests you the most and something else, like what
pays the best.

And indeed if your main goal is to make money, you can't usually afford
to work on what interests you the most. People pay you for doing what
they want, not what you want. But there's an obvious exception: when you
both want the same thing. For example, if you love football, and you're
good enough at it, you can get paid a lot to play it.

Of course the odds are against you in a case like football, because so
many other people like playing it too. This is not to say you shouldn't
try though. It depends how much ability you have and how hard you're
willing to work.

The odds are better when you have strange tastes: when you like
something that pays well and that few other people like. For example,
it's clear that Bill Gates truly loved running a software company. He
didn't just love programming, which a lot of people do. He loved writing
software for customers. That is a very strange taste indeed, but if you
have it, you can make a lot by indulging it.

There are even some people who have a genuine intellectual interest in
making money. This is distinct from mere greed. They just can't help
noticing when something is mispriced, and can't help doing something
about it. It's like a puzzle for them. \footnote{These examples show why it's a mistake to assume that economic
inequality must be evidence of some kind of brokenness or unfairness.
It's obvious that different people have different interests, and that
some interests yield far more money than others, so how can it not be
obvious that some people will end up much richer than others? In a world
where some people like to write enterprise software and others like to
make studio pottery, economic inequality is the natural outcome.}

In fact there's an edge case here so spectacular that it turns all the
preceding advice on its head. If you want to make a really huge amount
of money --- hundreds of millions or even billions of dollars --- it
turns out to be very useful to work on what interests you the most. The
reason is not the extra motivation you get from doing this, but that the
way to make a really large amount of money is to start a startup, and
working on what interests you is an excellent way to discover
\href{startupideas.html}{startup ideas}.

Many if not most of the biggest startups began as projects the founders
were doing for fun. Apple, Google, and Facebook all began that way. Why
is this pattern so common? Because the best ideas tend to be such
outliers that you'd overlook them if you were consciously looking for
ways to make money. Whereas if you're young and good at technology, your
unconscious instincts about what would be interesting to work on are
very well aligned with what needs to be built.

So there's something like a midwit peak for making money. If you don't
need to make much, you can work on whatever you're most interested in;
if you want to become moderately rich, you can't usually afford to; but
if you want to become super rich, and you're young and good at
technology, working on what you're most interested in becomes a good
idea again.

What if you're not sure what you want? What if you're attracted to the
idea of making money and more attracted to some kinds of work than
others, but neither attraction predominates? How do you break ties?

The key here is to understand that such ties are only apparent. When you
have trouble choosing between following your interests and making money,
it's never because you have complete knowledge of yourself and of the
types of work you're choosing between, and the options are perfectly
balanced. When you can't decide which path to take, it's almost always
due to ignorance. In fact you're usually suffering from three kinds of
ignorance simultaneously: you don't know what makes you happy, what the
various kinds of work are really like, or how well you could do them.
\footnote{Difficulty choosing between interests is a different matter.
That's not always due to ignorance. It's often intrinsically difficult.
I still have trouble doing it.}

In a way this ignorance is excusable. It's often hard to predict these
things, and no one even tells you that you need to. If you're ambitious
you're told you should go to college, and this is good advice so far as
it goes, but that's where it usually ends. No one tells you how to
figure out what to work on, or how hard this can be.

What do you do in the face of uncertainty? Get more certainty. And
probably the best way to do that is to try working on things you're
interested in. That will get you more information about how interested
you are in them, how good you are at them, and how much scope they offer
for ambition.

Don't wait. Don't wait till the end of college to figure out what to
work on. Don't even wait for internships during college. You don't
necessarily need a job doing x in order to work on x; often you can just
start doing it in some form yourself. And since figuring out what to
work on is a problem that could take years to solve, the sooner you
start, the better.

One useful trick for judging different kinds of work is to look at who
your colleagues will be. You'll become like whoever you work with. Do
you want to become like these people?

Indeed, the difference in character between different kinds of work is
magnified by the fact that everyone else is facing the same decisions as
you. If you choose a kind of work mainly for how well it pays, you'll be
surrounded by other people who chose it for the same reason, and that
will make it even more soul-sucking than it seems from the outside.
Whereas if you choose work you're genuinely interested in, you'll be
surrounded mostly by other people who are genuinely interested in it,
and that will make it extra inspiring. \footnote{You can't always take people at their word on this. Since it's
more prestigious to work on things you're interested in than to be
driven by money, people who are driven mainly by money will often claim
to be more interested in their work than they actually are. One way to
test such claims is by doing the following thought experiment: if their
work didn't pay well, would they take day jobs doing something else in
order to do it in their spare time? Lots of mathematicians and
scientists and engineers would. Historically lots \emph{have}. But I
don't think as many investment bankers would.

This thought experiment is also useful for distinguishing between
university departments.}

The other thing you do in the face of uncertainty is to make choices
that are uncertainty-proof. The less sure you are about what to do, the
more important it is to choose options that give you more options in the
future. I call this ``staying upwind.'' If you're unsure whether to
major in math or economics, for example, choose math; math is upwind of
economics in the sense that it will be easier to switch later from math
to economics than from economics to math.

There's one case, though, where it's easy to say whether you should work
on what interests you the most: if you want to do
\href{greatwork.html}{great work}. This is not a sufficient condition
for doing great work, but it is a necessary one.

There's a lot of selection bias in advice about whether to ``follow your
passion,'' and this is the reason. Most such advice comes from people
who are famously successful, and if you ask someone who's famously
successful how to do what they did, most will tell you that you have to
work on what you're most interested in. And this is in fact true.

That doesn't mean it's the right advice for everyone. Not everyone can
do great work, or wants to. But if you do want to, the complicated
question of whether or not to work on what interests you the most
becomes simple. The answer is yes. The root of great work is a sort of
ambitious curiosity, and you can't manufacture that.


\textbf{Thanks} to Trevor Blackwell, Paul Buchheit, Jessica Livingston,
Robert Morris, Harj Taggar, and Garry Tan for reading drafts of this.
\chapter{Writes and Write-Nots}

October 2024

I'm usually reluctant to make predictions about technology, but I feel
fairly confident about this one: in a couple decades there won't be many
people who can write.

One of the strangest things you learn if you're a writer is how many
people have trouble writing. Doctors know how many people have a mole
they're worried about; people who are good at setting up computers know
how many people aren't; writers know how many people need help writing.

The reason so many people have trouble writing is that it's
fundamentally difficult. To write well you have to think clearly, and
thinking clearly is hard.

And yet writing pervades many jobs, and the more prestigious the job,
the more writing it tends to require.

These two powerful opposing forces, the pervasive expectation of writing
and the irreducible difficulty of doing it, create enormous pressure.
This is why eminent professors often turn out to have resorted to
plagiarism. The most striking thing to me about these cases is the
pettiness of the thefts. The stuff they steal is usually the most
mundane boilerplate --- the sort of thing that anyone who was even
halfway decent at writing could turn out with no effort at all. Which
means they're not even halfway decent at writing.

Till recently there was no convenient escape valve for the pressure
created by these opposing forces. You could pay someone to write for
you, like JFK, or plagiarize, like MLK, but if you couldn't buy or steal
words, you had to write them yourself. And as a result nearly everyone
who was expected to write had to learn how.

Not anymore. AI has blown this world open. Almost all pressure to write
has dissipated. You can have AI do it for you, both in school and at
work.

The result will be a world divided into writes and write-nots. There
will still be some people who can write. Some of us like it. But the
middle ground between those who are good at writing and those who can't
write at all will disappear. Instead of good writers, ok writers, and
people who can't write, there will just be good writers and people who
can't write.

Is that so bad? Isn't it common for skills to disappear when technology
makes them obsolete? There aren't many blacksmiths left, and it doesn't
seem to be a problem.

Yes, it's bad. The reason is something I mentioned earlier: writing is
thinking. In fact there's a kind of thinking that can only be done by
writing. You can't make this point better than Leslie Lamport did:

\begin{quote}
If you're thinking without writing, you only think you're thinking.
\end{quote}

So a world divided into writes and write-nots is more dangerous than it
sounds. It will be a world of thinks and think-nots. I know which half I
want to be in, and I bet you do too.

This situation is not unprecedented. In preindustrial times most
people's jobs made them strong. Now if you want to be strong, you work
out. So there are still strong people, but only those who choose to be.

It will be the same with writing. There will still be smart people, but
only those who choose to be.

\textbf{Thanks} to Jessica Livingston, Ben Miller, and Robert Morris for
reading drafts of this.
\chapter{The Origins of Wokeness}

January 2025

The word ``prig'' isn't very common now, but if you look up the
definition, it will sound familiar. Google's isn't bad:

\begin{quote}
A self-righteously moralistic person who behaves as if superior to
others.
\end{quote}

This sense of the word originated in the 18th century, and its age is an
important clue: it shows that although wokeness is a comparatively
recent phenomenon, it's an instance of a much older one.

There's a certain kind of person who's attracted to a shallow, exacting
kind of moral purity, and who demonstrates his purity by attacking
anyone who breaks the rules. Every society has these people. All that
changes is the rules they enforce. In Victorian England it was Christian
virtue. In Stalin's Russia it was orthodox Marxism-Leninism. For the
woke, it's social justice.

So if you want to understand wokeness, the question to ask is not why
people behave this way. Every society has prigs. The question to ask is
why our prigs are priggish about these ideas, at this moment. And to
answer that we have to ask when and where wokeness began.

The answer to the first question is the 1980s. Wokeness is a second,
more aggressive wave of political correctness, which started in the late
1980s, died down in the late 1990s, and then returned with a vengeance
in the early 2010s, finally peaking after the riots of 2020.

What was political correctness, exactly? I'm often asked to define both
this term and wokeness by people who think they're meaningless labels,
so I will. They both have the same definition:

\begin{quote}
An aggressively performative focus on social justice.
\end{quote}

In other words, it's people being prigs about social justice. And that's
the real problem --- the performativeness, not the social justice.
\footnote{This was not the original meaning of ``woke,'' but it's rarely
used in the original sense now. Now the pejorative sense is the dominant
one.}

Racism, for example, is a genuine problem. Not a problem on the scale
that the woke believe it to be, but a genuine one. I don't think any
reasonable person would deny that. The problem with political
correctness was not that it focused on marginalized groups, but the
shallow, aggressive way in which it did so. Instead of going out into
the world and quietly helping members of marginalized groups, the
politically correct focused on getting people in trouble for using the
wrong words to talk about them.

As for where political correctness began, if you think about it, you
probably already know the answer. Did it begin outside universities and
spread to them from this external source? Obviously not; it has always
been most extreme in universities. So where in universities did it
begin? Did it begin in math, or the hard sciences, or engineering, and
spread from there to the humanities and social sciences? Those are
amusing images, but no, obviously it began in the humanities and social
sciences.

Why there? And why then? What happened in the humanities and social
sciences in the 1980s?

A successful theory of the origin of political correctness has to be
able to explain why it didn't happen earlier. Why didn't it happen
during the protest movements of the 1960s, for example? They were
concerned with much the same issues. \footnote{Why did 1960s radicals focus on the causes they did? One of the
people who reviewed drafts of this essay explained this so well that I
asked if I could quote him:

\begin{quote}
The middle-class student protestors of the New Left rejected the
socialist/Marxist left as unhip. They were interested in sexier forms of
oppression uncovered by cultural analysis (Marcuse) and abstruse
``Theory''. Labor politics became stodgy and old-fashioned. This took a
couple generations to work through. The woke ideology's conspicuous lack
of interest in the working class is the tell-tale sign. Such fragments
as are, er, left of the old left are anti-woke, and meanwhile the actual
working class shifted to the populist right and gave us Trump. Trump and
wokeness are cousins.

The middle-class origins of wokeness smoothed its way through the
institutions because it had no interest in ``seizing the means of
production'' (how quaint such phrases seem now), which would quickly
have run up against hard state and corporate power. The fact that
wokeness only expressed interest in other kinds of class (race, sex,
etc) signalled compromise with existing power: give us power within your
system and we'll bestow the resource we control --- moral rectitude ---
upon you. As an ideological stalking horse for gaining control over
discourse and institutions, this succeeded where a more ambitious
revolutionary program would not have.
\end{quote}}

The reason the student protests of the 1960s didn't lead to political
correctness was precisely that --- they were student movements. They
didn't have any real power. The students may have been talking a lot
about women's liberation and black power, but it was not what they were
being taught in their classes. Not yet.

But in the early 1970s the student protestors of the 1960s began to
finish their dissertations and get hired as professors. At first they
were neither powerful nor numerous. But as more of their peers joined
them and the previous generation of professors started to retire, they
gradually became both.

The reason political correctness began in the humanities and social
sciences was that these fields offered more scope for the injection of
politics. A 1960s radical who got a job as a physics professor could
still attend protests, but his political beliefs wouldn't affect his
work. Whereas research in sociology and modern literature can be made as
political as you like. \footnote{It helped that the humanities and social sciences also included
some of the biggest and easiest undergrad majors. If a political
movement had to start with physics students, it could never get off the
ground; there would be too few of them, and they wouldn't have the time
to spare.

At the top universities these majors are not as big as they used to be,
though. A
\href{https://www.newyorker.com/magazine/2023/03/06/the-end-of-the-english-major}{2022
survey} found that only 7\% of Harvard undergrads plan to major in the
humanities, vs nearly 30\% during the 1970s. I expect wokeness is at
least part of the reason; when undergrads consider majoring in English,
it's presumably because they love the written word and not because they
want to listen to lectures about racism.}

I saw political correctness arise. When I started college in 1982 it was
not yet a thing. Female students might object if someone said something
they considered sexist, but no one was getting \emph{reported} for it.
It was still not a thing when I started grad school in 1986. It was
definitely a thing in 1988 though, and by the early 1990s it seemed to
pervade campus life.

What happened? How did protest become punishment? Why were the late
1980s the point at which protests against male chauvinism (as it used to
be called) morphed into formal complaints to university authorities
about sexism? Basically, the 1960s radicals got tenure. They became the
Establishment they'd protested against two decades before. Now they were
in a position not just to speak out about their ideas, but to enforce
them.

A new set of moral rules to enforce was exciting news to a certain kind
of student. What made it particularly exciting was that they were
allowed to attack professors. I remember noticing that aspect of
political correctness at the time. It wasn't simply a grass-roots
student movement. It was faculty members encouraging students to attack
other faculty members. In that respect it was like the Cultural
Revolution. That wasn't a grass-roots movement either; that was Mao
unleashing the younger generation on his political opponents. And in
fact when Roderick MacFarquhar started teaching a class on the Cultural
Revolution at Harvard in the late 1980s, many saw it as a comment on
current events. I don't know if it actually was, but people thought it
was, and that means the similarities were obvious.
\footnote{The puppet-master-and-puppet character of political correctness
became clearly visible when a bakery near Oberlin College was falsely
accused of race discrimination in 2016. In the subsequent civil trial,
lawyers for the bakery produced a text message from Oberlin Dean of
Students Meredith Raimondo that read ``I'd say unleash the students if I
wasn't convinced this needs to be put behind us.''}

College students larp. It's their nature. It's usually harmless. But
larping morality turned out to be a poisonous combination. The result
was a kind of moral etiquette, superficial but very complicated. Imagine
having to explain to a well-meaning visitor from another planet why
using the phrase ``people of color'' is considered particularly
enlightened, but saying ``colored people'' gets you fired. And why
exactly one isn't supposed to use the word ``negro'' now, even though
Martin Luther King used it constantly in his speeches. There are no
underlying principles. You'd just have to give him a long list of rules
to memorize. \footnote{The woke sometimes claim that wokeness is simply treating people
with respect. But if it were, that would be the only rule you'd have to
remember, and this is comically far from being the case. My younger son
likes to imitate voices, and at one point when he was about seven I had
to explain which accents it was currently safe to imitate publicly and
which not. It took about ten minutes, and I still hadn't covered all the
cases.}

The danger of these rules was not just that they created land mines for
the unwary, but that their elaborateness made them an effective
substitute for virtue. Whenever a society has a concept of heresy and
orthodoxy, orthodoxy becomes a substitute for virtue. You can be the
worst person in the world, but as long as you're orthodox you're better
than everyone who isn't. This makes orthodoxy very attractive to bad
people.

But for it to work as a substitute for virtue, orthodoxy must be
difficult. If all you have to do to be orthodox is wear some garment or
avoid saying some word, everyone knows to do it, and the only way to
seem more virtuous than other people is to actually be virtuous. The
shallow, complicated, and frequently changing rules of political
correctness made it the perfect substitute for actual virtue. And the
result was a world in which good people who weren't up to date on
current moral fashions were brought down by people whose characters
would make you recoil in horror if you could see them.

One big contributing factor in the rise of political correctness was the
lack of other things to be morally pure about. Previous generations of
prigs had been prigs mostly about religion and sex. But among the
cultural elite these were the deadest of dead letters by the 1980s; if
you were religious, or a virgin, this was something you tended to
conceal rather than advertise. So the sort of people who enjoy being
moral enforcers had become starved of things to enforce. A new set of
rules was just what they'd been waiting for.

Curiously enough, the tolerant side of the 1960s left helped create the
conditions in which the intolerant side prevailed. The relaxed social
rules advocated by the old, easy-going hippy left became the dominant
ones, at least among the elite, and this left nothing for the naturally
intolerant to be intolerant about.

Another possibly contributing factor was the fall of the Soviet empire.
Marxism had been a popular focus of moral purity on the left before
political correctness emerged as a competitor, but the pro-democracy
movements in Eastern Bloc countries took most of the shine off it.
Especially the fall of the Berlin Wall in 1989. You couldn't be on the
side of the Stasi. I remember looking at the moribund Soviet Studies
section of a used bookshop in Cambridge in the late 1980s and thinking
``what will those people go on about now?'' As it turned out the answer
was right under my nose.

One thing I noticed at the time about the first phase of political
correctness was that it was more popular with women than men. As many
writers (perhaps most eloquently George Orwell) have observed, women
seem more attracted than men to the idea of being moral enforcers. But
there was another more specific reason women tended to be the enforcers
of political correctness. There was at this time a great backlash
against sexual harassment; the mid 1980s were the point when the
definition of sexual harassment was expanded from explicit sexual
advances to creating a ``hostile environment.'' Within universities the
classic form of accusation was for a (female) student to say that a
professor made her ``feel uncomfortable.'' But the vagueness of this
accusation allowed the radius of forbidden behavior to expand to include
talking about heterodox ideas. Those make people uncomfortable too.
\footnote{In 1986 the Supreme Court ruled that creating a hostile work
environment could constitute sex discrimination, which in turn affected
universities via Title IX. The court specified that the test of a
hostile environment was whether it would bother a reasonable person, but
since for a professor merely being the subject of a sexual harassment
complaint would be a disaster whether the complainant was reasonable or
not, in practice any joke or remark remotely connected with sex was now
effectively forbidden. Which meant we'd now come full circle to
Victorian codes of behavior, when there was a large class of things that
might not be said ``with ladies present.''}

Was it sexist to propose that Darwin's greater male variability
hypothesis might explain some variation in human performance? Sexist
enough to get Larry Summers pushed out as president of Harvard,
apparently. One woman who heard the talk in which he mentioned this idea
said it made her feel ``physically ill'' and that she had to leave
halfway through. If the test of a hostile environment is how it makes
people feel, this certainly sounds like one. And yet it does seem
plausible that greater male variability explains some of the variation
in human performance. So which should prevail, comfort or truth? Surely
if truth should prevail anywhere, it should be in universities; that's
supposed to be their specialty; but for decades starting in the late
1980s the politically correct tried to pretend this conflict didn't
exist. \footnote{Much as they tried to pretend there was no conflict between
diversity and quality. But you can't simultaneously optimize for two
things that aren't identical. What diversity actually means, judging
from the way the term is used, is proportional representation, and
unless you're selecting a group whose purpose is to be representative,
like poll respondents, optimizing for proportional representation has to
come at the expense of quality. This is not because of anything about
representation; it's the nature of optimization; optimizing for x has to
come at the expense of y unless x and y are identical.}

Political correctness seemed to burn out in the second half of the
1990s. One reason, perhaps the main reason, was that it literally became
a joke. It offered rich material for comedians, who performed their
usual disinfectant action upon it. Humor is one of the most powerful
weapons against priggishness of any sort, because prigs, being
humorless, can't respond in kind. Humor was what defeated Victorian
prudishness, and by 2000 it seemed to have done the same thing to
political correctness.

Unfortunately this was an illusion. Within universities the embers of
political correctness were still glowing brightly. After all, the forces
that created it were still there. The professors who started it were now
becoming deans and department heads. And in addition to their
departments there were now a bunch of new ones explicitly focused on
social justice. Students were still hungry for things to be morally pure
about. And there had been an explosion in the number of university
administrators, many of whose jobs involved enforcing various forms of
political correctness.

In the early 2010s the embers of political correctness burst into flame
anew. There were several differences between this new phase and the
original one. It was more virulent. It spread further into the real
world, although it still burned hottest within universities. And it was
concerned with a wider variety of sins. In the first phase of political
correctness there were really only three things people got accused of:
sexism, racism, and homophobia (which at the time was a neologism
invented for the purpose). But between then and 2010 a lot of people had
spent a lot of time trying to invent new kinds of -isms and -phobias and
seeing which could be made to stick.

The second phase was, in multiple senses, political correctness
metastasized. Why did it happen when it did? My guess is that it was due
to the rise of social media, particularly Tumblr and Twitter, because
one of the most distinctive features of the second wave of political
correctness was the \emph{cancel mob}: a mob of angry people uniting on
social media to get someone ostracized or fired. Indeed this second wave
of political correctness was originally called ``cancel culture''; it
didn't start to be called ``wokeness'' till the 2020s.

One aspect of social media that surprised almost everyone at first was
the popularity of outrage. Users seemed to \emph{like} being outraged.
We're so used to this idea now that we take it for granted, but really
it's pretty strange. Being outraged is not a pleasant feeling. You
wouldn't expect people to seek it out. But they do. And above all, they
want to share it. I happened to be running a forum from 2007 to 2014, so
I can actually quantify how much they want to share it: our users were
about three times more likely to upvote something if it outraged them.

This tilt toward outrage wasn't due to wokeness. It's an inherent
feature of social media, or at least this generation of it. But it did
make social media the perfect mechanism for fanning the flames of
wokeness. \footnote{Maybe societies will eventually develop antibodies to viral
outrage. Maybe we were just the first to be exposed to it, so it tore
through us like an epidemic through a previously isolated population.
I'm fairly confident that it would be possible to create new social
media apps that were less driven by outrage, and an app of this type
would have a good chance of stealing users from existing ones, because
the smartest people would tend to migrate to it.}

It wasn't just public social networks that drove the rise of wokeness
though. Group chat apps were also critical, especially in the final
step, cancellation. Imagine if a group of employees trying to get
someone fired had to do it using only email. It would be hard to
organize a mob. But once you have group chat, mobs form naturally.

Another contributing factor in this second wave of political correctness
was the dramatic increase in the polarization of the press. In the print
era, newspapers were constrained to be, or at least seem, politically
neutral. The department stores that ran ads in the New York Times wanted
to reach everyone in the region, both liberal and conservative, so the
Times had to serve both. But the Times didn't regard this neutrality as
something forced upon them. They embraced it as their duty as a
\emph{paper of record} --- as one of the big newspapers that aimed to be
chronicles of their times, reporting every sufficiently important story
from a neutral point of view.

When I grew up the papers of record seemed timeless, almost sacred
institutions. Papers like the New York Times and Washington Post had
immense prestige, partly because other sources of news were limited, but
also because they did make some effort to be neutral.

Unfortunately it turned out that the paper of record was mostly an
artifact of the constraints imposed by print. \footnote{I say ``mostly'' because I have hopes that journalistic
neutrality will return in some form. There is some market for unbiased
news, and while it may be small, it's valuable. The rich and powerful
want to know what's really going on; that's how they became rich and
powerful.}
When your market was determined by geography, you had to be neutral. But
publishing online enabled --- in fact probably forced --- newspapers to
switch to serving markets defined by ideology instead of geography. Most
that remained in business fell in the direction they'd already been
leaning: left. On October 11, 2020 the New York Times announced that
``The paper is in the midst of an evolution from the stodgy paper of
record into a juicy collection of great narratives.''
\footnote{The Times made this momentous announcement very informally, in
passing in the middle of an
\href{https://www.nytimes.com/2020/10/11/business/media/new-york-times-rukmini-callimachi-caliphate.html}{article}
about a Times reporter who'd been criticized for inaccuracy. It's quite
possible no senior editor even approved it. But it's somehow appropriate
that this particular universe ended with a whimper rather than a bang.} Meanwhile journalists, of a sort, had arisen to
serve the right as well. And so journalism, which in the previous era
had been one of the great centralizing forces, now became one of the
great polarizing ones.

The rise of social media and the increasing polarization of journalism
reinforced one another. In fact there arose a new variety of journalism
involving a loop through social media. Someone would say something
controversial on social media. Within hours it would become a news
story. Outraged readers would then post links to the story on social
media, driving further arguments online. It was the cheapest source of
clicks imaginable. You didn't have to maintain overseas news bureaus or
pay for month-long investigations. All you had to do was watch Twitter
for controversial remarks and repost them on your site, with some
additional comments to inflame readers further.

For the press there was money in wokeness. But they weren't the only
ones. That was one of the biggest differences between the two waves of
political correctness: the first was driven almost entirely by amateurs,
but the second was often driven by professionals. For some it was their
whole job. By 2010 a new class of administrators had arisen whose job
was basically to enforce wokeness. They played a role similar to that of
the political commissars who got attached to military and industrial
organizations in the USSR: they weren't directly in the flow of the
organization's work, but watched from the side to ensure that nothing
improper happened in the doing of it. These new administrators could
often be recognized by the word ``inclusion'' in their titles. Within
institutions this was the preferred euphemism for wokeness; a new list
of banned words, for example, would usually be called an ``inclusive
language guide.'' \footnote{As the acronym DEI goes out of fashion, many of these
bureaucrats will try to go underground by changing their titles. It
looks like ``belonging'' will be a popular option.}

This new class of bureaucrats pursued a woke agenda as if their jobs
depended on it, because they did. If you hire people to keep watch for a
particular type of problem, they're going to find it, because otherwise
there's no justification for their existence. \footnote{If you've ever wondered why our legal system includes
protections like the separation of prosecutor, judge, and jury, the
right to examine evidence and cross-examine witnesses, and the right to
be represented by legal counsel, the de facto
\href{https://www.nytimes.com/2020/03/18/magazine/title-ix-sexual-harassment-accusations.html}{parallel
legal system} established by Title IX makes that all too clear.}
But these bureaucrats also represented a second and possibly even
greater danger. Many were involved in hiring, and when possible they
tried to ensure their employers hired only people who shared their
political beliefs. The most egregious cases were the new ``DEI
statements'' that some universities started to require from faculty
candidates, proving their commitment to wokeness. Some universities used
these statements as the initial filter and only even considered
candidates who scored high enough on them. You're not hiring Einstein
that way; imagine what you get instead.

Another factor in the rise of wokeness was the Black Lives Matter
movement, which started in 2013 when a white man was acquitted after
killing a black teenager in Florida. But this didn't launch wokeness; it
was well underway by 2013.

Similarly for the Me Too Movement, which took off in 2017 after the
first news stories about Harvey Weinstein's history of raping women. It
accelerated wokeness, but didn't play the same role in launching it that
the 80s version did in launching political correctness.

The election of Donald Trump in 2016 also accelerated wokeness,
particularly in the press, where outrage now meant traffic. Trump made
the New York Times a lot of money: headlines during his first
administration mentioned his name at about four times the rate of
previous presidents.

In 2020 we saw the biggest accelerant of all, after a white police
officer asphyxiated a black suspect on video. At this point the
metaphorical fire became a literal one, as violent protests broke out
across America. But in retrospect this turned out to be peak woke, or
close to it. By every measure I've seen, wokeness peaked in 2020 or
2021.

Wokeness is sometimes described as a mind-virus. What makes it viral is
that it defines new types of impropriety. Most people are afraid of
impropriety; they're never exactly sure what the social rules are or
which ones they might be breaking. Especially if the rules change
rapidly. And since most people already worry that they might be breaking
rules they don't know about, if you tell them they're breaking a rule,
their default reaction is to believe you. Especially if multiple people
tell them. Which in turn is a recipe for exponential growth. Zealots
invent some new impropriety to avoid. The first people to adopt it are
fellow zealots, eager for new ways to signal their virtue. If there are
enough of these, the initial group of zealots is followed by a much
larger group, motivated by fear. They're not trying to signal virtue;
they're just trying to avoid getting in trouble. At this point the new
impropriety is now firmly established. Plus its success has increased
the rate of change in social rules, which, remember, is one of the
reasons people are nervous about which rules they might be breaking. So
the cycle accelerates. \footnote{The invention of new improprieties is most visible in the rapid
evolution of woke nomenclature. This is particularly annoying to me as a
writer, because the new names are always worse. Any religious observance
has to be inconvenient and slightly absurd; otherwise gentiles would do
it too. So ``slaves'' becomes ``enslaved individuals.'' But web search
can show us the leading edge of moral growth in real time: if you search
for ``individuals experiencing slavery'' you will as of this writing
find five legit attempts to use the phrase, and you'll even find two for
``individuals experiencing enslavement.''}

What's true of individuals is even more true of organizations.
Especially organizations without a powerful leader. Such organizations
do everything based on ``best practices.'' There's no higher authority;
if some new ``best practice'' achieves critical mass, they \emph{must}
adopt it. And in this case the organization can't do what it usually
does when it's uncertain: delay. It might be committing improprieties
right now! So it's surprisingly easy for a small group of zealots to
capture this type of organization by describing new improprieties it
might be guilty of. \footnote{Organizations that do dubious things are particularly concerned
with propriety, which is how you end up with absurdities like tobacco
and oil companies having higher ESG ratings than Tesla.}

How does this kind of cycle ever end? Eventually it leads to disaster,
and people start to say enough is enough. The excesses of 2020 made a
lot of people say that.

Since then wokeness has been in gradual but continual retreat. Corporate
CEOs, starting with Brian Armstrong, have openly rejected it.
Universities, led by the University of Chicago and MIT, have explicitly
confirmed their commitment to free speech. Twitter, which was arguably
the hub of wokeness, was bought by Elon Musk in order to neutralize it,
and he seems to have succeeded --- and not, incidentally, by censoring
left-wing users the way Twitter used to censor right-wing ones, but
without censoring either. \footnote{Elon did something else that tilted Twitter rightward though:
he gave more visibility to paying users. Paying users lean right on
average, because people on the far left dislike Elon and don't want to
give him money. Elon presumably knew this would happen. On the other
hand, the people on the far left have only themselves to blame; they
could tilt Twitter back to the left tomorrow if they wanted to.} Consumers have
emphatically rejected brands that ventured too far into wokeness. The
Bud Light brand may have been permanently damaged by it. I'm not going
to claim Trump's second victory in 2024 was a referendum on wokeness; I
think he won, as presidential candidates always do, because he was more
\href{charisma.html}{charismatic}; but voters' disgust with wokeness
must have helped.

So what do we do now? Wokeness is already in retreat. Obviously we
should help it along. What's the best way to do that? And more
importantly, how do we avoid a third outbreak? After all, it seemed to
be dead once, but came back worse than ever.

In fact there's an even more ambitious goal: is there a way to prevent
any similar outbreak of aggressively performative moralism in the future
--- not just a third outbreak of political correctness, but the next
thing like it? Because there will be a next thing. Prigs are prigs by
nature. They need rules to obey and enforce, and now that Darwin has cut
off their traditional supply of rules, they're constantly hungry for new
ones. All they need is someone to meet them halfway by defining a new
way to be morally pure, and we'll see the same phenomenon again.

Let's start with the easier problem. Is there a simple, principled way
to deal with wokeness? I think there is: to use the customs we already
have for dealing with religion. Wokeness is effectively a religion, just
with God replaced by protected classes. It's not even the first religion
of this kind; Marxism had a similar form, with God replaced by the
masses. \footnote{It even, as James Lindsay and Peter Boghossian pointed out, has
a concept of original sin: privilege. Which means unlike Christianity's
egalitarian version, people have varying degrees of it. An able-bodied
straight white American male is born with such a load of sin that only
by the most abject repentance can he be saved.

Wokeness also shares something rather funny with many actual versions of
Christianity: like God, the people for whose sake wokeness purports to
act are often revolted by the things done in their name.} And we already have well-established
customs for dealing with religion within organizations. You can express
your own religious identity and explain your beliefs, but you can't call
your coworkers infidels if they disagree, or try to ban them from saying
things that contradict its doctrines, or insist that the organization
adopt yours as its official religion.

If we're not sure what to do about any particular manifestation of
wokeness, imagine we were dealing with some other religion, like
Christianity. Should we have people within organizations whose jobs are
to enforce woke orthodoxy? No, because we wouldn't have people whose
jobs were to enforce Christian orthodoxy. Should we censor
\href{https://www.telegraph.co.uk/news/2023/02/17/roald-dahl-books-rewritten-offensive-matilda-witches-twits/}{writers}
or
\href{https://x.com/KoustaStavroula/status/1562034502897106946}{scientists}
whose work contradicts woke doctrines? No, because we wouldn't do this
to people whose work contradicted Christian teachings. Should job
candidates be required to write DEI statements? Of course not; imagine
an employer requiring proof of one's religious beliefs. Should students
and employees have to participate in woke indoctrination sessions in
which they're required to answer questions about their beliefs to ensure
compliance? No, because we wouldn't dream of catechizing people in this
way about their religion. \footnote{There is one exception to most of these rules: actual religious
organizations. It's reasonable for them to insist on orthodoxy. But they
in turn should declare that they're religious organizations. It's
rightly considered shady when something that appears to be an ordinary
business or publication turns out to be a religious organization.}

One shouldn't feel bad about not wanting to watch woke movies any more
than one would feel bad about not wanting to listen to Christian rock.
In my twenties I drove across America several times, listening to local
radio stations. Occasionally I'd turn the dial and hear some new song.
But the moment anyone mentioned Jesus I'd turn the dial again. Even the
tiniest bit of being preached to was enough to make me lose interest.

But by the same token we should not automatically reject everything the
woke believe. I'm not a Christian, but I can see that many Christian
principles are good ones. It would be a mistake to discard them all just
because one didn't share the religion that espoused them. It would be
the sort of thing a religious zealot would do.

If we have genuine pluralism, I think we'll be safe from future
outbreaks of woke intolerance. Wokeness itself won't go away. There will
for the foreseeable future continue to be pockets of woke zealots
inventing new moral fashions. The key is not to let them treat their
fashions as normative. They can change what their coreligionists are
allowed to say every few months if they like, but they mustn't be
allowed to change what we're allowed to say. \footnote{I don't want to give the impression that it will be simple to
roll back wokeness. There will be places where the fight inevitably gets
messy --- particularly within universities, which everyone has to share,
yet which are currently the most pervaded by wokeness of any
institutions.}

The more general problem --- how to prevent similar outbreaks of
aggressively performative moralism --- is of course harder. Here we're
up against human nature. There will always be prigs. And in particular
there will always be the enforcers among them, the
\href{conformism.html}{aggressively conventional-minded}. These people
are born that way. Every society has them. So the best we can do is to
keep them bottled up.

The aggressively conventional-minded aren't always on the rampage.
Usually they just enforce whatever random rules are nearest to hand.
They only become dangerous when some new ideology gets a lot of them
pointed in the same direction at once. That's what happened during the
Cultural Revolution, and to a lesser extent (thank God) in the two waves
of political correctness we've experienced.

We can't get rid of the aggressively conventional-minded.
\footnote{You can however get rid of aggressively conventional-minded
people within an organization, and in many if not most organizations
this would be an excellent idea. Even a handful of them can do a lot of
damage. I bet you'd feel a noticeable improvement going from a handful
to none.} And we couldn't prevent people from creating
new ideologies that appealed to them even if we wanted to. So if we want
to keep them bottled up, we have to do it one step downstream.
Fortunately when the aggressively conventional-minded go on the rampage
they always do one thing that gives them away: they define new
\href{heresy.html}{heresies} to punish people for. So the best way to
protect ourselves from future outbreaks of things like wokeness is to
have powerful antibodies against the concept of heresy.

We should have a conscious bias against defining new forms of heresy.
Whenever anyone tries to ban saying something that we'd previously been
able to say, our initial assumption should be that they're wrong. Only
our initial assumption of course. If they can prove we should stop
saying it, then we should. But the burden of proof is on them. In
liberal democracies, people trying to prevent something from being said
will usually claim they're not merely engaging in censorship, but trying
to prevent some form of ``harm''. And maybe they're right. But once
again, the burden of proof is on them. It's not enough to claim harm;
they have to prove it.

As long as the aggressively conventional-minded continue to give
themselves away by banning heresies, we'll always be able to notice when
they become aligned behind some new ideology. And if we always fight
back at that point, with any luck we can stop them in their tracks.

The number of true things we \href{say.html}{can't say} should not
increase. If it does, something is wrong.


\textbf{Thanks} to Sam Altman, Ben Miller, Daniel Gackle, Robin Hanson,
Jessica Livingston, Greg Lukianoff, Harj Taggar, Garry Tan, and Tim
Urban for reading drafts of this.
\chapter{What to Do}

March 2025

What should one do? That may seem a strange question, but it's not
meaningless or unanswerable. It's the sort of question kids ask before
they learn not to ask big questions. I only came across it myself in the
process of investigating something else. But once I did, I thought I
should at least try to answer it.

So what \emph{should} one do? One should help people, and take care of
the world. Those two are obvious. But is there anything else? When I ask
that, the answer that pops up is \emph{Make good new things}.

I can't prove that one should do this, any more than I can prove that
one should help people or take care of the world. We're talking about
first principles here. But I can explain why this principle makes sense.
The most impressive thing humans can do is to think. It may be the most
impressive thing that can be done. And the best kind of thinking, or
more precisely the best proof that one has thought well, is to make good
new things.

I mean new things in a very general sense. Newton's physics was a good
new thing. Indeed, the first version of this principle was to have good
new ideas. But that didn't seem general enough: it didn't include making
art or music, for example, except insofar as they embody new ideas. And
while they may embody new ideas, that's not all they embody, unless you
stretch the word ``idea'' so uselessly thin that it includes everything
that goes through your nervous system.

Even for ideas that one has consciously, though, I prefer the phrasing
``make good new things.'' There are other ways to describe the best kind
of thinking. To make discoveries, for example, or to understand
something more deeply than others have. But how well do you understand
something if you can't make a model of it, or write about it? Indeed,
trying to express what you understand is not just a way to prove that
you understand it, but a way to understand it better.

Another reason I like this phrasing is that it biases us toward
creation. It causes us to prefer the kind of ideas that are naturally
seen as making things rather than, say, making critical observations
about things other people have made. Those are ideas too, and sometimes
valuable ones, but it's easy to trick oneself into believing they're
more valuable than they are. Criticism seems sophisticated, and making
new things often seems awkward, especially at first; and yet it's
precisely those first steps that are most rare and valuable.

Is newness essential? I think so. Obviously it's essential in science.
If you copied a paper of someone else's and published it as your own, it
would seem not merely unimpressive but dishonest. And it's similar in
the arts. A copy of a good painting can be a pleasing thing, but it's
not impressive in the way the original was. Which in turn implies it's
not impressive to make the same thing over and over, however well;
you're just copying yourself.

Note though that we're talking about a different kind of should with
this principle. Taking care of people and the world are shoulds in the
sense that they're one's duty, but making good new things is a should in
the sense that this is how to live to one's full potential. Historically
most rules about how to live have been a mix of both kinds of should,
though usually with more of the former than the latter.
\footnote{We could treat all three as the same kind of should by saying
that it's one's duty to live well --- for example by saying, as some
Christians have, that it's one's duty to make the most of one's
God-given gifts. But this seems one of those casuistries people invented
to evade the stern requirements of religion: it was permissible to spend
time studying math instead of praying or performing acts of charity
because otherwise you were rejecting a gift God had given you. A useful
casuistry no doubt, but we don't need it.

We could also combine the first two principles, since people are part of
the world. Why should our species get special treatment? I won't try to
justify this choice, but I'm skeptical that anyone who claims to think
differently actually lives according to their principles.}

For most of history the question ``What should one do?'' got much the
same answer everywhere, whether you asked Cicero or Confucius. You
should be wise, brave, honest, temperate, and just, uphold tradition,
and serve the public interest. There was a long stretch where in some
parts of the world the answer became ``Serve God,'' but in practice it
was still considered good to be wise, brave, honest, temperate, and
just, uphold tradition, and serve the public interest. And indeed this
recipe would have seemed right to most Victorians. But there's nothing
in it about taking care of the world or making new things, and that's a
bit worrying, because it seems like this question should be a timeless
one. The answer shouldn't change much.

I'm not too worried that the traditional answers don't mention taking
care of the world. Obviously people only started to care about that once
it became clear we could ruin it. But how can making good new things be
important if the traditional answers don't mention it?

The traditional answers were answers to a slightly different question.
They were answers to the question of how to be, rather than what to do.
The audience didn't have a lot of choice about what to do. The audience
up till recent centuries was the landowning class, which was also the
political class. They weren't choosing between doing physics and writing
novels. Their work was foreordained: manage their estates, participate
in politics, fight when necessary. It was ok to do certain other kinds
of work in one's spare time, but ideally one didn't have any. Cicero's
\emph{De Officiis} is one of the great classical answers to the question
of how to live, and in it he explicitly says that he wouldn't even be
writing it if he hadn't been excluded from public life by recent
political upheavals. \footnote{Confucius was also excluded from public life after ending up on
the losing end of a power struggle, and presumably he too would not be
so famous now if it hadn't been for this long stretch of enforced
leisure.}

There were of course people doing what we would now call ``original
work,'' and they were often admired for it, but they weren't seen as
models. Archimedes knew that he was the first to prove that a sphere has
2/3 the volume of the smallest enclosing cylinder and was very pleased
about it. But you don't find ancient writers urging their readers to
emulate him. They regarded him more as a prodigy than a model.

Now many more of us can follow Archimedes's example and devote most of
our attention to one kind of work. He turned out to be a model after
all, along with a collection of other people that his contemporaries
would have found it strange to treat as a distinct group, because the
vein of people making new things ran at right angles to the social
hierarchy.

What kinds of new things count? I'd rather leave that question to the
makers of them. It would be a risky business to try to define any kind
of threshold, because new kinds of work are often despised at first.
Raymond Chandler was writing literal pulp fiction, and he's now
recognized as one of the best writers of the twentieth century. Indeed
this pattern is so common that you can use it as a recipe: if you're
excited about some kind of work that's not considered prestigious and
you can explain what everyone else is overlooking about it, then this is
not merely a kind of work that's ok to do, but one to seek out.

The other reason I wouldn't want to define any thresholds is that we
don't need them. The kind of people who make good new things don't need
rules to keep them honest.

So there's my guess at a set of principles to live by: take care of
people and the world, and make good new things. Different people will do
these to varying degrees. There will presumably be lots who focus
entirely on taking care of people. There will be a few who focus mostly
on making new things. But even if you're one of those, you should at
least make sure that the new things you make don't net \emph{harm}
people or the world. And if you go a step further and try to make things
that help them, you may find you're ahead on the trade. You'll be more
constrained in what you can make, but you'll make it with more energy.

On the other hand, if you make something amazing, you'll often be
helping people or the world even if you didn't mean to. Newton was
driven by curiosity and ambition, not by any practical effect his work
might have, and yet the practical effect of his work has been enormous.
And this seems the rule rather than the exception. So if you think you
can make something amazing, you should probably just go ahead and do it.


\textbf{Thanks} to Trevor Blackwell, Jessica Livingston, and Robert
Morris for reading drafts of this.
\chapter{Good Writing}

May 2025

There are two senses in which writing can be good: it can sound good,
and the ideas can be right. It can have nice, flowing sentences, and it
can draw correct conclusions about important things. It might seem as if
these two kinds of good would be unrelated, like the speed of a car and
the color it's painted. And yet I don't think they are. I think writing
that sounds good is more likely to be right.

So here we have the most exciting kind of idea: one that seems both
preposterous and true. Let's examine it. How can this possibly be true?

I know it's true from writing. You can't simultaneously optimize two
unrelated things; when you push one far enough, you always end up
sacrificing the other. And yet no matter how hard I push, I never find
myself having to choose between the sentence that sounds best and the
one that expresses an idea best. If I did, it would be frivolous to care
how sentences sound. But in practice it feels the opposite of frivolous.
Fixing sentences that sound bad seems to help get the ideas right.
\footnote{The closest thing to an exception is when you have to go back
and insert a new point into the middle of something you've written. This
often messes up the flow, sometimes in ways you can never quite repair.
But I think the ultimate source of this problem is that ideas are
tree-shaped and essays are linear. You inevitably run into difficulties
when you try to cram the former into the latter. Frankly it's surprising
how much you can get away with. But even so you sometimes have to resort
to an endnote.}

By right I mean more than just true. Getting the ideas right means
developing them well --- drawing the conclusions that matter most, and
exploring each one to the right level of detail. So getting the ideas
right is not just a matter of saying true things, but saying the right
true things.

How could trying to make sentences sound good help you do that? The clue
to the answer is something I noticed 30 years ago when I was doing the
layout for my first book. Sometimes when you're laying out text you have
bad luck. For example, you get a section that runs one line longer than
the page. I don't know what ordinary typesetters do in this situation,
but what I did was rewrite the section to make it a line shorter. You'd
expect such an arbitrary constraint to make the writing worse. But I
found, to my surprise, that it never did. I always ended up with
something I liked better.

I don't think this was because my writing was especially careless. I
think if you pointed to a random paragraph in anything written by anyone
and told them to make it slightly shorter (or longer), they'd probably
be able to come up with something better.

The best analogy for this phenomenon is when you shake a bin full of
different objects. The shakes are arbitrary motions. Or more precisely,
they're not calculated to make any two specific objects fit more closely
together. And yet repeated shaking inevitably makes the objects discover
brilliantly clever ways of packing themselves. Gravity won't let them
become less tightly packed, so any change has to be a change for the
better. \footnote{Obviously if you shake the bin hard enough the objects in it can
become less tightly packed. And similarly, if you imposed some huge
external constraint on your writing, like using alternating one and two
syllable words, the ideas would start to suffer.}

So it is with writing. If you have to rewrite an awkward passage, you'll
never do it in a way that makes it \emph{less} true. You couldn't bear
it, any more than gravity could bear things floating upward. So any
change in the ideas has to be a change for the better.

It's obvious once you think about it. Writing that sounds good is more
likely to be right for the same reason that a well-shaken bin is more
likely to be tightly packed. But there's something else going on as
well. Sounding good isn't just a random external force that leaves the
ideas in an essay better off. It actually helps you to get them right.

The reason is that it makes the essay easier to read. It's less work to
read writing that flows well. How does that help the writer?
\emph{Because the writer is the first reader.} When I'm working on an
essay, I spend far more time reading than writing. I'll reread some
parts 50 or 100 times, replaying the thoughts in them and asking myself,
like someone sanding a piece of wood, does anything catch? Does anything
feel wrong? And the easier the essay is to read, the easier it is to
notice if something catches.

So yes, the two senses of good writing are connected in at least two
ways. Trying to make writing sound good makes you fix mistakes
unconsciously, and also helps you fix them consciously; it shakes the
bin of ideas, and also makes mistakes easier to see. But now that we've
dissolved one layer of preposterousness, I can't resist adding another.
Does sounding good do more than just help you get the ideas right? Is
writing that sounds good \emph{inherently} more likely to be right?
Crazy as it may seem, I think that's true too.

Obviously there's a connection at the level of individual words. There
are lots of words in English that sound like what they mean, often in
wonderfully subtle ways. Glitter. Round. Scrape. Prim. Cavalcade. But
the sound of good writing depends even more on the way you put words
together, and there's a connection at that level too.

When writing sounds good, it's mostly because it has good rhythm. But
the rhythm of good writing is not the rhythm of music, or the meter of
verse. It's not so regular. If it were, it wouldn't be good, because the
rhythm of good writing has to match the ideas in it, and ideas have all
kinds of different shapes. Sometimes they're simple and you just state
them. But other times they're more subtle, and you need longer, more
complicated sentences to tease out all the implications.

An essay is a cleaned up train of thought, in the same way dialogue is
cleaned up conversation, and a train of thought has a natural rhythm. So
when an essay sounds good, it's not merely because it has a pleasing
rhythm, but because it has its natural one. Which means you can use
getting the rhythm right as a heuristic for getting the ideas right. And
not just in principle: good writers do both simultaneously as a matter
of course. Often I don't even distinguish between the two problems. I
just think Ugh, this doesn't sound right; what do I mean to say here?
\footnote{Bizarrely enough, this happened in the writing of this very
paragraph. An earlier version shared several phrases in common with the
preceding paragraph, and the repetition bugged me each time I reread it.
When I got annoyed enough to fix it, I discovered that the repetition
reflected a problem in the underlying ideas, and I fixed both
simultaneously.}

The sound of writing turns out to be more like the shape of a plane than
the color of a car. If it looks good, as Kelly Johnson used to say, it
will fly well.

This is only true of writing that's used to develop ideas, though. It
doesn't apply when you have ideas in some other way and then write about
them afterward --- for example, if you build something, or conduct an
experiment, and then write a paper about it. In such cases the ideas
often live more in the work than the writing, so the writing can be bad
even though the ideas are good. The writing in textbooks and popular
surveys can be bad for the same reason: the author isn't developing the
ideas, merely describing other people's. It's only when you're writing
to develop ideas that there's such a close connection between the two
senses of doing it well.

Ok, many people will be thinking, this seems plausible so far, but what
about liars? Is it not notoriously possible for a smooth-tongued liar to
write something beautiful that's completely false?

It is, of course. But not without method acting. The way to write
something beautiful and false is to begin by making yourself almost
believe it. So just like someone writing something beautiful and true,
you're presenting a perfectly-formed train of thought. The difference is
the point where it attaches to the world. You're saying something that
would be true if certain false premises were. If for some bizarre reason
the number of jobs in a country were fixed, then immigrants really would
be taking our jobs.

So it's not quite right to say that better sounding writing is more
likely to be true. Better sounding writing is more likely to be
internally consistent. If the writer is honest, internal consistency and
truth converge.

But while we can't safely conclude that beautiful writing is true, it's
usually safe to conclude the converse: something that seems clumsily
written will usually have gotten the ideas wrong too.

Indeed, the two senses of good writing are more like two ends of the
same thing. The connection between them is not a rigid one; the goodness
of good writing is not a rod but a rope, with multiple overlapping
connections running through it. But it's hard to move one end without
moving the other. It's hard to be right without sounding right.


\textbf{Thanks} to Jessica Livingston and Courtenay Pipkin for reading
drafts of this.
\chapter{The Shape of the Essay Field}

June 2025

An essay has to tell people something they don't already know. But there
are three different reasons people might not know something, and they
yield three very different kinds of essays.

One reason people won't know something is if it's not important to know.
That doesn't mean it will make a bad essay. For example, you might write
a good essay about a particular model of car. Readers would learn
something from it. It would add to their picture of the world. For a
handful of readers it might even spur some kind of epiphany. But unless
this is a very unusual car it's not critical for everyone to know about
it. \footnote{It's hard to write a really good essay about an unimportant
topic, though, because a really good essayist will inevitably draw the
topic into deeper waters. E. B. White could write an essay about how to
boil potatoes that ended up being full of timeless wisdom. In which
case, of course, it wouldn't really be about how to boil potatoes; that
would just have been the starting point.}

If something isn't important to know, there's no answer to the question
of why people don't know it. Not knowing random facts is the default.
But if you're going to write about things that \emph{are} important to
know, you have to ask why your readers don't already know them. Is it
because they're smart but inexperienced, or because they're obtuse?

So the three reasons readers might not already know what you tell them
are (a) that it's not important, (b) that they're obtuse, or (c) that
they're inexperienced.

The reason I did this breakdown was to get at the following fact, which
might have seemed controversial if I'd led with it, but should be
obvious now. If you're writing for smart people about important things,
you're writing for the young.

Or more precisely, that's where you'll have the most effect. Whatever
you say should also be at least somewhat novel to you, however old you
are. It's not an essay otherwise, because an essay is something you
write to figure something out. But whatever you figure out will
presumably be more of a surprise to younger readers than it is to you.

There's a continuum of surprise. At one extreme, something you read can
change your whole way of thinking. \emph{The Selfish Gene} did this to
me. It was like suddenly seeing the other interpretation of an ambiguous
image: you can treat genes rather than organisms as the protagonists,
and evolution becomes easier to understand when you do. At the other
extreme, writing merely puts into words something readers were already
thinking --- or thought they were.

The impact of an essay is how much it changes readers' thinking
multiplied by the importance of the topic. But it's hard to do well at
both. It's hard to have big new ideas about important topics. So in
practice there's a tradeoff: you can change readers' thinking a lot
about moderately important things, or change it a little about very
important ones. But with younger readers the tradeoff shifts. There's
more room to change their thinking, so there's a bigger payoff for
writing about important things.

The tradeoff isn't a conscious one, at least not for me. It's more like
a kind of gravitational field that writers work in. But every essayist
works in it, whether they realize it or not.

This seems obvious once you state it, but it took me a long time to
understand. I knew I wanted to write for smart people about important
topics. I noticed empirically that I seemed to be writing for the young.
But it took me years to understand that the latter was an automatic
consequence of the former. In fact I only really figured it out as I was
writing this essay.

Now that I know it, should I change anything? I don't think so. In fact
seeing the shape of the field that writers work in has reminded me that
I'm not optimizing for returns in it. I'm not trying to surprise readers
of any particular age; I'm trying to surprise myself.

The way I usually decide what to write about is by following curiosity.
I notice something new and dig into it. It would probably be a mistake
to change that. But seeing the shape of the essay field has set me
thinking. What would surprise young readers? Which important things do
people tend to learn late? Interesting question. I should think about
that.

\textbf{Note}

\textbf{Thanks} to Jessica Livingston and Michael Nielsen for reading
drafts of this.

\end{document}
